\documentclass[11pt,letterpaper]{report}
\usepackage{xr-hyper}
\usepackage{scientific_report}
\usepackage{mathtools}
\usepackage{bm}
\usepackage{mathrsfs}
\usepackage{microtype}
\usepackage{longtable}
\usepackage{cleveref}
\usepackage{expl3,xparse}

% The technical documents. Labels are imported with a prefix; links open the
% technical PDFs next to this folder at the cited statement.
\externaldocument[MT-]{main}[../MultiAgentELBO_Theory/main.pdf]
\externaldocument[SG-]{supplement_gaussian}[../MultiAgentELBO_Theory/supplement_gaussian.pdf]
\externaldocument[SB-]{supplement_base}[../MultiAgentELBO_Theory/supplement_base.pdf]

\setlength{\headheight}{22pt}
\renewcommand{\chaptermark}[1]{\markboth{\thechapter.\ #1}{}}
\makeatletter
\renewcommand{\@pnumwidth}{2em}
\renewcommand{\@tocrmarg}{3em}
\makeatother
\usetikzlibrary{arrows.meta}
\fancyhead[R]{\small\textit{A Physicist's Companion}}

\definecolor{techlink}{RGB}{0,110,110}
\hypersetup{
  filecolor=techlink,
  pdftitle={A Physicist's Companion to Gauge-Covariant Variational Free Energy and Exact Coarse-Graining},
  pdfauthor={Robert C. Dennis},
  pdfsubject={Intuitive explanations, notation and jargon decoding, and a status and delivery assessment of the single-point theory and its two supplements},
}

% ---- shared macros of the technical documents (identical definitions) ----
\newcommand{\KL}{D_{\mathrm{KL}}}
\newcommand{\E}{\mathbb E}
\newcommand{\R}{\mathbb R}
\newcommand{\given}{\mid}
\newcommand{\Tr}{\operatorname{Tr}}
\newcommand{\Sym}{\mathbb S}
\newcommand{\PSD}{\mathbb S_{+}}
\newcommand{\GL}{\mathrm{GL}}
\newcommand{\Fenergy}{\mathcal F}
\newcommand{\Lelbo}{\mathcal L}

% ---- cleveref names for the statement types of the technical documents ----
\crefname{definition}{definition}{definitions}
\crefname{lemma}{lemma}{lemmas}
\crefname{proposition}{proposition}{propositions}
\crefname{theorem}{theorem}{theorems}
\crefname{corollary}{corollary}{corollaries}
\crefname{conjecture}{conjecture}{conjectures}
\crefname{openproblem}{open problem}{open problems}
\crefname{hypothesis}{hypothesis}{hypotheses}
\crefname{remark}{remark}{remarks}
\Crefname{definition}{Definition}{Definitions}
\Crefname{lemma}{Lemma}{Lemmas}
\Crefname{proposition}{Proposition}{Propositions}
\Crefname{theorem}{Theorem}{Theorems}
\Crefname{corollary}{Corollary}{Corollaries}
\Crefname{conjecture}{Conjecture}{Conjectures}
\Crefname{openproblem}{Open Problem}{Open Problems}
\Crefname{hypothesis}{Hypothesis}{Hypotheses}
\Crefname{remark}{Remark}{Remarks}
\crefname{subappendix}{appendix}{appendices}
\Crefname{subappendix}{Appendix}{Appendices}

% ---- references into the technical documents ----
% \MT{label[,label...]}  main manuscript, printed "MT Theorem 5.3"
% \SG{label[,...]}       Supplement G, printed "Proposition G4.2"
% \SB{label[,...]}       Supplement B, printed "Theorem B6.1"
% \MTeq, \SGeq, \SBeq    equation numbers, printed "MT (5.3)"
% Lowercase \mtc, \sgc, \sbc give lowercase type names mid-sentence.
\ExplSyntaxOn
\clist_new:N \l__comp_refs_clist
\cs_new_protected:Npn \__comp_prefix:nn #1#2
  {
    \clist_clear:N \l__comp_refs_clist
    \clist_map_inline:nn {#2} { \clist_put_right:Nn \l__comp_refs_clist {#1##1} }
  }
\NewDocumentCommand{\MT}{m}{ \__comp_prefix:nn {MT-} {#1} \textsc{mt}~\exp_args:NV \Cref \l__comp_refs_clist }
\NewDocumentCommand{\mtc}{m}{ \__comp_prefix:nn {MT-} {#1} \textsc{mt}~\exp_args:NV \cref \l__comp_refs_clist }
\NewDocumentCommand{\SG}{m}{ \__comp_prefix:nn {SG-} {#1} \exp_args:NV \Cref \l__comp_refs_clist }
\NewDocumentCommand{\sgc}{m}{ \__comp_prefix:nn {SG-} {#1} \exp_args:NV \cref \l__comp_refs_clist }
\NewDocumentCommand{\SB}{m}{ \__comp_prefix:nn {SB-} {#1} \exp_args:NV \Cref \l__comp_refs_clist }
\NewDocumentCommand{\sbc}{m}{ \__comp_prefix:nn {SB-} {#1} \exp_args:NV \cref \l__comp_refs_clist }
\ExplSyntaxOff
\newcommand{\MTeq}[1]{\textsc{mt}~(\ref{MT-#1})}
\newcommand{\SGeq}[1]{(\ref{SG-#1})}
\newcommand{\SBeq}[1]{(\ref{SB-#1})}

% ---- status badges (what kind of knowledge a result is) ----
\newcommand{\badge}[2]{\tikz[baseline=(b.base)]\node[fill=#1,text=white,rounded corners=1.2pt,inner xsep=2pt,inner ysep=1.4pt,font=\scriptsize\bfseries\sffamily](b){#2};}
\newcommand{\SOLID}{\badge{darkgreen}{SOLID}}
\newcommand{\TOY}{\badge{cautionorange}{TOY}}
\newcommand{\COND}{\badge{accentblue}{COND}}
\newcommand{\NUM}{\badge{recommendpurple}{NUM}}
\newcommand{\OPEN}{\badge{criticalred}{OPEN}}
\newcommand{\VAGUE}{\badge{mediumgray}{VAGUE}}
\newcommand{\CHOICE}{\badge{primaryblue}{CHOICE}}

% ---- delivery verdicts (does the theory do what it says) ----
\newcommand{\DELIVERED}{\badge{darkgreen}{DELIVERED}}
\newcommand{\NARROWER}{\badge{cautionorange}{NARROWER THAN IT SOUNDS}}
% two-line form for narrow table columns
\newcommand{\NARROWERstack}{\tikz[baseline={([yshift=-6.6pt]b.north)}]\node[fill=cautionorange,text=white,rounded corners=1.2pt,inner xsep=2pt,inner ysep=1.4pt,font=\scriptsize\bfseries\sffamily,align=center](b){NARROWER THAN\\IT SOUNDS};}
\newcommand{\PARTIAL}{\badge{cautionorange}{PARTLY DELIVERED}}
\newcommand{\NOTDELIVERED}{\badge{criticalred}{NOT DELIVERED}}

% ---- boxes ----
\newtcolorbox{verdictbox}[1]{enhanced,breakable,colback=lightgray,colframe=primaryblue,
  boxrule=0.6pt,arc=2pt,left=8pt,right=8pt,top=6pt,bottom=6pt,
  fonttitle=\bfseries\sffamily,coltitle=white,
  title={Does it deliver what it claims?\hfill #1}}
\newtcolorbox{analogy}[1]{enhanced,breakable,colback=lightblue,colframe=secondaryblue,
  boxrule=0.4pt,arc=2pt,left=8pt,right=8pt,top=5pt,bottom=5pt,
  fonttitle=\bfseries\sffamily\small,coltitle=primaryblue,colbacktitle=lightblue,
  title={Physics analogue (#1)}}
\newtcolorbox{hygiene}{enhanced,breakable,colback=white,colframe=tableborder,
  boxrule=0.4pt,arc=2pt,left=8pt,right=8pt,top=5pt,bottom=5pt,
  fonttitle=\bfseries\sffamily\small,coltitle=darkgray,colbacktitle=lightgray,
  title={Conditions: which ones matter}}

% ---- scorecard table: Result | What it says | Status | Scope and caveats ----
\newenvironment{scorecard}{%
  \small\renewcommand{\arraystretch}{1.25}%
  \begin{longtable}{@{}>{\raggedright\arraybackslash}p{0.185\textwidth}>{\raggedright\arraybackslash}p{0.365\textwidth}>{\raggedright\arraybackslash}p{0.09\textwidth}>{\raggedright\arraybackslash}p{0.26\textwidth}@{}}
  \toprule
  \textbf{Result} & \textbf{What it says} & \textbf{Status} & \textbf{Scope and caveats}\\
  \midrule\endhead
  \bottomrule\endfoot}
  {\end{longtable}}

% ---- decoder table: Term | Plain meaning | Physics analogue ----
\newenvironment{decoder}{%
  \small\renewcommand{\arraystretch}{1.2}%
  \begin{longtable}{@{}>{\raggedright\arraybackslash}p{0.22\textwidth}>{\raggedright\arraybackslash}p{0.45\textwidth}>{\raggedright\arraybackslash}p{0.25\textwidth}@{}}
  \toprule
  \textbf{Term or symbol} & \textbf{Plain meaning} & \textbf{Physics analogue}\\
  \midrule\endhead
  \bottomrule\endfoot}
  {\end{longtable}}

\begin{document}
\hypersetup{pageanchor=false}
\makereporttitle
  {A Physicist's Companion to\\Gauge-Covariant Variational Free Energy}
  {Intuition, Notation, and an Honest Status Map for the Single-Point Theory and Its Supplements}
  {Robert C. Dennis}
  {Independent Researcher}
  {September 2026}

\hypersetup{pageanchor=true}
\pagenumbering{roman}
\tableofcontents
\clearpage
\pagenumbering{arabic}

\chapter*{How to Read This Companion}
\phantomsection\addcontentsline{toc}{chapter}{How to Read This Companion}
\markboth{How to Read This Companion}{}
\label{comp:guide}

\section*{What this companion is}
\phantomsection\addcontentsline{toc}{section}{What this companion is}

The theory lives in three technical documents: the main manuscript, cited here as \textsc{mt}, which
develops everything at one base point $c_*$; Supplement~G, the multivariate-Gaussian realization; and
Supplement~B, which holds base transport, infinite networks, and the lattice cluster, all deferred
until the single-point theory is settled. Together they run to more than six hundred pages. They are
written for correctness under adversarial checking, not for a first reading, and they are the only
authority. This companion proves nothing. It explains what the technical documents say, in the
language of a physicist, and it says plainly what kind of knowledge each main result is and whether
the theory delivers what its own words promise.

When the technical documents prove something, the companion states it and cites it. A citation such
as \MT{thm:elbo-extended-gap} refers to the main manuscript, and statements of the supplements carry
their own letters (Proposition~G4.2, Theorem~B6.1). In the PDF these references are links that open
the technical document at the cited statement, provided the technical PDFs sit in the neighboring
folder where they are built. The companion's main part mirrors the sixteen chapters of the main
manuscript one to one, including the latent-link chapter (\cref{comp:lie-links}). Three chapters cover
Supplement~G (\cref{comp:sg1,comp:sg2,comp:gaussian-latent-links}), and \cref{comp:sb} covers Supplement~B. The
companion has no numbered statements of its own, so every theorem, proposition, or lemma number you
meet here points into a technical document.

Each companion chapter follows one template: the question the technical chapter answers and its
answer in plain language; the physical picture, with the nearest physics analogue graded as exact,
structural, or loose; a decoder for the chapter's notation and jargon; the key formulas explained
term by term with their limiting cases; a scorecard of the main results; a box that sorts the
chapter's hypotheses into those that matter and those that are hygiene; the open ends; and a
verdict. Where a proof would help intuition, the companion gives a sentence or two of informal
reason and labels it as such.

Highlighted equation boxes mark the central constructions and balances of the dynamic-link model.
Their titles identify the role of each formula; the surrounding text states its hypotheses and
points to the technical result. The Proca-type equation in \cref{comp:lie-links} is a linearized
Euclidean stationarity equation across inference episodes, whereas the recognition flow is first order.

\section*{Two questions about every result}
\phantomsection\addcontentsline{toc}{section}{Two questions about every result}

The first question is what kind of knowledge a result is. Nearly every theorem-type statement in the
main manuscript carries the tag \textsc{established}, meaning it is proved there. The tag says
nothing about how general the result is, how much it depends on a special instance, or whether it
rests on an assumption that might fail in the cases the program cares about. The companion adds that
information with one badge per result.

\begin{sciencetable}{Status badges used in this companion.}
\label{comp:tab-badges}
\begin{tabularx}{\textwidth}{>{\raggedright\arraybackslash}p{0.13\textwidth}>{\raggedright\arraybackslash}X}
\toprule
\textbf{Badge} & \textbf{Meaning} \\
\midrule
\SOLID & Proved in the technical documents, or a checked textbook result, at the generality the theory relies on. Established negative results (refutations and no-go statements) are solid. When the mathematics is textbook and the novelty is only the typing, the companion says ``standard in substance''. \\
\TOY & Proved, but only on a named small or special instance: the four-cycle, a pair, three binary agents, a scalar or abelian case, a particular constant schedule. The mechanism is not established beyond that instance. \\
\COND & Proved under a hypothesis that could fail in the models the program cares about and is not shown to hold there. The companion names the hypothesis. \\
\NUM & Supported only by a named computation. \\
\OPEN & An open problem or a conjecture. \\
\VAGUE & A reading, an interpretation, or a programmatic statement with no precise mathematical statement behind it, or a construction whose meaning is not pinned down. \\
\CHOICE & A definition or standing hypothesis that is a substantive modeling commitment, the place where physics enters, as opposed to notation. \\
\bottomrule
\end{tabularx}
\end{sciencetable}

The second question is whether the theory does what it says. Each companion chapter ends with a
verdict on the technical chapter's own claims, taken from its opening paragraph and from the reading
map of the introduction, measured against what its statements actually prove.

\begin{sciencetable}{Delivery verdicts.}
\label{comp:tab-verdicts}
\begin{tabularx}{\textwidth}{>{\raggedright\arraybackslash}p{0.27\textwidth}>{\raggedright\arraybackslash}X}
\toprule
\textbf{Verdict} & \textbf{Meaning} \\
\midrule
\DELIVERED & The words match the mathematics. \\
\NARROWER & The claim is proved, but in a setting narrower than the headline wording suggests: a supplied hierarchy, a single small network, a particular family. \\
\PARTIAL & Part of the claim is proved and part is open. \\
\NOTDELIVERED & The claim is posed but not proved; the technical document usually says so itself. \\
\bottomrule
\end{tabularx}
\end{sciencetable}

A result that is correct and standard is not a defect, and the companion says ``standard in
substance'' as information, not as criticism: a physicist gains from knowing that an identity is the
Gibbs variational principle in new notation, because then the physicist already knows what it can
and cannot do.

\section*{Reading the technical documents: the house style decoded}
\phantomsection\addcontentsline{toc}{section}{Reading the technical documents: the house style decoded}

The technical documents are written under an authoring specification whose rules explain most of
what makes them hard to read. Knowing the rules lets you skip what is bookkeeping and slow down where
a qualifier carries content.

\paragraph{Status tags and scope paragraphs.}
Every lemma, proposition, theorem, corollary, and claim-bearing remark carries exactly one tag:
\textsc{established} (proved in the document), \textsc{standard} (a known result, cited with a
locator), or \textsc{numerical} (supported only by a named computation). Definitions, standing
hypotheses, conjectures, and open problems carry their status in the heading. After its opening, a
chapter may add one scope paragraph, tagged \textsc{not-claimed}, listing what it declines to claim;
in practice every chapter has one. Declining is
not refuting: a declined statement may be true, and a refuted statement appears instead as an
established negative result. The scope paragraphs are the fastest way to find out what a chapter does
not do, and the companion's verdicts lean on them.

\paragraph{Defaults, and why the prose is full of qualifiers.}
The introduction of the main manuscript fixes defaults once for the whole document
(\MT{sec:intro-conventions}): measurable spaces are standard Borel, laws and kernels are normalized
and measurable, a coarse channel is a Markov kernel that does not read the recognition law and
leaves observations alone, admitted observations have finite positive evidence, structural data are
parameters unless a law for them is declared, and the interaction multigraph contains the reversal
closure of the union of the source-prior supports. A receiving row's support is complete unless
masked and need not be symmetric. In the latent-link model the multigraph and frames remain
structural, while Lie-algebra generators are inferred under a prior
(\MT{def:link-latent-model}). A statement mentions only departures from these defaults, so a
hypothesis that appears inside a statement is a departure and worth reading. Three other sources of
qualifiers are worth knowing. The documents work with arbitrary laws on arbitrary spaces rather than
Gaussians, which forces extended-real values (a relative entropy may be $+\infty$) and almost-sure
qualifiers. They keep apart objects that coincide in the Gaussian or flat cases but not in general:
three description levels (sample, law point, configuration), three transports (frame comparison,
graph link, base parallel transport), three kinds of agreement between agents, three coarse
operations (Markov pushforward, energy precomposition, Galerkin restriction), and five attention objects
(the source prior, the generative and the recognition attention rows, the receiver occupancy, and the
event law).
Keeping them apart is deliberate: each pair can differ outside the Gaussian or flat case, and a
theorem about one does not transfer to the other. And the scope rules described next constrain what
may be called formation, selection, or recursion.

\paragraph{The Curie rule (R13): what may be called selection by learning.}
Curie's principle says that symmetric causes cannot produce less symmetric effects when the effect
is uniquely determined by the cause. The rule applies it to learning. Let $H$ be the symmetries shared
by the generative model with its records, the parameterization, and the kinetics, including the
kinetic metric and the proposal mechanism ($H$ may be a group of graph automorphisms or a residual
gauge group). A claim that learning creates or selects a structure that $H$ does not fix must name
the $H$-breaking datum it amplifies (the initial law, an initial parameter, the model, or the
records) and must compare with the frozen-learning arm on the same data. A claim that learning
selects a structure that $H$ does fix must name the symmetry it restores, with the frozen arm at the
non-invariant initial value as its null. A claim of selection with no breaking datum must name the
hypothesis its mechanism violates and declare the mechanism as part of the model: a rule that reads
realized paths, records, or a finite ensemble; selection among non-unique optimizers, critical
points, or branches; non-uniqueness at masked entries, which amounts to edge creation; or a
non-equivariant metric, mobility, or proposal mechanism. This is substantive, not bookkeeping: an
equivariant law-level learning rule whose coupled law and parameter dynamics is unique cannot break a
symmetry that its initial law and initial parameter share. In physics terms (the companion's
reading, not part of the rule), spontaneous symmetry breaking needs non-uniqueness, which in
statistical mechanics arrives with the thermodynamic limit. The four-cycle results of
\cref{comp:ch14} are written under this rule, which is why they speak of amplifying an initial
asymmetry rather than of spontaneous selection.

\paragraph{Recursion versus a supplied-scale sequence (R14).}
A two-stage claim declares a handoff time, the first- and second-stage events with their blocks, one
statistic of the handoff state, and a class of handoffs fixed in advance that also contains, with
probability bounded below, handoffs of histories in which the first stage did not form. If the
second-stage event occurs with probability tending to one uniformly over the handoff class, the
sequence is a supplied-scale sequence, not recursion, whatever differences appear earlier. Otherwise
the second coalescence counts as recursive only if, after first-stage formation, the second-stage
event has probability bounded below, its block is a union of at least two first-stage blocks, and
the declared statistic exceeds its value after a counterfactual handoff by a margin bounded below. A
margin in the other direction, where the first stage delays or suppresses the second, is recorded as
first-stage dependence, not recursion. In physics terms (again the companion's reading), a
supplied-scale sequence is a hierarchy built into the interaction strengths through which the
dynamics relaxes in order, and the rule separates a genuine cascade from the successive freezing of
modes whose stiffnesses were prescribed. It matters most in \cref{comp:ch13}.

\paragraph{Coherence episodes (R15): what counts as formation.}
A block counts as having formed only if two clauses hold. The restoring clause names the interaction
that holds the block coherent against thermal noise at the stated resolution, with restoring rate
times squared resolution divided by temperature tending to infinity, and the rate is evaluated on the
actual law of the assignment sheets or on the landscape marginalized over sources, not on a sheet
held fixed by hand. The horizon clause requires the
coherence to be a dynamical event rather than a property of equilibrium: either the observation
horizon is short compared with the stationary recurrence time of the event, or the event is more
probable than it is in the stationary state by a margin. Coherence that is restored but merely
reflects equilibrium (a strongly bound pair in thermal equilibrium is always coherent) is called
restored equilibrium coherence; coherence with no restoring interaction is coincidental. In physics
terms, the rule separates ordering produced by the dynamics from equilibrium order and from
fluctuations.

\paragraph{Settling conditions (R16).}
Every open problem and conjecture states a class of models and two outcomes, a proof or a
counterexample within that class, and neither outcome may already be met by an existing statement or
by a trivial construction. This is why the open problems read like contracts. It keeps open items
falsifiable, at the price of posing them in narrower classes than the questions that motivate them.

\paragraph{No Wilsonian vocabulary in the main manuscript.}
The main manuscript constructs exact finite-step coarse-graining with no rescaling, so it does not
use the words relevant, irrelevant, fixed point, or beta function except in one pointer to
Supplement~B and when describing other authors' work. The asymptotic apparatus (rescalings, spectra, fixed points, universality) lives in
Supplement~B. When the companion uses renormalization-group language, it says which of the two it
means.

\paragraph{Labels and names.}
Labels are frozen when a statement changes type, so a label beginning \texttt{thm:} may belong to a
corollary and \texttt{conj:phil-contrast-formation} is a theorem. Trust the printed heading. Words
such as agent, meta-agent, parent, and block are naming decisions: the documents make no claim of
agency, and population laws are not agents.

\section*{A notation survival kit}
\phantomsection\addcontentsline{toc}{section}{A notation survival kit}

The notation table of the main manuscript (\MT{app:notation-contract}) is complete but long. The
following symbols carry most of the load; the glossary at the end of the companion
(\hyperref[comp:glossary]{Glossary}) covers the jargon.

\begin{decoder}
$c_*$ & The single base point (context) at which everything happens. & A zero-dimensional base; one spacetime point. \\
$G$, $\mathscr P_G$ & Structure group and principal bundle; at $c_*$ the fiber is a $G$-torsor (a copy of $G$ with no preferred identity). & Gauge group; the space of frames at a point. \\
$\rho_b$, $\rho_m$ & Representations of $G$ on the belief and model sample spaces; they may differ. & Matter representations. \\
$q_i^b$, $q_i^m$ & Agent $i$'s belief law and model law (probability laws on the fibers). & A mixed state of a matter field at site $i$. \\
$u_i$ & Agent $i$'s frame, one for both channels at $c_*$ (a gauge choice). & A local gauge fixing at site $i$. \\
$\mathfrak I$, $e$, $\bar e$ & Interaction multigraph, an oriented edge copy, its reversal. & The lattice, with multiple links allowed. \\
$\Theta^x_e$ & Independent group-valued link on edge copy $e$ for channel $x\in\{b,m\}$, with $\Theta_{\bar e}=\Theta_e^{-1}$. & Link variable $U_\ell$ of lattice gauge theory. \\
$\mathring\omega^x_e$, $R$ & Generator of the reference twist, $\Theta^x_e=U_i\exp(\mathring\omega^x_e)U_j^{-1}$; recognition law of the generators. & Lie-algebra field and its variational distribution. \\
$\mathring\alpha$, $\mathring\Sigma$ & Anchor and covariance of the Gaussian base prior; generative parameters. & Preferred field and bare restoring scale. \\
$\mathscr S$, $\mathbf m$, $n$ & Loop action, history increment precision, and learning-episode index. & Spatial loop penalty, Euclidean history mass, and an inference index. \\
$\mathrm{Hol}^x(p)$ & Holonomy: product of links around a closed walk $p$. & Wilson loop before the trace. \\
$K_a$, $\mathfrak A$ & A normalized record kernel (an interaction or observation factor) and the set of them. & A factor of the Boltzmann weight; a local term of the action. \\
$\mathbb P_{\theta,V}$ & The one normalized joint generative law of all agents and records. & The Gibbs measure $e^{-S}/Z$. \\
$\mathbb Q_{V,o,X}$ & The one recognition law of the whole population (a coupling of the agents' local laws). & The variational trial state. \\
$\boldsymbol\Pi_{\theta,V,o,X}$ & The exact posterior given observation $o$ and structure $X$. & The exact Gibbs state conditioned on the data. \\
$\Lelbo_V$, $\Fenergy_V=-\Lelbo_V$ & Evidence lower bound and variational free energy of the population. & Minus the variational free energy; the Gibbs--Bogoliubov functional. \\
$\pi_{ij}$, $\beta_{ij}$, $\tau_i$ & Source prior, attention weight of receiver $i$ on source $j$, attention temperature. & Prior weight, Boltzmann weight over sources, temperature. \\
$S^Q_e$, $S^K_e$, $\mathsf L_e$ & Query and key comparison maps and the composite communication map of an edge. & None standard; linear maps between internal spaces. \\
$C_A$, $\ell$ & Coarse channel that makes parent $A$ from its child block; coarse-graining depth. & Blocking kernel; step index of a renormalization without rescaling. \\
$\mathrm T$, $\varepsilon$ & Temperature and small parameters of the formation families (\cref{comp:ch13,comp:ch14}). & Temperature; small expansion parameters. \\
\end{decoder}

\section*{Reading paths}
\phantomsection\addcontentsline{toc}{section}{Reading paths}

For a first pass, read \hyperref[comp:glance]{The Theory at a Glance} and the verdict box at the end of each chapter. For the
physical core of the theory, read \cref{comp:ch02,comp:ch05,comp:ch06,comp:lie-links,comp:ch09,comp:ch11,comp:ch13,comp:ch14}
in that order. For what the theory says about Gaussian networks with $\GL(K)$ links at large size,
read \cref{comp:ch02,comp:lie-links,comp:ch10,comp:sg1,comp:sg2,comp:gaussian-latent-links}. For the state of the renormalization program at one
base point, read \cref{comp:ch09,comp:ch11,comp:ch12,comp:sb}. The open ends of the whole theory are
collected in \hyperref[comp:open-ends]{Open Ends of the Theory}.

\chapter*{The Theory at a Glance}
\phantomsection\addcontentsline{toc}{chapter}{The Theory at a Glance}
\markboth{The Theory at a Glance}{}
\label{comp:glance}

\section*{The theory in one page}
\phantomsection\addcontentsline{toc}{section}{The theory in one page}

At one base point $c_*$, a population of agents is a lattice gauge theory on a graph whose site
variables are probability laws. Each agent carries a belief law and a model law on spaces where a
structure group $G$ acts, and a frame saying in which gauge those laws are written; each edge carries
an independent group-valued link, so products of links around loops carry holonomy. Agents interact
through normalized record kernels. From these local data the theory builds one normalized joint law,
one recognition law for the whole population, and one evidence bound, which is the Gibbs variational
principle applied to that joint. Attention is a posterior over which neighbor a receiver listens to,
and selective communication adds query, key, and value maps between the agents' internal spaces.
Fisher geometry turns free-energy descent into oriented histories with an information clock, once a
metric is chosen, and a fast-slow split separates beliefs from models.

Coarse-graining is a Markov kernel from a block of agents to a parent. Applied to the joint, the
posterior, and the recognition law at once, it preserves the evidence and splits the free energy
exactly into the parent's free energy plus the information discarded inside the block. The parent's
law therefore always exists, with an exact effective action. What coarse-graining does not give is a
closed flow on a finite set of couplings, a thermodynamic limit, or a proof that the parent evolves as
an autonomous process. Formation of meta-agents is proved in two settings. In a family with a
hierarchy of interaction strengths supplied by hand, exact parents track a prepared state as it
relaxes through the hierarchy. On a four-agent ring, a pair forms by feedback between attention and
belief relaxation, and slow learning selects one of the two pairings by amplifying an asymmetry of
the initial law. Formation that the agents' coupled inference produces on its own, beyond one level
and beyond four agents, is open.

\section*{Where the chapters stand}
\phantomsection\addcontentsline{toc}{section}{Where the chapters stand}

\colorlet{VCOLch02}{lightgreen}\colorlet{VCOLch03}{lightgreen}\colorlet{VCOLch04}{lightgreen}
\colorlet{VCOLch05}{lightgreen}\colorlet{VCOLch06}{lightgreen}\colorlet{VCOLch07}{lightorange}
\colorlet{VCOLch08}{lightorange}\colorlet{VCOLch09}{lightgreen}\colorlet{VCOLch10}{lightgreen}
\colorlet{VCOLch11}{lightorange}\colorlet{VCOLch12}{lightgreen}\colorlet{VCOLch13}{lightorange}
\colorlet{VCOLch14}{lightorange}\colorlet{VCOLsg}{lightgreen}\colorlet{VCOLsb}{lightorange}
\begin{figure}[htbp]
\centering
\begin{tikzpicture}[
  every node/.style={font=\sffamily\scriptsize,align=center},
  box/.style={draw=primaryblue,rounded corners=2pt,minimum width=3.1cm,minimum height=0.95cm,inner sep=3pt,line width=0.6pt},
  arr/.style={-{Latex[length=2mm]},draw=darkgray,line width=0.5pt},
  sup/.style={box,draw=mediumgray,dashed},
]
\node[box,fill=VCOLch02] (c2) at (0,0) {2\enspace Geometry at $c_*$\\bundles, links, holonomy};
\node[box,fill=VCOLch03] (c3) at (0,-1.35) {3\enspace Probability semantics\\kernels, levels, chain rule};
\node[box,fill=VCOLch04] (c4) at (0,-2.7) {4\enspace Normalized joint law\\typing prohibition};
\node[box,fill=VCOLch05] (c5) at (0,-4.05) {5\enspace One evidence bound\\local--global identity};
\node[box,fill=VCOLch06] (c6) at (4.3,-2.7) {6\enspace Attention and\\selective communication};
\node[box,fill=VCOLch07] (c7) at (4.3,-4.05) {7\enspace Fiber information\\geometry, histories};
\node[box,fill=VCOLch08] (c8) at (4.3,-5.4) {8\enspace Fast beliefs,\\slow models};
\node[box,fill=VCOLch09] (c9) at (0,-6.75) {9\enspace Coarse channels\\exact coarse free energy};
\node[box,fill=VCOLch10] (c10) at (0,-8.1) {10\enspace Network state\\fixed sectors, holonomy};
\node[box,fill=VCOLch11] (c11) at (4.3,-8.1) {11\enspace Exact network\\coarse-graining};
\node[box,fill=VCOLch12] (c12) at (4.3,-9.45) {12\enspace Multiscale VFE,\\closing the tower};
\node[box,fill=VCOLch13] (c13) at (8.6,-9.45) {13\enspace Formation under a\\supplied hierarchy};
\node[box,fill=VCOLch14] (c14) at (8.6,-8.1) {14\enspace Four-cycle selection\\and parent activation};
\node[sup,fill=VCOLsg] (sg) at (8.6,-2.7) {Supplement G\\Gaussian realization};
\node[sup,fill=VCOLsb] (sb) at (8.6,-5.4) {Supplement B\\base, infinite networks,\\lattice cluster};
\draw[arr] (c2) -- (c3);
\draw[arr] (c3) -- (c4);
\draw[arr] (c4) -- (c5);
\draw[arr] (c5) -- (c6);
\draw[arr] (c5) -- (c7);
\draw[arr] (c7) -- (c8);
\draw[arr] (c5) -- (c9);
\draw[arr] (c9) -- (c10);
\draw[arr] (c10) -- (c11);
\draw[arr] (c9) -- (c11);
\draw[arr] (c11) -- (c12);
\draw[arr] (c12) -- (c13);
\draw[arr] (c8) -- (c13);
\draw[arr] (c13) -- (c14);
\draw[arr,dashed,draw=mediumgray] (sg) -- (c6);
\draw[arr,dashed,draw=mediumgray] (sg.south) -- (sb.north);
\draw[arr,dashed,draw=mediumgray] (sb) -- (c11);
\end{tikzpicture}
\caption{How the chapters depend on one another. Fill color gives the companion's verdict on each
technical chapter's own claims: green, delivered as worded; orange, delivered but narrower than the
wording. Dashed boxes are the supplements, on which no main-manuscript result depends. Chapter~1 is
the program statement, whose verdict is ``partly delivered'' (\cref{comp:ch01}), and chapter~15 is
related work.}
\label{comp:fig-map}
\end{figure}

\Cref{comp:fig-map} shows the chapters in dependency order, colored by verdict. No chapter is ``not
delivered'': every chapter proves what its statements say, and the open problems are posed as open.
Where a verdict is ``narrower than it sounds'', the gap is between an opening paragraph or a title
and the theorems, and the companion chapter names the sentence. The theory's unproved ambitions sit
in its open problems, which \hyperref[comp:open-ends]{Open Ends of the Theory} collects; the two
``not delivered'' rows of the table below are two of those ambitions, which the technical documents
themselves pose as open.

\clearpage
\section*{The headline claims}
\phantomsection\addcontentsline{toc}{section}{The headline claims}

\begingroup
\small\renewcommand{\arraystretch}{1.25}
\begin{longtable}{@{}>{\raggedright\arraybackslash}p{0.34\textwidth}>{\raggedright\arraybackslash}p{0.19\textwidth}>{\raggedright\arraybackslash}p{0.39\textwidth}@{}}
\toprule
\textbf{Claim (where)} & \textbf{Verdict} & \textbf{In one sentence}\\
\midrule\endhead
\bottomrule\endfoot
One evidence bound for the whole population, tight exactly at the posterior (\cref{comp:ch05}) & \DELIVERED & The Gibbs variational principle for a correctly normalized joint, written in the extended reals; standard in substance, new in its typing. \\
Local objectives are exact conditional pieces of the collective one (\cref{comp:ch05}) & \DELIVERED & The chain rule for relative entropy with the rest of the population held fixed; the ``exact potential game'' needs singleton blocks with product profiles. \\
Everything is gauge covariant at one base point (\cref{comp:ch02,comp:ch04}) & \DELIVERED & Evidence and bound are invariant under passive frame changes when kernels are covariant; links carry per-edge twist classes and loop holonomy. \\
Attention is inference (\cref{comp:ch06}) & \DELIVERED & The softmax is Bayes' rule for a declared latent pointer; the transformer's free value and output maps are replaced by a tied value map in the exact update. \\
Information geometry and inference histories (\cref{comp:ch07}) & \NARROWERstack & Fisher contraction and the information clock are proved; they cover a meta-agent only when its law is a parameter-independent Markov image, and every orbit rests on a chosen metric. \\
Collective free energy does not factor through marginals (\cref{comp:ch07}) & \DELIVERED & An exact refutation by a three-agent parity witness. \\
Fast beliefs, slow models (\cref{comp:ch08}) & \NARROWERstack & The reduced slow flow is the kinematics of adiabatic elimination; that the actual flow tracks it is proved only on the four-cycle's held sheets. \\
Coarse channels preserve evidence and split the free energy exactly; a parent law always exists (\cref{comp:ch09}) & \DELIVERED & The relative-entropy chain rule for one channel pushing three laws; exact, static, and with no closure or selected channel. \\
Holonomy decides what a coarse description can retain (\cref{comp:ch10}) & \DELIVERED & Fixed sectors and a holonomy-conditioned parent score; for generic $\GL(K)$ links the Gaussian score is infinite on any block with a loop. \\
Exact coarse-graining of the agent network (\cref{comp:ch11}) & \NARROWERstack & The effective action exists for every channel, but closure keeps the complete laws, the coupling calculus excludes Gaussian $\GL(K)$ couplings, everything is finite-network, and there is no thermodynamic limit. \\
One free energy for the multiscale tower; parents influence children (\cref{comp:ch12}) & \DELIVERED & As accounting and typing; parent influence is proved as a typed possibility, not as dynamics. \\
Successive meta-agents form under free-energy descent with learning (\cref{comp:ch13}) & \NARROWERstack & Exact parents track linear relaxation through a hierarchy supplied by hand; attention and learning play no role in the coherence conclusion. \\
A pair forms endogenously on a four-cycle (\cref{comp:ch13}) & \DELIVERED{} \TOY & One pairing episode on four agents, seeded by the random initial assignment, with no learning. \\
Slow learning selects a pairing (\cref{comp:ch14}) & \DELIVERED{} \TOY & A law-level theorem on the four-cycle that amplifies an asymmetry of the initial law; not spontaneous symmetry breaking. \\
Parents activated after formation behave well (\cref{comp:ch14}) & \DELIVERED{} \TOY & On the four-cycle, activation at a deterministic time after the fast modes thermalize works; activation at the first completion is refuted. \\
A parent evolves as an autonomous process (\cref{comp:ch11}) & \NOTDELIVERED & Posed as the defining open problem of the phase, \MT{open:rg-pointwise-parent-dynamics}. \\
Collectives form through the agents' coupled inference, recursively (\cref{comp:ch01,comp:ch13}) & \NOTDELIVERED & The program aim; posed as \MT{open:rg-formation-general} and \MT{open:full-graph-hierarchy}. \\
Gaussian realization: aggregation, Kron reduction, holonomy (\cref{comp:sg1,comp:sg2}) & \DELIVERED & Aggregation is the exact algebraic-multigrid coarse operator; Kron closure is classified; the interaction family is declared, not derived, and complex Gaussians are not treated. \\
Asymptotic renormalization and the lattice cluster (\cref{comp:sb}) & \NARROWERstack & Finite depth gives no asymptotics; the $\mathbb Z_2$ lattice cluster is a genuine constructive result whose sites are not yet typed as agents. \\
\end{longtable}
\endgroup

\section*{What a physicist should take away}
\phantomsection\addcontentsline{toc}{section}{What a physicist should take away}

\paragraph{The foundations are right, and mostly standard in substance.}
The evidence bound, the local objectives, gauge covariance, and exact coarse-graining are proved at
full generality, and most of them are known identities (the Gibbs variational principle, the chain
rule for relative entropy, change of variables, pushforward densities) written in a typed language.
The typing is where the value lies: it separates objects that coincide only in Gaussian or flat
cases, and it prevents collective objects from being sums that merely look global. The one
foundational gap with consequences is that belief-alignment divergences, the attention energies the
program relies on, are law-point energies without a normalized prior over law points
(\MT{open:prob-law-point-prior}), so they are not yet terms of one evidence bound at that level
(\cref{comp:ch03}).

\paragraph{The renormalization at one base point is exact but not closed.}
Every coarse channel produces an exact parent law and an exact effective action. There is no closed
truncation, no flow on a finite set of couplings (a finite-network obstruction rules one out across
reducing steps), no thermodynamic limit, and no autonomous parent dynamics. None of the main
manuscript's open problems takes the limit of many agents, so the thermodynamic-limit view of
renormalization at $c_*$ is not yet posed there (\cref{comp:ch11}, \hyperref[comp:open-ends]{Open
Ends of the Theory}).

\paragraph{Formation is proved only where it is supplied or small.}
The finite-depth theorem shows that exact parents survive a relaxation through a supplied hierarchy
of stiffnesses; it is a statement about tracking, not about emergence. The endogenous results live
on four agents, where selection by learning amplifies an initial asymmetry, as Curie's principle
requires of a unique law-level dynamics. By the small-network rule these results ground the
mechanisms; they neither establish nor exclude formation at scale, and they say nothing for or
against phase transitions, which need the many-agent limit (\cref{comp:ch13,comp:ch14}).

\paragraph{Noncompact groups meet specific friction.}
For the program's target of hundreds of agents with $\GL(K,\R)$, $\GL(K,\mathbb C)$, or $\mathrm{SU}(N)$
links, the documents record several boundaries. A normalized law on noncompact links is not
constructed (\cref{comp:ch02}). The holonomy-conditioned Gaussian parent score is infinite for
generic $\GL(K)$ loops (\cref{comp:ch10}). The coupling calculus of exact network coarse-graining
is proved only in a bounded tier that excludes Gaussian $\GL(K)$ couplings (\cref{comp:ch11}). Gauge
absorption of Gaussian gains is proved on forests and pairs, with only sign conditions on cycles, and
complex Gaussians are not treated (\cref{comp:sg1}). No symmetric rule that ignores a reference point
averages frames in $\GL^+(K)$, so a meta-agent's frame must be chosen with a reference
(\cref{comp:ch10}). $\mathrm{SU}(N)$, being compact, avoids the first two.

\paragraph{What is new.}
Beyond the typing, the new mathematics is concentrated in a few places: the per-edge twist invariants
that a shared fiber gives links at one point (\cref{comp:ch02}), the holonomy-conditioned fixed
sectors and parent score (\cref{comp:ch10}), the parity no-go for factoring collective free energy
(\cref{comp:ch07}), the Kron closure classification (\cref{comp:sg2}), the formation machinery that
carries exact parents through a relaxation (\cref{comp:ch13}), and the four-cycle selection and
activation-clock results (\cref{comp:ch14}).

\section*{How far this companion was checked}
\phantomsection\addcontentsline{toc}{section}{How far this companion was checked}

The companion's chapters were drafted by Claude Opus 5.5 against the technical statements. A second
model, Claude Fable 5.1, checked the reading guide, chapters~1, 5, 6, 9, 10, and~18, and the open
ends against the technical documents, and its findings were applied: none of them changed a verdict,
and several corrected paraphrases that were stronger than their sources. The remaining chapters
(2~to~4, 7, 8, 11~to~17), this overview, and the glossary were not checked by a second model. Their
drafters verified every cited label and recomputed many numerical witnesses, and the central claim
of \cref{comp:ch13}, that attention and learning play no role in the coherence conclusion of the
formation theorem, was checked against its source. Read the unchecked chapters with that in mind,
and trust the technical statement over the companion wherever they differ.


\part{The Main Manuscript}
\chapter{Program, Scope, and Standing Conventions}
\label{comp:ch01}

\section{The question and the answer}
\label{comp:ch01-question}

The main manuscript asks two questions. What collective inference is available to a population of
agents, each a local section of a belief bundle and a model bundle of one principal bundle, all
placed at one base point $c_*$ and interacting through independent links under one normalized joint
law? And which parts of that inference survive when agents are coarse-grained exactly into
meta-agents? Its answer, announced in the opening of \MT{ch:introduction}, has four parts. The whole
population has one evidence bound, whose gap vanishes exactly when the recognition law equals the
posterior (\MT{thm:elbo-extended-gap}, \MT{cor:elbo-bound-tightness}). Each block of agents has a
local objective, a conditional free energy, such that when the block's conditional law is replaced
with the rest of the population held fixed, the change of the collective free energy equals the
outside-averaged change of the block's local objective (\MT{thm:obs-local-global-potential}); holding
the outside marginal fixed is load-bearing. Coarse channels preserve the evidence and split the free
energy exactly (\MT{thm:rg-exact-coarse-vfe}), with an exact effective action for what they retain
(\MT{thm:rg-effective-action}). And a finite-depth family of models forms successive meta-agents
under free-energy descent with learning, in a stated scope (\MT{thm:finite-depth-adaptive-formation}).
The program-aim section announces two further results on a homogeneous four-cycle: seeded formation
from every Gaussian initial law with source learning idle (\MT{prop:phil-seeded-contrast-formation}),
and law-level selection of one pairing by slow source learning for almost every Gaussian initial
mean (\MT{conj:phil-contrast-formation}, a theorem despite its label). The chapter also names what is
not done: how a parent's own dynamics are determined by the network it summarizes is
\MT{open:rg-pointwise-parent-dynamics}, which it calls the defining open problem of the present
phase.

\section{The picture}
\label{comp:ch01-picture}

Think of a lattice gauge theory on a graph whose sites all sit at one spacetime point. Each site
carries an agent. What the agent carries is not a field value but two probability laws, a belief
$q_i^b$ and a model $q_i^m$, on spaces on which the structure group $G$ acts through representations
$\rho_b$ and $\rho_m$; the agent also carries a frame $u_i$, a point of the one $G$-torsor over $c_*$
that says in which gauge the agent's laws are written. Each edge of the interaction graph carries a
group element $\Theta_e$ for each channel, an independent variable rather than something computed from
the frames, so products of links around loops (holonomies) can be nontrivial. Agents interact through
normalized record kernels $K_a$, which play the role of local factors of a Boltzmann weight.

From these local data the manuscript builds three global objects, in this order: one normalized joint
generative law $\mathbb P_{\theta,V}$ of everything, one recognition law $\mathbb Q_{V,o,X}$ of the
whole population, and one evidence bound $\Lelbo_V$. The order matters to the author, and it is the
first commitment of the program: every collective object is a complete probability object,
constructed from the agents' mechanisms and records or selected explicitly, never an extra agent added
to the population and never a sum of local objectives that happens to look like a global one.

\begin{analogy}{exact for the variational principle, structural for the gauge part}
The evidence identity below is the Gibbs variational principle of statistical mechanics: for any trial
state $q$, $\log Z=-F[q]+\KL(q\Vert p)$ with $F[q]=\E_q[\text{energy}]-\text{entropy}(q)$. Here the
``energy'' is $-\log$ of the joint generative density and $Z$ is the evidence. The gauge part is lattice
gauge theory over a zero-dimensional base: frames are local gauge choices, links are link variables,
holonomies are untraced Wilson loops. Two differences matter. The site variables are probability laws
on representation spaces, not field values, and the links need not carry a kinetic (plaquette) term;
whether they are fixed parameters or dynamical variables is declared model by model.
\end{analogy}

\section{Objects and notation}
\label{comp:ch01-objects}

\begin{decoder}
Agent $\mathcal A^i=(\mathcal C_i;q_i^b,q_i^m,u_i)$ & A belief section, a model section, and one common frame (\MT{def:geo-agent}); at one base point, two laws and a frame. & A site with its matter content and gauge choice. \\
Common frame $u_i$ & One frame for both channels; a gauge choice, not an assumption (\MT{prop:geo-common-frame-gauge}). & Gauge fixing at a site. \\
Cross morphism $\Phi$ & A map from belief fibers to model fibers (or back); extra structure, not implied by the bundles (\MT{def:geo-cross-morphisms}). & A coupling between two matter species. \\
Record kernel $K_a$ & A normalized conditional law on the scope of an interaction or observation. & A local factor of the Boltzmann weight. \\
Evaluator & The map that turns a model presentation $m_i$ into a generative kernel; a model is not itself a law. & The map from coupling constants to a Hamiltonian. \\
Constructed joint $\mathbb P_{\theta,V}$ & The normalized product of all kernels (\MT{prop:gen-exact-normalization}). & The Gibbs measure. \\
Selected coupling $\mathbb Q_{V,o,X}$ & One joint recognition law whose marginals are the agents' local laws, chosen by a declared measurable rule (\MT{hyp:prob-measurable-coupling-selector}). & A correlated trial state; not a product state unless chosen so. \\
Posterior slice $\boldsymbol\Pi_{\theta,V,o,X}$ & The exact posterior given the observation $o$ and the structural data $X$. & The exact conditioned Gibbs state. \\
Structural data $X$ & Graph, masks, and other structure, held as parameters unless a law for them is declared. & Quenched background. \\
\end{decoder}

A notation trap: the manuscript's $T_{ij}$ is a frame comparison, not a temperature or a transfer
matrix, and the formation chapters write their temperature as $\mathrm T$. The letter $\pi$ is both
the bundle projection and, with indices, the source prior of attention.

\section{The formulas, explained}
\label{comp:ch01-formulas}

The architecture display \MTeq{eq:intro-local-to-global-architecture} is a flow chart:
\[
\{\mathcal A^i,G_{\theta,i},\mathbb Q_{i,o,X},K_a\}_{i\in V,a\in\mathfrak A}
\longrightarrow
\bigl(\mathbb P_{\theta,V}\text{ by construction},\ \mathbb Q_{V,o,X}\in\operatorname{Cpl}(\mathbb Q_{\bullet,o,X})\text{ by selection}\bigr)
\longrightarrow\Lelbo_V .
\]
On the left are the local data: agents, their generative mechanisms, their local recognition laws,
and the record kernels. The joint law is built by multiplying normalized kernels; the population's
recognition law is chosen among the couplings $\operatorname{Cpl}$ of the local laws, the joint laws
with those marginals. The selection step is not trivial. Once two agents have non-degenerate laws, the
marginals do not determine the joint, and the correlations the selection puts in are exactly what a
mean-field (product) choice leaves out.

The one identity \MTeq{eq:elbo-intro-identity},
\[
\log p_{\theta,V}(o\given X)=\Lelbo_V(\mathbb Q_{V,o,X};X,o)+\KL\bigl(\mathbb Q_{V,o,X}\big\Vert\boldsymbol\Pi_{\theta,V,o,X}\bigr),
\]
reads as follows. The left side is the log evidence, the log partition function of the system
conditioned on the observed data. The first term on the right is the evidence lower bound, minus the
variational free energy of the population's recognition law. The second term is the relative entropy
from that recognition law to the exact posterior: nonnegative, zero exactly when the two agree as
measures. Two consequences are spelled out in the chapter. Getting every single-agent marginal right
does not close the gap when the posterior is correlated (\MT{cor:elbo-bound-tightness}); the missing
piece is the correlation structure. And the identity does not change when agents re-express their
laws in other frames, provided the kernels transform covariantly
(\MT{prop:gen-product-evidence-invariance}, \MT{cor:elbo-independent-frame-invariance}): the evidence,
like a partition function, is gauge invariant. A subtler point concerns restricted families: if the
posterior is excluded from the recognition family, every individual gap is positive, but the best gap
over the family may still be zero, approached without being attained (\MT{sec:restrict-principle}).

\section{The standing conventions, and which of them matter}
\label{comp:ch01-conventions}

\MT{sec:intro-conventions} fixes defaults for the whole manuscript, so that later statements mention
only departures. The reading guide lists them. Three conventions are physics rather than hygiene.

The first is that the interaction graph at $c_*$ is the support of the source prior of attention, a
structural parameter that is the complete graph unless masked. Gibbs attention weights are then
strictly positive on every allowed edge whose comparison cost is finite (the extended-cost rows of
\MT{prop:full-graph-gibbs-rows} may put zero weight on an infinite-cost edge), and all dynamic structure
is carried by attention weights and link values. Creating or deleting edges is not part of the model;
it is \MT{open:obs-edge-creation}. Loosely, this fixes the set of possible couplings the way a
canonical ensemble fixes the particle number: only coupling strengths and link values move, never
the graph. Induced coarse multigraphs are derived objects.

The second is that three transports must stay distinct. A frame comparison $T_{ij}$ relates two frames
in the same fiber: because all agents sit over the same point, their frames are points of one torsor
and differ by a unique group element, with no path involved (\MT{prop:geo-frame-comparisons}). A link
$\Theta_e$ is an independent variable on an edge (\MT{def:geo-graph-links}); setting every link equal to
the frame comparison gives frame-aligned, holonomy-free links, and general links carry holonomy that no
frame change removes. Parallel transport along a curve in the base is not used at all. Graph links are
therefore not base transport, and graph holonomy is neither the holonomy of a base connection nor a
topological property of the bundle.

The third is that ``coarse'' has distinct meanings that the manuscript keeps apart: a Markov channel,
which discards sample information while preserving normalization; a hard restriction, which
substitutes a lower-dimensional configuration into an energy and defines a new model; and parameter
aggregation inside a closed family. \MT{sec:cg-three-operations}, where the three are defined, calls
them Markov pushforward, energy precomposition, and Galerkin restriction (the compression of an
operator, which is how aggregation acts on a Gaussian precision). They can agree in special
realizations (Gaussian ones are worked out in Supplement~G), and a theorem about one does not
transfer to another.

\begin{hygiene}
Standard Borel spaces, measurability of laws and kernels, and the finite positive evidence of admitted
observations are hygiene; every model the program uses satisfies them. Absolute continuity of the
recognition law with respect to the posterior is deliberately not a default: it is clause (H3) of
\MT{hyp:elbo-evidence-domain}, and when it fails the bound is $-\infty$. This matters in practice,
because a recognition law supported on a lower-dimensional set (for example, one that identifies
agents exactly) is singular and costs an infinite free energy. The decision that structural data are
parameters unless a law is declared matters too: every statement about random graphs, random
partitions, or random memberships needs that law to be named.
\end{hygiene}

\section{What the manuscript declines to claim}
\label{comp:ch01-scope}

The scope paragraph of \MT{sec:intro-conventions} declines several claims for the whole manuscript. No
index of the theory is identified with physical time: not the inference step, the flow parameter, the
coarse-graining depth, or the Fisher duration. The chapter states what such an identification would
need: a named target clock, a calibration map with an uncertainty model, causal and compositional
consistency, and evidence against alternative parameterizations. The base is not spacetime and
carries no Lorentzian structure. Agents, laws, and meta-agents are mathematical objects, and calling a
node an agent is a naming decision. No canonical rule for choosing a recognition coupling, a
partition, or a membership is asserted. No renormalization semigroup with rescaling, fixed points, or
beta functions is constructed in the main manuscript; that apparatus is in Supplement~B
(\cref{comp:sb}). Identifying any limiting object with a physical law is posed as
\MT{claim:physical-law-identification}. Each of these is declined, not refuted.

\section{The program's commitments}
\label{comp:ch01-commitments}

\begin{scorecard}
\MT{sec:intro-program-aim} & Every collective object is a complete probability object built from the agents' mechanisms and records or selected explicitly. & \CHOICE & Methodological; it excludes mean-field sums that only look global. \\
\MT{sec:intro-program-aim} & Every statement about a parent holds for the normalized parent law that coarse-graining produces. & \CHOICE & Rules out ``effective'' parents defined by hand. \\
\MT{sec:intro-program-aim} & Formation of a parent is a theorem about a declared family, with its scope stated. & \CHOICE & The formation results are held to this; see \cref{comp:ch13,comp:ch14}. \\
\MT{sec:intro-conventions} & One base point; no base maps, neighborhoods, or connection descent. & \CHOICE & A phase decision; base transport is deferred to Supplement~B. \\
\MT{sec:intro-conventions} & Interaction graph is the support of the source prior, complete unless masked; no edge creation. & \CHOICE & Edge creation is \MT{open:obs-edge-creation}. \\
\MT{sec:intro-conventions} & Frame comparison, graph link, and base transport are distinct. & \CHOICE & Links are independent variables; holonomy is network-loop holonomy. \\
\end{scorecard}

\section{Open ends}
\label{comp:ch01-open}

The chapter names one open problem as the defining one of the phase: how a parent's dynamics follow
from the network it summarizes (\MT{open:rg-pointwise-parent-dynamics}). The coarse-graining of
\cref{comp:ch09,comp:ch11} produces the parent's law at each instant exactly; what is missing is a
theorem that this law evolves as an autonomous process of its own. The other open ends of the program
are collected in \hyperref[comp:open-ends]{Open Ends of the Theory}.

\begin{verdictbox}{\PARTIAL}
The chapter claims four results and names one open problem. Three of the four are proved as stated:
the one evidence bound, the local-global potential identity, and exact coarse-graining with an
effective action. ``Exact'' there means an identity for one channel, not a closed renormalization
flow; the chapter does not claim a flow, and \cref{comp:ch11} explains what is and is not closed. The
fourth, formation, is proved in a narrower sense than the program aim suggests: the finite-depth
theorem relaxes a prepared state through a hierarchy of interaction strengths supplied by hand, and
the endogenous witness, seeded formation, and the pairing-selection theorem all live on the
four-cycle (\cref{comp:ch13,comp:ch14}); the selection theorem is the strongest learning result, a
law-level statement for almost every initial mean, and it amplifies an asymmetry of the initial law
rather than creating one. The program aim
that parents act back on their children through normalized laws is set up (\cref{comp:ch12}) but not
closed, because the parent's own dynamics is the open problem the chapter names. The chapter is
candid about this; its opening states the scope of the formation theorem and names the open problem.
The program-aim paragraph, which speaks of collectives forming through the agents' coupled inference,
describes the goal rather than the delivered result.
\end{verdictbox}

\chapter{Geometry at One Base Point}
\label{comp:ch02}

\section{The question and the answer}
\label{comp:ch02-question}

The technical chapter asks what geometric data an agent carries at the single base point $c_*$ at which the
whole main manuscript works, and which of those data are invariant, meaning unchanged by any change in how
the data are written down. It gives three answers. The group of genuine symmetries at $c_*$, the active
gauge group, is one copy of the structure group $G$ sitting inside the much larger group $\prod_{i\in V}G$
of independent per-agent frame changes, which the chapter calls passive (\MT{prop:geo-active-image-point}).
Every independent link between two agents carries a conjugacy class that neither passive nor active changes
can alter, and the holonomy of a closed loop of agents is the telescoped product of the per-edge twists whose
classes these are (\MT{prop:geo-edge-twist-invariance}); the classes alone do not determine the holonomy.
Three notions of agreement on a block of agents, frame agreement, agreement of an equivariant statistic
such as the mean, and law coherence, are logically distinct, and the chapter proves exactly which
implications hold among them (\MT{prop:geo-three-agreements}). These objects are used downstream: passive
re-expression is the transformation under which the generative law must be covariant (\cref{comp:ch04}),
link-transported comparisons are what attention compares (\cref{comp:ch06}), and holonomy together with the
trivialization criterion becomes the fixed sectors of coarse channels (\cref{comp:ch10}).

Is the principal- and associated-bundle machinery doing real work at one base point? Only as typing. At
$c_*$ the principal bundle enters only through its fiber $(\mathscr P_G)_{c_*}$, which is a $G$-torsor,
that is, a copy of $G$ with no preferred identity element; the belief and model bundles enter only through
their fibers, which are $G$-spaces of probability laws; and a section enters only through its value at
$c_*$, which is a law. Nothing in the chapter uses a connection, a base curve, the topology of
$\mathscr P_G$, or the manifold structure of the context space $\mathcal C$, and the manuscript says so
itself: graph flatness is a property of the links alone, whatever the topology of $\mathscr P_G$
(\MT{rem:geo-flatness-strictness}), and graph-link holonomy is not the holonomy of any base connection
(\MT{sec:intro-conventions}). What the bundle language contributes is a clean separation between coordinates
(frames, link coordinates, sample coordinates) and frame-independent objects (link maps, points of the
associated fibers), and both unfamiliar statements of the chapter come out of that separation. The link sector
itself is lattice gauge theory on the interaction multigraph, with the lattice vertex law, Wilson loops, and the
lattice criterion for pure gauge; what is added is a group-valued frame at every agent, all frames lying in one
shared fiber. Transport along the base, which would make the bundle language indispensable, is deferred to
Supplement B.

\section{The picture}
\label{comp:ch02-picture}

\subsection{One fiber, many frames}

All agents sit at the same base point, so their frames $u_i$ are points of one and the same fiber. Pick any
reference point $\sigma$ of that fiber and each frame becomes a group element $U_i$ through
$u_i=\sigma\cdot U_i^{-1}$; any two frames are then compared by $T_{ij}=U_iU_j^{-1}$, which does not depend
on $\sigma$ (\MT{prop:geo-frame-comparisons}). A useful picture is a set of observers in one laboratory,
each with private coordinate axes. An observer who changes axes, a passive re-expression $g_i$, changes
every description: sample coordinates are multiplied by $R_i^x=\rho_x(g_i)$, laws are pushed forward, the
comparisons become $g_iT_{ij}g_j^{-1}$, and nothing in the laboratory has moved. Rotating the laboratory
itself, an active transformation, that is, an automorphism of the fiber that commutes with the right action
of $G$, is a single operation, which observer $i$ records as $\varphi_i$ and observer $j$ as the conjugate
$\varphi_j=T_{ij}^{-1}\varphi_iT_{ij}$. The active transformations therefore form a copy of $G$, of
dimension $\dim G$ rather than $|V|\dim G$, and a tuple that changes the frame of one agent alone is not
active when $G$ is nontrivial and there are at least two agents (\MT{prop:geo-active-image-point}).

\begin{analogy}{structural}
In lattice gauge theory a gauge transformation has two readings that give the same group: a change of local
trivialization at every site, and a vertical automorphism of the bundle over the lattice. Both are
$\prod_xG$, because every site $x$ has its own fiber and nothing but the link variable relates the fibers at
two sites; Elitzur's theorem \citep{elitzur1975impossibility} concerns this local group. At $c_*$ the two
readings come apart. The passive group of frame changes is still $\prod_{i\in V}G$, and it acts on link
coordinates by exactly the lattice formula $\Theta_e\mapsto g_i\Theta_eg_j^{-1}$
(\MTeq{eq:geo-regime-two-gauge-law}). The automorphisms form only one copy of $G$, because the agents are not
sites with fibers of their own: they are frames in one shared fiber, and every automorphism of the whole
bundle restricts at $c_*$ to one element of that copy (\MT{prop:geo-active-image-point}). Lattice sites
correspond to base points, not to agents at one base point; which reading a lattice model of the supplements
has is left open there (\SB{open:rg-lattice-typing}). The correspondence is structural: the same
transformation formula for the links, a different automorphism group.
\end{analogy}

\subsection{A link is an automorphism of the shared fiber}

A link $\Theta_e^x\in G$ on the oriented edge copy $e:j\to i$ defines the link map
$\mathsf t_e^x(u_j\cdot g)=u_i\cdot\Theta_e^xg$, a right-equivariant bijection of the one fiber onto itself
that does not depend on the frames used to write it (\MT{prop:geo-edge-twist-invariance}). An equivariant
self-map of a torsor is, once a reference point is chosen, left multiplication by one group element, here
the per-edge twist $\mathring\Theta_e^x=U_i^{-1}\Theta_e^xU_j$. Passive re-expressions leave every twist
unchanged, a change of reference point conjugates all twists by one common element, and an active
transformation does the same. Each edge copy therefore carries the invariant conjugacy class
$[\mathring\Theta_e^x]$, and the full invariant content of the links is the twist tuple modulo one
simultaneous conjugation. In lattice gauge theory a link is a map between two different fibers and has no
conjugacy class; here all agents share a fiber, so a link is a symmetry of that fiber and has one.

In lattice language, and informally, the twist is the link measured against the direct frame comparison.
The comparison $T_{ij}$ transforms under passive re-expression exactly like a link from $j$ to $i$, so
$T_{ij}^{-1}\Theta_e^x$ is the holonomy of a two-step loop based at $j$: out to $i$ along the link, back
along the frame comparison. It is conjugate to the twist. The per-edge class is thus a Wilson loop, of a
loop that closes through the frame comparison instead of through another link, and pure lattice gauge
theory has no such loops because it has no frame comparison between sites. The same reading removes an
apparent clash with lattice intuition. A single link coordinate between two distinct agents carries no
invariant by itself, since the re-expression $g_i=e_G$, $g_j=\Theta_e^x$ sends it to $e_G$
(\MT{rem:geo-flatness-strictness}), exactly as in lattice gauge theory. The invariant belongs to the pair
formed by the link and the frames.

Whether a model can see these extra invariants depends on whether it reads the frames. Informally: if a model
satisfies the kernel-covariance hypothesis of \cref{comp:ch04} with parameters that the re-expression leaves
unchanged, and none of its kernels reads the frames, then two link configurations related by the lattice
vertex law, with the frames held fixed, have the same evidence, by the covariance of the joint
(\MT{prop:gen-product-evidence-invariance}). Such a model sees only the Wilson-loop data, and the running
example below is of this kind. A mechanism that contains a map shared by all agents and referred to their
frames reads the twists directly; the shared query and key reference maps of attention are of this kind
(\MT{tab:obs-transformer-correspondence}). The per-edge classes are physical to the extent that the model
contains frame-referenced structure.

\begin{analogy}{structural}
Write the frames as group elements $U_i$ against one reference point. A passive re-expression acts by
$U_i\mapsto g_iU_i$, a change of reference point by $U_i\mapsto U_ig_0$, and the twist $U_i^{-1}\Theta_eU_j$
is invariant under the first and conjugated by the second. These are the field content and the symmetry
pattern of a lattice gauge--Higgs model whose Higgs field is group valued (for $SU(2)$, a fundamental doublet
of fixed length written as a matrix): the twist is the link in unitary gauge, and the single remaining global
group acts on unitary-gauge links by conjugation \citep{FradkinShenker1979,OsterwalderSeiler1978}.
Supplement B spells out this reading for $G=\mathbb Z_2$, identifying the matter sign with the frame
coordinate and the dressed link sign with the twist (\SB{open:rg-lattice-typing}). The correspondence is
structural and kinematic only. The frames are structural data at $c_*$, not a field that is summed over, so
nothing about gauge--Higgs dynamics, phase structure, or complementarity carries over.
\end{analogy}

\subsection{Pure gauge, frame alignment, and Wilson loops}

Links are graph flat on a block $I$ when every closed internal walk has trivial product, equivalently when they
are a coboundary $\Theta_e^x=V_iV_j^{-1}$ for vertex potentials $V_i$, unique up to one right multiplication per
connected component (\MT{prop:geo-trivializing-criterion}, \MT{def:geo-graph-flat-links}). This is pure gauge.
Links are frame aligned when $\Theta_e^x=T_{ij}$, the canonical pure-gauge configuration supplied by the frames
themselves; then every twist is $e_G$ and every link map is the identity of the fiber. Frame alignment implies
graph flatness and not conversely, and on a tree every configuration is graph flat while its per-edge classes can
all be nontrivial (\MT{rem:geo-flatness-strictness}). In lattice terms: on a tree all links are pure gauge, as in
the maximal-tree gauge fixing \citep{Creutz1977}, and here they still carry twists relative to the frames.

\subsection{Three kinds of agreement}

On a connected block $I$ the chapter separates three conditions (\MT{def:geo-law-coherence}). Frame
agreement is graph flatness of the links on $I$, a condition on the links alone. Agreement of an equivariant
statistic, with the mean as the basic case, asks that transported statistics match on every internal edge,
$\mu_i=\rho_x(\Theta_e^x)\mu_j$, a condition on links and means. Law coherence asks that whole laws match,
$q_i=\widehat\rho_x(\Theta_e^x)q_j$, a condition on links and laws. The relations are proved in
\MT{prop:geo-three-agreements}. Law coherence holds exactly when the base law $q_{i_0}$ is fixed by every
loop holonomy based at $i_0$ and the other laws are its transports, and mean agreement holds exactly when
the base mean lies in the common fixed space of the represented holonomies and the other means are its
transports. Law coherence implies agreement
of every equivariant statistic, and the converse fails for the mean. Frame agreement neither implies nor is
implied by either of the other two. With several independent cycles, a condition checked on one loop is
necessary and not sufficient.

The distinctions matter because each notion is a candidate criterion for when agents form a collective,
and the three give different verdicts. Frame agreement without law coherence is a block with no flux whose
agents simply disagree. Law coherence without frame agreement is a block with nontrivial flux whose laws
agree anyway, because the flux lies in the stabilizer of the common law: isotropic Gaussians around the
running example's loop, or, for $G=\GL(2,\R)$, the law $\mathcal N(0,\Sigma)$ transported around a loop
whose holonomy $\Sigma^{1/2}R(0.7)\Sigma^{-1/2}$ is not orthogonal and fixes no nonzero vector. Mean
agreement without law coherence is a block whose covariances disagree. The coarse-graining chapters use the
third notion: they score a block by its divergence from holonomy-fixed parent laws, and an attained zero score
means exactly that the transported member laws coincide in one such law, which makes them law coherent
(\MT{thm:cg-holonomy-kl-marginal}). A collective criterion built on frame agreement would be too strong for that
purpose, and one built on means too weak. None of the three selects a partition into meta-agents,
which is part of \MT{open:cg-natural-selector}.

\begin{analogy}{exact}
Part (i) of \MT{prop:geo-three-agreements} is the graph version of the holonomy principle of differential
geometry: a covariantly constant section through a given fiber element exists exactly when that element is
fixed by the holonomy group. Law coherence is a covariantly constant section of the law bundle over the
block, with the links playing the connection.
\end{analogy}

\begin{analogy}{structural}
A covariantly constant Higgs field forces the holonomy of the gauge field into the stabilizer of the Higgs
value, the unbroken subgroup, while flux inside that subgroup is allowed. Law coherence with nontrivial
holonomy has the same pattern with a probability law in place of the Higgs value: flux is allowed exactly
inside the law's stabilizer. No symmetry breaking is claimed by the chapter.
\end{analogy}

\subsection{Compact and noncompact structure groups}

Nearly everything in the chapter holds for an arbitrary Lie group $G$, and the trivialization criterion
holds for an arbitrary group, so the frame comparisons, the single active copy of $G$, the per-edge classes,
and the three agreements hold verbatim for $\GL(K,\R)$, $\GL(K,\mathbb C)$, and $SU(N)$. The exception is
the construction of a probability law on the links from loop penalties (\MT{rem:geo-loop-penalties}). For a
unitary representation the penalty $1-\operatorname{Re}\Tr\mathsf U/n$ of a represented loop holonomy
$\mathsf U$ is Wilson's plaquette action \citep{WilsonConfinement1974} normalized per dimension; it lies in
$[0,2]$, and for compact $G$ the Boltzmann weight of finitely many such penalties against the product Haar
probability measure on the links has a normalizer between $e^{-2\sum_p\mathrm c_p}$ and $1$. This covers
$SU(N)$ and $SO(K)$ with unitary representations. For $\GL(n,\R)$ with $n\geq2$ the same remark proves a
no-go: no norm on $\mathfrak{gl}(n,\R)$ is invariant under conjugation by all of $\GL(n,\R)$. An invariant
nonnegative penalty vanishing exactly at the identity exists once a symmetric positive-definite fiber metric
at the base vertex is declared and made to transform covariantly, which is additional data, playing the role
a metric plays in reducing a general linear group to an orthogonal one. Whether any such penalty gives a
normalizable law on $\GL$-valued links is declined in the chapter's scope paragraph.

The gap is not a technicality. Informally, the Haar measure of a noncompact group is infinite, and a penalty
that depends only on loop holonomies leaves every tree direction of the link configuration unpenalized, so
the naive Boltzmann weight already fails to integrate on a graph with a single edge. A proper reference that
is not Haar would normalize such a law at the cost of the invariance the penalty was built to have; the
manuscript makes that trade once, for the Gaussian presentations of its formation model, as recorded in
\MT{open:prob-law-point-prior}. Two further facts bear on the target group. Links are structural parameters
by default, not random variables (\MT{sec:intro-conventions}), so the missing $\GL$ link law blocks only models
that integrate over links. And for nondegenerate Gaussian laws under independent $\GL(K)$ links, law coherence
around a loop requires the holonomy to be conjugate into $\mathrm O(K)$, which forces $\lvert\det\rvert=1$ on
the loop and so holds only on a Lebesgue-null set of link configurations
(\MT{cor:cg-holonomy-noncompact-gaussian}, discussed in \cref{comp:ch10}); for the target group, law coherence of
Gaussian beliefs around a loop is nongeneric.

\section{Objects and notation}
\label{comp:ch02-objects}

\begin{decoder}
Base point $c_*$; context space $\mathcal C$ & The single point of the context manifold at which everything happens; $\mathcal C$ itself is idle in this chapter. & One spacetime point, or one lattice site \\
Principal bundle $\mathscr P_G$; its fiber at $c_*$ & At $c_*$ only the fiber matters: a $G$-torsor, a copy of $G$ without a preferred identity. & Space of frames at one point \\
Frame $u_i$; reference point $\sigma$; $U_i$ & Agent $i$'s chosen point of the shared fiber; against a reference $\sigma$ it is the group element $U_i$ with $u_i=\sigma\cdot U_i^{-1}$. & Local frame; group-valued Higgs field in the analogy above \\
Passive re-expression $g_i$; passive vertex group $\prod_{i\in V}G$ & Agent $i$ replaces its frame by $u_i\cdot g_i^{-1}$; coordinates change, nothing else does. & Local gauge transformation read as a change of trivialization \\
Active gauge group $\mathcal A_{c_*}$; local functions $\varphi_i$ & Automorphisms of the shared fiber; one copy of $G$, written in agent $i$'s frame as $\varphi_i$. & One global symmetry transformation \\
Frame comparison $T_{ij}$ & The unique group element with $u_j=u_i\cdot T_{ij}$, equal to $U_iU_j^{-1}$. & Canonical pure-gauge link on the complete graph \\
Sample fibers $\mathsf K$, $\mathsf M$; representations $\rho_b$, $\rho_m$ & Spaces of belief samples and model presentations, with $G$ acting; the two actions may be inequivalent, unfaithful, and of different dimension. & Matter representations \\
Law fibers $\mathcal B_b$, $\mathcal B_m$; pushforward action $\widehat\rho_x$ & Families of probability laws on the sample fibers, on which $G$ acts by pushforward. & A nonlinear target space whose points are laws \\
Associated bundles $\mathcal E_b$, $\mathcal E_m$ & Bundles whose fiber at $c_*$ is a space of laws, identified with $\mathcal B_x$ by each frame. & Matter bundle, with laws in place of field values \\
Agent $\mathcal A^i=(\mathcal C_i;q_i^b,q_i^m,u_i)$ & A frame plus a belief law and a model law; at $c_*$ only their values matter. & None; a naming decision \\
Cross morphisms $\Phi$, $\widetilde\Phi$ & Fiberwise maps from belief laws to model laws and back; extra structure, not supplied by the bundle. & None \\
Multigraph $\mathfrak I=(V,E,\bar{\ })$; edge copy $e:j\to i$; reversal $\bar e$ & Oriented edges between agents, parallel copies allowed; every copy has a reversed partner. & Lattice with multiple bonds \\
Link $\Theta_e^x$; vertex law & A group element per edge copy and channel, independent of the frames; transforms by $g_i\Theta_e^xg_j^{-1}$. & Link variable $U_{xy}$ \\
Link map $\mathsf t_e^x$ & The frame-independent automorphism of the fiber that the link defines. & None in pure lattice gauge theory \\
Per-edge twist $\mathring\Theta_e^x$ & $U_i^{-1}\Theta_e^xU_j$, the link written against the single reference point; its conjugacy class is invariant. & Link in unitary gauge \\
Path product $\Theta_p^x$; holonomy $\mathrm{Hol}^x(p)$ & Ordered product of links along a walk; for a closed walk, a group element covariant under conjugation. & Wilson loop, before taking the class \\
Graph flat; jointly trivializing & Trivial holonomy on every internal loop; in both channels at once. & Pure gauge; zero field strength \\
Frame aligned & Links equal the frame comparisons, so all twists are trivial. & The canonical pure-gauge background \\
Law coherent & Laws on a block equal each other's transports along every internal edge. & Covariantly constant matter field \\
$s$-agreeing; mean agreement & An equivariant statistic, such as the mean, agrees after transport. & Agreement of an order parameter \\
\end{decoder}

Several words and glyphs trip a physicist. "Gauge" appears with two qualifiers: the passive group is the local
one and the active group is global, and "gauge invariant" without qualifier means invariant under both.
"Holonomy", "flat", and "twist" are graph-link notions: a flat link configuration says nothing about the
curvature of any connection on $\mathcal C$, and a twist is not a twist of the bundle, which over one point is
always trivial. "Common frame" means one frame shared by an agent's belief and model channels
(\MT{prop:geo-common-frame-gauge}), not one frame shared by all agents. The orientation runs from source to
receiver: $e:j\to i$ carries $j$'s coordinates into $i$'s, a path product is written so that its rightmost link
acts first, and $T_{ij}$ transforms like a link on $e:j\to i$. $T_{ij}$ is never a temperature or a transfer
matrix, a hat on $\rho$ means pushforward, and a sans-serif $\mathsf G$ is a metric, not the group. Calling a
section-bearing object an agent is a naming decision with no claim of agency attached
(\MT{sec:intro-conventions}).

\section{The formulas, explained}
\label{comp:ch02-formulas}

\paragraph{Frames and comparisons.}
The fiber at $c_*$ is a right $G$-torsor, and the chapter's first relations are
\[
u_i=\sigma\cdot U_i^{-1},
\qquad
u_j=u_i\cdot T_{ij},
\qquad
T_{ij}=U_iU_j^{-1},
\qquad
T_{ij}'=g_iT_{ij}g_j^{-1},
\qquad
U_i'=g_iU_i
\]
(\MT{prop:geo-frame-comparisons}). The reference point $\sigma$ is an arbitrary origin, and $T_{ij}$ does not
depend on it. The identities $T_{ii}=e_G$, $T_{ji}=T_{ij}^{-1}$, $T_{ij}T_{jl}=T_{il}$ make $(T_{ij})$ a pure-gauge
configuration on the complete graph. The re-expression $g_i=U_i^{-1}$ gives every agent the frame $\sigma$ and
all $T_{ij}=e_G$; this equal-frame gauge is always reachable at $c_*$, and the running example uses it.

\paragraph{The active image.}
An automorphism $F$ of the fiber, written in agent $i$'s frame as $F(u_i)=u_i\cdot\varphi_i$, satisfies
$\varphi_j=T_{ij}^{-1}\varphi_iT_{ij}$ (\MTeq{eq:geo-active-image-point}), so the whole tuple is fixed by any one
entry. In coordinates it acts like the passive re-expression $g_i=\varphi_i$; the difference is that the active
map moves the configuration while the passive one relabels it.

\paragraph{Links, the vertex law, and holonomy.}
For each edge copy $e:j\to i$ and channel $x\in\{b,m\}$,
\[
\Theta_{\bar e}^x=(\Theta_e^x)^{-1},
\qquad
\Theta_e^{x\prime}=g_i\Theta_e^xg_j^{-1},
\qquad
\mathrm{Hol}^x(p)=\Theta_{e_0}^x\Theta_{e_1}^x\cdots\Theta_{e_{r-1}}^x,
\qquad
\mathrm{Hol}^x(p)'=g_{i_0}\mathrm{Hol}^x(p)g_{i_0}^{-1},
\]
the last two for a closed walk based at $i_0$ (\MT{def:geo-graph-links}). These are the lattice formulas term
for term. The two channels carry separate link families, and the common-link case $\Theta_e^b=\Theta_e^m$ is a
declared specialization.

\paragraph{Twists.}
With $u_i=\sigma\cdot U_i^{-1}$,
\[
\mathring\Theta_e^x=U_i^{-1}\Theta_e^xU_j,
\qquad
\mathrm{Hol}^x(p)=U_{i_0}\mathring\Theta_{e_0}^x\mathring\Theta_{e_1}^x\cdots\mathring\Theta_{e_{r-1}}^xU_{i_0}^{-1}
\]
(\MTeq{eq:geo-edge-twist}, \MTeq{eq:geo-edge-twist-holonomy}). The twist is the link rewritten against the
single reference point $\sigma$; frame factors cancel in pairs along a closed walk. For abelian $G$ each twist is
itself invariant (Supplement B's $\mathbb Z_2$ reading calls it the dressed link sign), and on a tree the twists
are the entire invariant content.

\paragraph{The trivialization criterion.}
For group elements $\Xi_e$ on the internal edge copies of a block,
\[
\begin{aligned}
&\Xi_e=V_iV_j^{-1}\ \text{for every internal }e:j\to i\\
\Longleftrightarrow\quad&\Xi_{e_0}\Xi_{e_1}\cdots\Xi_{e_{r-1}}=e_G\ \text{for every closed internal walk}
\end{aligned}
\]
(\MTeq{eq:geo-trivializing-criterion}). The first line is pure gauge and the second is triviality of all Wilson
loops. The proof uses only the group law, so it holds for $\GL(K)$, for its subgroups, and for represented
links, with no compactness or topology.

\paragraph{Loop penalties.}
For a represented loop holonomy $\mathsf U=\rho_x(\mathrm{Hol}^x(p))$ of dimension $n$,
\begin{gather*}
\frac{1}{2n}\Tr\bigl((I_n-\mathsf U)^*(I_n-\mathsf U)\bigr)=1-\frac1n\operatorname{Re}\Tr\mathsf U\in[0,2]
\qquad(\mathsf U\text{ unitary}),\\
\Tr\bigl((\mathsf G^{\mathrm{fib}})^{-1}(I_n-\mathsf U)^\top\mathsf G^{\mathrm{fib}}(I_n-\mathsf U)\bigr)
\qquad(\mathsf U\text{ real}).
\end{gather*}
The first is Wilson's action per loop, vanishing exactly at $\mathsf U=I_n$; for $U(1)$ it is $1-\cos\phi$. It
is invariant because a passive re-expression conjugates $\mathsf U$ by the unitary $\rho_x(g_{i_0})$. The second
is the general linear substitute (\MT{rem:geo-loop-penalties}); it is invariant only because the fiber metric
$\mathsf G^{\mathrm{fib}}$ is declared to transform by
$(R_{i_0}^x)^{-\top}\mathsf G^{\mathrm{fib}}(R_{i_0}^x)^{-1}$, so it is a covariant field at the base vertex,
not a constant.

\paragraph{Agreement.}
Law coherence and its mean-level shadow read
\[
q_i=\widehat\rho_x(\Theta_e^x)q_j\ \text{on each internal }e:j\to i,
\qquad
\mu_{i_0}\in\bigcap_{p\text{ closed at }i_0}\ker\bigl(\rho_x(\mathrm{Hol}^x(p))-\operatorname{id}\bigr)
\]
(\MTeq{eq:phil-law-coherence}, \MTeq{eq:phil-agreement-kernel}). The first is a parallel section of the law
bundle over the block; the second says the base mean survives every loop. For a nontrivial planar rotation loop
the kernel is $\{0\}$ and every agreeing mean is zero.

\paragraph{Cross morphisms.}
A cross morphism has representatives transforming by
$\Phi_i'=\widehat\rho_m(g_i)\circ\Phi_i\circ\widehat\rho_b(g_i)^{-1}$, and one map $\Phi_0$ serves in every frame
exactly when it intertwines the two actions (\MT{prop:geo-intertwining-cross-map}). In the running example's
setting, with $SO(3)$ acting on $\R^3$ in the belief channel and trivially on $\R$ in the model channel, the only
equivariant linear map $\R^3\to\R$ is zero by Schur's lemma, while pushforward of laws by the norm,
$\beta\mapsto(\lvert\cdot\rvert)_\#\beta$, is an equivariant map of laws: law-level equivariance is weaker than
linear equivariance.

\paragraph{The running example.}
Take $G=SO(3)$, $\rho_b$ the defining representation on $\R^3$, $\rho_m$ trivial on $\R$, three agents on a
triangle with $e_1:2\to1$, $e_2:3\to2$, $e_3:1\to3$, all frames equal, and belief links
\[
\Theta_{e_1}^b=R_z(\tfrac\pi2),
\qquad
\Theta_{e_2}^b=R_x(\tfrac\pi2),
\qquad
\Theta_{e_3}^b=I_3,
\qquad
\mathrm{Hol}^b(p)=R_z(\tfrac\pi2)R_x(\tfrac\pi2),
\qquad
\Tr\mathrm{Hol}^b(p)=0
\]
(\MT{sec:geo-running-example}, \MTeq{eq:geo-running-holonomy}). The per-edge classes are the rotation angles
$\pi/2$, $\pi/2$, $0$, and the holonomy is a rotation by $2\pi/3$, here about the axis $(1,1,1)$. Replacing the
second quarter turn by one about the same axis as the first gives a rotation by $\pi$ with the same three
classes, which is the concrete reason the class tuple alone does not determine the holonomy while the tuple
modulo one simultaneous conjugation does. The loop is not graph flat, yet the isotropic laws
$\mathcal N(0,s^2I_3)$ are law coherent on it, because every rotation fixes them. The one record of the
example has success probability $\exp(-\tfrac12\lvert k_1-\rho_b(\Theta_{e_1}^b)k_2\rvert^2)$, a transported
discrepancy whose length no passive re-expression changes (\MTeq{eq:geo-running-record}). In lattice terms,
this is a frustrated plaquette with flux $2\pi/3$ and a matter law that does not feel the flux.

\section{What is proved, and how solid it is}
\label{comp:ch02-status}

\begin{scorecard}
Agents and law bundles: \MT{def:geo-associated-bundles}, \MT{def:geo-agent} & An agent is a frame in the one fiber plus a belief law and a model law; the channels may carry inequivalent representations. & \CHOICE & Standard construction; the commitment is laws valued in $G$-spaces at one fiber. \\
Frame comparisons: \MT{prop:geo-frame-comparisons} & $T_{ij}$ is unique, a cocycle, and transforms by $g_iT_{ij}g_j^{-1}$. & \SOLID & Any Lie group; torsor algebra. \\
Active image: \MT{prop:geo-active-image-point} & Fiber automorphisms form one copy of $G$ inside $\prod_iG$, proper when $G\neq\{e_G\}$ and $|V|\geq2$. & \SOLID & Any Lie group; standard in substance. \\
Common frame: \MT{prop:geo-common-frame-gauge} & Separate belief and model frames reduce to one by a re-expression. & \SOLID & Harmless because nothing reads the relative frame. \\
Cross morphisms: \MT{prop:geo-intertwining-cross-map} & Gauge law of representatives; intertwiners are exactly the frame-independent ones; not every cross morphism is induced; Schur vanishing does not remove cross morphisms. & \SOLID & Their existence is additional structure. \\
Independent links: \MT{def:geo-graph-links} & One group element per oriented edge copy and channel, independent of the frames. & \CHOICE & Structural parameters by default. \\
Twist invariance: \MT{prop:geo-edge-twist-invariance} & Each edge carries an invariant twist class; link sets are actively related exactly when their twist tuples are simultaneously conjugate. & \SOLID & Any Lie group; standard in substance; an invariant of links with frames. \\
Trivialization: \MT{prop:geo-trivializing-criterion} & Coboundary exactly when all closed-walk products are trivial; potentials unique up to one right factor per component. & \SOLID & Any group; pure gauge exactly when Wilson loops are trivial. \\
Strictness of flatness: \MT{rem:geo-flatness-strictness} & Frame alignment is strictly stronger than graph flatness; one link coordinate has no invariant; trees can carry nontrivial classes. & \SOLID & General non-implications on minimal graphs. \\
Compact loop penalty: \MT{rem:geo-loop-penalties} & For unitary $\rho_x$, $1-\operatorname{Re}\Tr\mathsf U/n\in[0,2]$ is invariant; for compact $G$ the penalty law against product Haar probability has normalizer in $[e^{-2\sum_p\mathrm c_p},1]$. & \COND & Needs $G$ compact, $\rho_x$ unitary: holds for $SU(N)$, $SO(K)$; fails for $\GL(K,\R)$, $\GL(K,\mathbb C)$. \\
No invariant norm for $\GL$: \MT{rem:geo-loop-penalties} & No norm on $\mathfrak{gl}(n,\R)$, $n\geq2$, is $\GL(n,\R)$-conjugation invariant; a covariant fiber metric restores an invariant penalty. & \SOLID & Negative result; the metric is extra data. \\
Link law for noncompact $G$ (scope paragraph) & Normalizability of a loop-penalty link law for noncompact or general linear $G$. & \OPEN & Declined, with no ledger entry; the program's target case. \\
Three agreements: \MT{prop:geo-three-agreements} & Holonomy characterizations; law coherence implies statistic agreement, not conversely; frame agreement independent; one loop insufficient; $\GL(2,\R)$ witness with non-orthogonal holonomy. & \SOLID & Any $G$; witnesses refute general implications. \\
Agreement and partitions (scope paragraph) & Whether law coherence bears on selecting a partition into meta-agents. & \OPEN & Part of \MT{open:cg-natural-selector}. \\
\end{scorecard}

\begin{hygiene}
Three conditions matter for the models the program cares about. Compactness of $G$ and unitarity of the
representation are what make the loop-penalty link law normalizable; they hold for $SU(N)$ and fail for
$\GL(K,\R)$ and $\GL(K,\mathbb C)$, and the chapter supplies no substitute. The existence of cross morphisms is
additional structure: the bundles do not supply maps between belief laws and model laws, and with irreducible
inequivalent linear representations no nonzero equivariant linear map exists. The single frame shared by the
two channels of an agent is justified only because no energy, kernel, or evaluator reads the relative frame
between them (\MT{prop:geo-common-frame-gauge}); a model that read it would need both frames back. The
common-link specialization $\Theta_e^b=\Theta_e^m$ is an optional modeling choice, not a consequence.

The rest is hygiene here: the context space is a second countable Hausdorff manifold (idle at $c_*$),
representations are measurable, sample fibers are standard Borel, and the smooth tier (\MT{hyp:geo-smooth-tier}),
which excludes infinite-dimensional law fibers, matters only for the fiber geometry of \cref{comp:ch07}.
\end{hygiene}

The chapter has no toy results: its small graphs and its $SO(2)$, $SO(3)$, and $\GL(2,\R)$ instances are
counterexamples to general implications, solid at full generality, or illustrations.

\section{Open ends}
\label{comp:ch02-open}

Two questions are left unsettled. The first is whether a normalizable law on links built from loop penalties
exists for a noncompact or general linear structure group. The chapter declines the claim and records no open
problem or settling condition for it; for the $\GL(K)$ target, links therefore remain structural parameters, and
the theory supplies no Wilson-type prior to integrate them against. The second is whether any of the three agreements bears on
choosing a partition of agents into meta-agents, part of \MT{open:cg-natural-selector}, which is settled by a
relabeling-natural, gauge-compatible functional of the laws, not monotone under refinement, with a nondegenerate
optimizer, or by a theorem that no such selector exists in a stated class. Connections, base transport, and
descent are out of scope by design rather than open.

\begin{verdictbox}{\DELIVERED}
The chapter says that it determines what geometric data an agent carries at $c_*$ and which are invariant: the
active gauge group is one copy of $G$ inside the passive vertex group, every link carries an invariant
conjugacy class with loop holonomy the telescoped product of per-edge twists, and three agreements are
distinct; the introduction's reading map repeats this. Every one of these statements is proved, for an
arbitrary Lie group, with short proofs. The mathematics is standard in substance: torsor algebra, automorphisms
of a torsor, degree-one nonabelian cohomology of a graph, and the holonomy principle on a graph. What is new is
the typing, and it does real work in one place: because all agents share one fiber, links become automorphisms
of that fiber and acquire per-edge invariants that pure lattice gauge theory lacks, visible to a model only
through frame-referenced structure. No connection, curvature, or bundle topology enters, and the text does not
claim otherwise. The gap that matters for the program is the declined normalizability of loop-penalty link laws
for noncompact $G$: for $\GL(K,\R)$ the chapter supplies an invariant penalty only after adding a covariant fiber
metric, and no normalized link law.
\end{verdictbox}


\chapter{Probability Semantics at \texorpdfstring{$c_*$}{c*}}
\label{comp:ch03}

\section{The question and the answer}
\label{comp:ch03-question}

The technical chapter asks on which measurable objects, and in which order, every probabilistic statement at
$c_*$ rests. Its answer is a discipline. Probability laws and Markov kernels are the primitive objects.
Densities are Radon--Nikodym derivatives against reference measures fixed once for a whole family of laws,
generative and recognition alike. Logarithms of density ratios are written only where the numerator is
absolutely continuous with respect to the denominator, and relative entropy is $+\infty$ everywhere else
(\MT{def:prob-log-density-ratio-convention}). The stated main results are that a family of mixed laws on the
real line has a common dominating measure exactly when the union of its atoms is countable
(\MT{thm:prob-mixed-family-common-domination}), the relative-entropy chain rule in $[0,+\infty]$ that every
later decomposition uses (\MT{lem:prob-kl-chain-rule}), and three levels of description, sample, law point,
and configuration, at which a divergence between laws is not an energy of samples
(\MT{prop:prob-law-point-energy}) and a point-mass recognition of a law point has infinite cost
(\MT{prop:prob-law-point-point-mass}). Along the way the chapter types the population recognition law as a
selected coupling of the agents' local laws and prices restricted access. The generative law of
\cref{comp:ch04} is built on these spaces, the evidence bound of \cref{comp:ch05} is an application of the
chain rule, and attention (\cref{comp:ch06}) and formation (\cref{comp:ch13}) are typed by the levels.

Which of this is substance and which is hygiene? Most of the chapter is measure theory done carefully so that
later identities hold without case distinctions, and a physicist can take it on trust. Three things are
substantive. The population recognition law is not determined by the agents' local beliefs and must be
selected, and that selection is where recognition-level correlations between agents live. Restricting what
each agent can see has an evidence price that no selection removes. And the level typing shows that a
relative entropy between two agents' beliefs, the building block of the program's attention free energies, is
not the expectation of any energy of samples, so its place inside one evidence bound depends on a prior over
beliefs that the theory has not constructed (\MT{open:prob-law-point-prior}). A fourth item sits on the border:
families of laws with moving atoms, and noiseless records with a continuum of values, have no density model,
and the evidence bound for such records is open (\MT{open:prob-nondominated-likelihood}).

\section{The picture}
\label{comp:ch03-picture}

\subsection{Laws first, because the theory needs degenerate laws}

A physicist's default is to write a probability as a density and a free energy as an integral of
$q\log(q/p)$. The chapter refuses that default for a concrete reason: the models it must cover include
point-mass model presentations (the definite case $q_i^m=\delta_{m_i}$ of \MT{def:prob-model-evaluation}),
laws carried by lower-dimensional sets such as proper affine subspaces, and laws with atoms whose locations
move with a parameter. None of these has a density against Lebesgue measure, and some families of them have
no common reference at all. A sigma-finite measure has at most countably many atoms, so the moving-atom family
$\mu_t=\tfrac12\mathcal N(0,1)+\tfrac12\delta_t$, $t\in\R$, has no common dominating measure, and neither does
any uncountable family of Dirac laws on an uncountable space (\MT{thm:prob-mixed-family-common-domination},
\MT{sec:exp-dominated-nondominated}). Working with measures and setting relative entropy to $+\infty$ off
absolute continuity keeps every identity meaningful in the extended reals: the evidence bound of a singular
recognition law is $-\infty$ rather than undefined. This is the whole content of "kernels before densities".
It costs nothing when densities exist, and it keeps the degenerate cases inside the theory when they do not.

\begin{analogy}{exact}
Conditioning on an event of probability zero is not determined by the joint law: the posterior at a single
observation depends on the version of the conditional one uses, which is the Borel--Kolmogorov paradox. The
chapter's regular-conditional convention (\MT{hyp:prob-regular-conditional-convention}) therefore asserts
every fixed-observation statement only for almost every observation, and its witness
(\MT{rem:prob-version-witness}) is an instance of the paradox: two versions of one joint density on $\R^2$ give
different posteriors at $o=0$.
\end{analogy}

\subsection{Local beliefs do not fix the population law}

Each agent has a local recognition law $\mathbb Q_{i,o,X}$ on its own block, computed from what it can access.
The population recognition law is a law on all blocks at once whose agent marginals are the local laws, that
is, a member of the coupling class $\operatorname{Cpl}(\mathbb Q_{\bullet,o,X})$. The class is never empty and
is convex (\MT{prop:prob-coupling-existence-convexity}), and it is a single point only when at most one local
law is non-Dirac (\MT{prop:prob-marginals-do-not-determine-joint}). As soon as two agents are uncertain, their
beliefs underdetermine the joint recognition law, and the theory selects one by a declared rule
(\MT{hyp:prob-measurable-coupling-selector}). The product coupling is one admissible selection, not the
default; copulas \citep{Nelsen2006}, shared latent variables, ordered conditionals, or graphical families
supply others. The coupling is a modeling input, not an output of the local beliefs, and the manuscript
asserts no canonical selector.

Restricted access sharpens the point (\MT{prop:prob-access-restriction}). An access map limits what each local
law may depend on, but a selector that reads all records can still make the population law depend on data that
no agent sees, so a claim of decentralized inference also needs a selector indexed by the agents' accesses
alone. And restricting access costs evidence. When an agent cannot distinguish two observations whose
posteriors differ, its single local law sits at a positive divergence from at least one of them; in the
chapter's two-agent binary witness the worst-case gap is at least the Chernoff information between the two
posterior marginals, about $0.0885$ nats \citep[Section~11.9]{CoverThomas2006}.

\begin{analogy}{structural}
The product coupling is the naive mean-field trial state of the Gibbs--Bogoliubov variational principle, and
correlated couplings are richer trial states with the same one-body marginals. The identification is exact for
the product case; for the others it is structural, since a coupling class fixes the one-body marginals and
leaves every higher correlation free.
\end{analogy}

\subsection{Three levels of description}

The chapter distinguishes three kinds of latent variable (\MT{def:prob-description-levels}). At the sample
level the latent state is the population sample $y=((k_i,m_i))_{i\in V}$, the agents' microstates, and an
agent's belief is the marginal on $k_i$ of the recognition law, a law over microstates. At the law-point level
the beliefs themselves are latent: the latent state of agent $i$ is a point $\mathfrak q_i^b$ of the
associated law fiber, a probability law treated as a random variable, and a recognition law is a law over laws.
At the configuration level the latent state is a presentation $w$, for instance a mean and a covariance,
together with a sample $y$ drawn from the law that $w$ presents, the two linked by a presentation kernel. The
levels define different optimization problems, and the manuscript requires every energy to name its level.

The reason is \MT{prop:prob-law-point-energy}. An energy $D$ of samples enters a free energy as $\E_QD$, which
is affine in the recognition law. A transported relative entropy $\KL(q_i\Vert\widehat\rho_x(\Theta_e^x)q)$ is
strictly convex in the source law $q$ along mixtures, because $-\log$ is strictly convex, so no sample energy
can match it at two source laws and one strict mixture of them. A divergence between two agents' beliefs is a
functional of the ensemble, like an entropy, and not a Hamiltonian term. The consequence runs through the
program. The belief-alignment terms $\KL(q_i\Vert(\Theta_e)_\#q_j)$ that weight attention rows
(\MTeq{eq:full-graph-transported-divergences}) are law-point energies, and they become terms of one evidence
bound only if beliefs are latent law points under a normalized prior, which is \MT{open:prob-law-point-prior}.

\MT{prop:prob-law-point-point-mass} closes the obvious shortcut. If law points are latent and their prior has
no atom at the law point one wants, a recognition law that puts all its mass on that law point has infinite
relative entropy to the posterior and an evidence bound of $-\infty$. A row or an update computed at fixed law
points is therefore not a recognition optimum at the law-point level. With law points latent it is a
generative posterior quantity conditional on them; with law points declared as parameters it is a function of
parameters whose change is an M-coordinate subject to the acceptance rule of
\MT{def:elbo-generalized-m-acceptance}. Both readings are legitimate, and neither is an agent's recognition of
its own belief.

\begin{analogy}{structural}
The entropy $-\int q\log q$ is a functional of the ensemble and not the average of any state-independent
microscopic energy: it is strictly concave in $q$, while every average $\int qH$ is linear. The cross-entropy
between two agents' beliefs has the same character in its second argument, and that is all
\MT{prop:prob-law-point-energy} uses. The point-mass result has an exact counterpart: a trial state
concentrated on one point of a continuum has entropy $-\infty$ relative to a continuous reference.
\end{analogy}

\subsection{Auxiliary variables are free only when exact}

The configuration level is reached from the sample level by a lift
(\MT{prop:prob-configuration-lift-identity}). Draw configuration coordinates from a kernel
$\kappa^{\mathrm{cfg}}(y,dw)$ given the sample: the evidence does not change, and a lifted recognition law pays,
relative to the sample-level bound of its sample marginal, the expected divergence between its conditional law
of $w$ given $y$ and $\kappa^{\mathrm{cfg}}$. The lift never raises the bound, and it leaves a finite bound
unchanged exactly when recognition uses $\kappa^{\mathrm{cfg}}$ itself as its conditional.

\begin{analogy}{exact}
Introducing an auxiliary variable with its exact conditional law, as in a Hubbard--Stratonovich or
data-augmentation step, leaves the partition function unchanged, and a variational ansatz on the enlarged
space pays the conditional relative entropy for any mismatch. Both statements are the chain rule of
\MT{lem:prob-kl-chain-rule}.
\end{analogy}

\section{Objects and notation}
\label{comp:ch03-objects}

\begin{decoder}
Sample blocks $\mathsf Y_i=\mathsf K_i\times\mathsf M_i$; latent space $\mathsf Y$ & Agent $i$'s belief-sample and model-presentation coordinates; their product over agents. & Single-site and full configuration space \\
Record $a\in\mathfrak A$; scope $\partial a$; record space $\mathsf O_a$ & An observation factor attached to a set of agents; a single-agent observation has singleton scope. & Measured quantity on a bond or plaquette \\
Structural configuration $X\in\mathsf X$ & The multigraph, links, frames, and mechanism parameters, held fixed by declaration. & Lattice, couplings, background fields \\
Access map $\operatorname{Acc}_i$ & What agent $i$ may use from the records and the structure. & A local observer's data \\
Local recognition kernel $\mathbb Q_i$; law $\mathbb Q_{i,o,X}$ & Agent $i$'s approximate posterior on its own block, computed from its access. & One-body trial state \\
Population recognition law $\mathbb Q_{V,o,X}$ & A joint approximate posterior on all blocks, selected to have the local laws as marginals. & Many-body trial state \\
Coupling class $\operatorname{Cpl}(\mathbb Q_{\bullet,o,X})$ & All joint laws with the given agent marginals. & Fr\'echet class of joint laws \\
Selection rule $\mathsf S$ & A declared measurable rule picking one coupling. & Choice of variational ansatz \\
Recognition marginals $q_i^{b;o,X}$, $q_i^{m;o,X}$ & The belief and model parts of the local law; the values of agent $i$'s sections at $c_*$. & One-body reduced laws \\
Generative law $\mathbb P_{\theta,V}(\cdot\given X)$ & The constructed joint of latents and records (\cref{comp:ch04}). & Prior times likelihood; Boltzmann weight \\
Posterior slice $\boldsymbol\Pi_{\theta,V,o,X}$ & The exact conditional law of latents at observation $o$. & Gibbs state conditioned on data \\
Evidence $p_{\theta,V}(o\given X)$ & The marginal density of the records. & Partition function \\
References $\nu^{\mathsf Y}$, $\nu^{\mathsf O}$, $\nu_{\mathrm{mix}}$ & Fixed measures against which densities are taken; $\nu_{\mathrm{mix}}$ adds countably many atoms to Lebesgue measure. & Phase-space measure \\
Relative entropy $\KL(Q\Vert P)$ & Kullback--Leibler divergence, set to $+\infty$ off absolute continuity. & Free-energy excess over equilibrium \\
Presentation $m_i$; evaluator $\operatorname{ev}_i$ & A model coordinate that the evaluator turns into a generative kernel; not itself a law. & Coupling constants of a local Hamiltonian \\
Law point $\mathfrak q_i^x$; block $\mathsf Y_i^{\mathrm{law}}$ & A law treated as a latent random variable. & A one-body distribution treated as a field \\
Configuration space $\mathsf Y^{\mathrm{cfg}}$; $\operatorname{pres}$; $\kappa^{\mathrm{cfg}}$ & Presentations such as mean and covariance; the map to law points; the kernel linking them to samples. & Auxiliary-field representation \\
Regular set $\mathsf O^{\mathrm{reg}}_{\theta,V,X}$ & Observations of full probability with finite positive evidence, where the posterior formula is used. & None \\
Extended reals & Values in $[-\infty,+\infty]$; bounds may be $-\infty$ and divergences $+\infty$. & None \\
\end{decoder}

Several terms carry nonstandard meaning. "Configuration" is used twice, and the two uses must not be merged: the
structural configuration $X$ collects graph, links, frames, and parameters, while the configuration level of
description has latent presentations $w\in\mathsf Y^{\mathrm{cfg}}$. "Kernel" always means a Markov kernel, a
measurable family of probability laws, never an integral-operator kernel or a positive-definite kernel.
"Record" is an observation factor with a scope, a hyperedge on the agents, and never a memory. "Recognition
law" is the approximate posterior of variational inference, and the exact conditional is called the posterior
slice. "Evidence" is the marginal likelihood, which a physicist should read as a partition function.
"Presentation" is a parameter-like coordinate of a generative factor and has nothing to do with a group
presentation. "Version" means a representative up to a null set. "Selected" means chosen by a declared rule,
not derived.

\section{The formulas, explained}
\label{comp:ch03-formulas}

\paragraph{Relative entropy.}
The convention of \MTeq{eq:prob-kl-definition} is
\[
\KL(Q\Vert P)=
\begin{cases}
\displaystyle\int\log\Bigl(\frac{dQ}{dP}(y)\Bigr)Q(dy), & Q\ll P,\\[1.2ex]
+\infty, & \text{otherwise},
\end{cases}
\]
the divergence of \citet{Kullback1951} with the extension of \citet{Csiszar1967}. The negative part of the
integrand has integral at most $1/e$, so the first branch is a well-defined element of $[0,+\infty]$. The
typical infinite case in the program is a Gaussian with singular covariance measured against one with
nonsingular covariance; the typical finite case is two nondegenerate Gaussians. Whether $Q\ll P$ can be read
off from densities: against a common reference it holds exactly when $Q$ puts no mass where the density of $P$
vanishes (\MT{prop:prob-density-absolute-continuity}).

\paragraph{The mixed reference.}
When the atoms of a family of mixed laws on $\R$ form the countable set $\{a_1,a_2,\ldots\}$,
\[
\nu_{\mathrm{mix}}=\operatorname{Leb}+\sum_{n\geq1}2^{-n}\delta_{a_n}
\]
dominates the whole family with one parameter-independent measure (\MTeq{eq:prob-mixed-family-reference}). If the
atoms move over an uncountable set, no such measure exists, and the chapter then declares no density model on that
coordinate. For learning this is concrete: a law with an atom whose location is a learned parameter cannot have
its densities compared across parameter values against one reference.

\paragraph{The chain rule and its corollaries.}
For laws on a product with first marginals $Q_1,P_1$ and conditional kernels $Q_{2|1},P_{2|1}$,
\[
\KL(Q\Vert P)=\KL(Q_1\Vert P_1)+\int\KL\bigl(Q_{2|1}(y_1,\cdot)\big\Vert P_{2|1}(y_1,\cdot)\bigr)Q_1(dy_1)
\quad\text{in }[0,+\infty]
\]
(\MTeq{eq:elbo-kl-chain-rule}). This is the relative-entropy chain rule, standard in substance
\citep[Section~2.5]{CoverThomas2006}; the chapter's version needs no densities and no finiteness hypothesis, and
it does not depend on the chosen conditional versions. The same lemma gives four further tools
(\MT{lem:prob-kl-chain-rule}). Measured against a product reference, a joint law's divergence splits into its
total correlation $\KL(Q\Vert\bigotimes_bQ_b)$ plus the marginal divergences; physically, the connected
correlations separate from the one-body terms. With the first marginal fixed, the divergence is minimized by
copying the reference's conditional. Divergence cannot increase under a Markov kernel, which is the data
processing inequality (clause (d), also cited as \MT{thm:cg-kl-dpi-extended}). And divergence is unchanged by a
bimeasurable bijection, which is why frame changes preserve every divergence without any density hypothesis.

\paragraph{The evidence price of restricted access.}
In the witness of \MT{prop:prob-access-restriction}, agent $1$ sees only its own record, so its local law
$\operatorname{Ber}(t)$ is the same at the observations $(1,0)$ and $(1,1)$, whose posterior marginals for $y_1$
are $\Pi_1^{(0)}=(37/89,52/89)$ and $\Pi_1^{(1)}=(13/161,148/161)$. Then
\[
\max_{o_2\in\{0,1\}}\KL\bigl(\mathbb Q_{V,(1,o_2),X}\big\Vert\boldsymbol\Pi_{\theta,V,(1,o_2),X}\bigr)
\geq\max_{o_2}\KL\bigl(\operatorname{Ber}(t)\big\Vert\Pi_1^{(o_2)}\bigr)
\geq\sup_{\xi\in[0,1]}\bigl(-\log Z_\xi\bigr)=0.08849\ldots
\]
(\MTeq{eq:prob-access-gap-floor}), with $Z_\xi=\sum_{y_1}\Pi_1^{(0)}(y_1)^\xi\Pi_1^{(1)}(y_1)^{1-\xi}$. The first
step is data processing onto agent $1$'s coordinate, and the second is the Chernoff bound: a single law cannot
be close to two distinct targets at once, and the best compromise is the Chernoff point. The bound is attained.
The equality condition of the evidence bound, recognition equal to the posterior, lies outside the
access-restricted class at one of the two observations.

\paragraph{Transported divergences are not sample energies.}
For a receiver law $q_i$ and two source laws $p_1,p_2$ with finite transported divergences whose transported
densities differ on a set of positive $q_i$-measure,
\[
\KL\bigl(q_i\Vert\widehat\rho_x(\Theta_e^x)(sp_1+(1-s)p_2)\bigr)
<s\KL\bigl(q_i\Vert\widehat\rho_x(\Theta_e^x)p_1\bigr)+(1-s)\KL\bigl(q_i\Vert\widehat\rho_x(\Theta_e^x)p_2\bigr)
\]
for $0<s<1$ (\MT{prop:prob-law-point-energy}). The left side is what the divergence does along a mixture of
source laws, and the right side is what any energy averaged over samples must do. Besides finiteness, the
hypotheses exclude only the case in which the two transported laws have equal densities almost everywhere under
the receiver law.

\paragraph{Point masses on law points.}
At the law-point level, with a posterior dominated by the prior (true when every record kernel has a density
against a fixed reference) and a law point $\bar{\mathfrak q}$ that is not an atom of the prior,
\[
\KL\bigl(\delta_{\bar{\mathfrak q}}\big\Vert\boldsymbol\Pi^{\mathrm{law}}_{\theta,V,o,X}\bigr)=+\infty
\]
(\MTeq{eq:prob-law-point-point-mass}). The singleton $\{\bar{\mathfrak q}\}$ is a posterior-null set carrying all of
the recognition mass. With an atom at $\bar{\mathfrak q}$ the divergence can be finite; the proposition concerns
non-atomic priors, which is the situation for any continuous family of candidate beliefs.

\paragraph{The configuration lift.}
For a lifted recognition law $\mathbb J$ with sample marginal $\bar{\mathbb Q}$,
\[
\Lelbo^{\mathrm{ext,cfg}}_V(\mathbb J;X,o)
=\Lelbo^{\mathrm{ext}}_V(\bar{\mathbb Q};X,o)
-\E_{\bar{\mathbb Q}}\Bigl[\KL\bigl(\mathbb J(\cdot\given Y)\big\Vert\kappa^{\mathrm{cfg}}(Y,\cdot)\bigr)\Bigr]
\]
(\MTeq{eq:prob-configuration-lift-identity}). The first term is the bound of the sample-level problem, the second
is the price of any mismatch between recognition's conditional over configurations and the lift kernel, and both
live in the extended reals.

\section{What is proved, and how solid it is}
\label{comp:ch03-status}

\begin{scorecard}
Kernels first; levels: \MT{def:prob-normalized-kernels}, \MT{def:prob-description-levels} & Laws and kernels are primitive, densities are derivatives against family-wide references, and every energy names its level: sample, law point, or configuration. & \CHOICE & A typing discipline; it costs nothing when densities exist. \\
Common domination: \MT{thm:prob-mixed-family-common-domination} & A family of mixed laws on $\R$ without singular-continuous parts has a sigma-finite common dominating measure exactly when its atoms are countable; moving atoms are not dominated. & \SOLID & Standard in substance (a sigma-finite measure has countably many atoms). \\
Dirac families: \MT{sec:exp-dominated-nondominated} & An uncountable family of Dirac laws on an uncountable standard Borel space has no common sigma-finite dominating measure. & \SOLID & Standard. \\
Chain rule and corollaries: \MT{lem:prob-kl-chain-rule} & Chain rule, product-reference split, fixed-marginal minimization, data processing, and bijection invariance, all in $[0,+\infty]$. & \SOLID & Standard in substance; proved in \MT{app:measure-lemmas} without density or finiteness hypotheses. \\
Couplings exist: \MT{prop:prob-coupling-existence-convexity} & The coupling class is nonempty and convex; the product coupling is a measurable kernel. & \SOLID & Standard. \\
Marginals do not fix the joint: \MT{prop:prob-marginals-do-not-determine-joint} & The coupling class is a singleton exactly when at most one marginal is non-Dirac. & \SOLID & Standard in substance (Fr\'echet classes). \\
Coupling selection: \MT{hyp:prob-measurable-coupling-selector} & The population recognition law is chosen by a declared measurable rule inside the coupling class. & \CHOICE & No canonical selector; equivariance of the selector under frame changes is declined. \\
Access restriction, part (i): \MT{prop:prob-access-restriction} & Access maps constrain local laws, not the population law; decentralization also needs an access-indexed selector. & \SOLID & A general no-go shown by counterexample. \\
Access restriction, part (ii): \MT{prop:prob-access-restriction} & Every access-restricted recognition law has worst-case gap at least $0.08849\ldots$ nats, and the bound is attained. & \TOY & Two binary agents with one prior; the constant is instance specific. \\
Law-point energy: \MT{prop:prob-law-point-energy} & A transported relative entropy is strictly convex in the source along mixtures, so no sample energy reproduces it. & \SOLID & Needs finite divergences and transported laws that differ on the receiver's support. \\
Point-mass cost: \MT{prop:prob-law-point-point-mass} & Recognition by a point mass at a non-atomic law point has infinite divergence and bound $-\infty$; rows at fixed law points are posterior quantities or parameter functions. & \SOLID & Needs the posterior dominated by the prior, true for record kernels with densities. \\
Configuration lift: \MT{prop:prob-configuration-lift-identity} & Lifting leaves the evidence unchanged and lowers the bound by the expected conditional divergence. & \SOLID & Standard in substance (augmentation and chain rule). \\
Law-point prior: \MT{open:prob-law-point-prior} & A normalized generative law over law points for a general class. & \OPEN & Needed for transported divergences to be evidence-bound terms. \\
Nondominated likelihoods: \MT{open:prob-nondominated-likelihood} & Evidence bound and coarse channels for record families with no common dominating measure. & \OPEN & Covers noiseless records with a continuum of values. \\
\end{scorecard}

\begin{hygiene}
Four conditions can fail in the models the program cares about. Common domination fails for families with
moving atoms and for deterministic records with a continuum of values, and then the density formulas, and
everything that uses them, are not available (\MT{open:prob-nondominated-likelihood}). Absolute continuity of
recognition with respect to the posterior is not assumed; its failure is the singular branch on which the bound
is $-\infty$, and it genuinely occurs for degenerate Gaussian beliefs and point masses. The coupling selection is
a modeling commitment whose choice changes every collective quantity, and its equivariance under frame changes
is not claimed. The access structure, when imposed, carries an evidence price.

The rest is hygiene a physicist may skip on first reading: standard Borel spaces, sigma-finite references, joint
measurability of Radon--Nikodym derivatives and of relative entropies in their parameters (proved once in
\MT{app:measure-lemmas}), the almost-every-observation convention for posteriors, and bookkeeping in the extended
reals. The restriction of the common-domination theorem to laws without singular-continuous parts is also
hygiene: no model in the program uses such laws.
\end{hygiene}

\section{Toy versus general}
\label{comp:ch03-toy}

The only toy result is the quantitative half of the access-restriction proposition. Its mechanism is visible in
the proof and has a general form: data processing onto the coordinates an agent can see, followed by the Chernoff
bound for one law facing two targets. Nothing in that argument depends on binary variables or on two agents, so
the mechanism may well survive in large networks with many agents sharing an access class, but the manuscript
states and proves it only for the displayed instance, and the constant $0.08849\ldots$ and the optimizers
$\xi\approx0.548$ and $t\approx0.783$ are specific to that prior and those record accuracies. The qualitative
half, that access restriction constrains local laws and not the population law, is a general negative statement
proved by counterexample.

\section{Open ends}
\label{comp:ch03-open}

Two open problems are recorded. The first is a normalized prior over law points (\MT{open:prob-law-point-prior}).
Law fibers carry no canonical sigma-finite reference, since the space of all laws on a sample fiber is infinite
dimensional, and the Haar measure of a noncompact $G$ is infinite, so a Haar prior on a free orbit does not
normalize. The problem is settled, for a class of law points, by a sigma-finite reference on the law-point or
presentation space, a finite normalizer for the exponentiated law-point energies, the transformation law of the
prior under the active gauge action, record kernels that read law points, and a compatibility relation between
the law-point posterior and the sample-level recognition marginals. Until then, the transported divergences
\MTeq{eq:full-graph-transported-divergences} enter as quantities conditional on law points, as functions of
parameters, or at the configuration level under the one proper prior the formation model supplies. The second is the treatment of record families with no
common dominating measure (\MT{open:prob-nondominated-likelihood}): the joint law and its posterior exist, but the
likelihood and possibly the evidence density do not. It is settled by a treatment of the evidence bound, its
local disintegrations, and the coarse channels at the level of joint measures and regular conditional laws in
which the exact identities of the later chapters are proved, or by a record family for which they provably have
no measure-level counterpart.

The scope paragraph also records three limits. Injectivity of a model evaluator is not assumed, and no regularity
of the space of presentations modulo equal evaluated kernels is supplied, so presentations are kept as they are;
this is bookkeeping for overparameterized models. The selector's equivariance under frame re-expression is not
claimed; this is substantive, because a recognition procedure rerun in new frames need not return the
transformed coupling, and the invariance results of \cref{comp:ch04} are stated for a coupling transported as a
whole.

\begin{verdictbox}{\DELIVERED}
The chapter says that laws and kernels come first and densities follow against family-wide references, so that
the evidence bound becomes an identity in the extended reals, and names as its main results common domination
exactly when atoms are countable, the chain rule in $[0,+\infty]$, and three levels at which a law-point energy
is not a sample energy and point-mass recognition of law points has infinite cost; the reading map adds that it
proves the chain rule used everywhere and distinguishes the three levels. All of this is proved at the generality
the theory uses. Most of it is standard measure theory and information theory in careful form, and a physicist
loses nothing by treating it as such. The substantive content is elsewhere and is also delivered: local beliefs
do not determine the population recognition law, so a coupling must be selected; restricted access has a
provable evidence price; and a divergence between beliefs is not a sample energy. The last result exposes a gap
that the chapter names rather than closes: without a normalized law-point prior, the belief-alignment divergences
at the center of the program's attention free energies are not terms of one evidence bound at the law-point
level.
\end{verdictbox}


\chapter[The Fixed Normalized Generative Law]{The Fixed Normalized\\Generative Law}
\label{comp:ch04}

\section{The question and the answer}
\label{comp:ch04-question}

The technical chapter asks how one normalized joint law of latent samples and records is built at $c_*$, and
what a passive re-expression of the frames does to it. The construction is directed. Structural data fix a
directed acyclic graph $\Gamma$ on the agents. Each agent has a model factor, the law of its model presentation
given its parents' models, and a state factor, the law of its belief sample given its parents' beliefs and its
own model; composed, they form the agent's mechanism, and composing the mechanisms along a topological order of
$\Gamma$ gives the latent law (\MT{def:gen-parameter-split}, \MT{cons:gen-finite-directed-law}). Records are then
attached as conditionally independent children of the latents, each with an arbitrary scope. The chapter proves
that the result is normalized for every finite record hypergraph, cyclic or not (\MT{prop:gen-exact-normalization});
that locally normalized reciprocal couplings between latents can fail to give any joint
(\MT{prop:gen-gibbs-counterexample}); and that, under a covariance hypothesis on the kernels, a passive
re-expression pushes the whole joint forward, so the evidence and the bound of a transported recognition law are
unchanged (\MT{prop:gen-product-evidence-invariance}, \MT{cor:elbo-independent-frame-invariance}). It also states,
as a standing hypothesis, the typing prohibition: no generative factor may read a recognition law or the
posterior (\MT{req:gen-typing-prohibition}). The evidence of this joint is bounded in \cref{comp:ch05}, and
attention and coarse channels act on it in \cref{comp:ch06} and \cref{comp:ch09}.

What is nontrivial about normalizing on graphs with cycles? Mathematically, nothing: the proof integrates out the
records first, each has unit mass, and what remains is a directed composition of normalized kernels. This is the
standard fact that a Bayesian network is normalized whatever the undirected loops of its moral graph
\citep{Lauritzen1996}. The content lies in what the construction requires. Generative influence among latents
must be acyclic, and every interaction that closes a loop must be written as a record, a normalized law of an
observed quantity given the latents it reads. The alternative, reciprocal couplings written directly between
latents, has a parameter-dependent normalizer that can be infinite even when every local piece is normalized,
which is what the Gaussian counterexample shows. The chapter chooses the directed route because its
normalization does not depend on the parameters, and the price is that loop interactions enter only through
conditioning on records.

\section{The picture}
\label{comp:ch04-picture}

\subsection{Three graphs, not one}

A physicist will be tempted to see one interaction graph. There are three, with different roles. The directed
acyclic graph $\Gamma$ says who generates whom among latents. The interaction multigraph $\mathfrak I$ of
\cref{comp:ch02} carries the links. The record hypergraph, formed by the scopes $\partial a$, says which latents
each observation reads. The running example makes the distinction concrete: $\Gamma$ is a tree in which agent $1$
is the parent of agents $2$ and $3$, $\mathfrak I$ is the triangle, and two records with scopes $\{1,2\}$ and
$\{2,3\}$ close the latent tree into the triangle. Cycles live in the record hypergraph; $\Gamma$ never has any.

\subsection{Records close loops, and the evidence is the partition function}

When each record kernel has the form $K_a(X,y_{\partial a},do_a)=\exp(-E_{a,o}(y_{\partial a}))\nu_a(do_a)$, the
joint density is the latent density times $\exp(-\sum_aE_{a,o})$ (\MTeq{eq:gen-record-density}). Combined with
the posterior formula of \MT{hyp:prob-regular-conditional-convention}, this says that the posterior is a Gibbs
measure whose reference is the directed prior and whose Hamiltonian is the sum of the record energies, and that
its partition function is the evidence. A binary record whose success probability is $\exp(-E)$ with $E\geq0$,
observed at success, contributes exactly the Boltzmann factor $\exp(-E)$ to the posterior, and the running example
uses this device for its loop-closing records. Normalization of the joint is automatic; what can fail is
finiteness and positivity of the evidence, and the manuscript admits only observations with finite positive
evidence (\MT{sec:intro-conventions}). One physical consequence is easy to miss. Sampling the latents from the
generative model ignores the records altogether, because the records are children of the latents, not
constraints on them; the loop interaction shapes the latents only after conditioning on the observed records.

\begin{analogy}{exact}
Prior times likelihood, normalized by the evidence, is a Gibbs measure with the prior as reference and the
record energies as Hamiltonian; the evidence is its partition function. This is Bayes' rule and needs no further
argument.
\end{analogy}

\subsection{Why the loop is not a potential between latents}

The undirected alternative writes the joint as a product of potentials on agents, edges, and records, divided by
a partition function $Z_X$, and it is legitimate whenever $0<Z_X<\infty$ (\MT{def:gen-undirected-model}). The
trouble is that $Z_X$ depends on the parameters and nothing local controls it. Two scalar agents with normalized
Gaussian site potentials and the positive, finite, continuous edge potential $e^{\mathrm ck_1k_2}$ give the
quadratic form with matrix $\mathsf W_{\mathrm c}$, whose smaller eigenvalue $1-\mathrm c$ is not positive for
$\mathrm c\geq1$, and the integral diverges (\MT{prop:gen-gibbs-counterexample}). Each conditional
$k_1\given k_2$ and $k_2\given k_1$ is a perfectly normalized Gaussian, and still there is no joint.

\begin{analogy}{exact}
This is the Gaussian spin model beyond its critical coupling: once the ferromagnetic coupling exceeds the
single-site stiffness, the quadratic form loses positivity along the uniform mode and the partition function
diverges. The counterexample is the two-site case of the same integral.
\end{analogy}

A related temptation is to append an energy $U$ to the free energy of a joint that does not contain it. The
chapter's remark computes the result (\MT{rem:gen-appended-energy}): for $U$ bounded below,
$\Fenergy_V[Q;X,o]+\E_QU$ equals the free energy of the tilted joint $Z_U^{-1}e^{-U}\mathbb P_{\theta,V}$ minus
$\log Z_U$ (\MTeq{eq:gen-appended-energy}). Because $\log Z_U$ does not depend on $Q$, recognition updates under
the appended objective are those of the tilted model; because it depends on $(\theta,X)$, comparisons across
parameter values, and hence learning, are off by $\log Z_U$ unless it is restored.

\subsection{The typing prohibition}

The standing hypothesis \MT{req:gen-typing-prohibition} states that every generative factor and record kernel is a
function only of the parameters $\theta$, the structural data $X$, and its displayed latent coordinates. No
generative kernel takes as an argument the recognition law, one of its marginals or parameters, or the posterior.
In the two-argument form $\Fenergy_V(\mathbb Q;\theta)$ of the free energy \citep{Neal1998}, every evaluation holds
the generative argument fixed while the recognition argument varies, which excludes the composite
$\mathbb Q\mapsto\Fenergy_V(\mathbb Q;\theta(\mathbb Q))$ inside one evaluation. The chapter's example of what is
forbidden is a transition built from a live peer recognition marginal, such as agent $i$'s prior set to the
current transported belief of agent $j$. What is allowed is coupling through latents (agent $i$'s state factor
reads agent $j$'s sample), a parameter that depends on the observation or on an earlier recognition law provided
it is held fixed while the bound is evaluated, slow evolution of an agent's model as recognition dynamics of the
model channel or as accepted M-coordinates, each accepted value defining a new fixed joint, and loops of inference
updates under a fixed joint.

Why forbid it? The manuscript's answer is \MT{prop:gen-no-distinguished-target}: a family of joints indexed by
recognition laws satisfies the evidence identity for each member separately, but the change of the bound between
two members mixes a change of evidence with a change of gap, and the family designates no joint as the model. A
companion witness shows that the stronger failures one might expect, such as evidence that must vary along the
family, need not occur (\MT{prop:gen-moving-target-witness}). The prohibition is therefore a commitment to having
one target, not the avoidance of a contradiction.

\begin{analogy}{structural}
In naive mean-field theory for the Ising model, the Gibbs--Bogoliubov functional $\E_QH-TS[Q]$ over product trial
laws uses the fixed Hamiltonian $H=-J\sum_{\langle ij\rangle}s_is_j$, which reads samples. Scoring $Q$ against the
effective Hamiltonian $H_{\mathrm{eff}}(s;m)=-J\sum_i\bigl(\sum_{j\sim i}m_j\bigr)s_i$ with the current magnetizations
$m=\E_Qs$ inserted counts every bond twice, which is why the mean-field free energy carries its familiar correction
$+J\sum_{\langle ij\rangle}m_im_j$. Holding $m$ fixed during each variation and updating it between variations gives
the same stationarity equations. The first procedure is the composite the typing prohibition excludes; the second is
the admissible parameter held fixed while the bound is evaluated. The illustration is the companion's, not the
manuscript's.
\end{analogy}

\subsection{Passive re-expression moves the whole joint}

A passive re-expression acts on latent samples by the block map $R=\bigoplus_i(R_i^b\oplus R_i^m)$, leaves the
records alone, and transforms the frames and links in $X$ by the vertex law. The kernel-covariance hypothesis
(\MT{hyp:gen-kernel-covariance}) asks that re-expressed parameters $(\theta',X')$ exist for which every factor is
the pushforward of the old one and every record kernel reads the rotated latents as it read the old ones. Under
it, the new joint is the pushforward of the old one as a measure, with no density hypothesis, the observation
marginals agree, and the evidence is unchanged for almost every observation
(\MT{prop:gen-product-evidence-invariance}). Relative entropy is invariant under bijections, so a recognition law
transported as a whole keeps its divergence to the transported posterior and its extended bound, singular branch
included (\MT{cor:elbo-independent-frame-invariance}).

\begin{analogy}{exact}
This is a change of integration variables in a partition function with gauge-invariant observables: the records
are the observables, and the Jacobian $\lvert\det R\rvert^{-1}$ appears in densities but cancels from every ratio,
which is why the measure-level statement needs no density hypothesis.
\end{analogy}

\section{Objects and notation}
\label{comp:ch04-objects}

\begin{decoder}
Structural data $X$; parameters $\theta$ & $X$ holds the directed graph, the multigraph with links, and the frames; $\theta$ holds the numerical parameters that learning may move; the split is declared. & Lattice and background fields; couplings \\
Directed acyclic graph $\Gamma$; parents $\operatorname{pa}(i)$ & Who generates whom among latents. & A causal ordering; none in equilibrium statistical mechanics \\
Model factor $P^m_{\theta,i}$ & Law of agent $i$'s model presentation given its parents' models. & Hierarchical prior on couplings \\
State factor $P^k_{\theta,i}$; evaluated kernel $K^X_{i,m_i}$ & Law of agent $i$'s belief sample given its parents' beliefs and its own model; the same kernel indexed by the model sample. & Conditional Boltzmann weight at fixed couplings \\
Agent mechanism $G_{\theta,i}$ & Model factor then state factor, one normalized kernel from the parents' block to the agent's block. & Local stochastic update rule \\
Root; law bridge $\widetilde\Phi^{\mathrm{gen}}_i$ & An agent without parents; the belief law induced by a law over its model. & Posterior-predictive law \\
Directed latent law $\mathbb P^{\mathsf Y}_{\theta,V}$ & Ordered composition of the mechanisms along a topological order. & Bayesian-network prior \\
Record kernel $K_a$; energy $E_{a,o}$ & Normalized law of record $a$ given the latents in its scope; its negative log density at the observed value. & Likelihood factor; bond energy of the posterior \\
Joint $\mathbb P_{\theta,V}(dy,do\given X)$ & The directed latent law times the record kernels. & Prior times likelihood \\
Undirected model; potentials $\psi$; $Z_X$ & A Gibbs model on agents, edges, and records with its partition function. & Boltzmann distribution \\
Appended energy $U$; tilted joint; $Z_U$ & An energy added to a free energy without being part of the joint; the joint it would define; its normalizer. & A potential added without renormalizing \\
Typing prohibition & No generative kernel reads a recognition law, its marginals or parameters, or the posterior. & Hamiltonian independent of the trial state \\
Two-argument free energy $\Fenergy_V(\mathbb Q;\theta)$ & Free energy as a function of recognition law and parameters, varied one at a time. & Gibbs--Bogoliubov functional \\
M-coordinate; accepted value & A change of generative parameters accepted between evaluations. & M-step of expectation maximization \\
Maps $R$, $T$; gauge inert & The passive action on latents, and on latents and records; records are unchanged. & Change of variables; gauge-invariant observables \\
Kernel covariance & The model's coordinate form is compatible with re-expressed parameters. & Gauge covariance of the action \\
Reference factor $j_R$ & Radon--Nikodym factor of the reference under $R$; $\lvert\det R\rvert^{-1}$ for invertible linear blocks. & Jacobian \\
Transported coupling $R_\#\mathbb Q_{V,o,X}$ & A recognition law moved as a whole to the new frames. & Gauge-transformed trial state \\
\end{decoder}

The word "fixed" in the chapter title means held fixed while the evidence bound is evaluated; the joint changes
between evaluations when an M-coordinate is accepted, and each accepted value defines a new joint. "Model" has two
meanings: the agent's model channel ($m_i$ and its law $q_i^m$), and the generative model $\mathbb P_{\theta,V}$.
The model and state factors are generative factors and are not components of the agent of \cref{comp:ch02}, even
though they carry the agent's index; the agent's belief $q_i^b$ is a recognition object. "Cycle" refers to the
record hypergraph; $\Gamma$ is acyclic by definition. "Covariance" in the hypothesis name means gauge covariance,
not a second moment. The evidence invariance holds for almost every observation, not pointwise.

\section{The formulas, explained}
\label{comp:ch04-formulas}

\paragraph{The agent mechanism.}
\[
G_{\theta,i}\bigl(d(k_i,m_i)\given y_{\operatorname{pa}(i)},X\bigr)
=P_{\theta,i}^{m}(dm_i\given m_{\operatorname{pa}(i)},X)K^X_{i,m_i}(k_{\operatorname{pa}(i)},dk_i)
\]
(\MTeq{eq:gen-agent-mechanism}). The model presentation $m_i$ is drawn first from its parents' models, and the
belief sample $k_i$ then from its parents' beliefs using the kernel that $m_i$ evaluates to. At a root the first
factor is a proper law and the second a bridge from the agent's own model to its belief; averaging the bridge over
a model law $s$ gives the belief law $\widetilde\Phi^{\mathrm{gen}}_i(s)$ (\MTeq{eq:gen-root-law-bridge}), which may
be declared to represent the cross morphism of \cref{comp:ch02}. At a nonroot the evaluated kernel is indexed by
the parent beliefs and is not a map of law fibers.

\paragraph{The joint and its normalization.}
\[
\mathbb P_{\theta,V}(dy,do\given X)
=\mathbb P_{\theta,V}^{\mathsf Y}(dy\given X)\prod_{a\in\mathfrak A}K_a(X,y_{\partial a},do_a),
\qquad
\mathbb P_{\theta,V}(\mathsf Y\times\mathsf O\given X)=1
\]
(\MTeq{eq:gen-record-extension}, \MTeq{eq:gen-normalization}). The latent law is the ordered composition of the
mechanisms along a topological order (\MTeq{eq:gen-composition}), and each record kernel has unit mass in its own
coordinate, so integrating the records first removes the product whatever the overlaps of the scopes. Each record
is owned once by the joint, while the local free energies of \cref{comp:ch05} include every factor incident to a
coordinate, as a bond energy is counted once in the Hamiltonian but enters the local field of both of its sites.

\paragraph{The posterior as a Gibbs measure.}
Rewriting \MTeq{eq:gen-record-density} with the posterior formula of \MT{hyp:prob-regular-conditional-convention},
on the regular set of observations,
\begin{gather*}
\boldsymbol\Pi_{\theta,V,o,X}(dy)
=\frac{\exp\bigl(-\sum_{a}E_{a,o}(y_{\partial a})\bigr)\mathbb P^{\mathsf Y}_{\theta,V}(dy\given X)}{p_{\theta,V}(o\given X)},\\
p_{\theta,V}(o\given X)=\int\exp\Bigl(-\sum_{a\in\mathfrak A}E_{a,o}(y_{\partial a})\Bigr)\mathbb P^{\mathsf Y}_{\theta,V}(dy\given X).
\end{gather*}
With no records the posterior is the prior and the evidence is $1$. Cycles in the scopes make the posterior's
factor graph loopy, which matters for computing it and not at all for its normalization.

\paragraph{The running example.}
Agent $1$ is the root, with $k_1\given m_1\sim\mathcal N(0,e^{-m_1}I_3)$; agents $2$ and $3$ are its children, with
$k_2\given k_1,m_2\sim\mathcal N(\rho_b(\Theta^b_{\bar e_1})k_1,e^{-m_2}I_3)$ and
$k_3\given k_1,m_3\sim\mathcal N(\rho_b(\Theta^b_{e_3})k_1,e^{-m_3}I_3)$, and independent standard normal model
coordinates. The record on $\{2,3\}$ has success probability $\exp(-\tfrac12\lvert k_2-\rho_b(\Theta^b_{e_2})k_3\rvert^2)$.
At the conditional means of the children its residual is
\[
\rho_b(\Theta^b_{\bar e_1})k_1-\rho_b(\Theta^b_{e_2}\Theta^b_{e_3})k_1
=\rho_b(\Theta^b_{e_1})^{-1}\bigl(I_3-\rho_b(\mathrm{Hol}^b(p))\bigr)k_1,
\]
whose length is $\sqrt3$ times the component of $k_1$ orthogonal to the rotation axis of the holonomy (the
running-example paragraph after \MT{prop:gen-exact-normalization}). A record that closes the loop measures how far
going around the loop rotates the root's state; the records see the holonomy through an invariant, and with
trivial holonomy the residual vanishes.

\paragraph{The Gaussian counterexample.}
\begin{gather*}
\psi_i(k_i,m_i;X)=(2\pi)^{-1/2}e^{-k_i^2/2},
\qquad
\psi_{12}(y_1,y_2;X)=e^{\mathrm ck_1k_2},
\qquad
\mathrm c\geq1,\\
\mathsf W_{\mathrm c}=\begin{psmallmatrix}1&-\mathrm c\\-\mathrm c&1\end{psmallmatrix},
\qquad
\operatorname{spec}\mathsf W_{\mathrm c}=\{1-\mathrm c,1+\mathrm c\}
\end{gather*}
(\MTeq{eq:gen-gibbs-counterexample}, \MTeq{eq:gen-gibbs-counterexample-matrix}). The product of potentials is
$\exp(-\tfrac12k^\top\mathsf W_{\mathrm c}k)$ up to a constant; in the sum and difference coordinates the integral
factorizes, and the sum direction, with eigenvalue $1-\mathrm c\leq0$, diverges. For $0\leq\mathrm c<1$ the same
model is a legitimate joint.

\paragraph{The appended energy.}
\[
\Fenergy_V[Q;X,o]+\E_QU=\widetilde\Fenergy_V[Q;X,o]-\log Z_U,
\qquad
Z_U=\int e^{-U(y)}\mathbb P_{\theta,V}(dy,do\given X)\in(0,\infty)
\]
(\MTeq{eq:gen-appended-energy}). The left side is what one computes when an energy is simply added; the right side
is the free energy of the tilted joint, which does contain $U$, corrected by its normalizer. The identity holds on
the classical domain with $\E_Q\lvert U\rvert<\infty$.

\paragraph{Covariance.}
Under kernel covariance,
\begin{gather*}
\mathbb P_{\theta',V}(\cdot\given X')=T_\#\mathbb P_{\theta,V}(\cdot\given X),
\qquad
p_{\theta',V}(o\given X')=p_{\theta,V}(o\given X),\\
p_{\theta',V}(o,y'\given X')=\lvert\det R\rvert^{-1}p_{\theta,V}(o,R^{-1}y'\given X),
\end{gather*}
the second for almost every $o$ and the third for invertible linear blocks with Lebesgue references
(\MTeq{eq:gen-gauge-joint-pushforward}, \MTeq{eq:gen-gauge-evidence-invariance},
\MTeq{eq:gen-gauge-lebesgue-density}). For orthogonal representations $\lvert\det R\rvert=1$ and densities simply
move; for $\GL(K)$ they pick up the Jacobian, while measures, evidence, and divergences do not. The belief and model
channels may have different dimensions, and the reference must be the same within each channel. The root law
bridges transform as cross-morphism representatives do, which ties the generative construction to the geometry of
\cref{comp:ch02}. The corollary then gives
\[
\KL(\mathbb Q'_{V,o,X'}\Vert\boldsymbol\Pi_{\theta',V,o,X'})=\KL(\mathbb Q_{V,o,X}\Vert\boldsymbol\Pi_{\theta,V,o,X}),
\qquad
\Lelbo'^{\mathrm{ext}}_V(\mathbb Q'_{V,o,X'};X',o)=\Lelbo^{\mathrm{ext}}_V(\mathbb Q_{V,o,X};X,o)
\]
for $\mathbb Q'_{V,o,X'}=R_\#\mathbb Q_{V,o,X}$ (\MTeq{eq:elbo-gauge-invariance}), the recognition law being
selected before the re-expression and transported as a whole.

\section{What is proved, and how solid it is}
\label{comp:ch04-status}

\begin{scorecard}
Directed construction: \MT{def:gen-parameter-split}, \MT{cons:gen-finite-directed-law} & Latents are generated along a declared acyclic graph, model before belief within each agent; the split between structural data and parameters is declared. & \CHOICE & Excludes reciprocal latent influence; loops must be records. \\
Normalization: \MT{prop:gen-exact-normalization} & Latent composition is a probability kernel; with records conditionally independent given the latents, the joint is normalized for every finite record hypergraph, and its density is the latent density times $\exp(-\sum_aE_{a,o})$. & \SOLID & Standard in substance (Bayesian networks with observed children). \\
Local normalization fails: \MT{prop:gen-gibbs-counterexample} & Normalized agent potentials and a positive, finite, continuous edge potential need not give a finite partition function. & \SOLID & A general negative statement; the two-site Gaussian witness is textbook. \\
Undirected alternative: \MT{def:gen-undirected-model} & A Gibbs model with its partition function is a legitimate joint whenever $0<Z_X<\infty$. & \CHOICE & Admissible, but the finiteness check is parameter dependent and no general criterion is given. \\
Appended energy: \MT{rem:gen-appended-energy} & Adding a bounded-below energy to a free energy equals the free energy of the tilted joint minus $\log Z_U$. & \SOLID & Elementary; the consequence is for learning, since $\log Z_U$ depends on $(\theta,X)$. \\
Typing prohibition: \MT{req:gen-typing-prohibition} & No generative factor or record kernel reads a recognition law, its marginals or parameters, or the posterior; the composite $\Fenergy_V(\mathbb Q;\theta(\mathbb Q))$ is excluded inside one evaluation. & \CHOICE & A standing hypothesis; it admits fixed parameters updated between evaluations. \\
No distinguished target: \MT{prop:gen-no-distinguished-target}, \MT{prop:gen-moving-target-witness} & Recognition-indexed joints obey the identity member by member but designate no target; the stronger failures need not occur. & \SOLID & Justifies the prohibition as a choice of one target, not as avoidance of inconsistency. \\
Kernel covariance: \MT{hyp:gen-kernel-covariance} & Re-expressed parameters exist for which every factor and record kernel transforms covariantly. & \CHOICE & Excludes mechanisms compatible with no re-expression of their parameters. \\
Joint covariance: \MT{prop:gen-product-evidence-invariance} & The re-expressed joint is the pushforward of the old one; evidence is invariant for almost every observation; densities pick up $j_R$, equal to $\lvert\det R\rvert^{-1}$ for linear blocks; root bridges transform as cross-morphism representatives. & \SOLID & Standard in substance (change of variables). \\
Bound invariance: \MT{cor:elbo-independent-frame-invariance} & For a coupling transported as a whole, the posterior transforms by pushforward and the divergence and extended bound are unchanged. & \SOLID & A coupling reselected by a rule in the new frames is not covered; selector equivariance is declined. \\
\end{scorecard}

\begin{hygiene}
Conditions that matter: the latent graph $\Gamma$ must be acyclic, which excludes reciprocal generative influence
and forces loop interactions into records. Records must be conditionally independent given the latents; correlated
record noise has to be modeled by an added latent coordinate. Each record kernel must be a normalized law of its
record for every latent value, which is what makes normalization automatic and what an arbitrary pairwise
potential is not. The typing prohibition and kernel covariance are modeling commitments, the first excluding
peer-belief priors, the second excluding mechanisms whose coordinate form fixes a frame. Finite positive evidence
at an admitted observation is a standing default; it is automatic for binary records whose success probabilities
are positive on a set of positive prior mass, and must be checked for density records. For the undirected
alternative, finiteness of $Z_X$ is a parameter-dependent condition that can fail.

Hygiene: measurability of the kernels and of their composition, Tonelli's theorem, jointly measurable log
densities for the density formula, and the reference condition $R_\#\nu^{\mathsf Y}\ll\nu^{\mathsf Y}$ for the
density form of covariance, automatic for invertible linear blocks with Lebesgue reference.
\end{hygiene}

The chapter has no toy results: the two-site Gaussian counterexample refutes a general implication, and the
running example is an illustration.

\section{Open ends}
\label{comp:ch04-open}

The chapter records no open problem of its own, and its scope paragraph declines two things. Cross morphisms
supply no generative factor, and a model with factors in both directions between the belief and model channels
needs its own normalized joint, which the chapter does not construct; the directed order, model before belief, is
part of the construction rather than a convenience. No quotient of the parameter space by passive re-expression is
formed, and no parameter count subtracts the dimension of the vertex group, because its action can have
nontrivial stabilizers. A physicist used to counting physical degrees of freedom as total minus gauge should note
that the theory does not make this count. Two further limits follow from the construction without being flagged
as open. Reciprocal couplings between latents are admitted only through the undirected model with a verified
finite partition function, and no general criterion for that finiteness is supplied. And the one way around the
typing prohibition that keeps a single target, promoting beliefs to latent law points so that a generative factor
may read them, requires the normalized law-point prior that is \MT{open:prob-law-point-prior}.

\begin{verdictbox}{\DELIVERED}
The chapter says that it builds one normalized joint law of latent samples and records at $c_*$, normalized for
every finite record hypergraph, cyclic ones included; that local normalization of reciprocal latent conditionals
supplies no joint; and that under kernel covariance the evidence and the bound of a transported coupling are
invariant. The reading map adds that it states the typing prohibition. All of this is proved or stated as written.
The normalization and covariance results are standard in substance: a Bayesian network with observed children is
normalized whatever the loops of its moral graph, and a change of variables with gauge-invariant observables
preserves evidence and relative entropy. The substance is in the commitments the chapter makes explicit:
generative influence is acyclic, loop interactions enter as records and therefore act only after conditioning,
reciprocal latent couplings need a separate finiteness check that can fail, and no factor may read a recognition
law. The last commitment is justified as a choice of one target rather than by an inconsistency in the
alternative, which is what the cited propositions actually show. Two narrowings should be kept in view: "cycles"
means cycles of records, not of generative influence, and the invariance of the bound is for a recognition law
transported as a whole, not one recomputed by a selection rule in new frames.
\end{verdictbox}

% Companion chapters for MT chapter 5 (evidence bound) and MT chapter 6 (attention).
\chapter[One Evidence Bound and Local-Collective Free Energy]{One Evidence Bound and\\Local-Collective Free Energy}
\label{comp:ch05}

\section{The question and the answer}
\label{comp:ch05-question}

The foundations chapters (\cref{comp:ch03,comp:ch04}) built two objects and kept them apart:
the constructed joint $\mathbb P_{\theta,V}(dy,do\given X)$, one normalized law over every
latent coordinate and every record of the population, normalized on any record hypergraph,
cycles included (\MT{prop:gen-exact-normalization}); and the recognition law
$\mathbb Q_{V,o,X}$, one joint law over the same latents, selected as a coupling of the
agents' local recognition laws and correlated in general
(\MT{hyp:prob-measurable-coupling-selector}). The technical chapter, \mtc{ch:elbo}, asks what
exact relation binds the two, and how an agent's or a block's objective relates to the
population's.

The answer has two parts. The first is one number for the whole population,
\[
\Lelbo^{\mathrm{ext}}_V(Q;X,o)=\log p_{\theta,V}(o\given X)-\KL\bigl(Q\Vert\boldsymbol\Pi_{\theta,V,o,X}\bigr),
\]
defined for every probability law $Q$ on the latent space, equal to the log evidence exactly
when $Q$ is the posterior as a measure, and equal to $-\infty$ when $Q$ puts mass where the
posterior puts none (\MT{def:elbo-extended}, \MT{thm:elbo-extended-gap}). The second is the
local--global potential identity (\MT{thm:obs-local-global-potential}): hold the law of
everything outside a block $B$ fixed and change only the block's conditional recognition law;
the collective free energy then changes by exactly the outside average of the change of a local
free energy, the free energy of the block's own inference problem with its surroundings frozen.
For single agents and product recognition laws this makes the local objectives an exact
potential game with the collective free energy as potential (\MT{cor:obs-potential-game}). The
bound is a definition; every local object is the chain rule for relative entropy
(\MT{lem:prob-kl-chain-rule}) applied to it. The chapter also prices the error of summing local entropies (total correlation),
states what ELBO values cannot compare, types observations as records, counts shared records
correctly, treats links as random variables, and fixes how parameters may move. The bound is the
objective of \cref{comp:ch06,comp:ch07,comp:ch08}, and \cref{comp:ch09,comp:ch12} carry it along
evidence-preserving channels. The scope paragraph declines four claims: existence of a maximizer
over a restricted family, tightness for any restricted family, convergence of the alternating
scheme, and that a block with its local free energy is a meta-agent.

\paragraph{Is this more than the Gibbs variational principle?}
The bound is not. Reading $-\log p_{\theta,V}(o,y\given X)$ as an energy at unit temperature,
$\Lelbo^{\mathrm{ext}}_V=-\Fenergy$ and $\log Z=-\Fenergy[Q]+\KL(Q\Vert\boldsymbol\Pi)$ is the
Gibbs variational principle, the inequality behind the Gibbs--Bogoliubov--Feynman bound. What
the chapter adds is typing and bookkeeping: the target is one fixed joint, never a family
indexed by $Q$; singular trial laws get the value $-\infty$ instead of an undefined expression;
the recognition law is a correlated coupling, and the error of treating it as a product is
named and signed; and interactions enter as records, so cyclic interaction graphs sit inside a
normalized directed joint. Each is careful, and each is standard in substance.

\paragraph{Is the potential identity more than block coordinate ascent?}
In mechanism, no. With a product baseline and product laws, each agent's best response is the
textbook mean-field update \citep{Blei2017}, and the potential-game statement says these best
responses descend one functional. The identity is more general in one respect: the block law may
be any conditional kernel $r_B(dy_B\given y_{-B})$, so it covers correlated recognition and exact
block updates (a heat-bath step in law). It says nothing about simultaneous updates of several
correlated blocks, about equilibria, or about convergence.

\paragraph{What do observations as interaction records buy?}
A uniform type. Environmental observations and agent--agent interactions are both records with a
scope and a normalized kernel; pair interactions become likelihoods of observed records, which
puts a cyclic Hamiltonian inside a normalized joint and turns its partition function into the
evidence; and an environment can be rewritten as a boundary node emitting records
(\MT{thm:obs-agent-interaction-equivalence}). That rewriting is exact for the joint law, not for
variational inference over the node's internal randomness (\MT{rem:obs-environment-fiber}).

\paragraph{What converges?}
Evidence values, when each accepted parameter step follows a full exact posterior refresh
(\MT{prop:elbo-evidence-monotonicity}): classical EM monotonicity. Nothing is proved about the
parameters (\MT{open:elbo-alternating-convergence}), and finite natural-gradient steps can lower
the bound (\MT{prop:elbo-finite-step-nonmonotonicity}).

\section{The picture}
\label{comp:ch05-picture}

Read the constructed joint as a statistical-mechanics model on a hypergraph. The latent baseline
$\mathbb P^{\mathsf Y}_{\theta,V}$ is the a priori measure on configurations, correlated in general
by the directed latent model of \cref{comp:ch04}. Each record $a$ touches a set of agents, its scope
$\partial a$, and contributes a Boltzmann factor $e^{-E_{a,o}(y_{\partial a})}$, with $E_{a,o}$ the
negative log density of the record at its observed value. The sum $H_o=\sum_aE_{a,o}$ is a
Hamiltonian with one term per record, pair records being bonds. The posterior slice is the Gibbs
measure of $H_o$ and the evidence is its partition function (\MTeq{eq:obs-global-evidence},
\MTeq{eq:obs-global-posterior}):
\[
\boldsymbol\Pi_{\theta,V,o,X}(dy)=\frac{e^{-H_o(y)}\mathbb P^{\mathsf Y}_{\theta,V}(dy)}{p_{\theta,V}(o\given X)},
\qquad
p_{\theta,V}(o\given X)=\int e^{-H_o}d\mathbb P^{\mathsf Y}_{\theta,V}.
\]

\begin{analogy}{exact}
Gibbs variational principle. For every trial state
$\Fenergy^{\mathrm{ext}}_o(Q)\geq-\log p_{\theta,V}(o\given X)$, with gap $\KL(Q\Vert\boldsymbol\Pi_{\theta,V,o,X})$
and equality exactly at the Gibbs state; under absolute continuity to the baseline, finite relative
entropy, and an integrable negative energy part,
$\Fenergy^{\mathrm{ext}}_o(Q)=\E_QH_o+\KL(Q\Vert\mathbb P^{\mathsf Y}_{\theta,V})$. This is the variational
characterization of $-\log Z$ behind mean-field methods \citep{Wainwright2008}, with relative entropy to
the a priori measure in place of Shannon entropy. \MT{thm:elbo-extended-gap} is this statement for every
$Q$ in the extended reals. A point mass on a continuous configuration space
has free energy $+\infty$, as a delta-function state of a continuous system has entropy $-\infty$; that
is the singular branch.
\end{analogy}

Fix the state $y_{-B}$ outside a block $B$. Records wholly outside $B$ become constants; records
touching $B$ define a conditional Gibbs measure on the block with its own partition function
$Z_B(y_{-B})$.

\begin{analogy}{exact}
Finite-volume Gibbs specification. The local posterior
$\boldsymbol\Pi_{\theta,V,o,X,B}(dy_B\given y_{-B})\propto e^{-H_{B,o}(y_B;y_{-B})}\mathbb P^{\mathsf Y}_{\theta,V,B}(dy_B\given y_{-B})$
of \MT{lem:obs-block-posterior} is the Gibbs specification in volume $B$ with boundary condition
$y_{-B}$ \citep{Georgii1988}, with the baseline's conditional law as a priori measure; $Z_B$ is the
finite-volume partition function, and the local free energy of \MT{thm:obs-local-multiagent-elbo} is the
finite-volume free energy of a trial block state. The identification holds at every finite $N$.
\end{analogy}

Treat each agent as a player lowering its own local free energy, averaged over what the others
believe. Because each player's cost changes by exactly the change of one collective functional,
unilateral improvement is collective improvement.

\begin{analogy}{exact}
Mean-field coordinate ascent, exact on product profiles. With a product baseline
$\bigotimes_i\varrho_i$ and product laws, minimizing agent $i$'s averaged local free energy gives the
familiar $r_i\propto\varrho_i\exp(-\E_{Q_{-i}}H_{\{i\},o})$ \citep{Blei2017}, and the potential-game
corollary says these updates descend one functional. On product profiles the correspondence is exact;
the identity is strictly more general, because it also covers correlated block conditionals. The
update is the companion's orientation; the technical chapter does not display it.
\end{analogy}

\section{Objects and notation}
\label{comp:ch05-objects}

\begin{decoder}
$\mathbb P_{\theta,V}(dy,do\given X)$, constructed joint & One law for all latent states and records & A priori measure times Boltzmann factors\\
$\mathbb P^{\mathsf Y}_{\theta,V}$, latent baseline & Law of the latents before records; may correlate agents & A priori measure\\
Record $a$, scope $\partial a$, kernel $K_a$, energy $E_{a,o}$ & A datum attached to a set of agents, its likelihood, and its negative log density & Bond or plaquette term\\
$L_o$, $H_o=-\log L_o$ & Product of record likelihoods; total interaction energy & Boltzmann weight; Hamiltonian\\
$p_{\theta,V}(o\given X)$; $\boldsymbol\Pi_{\theta,V,o,X}$ & Evidence; posterior slice at the observed records & Partition function; Gibbs state\\
$\mathbb Q_{V,o,X}$, recognition coupling & Trial joint law over all agents' latents, correlated in general & Trial state\\
$\operatorname{Cpl}(\cdot)$, coupling & Set of joint laws with given marginals & None (not a coupling constant)\\
$\Lelbo^{\mathrm{ext}}_V$, $\Fenergy^{\mathrm{ext}}_V$ & ELBO and free energy for every $Q$, values $\pm\infty$ allowed & Minus the trial free energy\\
Singular branch & $Q$ charges a posterior-null set; bound $-\infty$ & Delta-function state of a continuum\\
(H1)--(H4) & Regular posterior version, finite positive evidence, absolute continuity, log-integrability & Hygiene, except (H3)\\
$\operatorname{TC}(Q)$ & Relative entropy of $Q$ from the product of its block marginals & Multi-information\\
Restriction to $\mathcal Q'$ & Best bound over an imposed recognition family & Variational ansatz\\
Block $B$, outside $-B$ & Any set of agents and its complement & Finite volume and exterior\\
$Z_B$, $\boldsymbol\Pi_{\theta,V,o,X,B}(\cdot\given y_{-B})$, $\Fenergy^{\mathrm{ext}}_{B,o}$ & Local evidence, local posterior, local conditional free energy & Finite-volume partition function, specification, free energy\\
E-coordinate; M-coordinate, accepted & Exact block posterior update; parameter step that does not lower the bound & Heat-bath update in law; EM M-step\\
Typing prohibition & No generative factor reads the recognition law or the posterior & A Hamiltonian independent of the trial state\\
\end{decoder}

Several words carry nonstandard meanings. A \emph{coupling} is a joint law with prescribed
marginals, never an interaction strength. A \emph{record} is any observed datum on a set of agents,
interactions included. \emph{Local} means block-conditional: at one base point there is no spatial
locality, and a block is any set of agents. The \emph{potential} of a potential game is
game-theoretic \citep{MondererShapley1996}. \emph{Free energy} is dimensionless, $\Fenergy=-\Lelbo$ at
unit temperature. \emph{Slice} means conditioning at the observed value, which is why posteriors hold
``for almost every $o$''. The posterior $\boldsymbol\Pi$ is unrelated to the attention prior $\pi_{ij}$
of \cref{comp:ch06}, and the Hamiltonian $H_o$ to its message spaces $H^\pm_e$.

\section{The formulas, explained}
\label{comp:ch05-formulas}

\subsection{The bound, its energy form, and the singular branch}
\label{comp:ch05-bound}

The extended evidence identity is
\begin{keyequation}{One collective evidence bound}
\[
\begin{aligned}
\Lelbo_V^{\mathrm{ext}}(Q)&=\log p_{\theta,V}(o\given X)-\KL(Q\Vert\boldsymbol\Pi),\\
\Fenergy_V^{\mathrm{ext}}(Q)&=-\log p_{\theta,V}(o\given X)+\KL(Q\Vert\boldsymbol\Pi).
\end{aligned}
\]
\end{keyequation}
The definition (\MTeq{eq:elbo-extended-definition}) never forms $\infty-\infty$: the evidence term is finite by (H2)
and relative entropy lies in $[0,+\infty]$. Under absolute continuity with respect to the baseline,
finite relative entropy, and an integrable negative part of the energy,
\[
\Fenergy^{\mathrm{ext}}_o(Q)=\KL\bigl(Q\Vert\mathbb P^{\mathsf Y}_{\theta,V}\bigr)+\E_QH_o
\]
(\MTeq{eq:obs-collective-energy-form}): a complexity cost measured against the prior plus the expected
misfit to the records. With no records the free energy is minimized at the prior with value zero; at
$Q=\boldsymbol\Pi$ it equals $-\log p_{\theta,V}(o\given X)$.

The singular branch is where (H3) fails: $Q$ charges a posterior-null set and the bound is $-\infty$,
for example a law on a proper affine subspace when the posterior has a Lebesgue density. A proper
closed support alone does not suffice: the uniform law on $[0,\tfrac12]$ is at divergence $\log2$ from the
uniform law on $[0,1]$ (\MT{prop:elbo-subspace-support-singular}). For this program the branch is not a
technicality: point-mass recognition of continuous law points lives there, which decides the status of the
softmax over current beliefs in \cref{comp:ch06}, and so does smooth recognition of an environment's
internal randomness (\cref{comp:ch05-records}).

The classical form $\log p=\E_{\mathbb Q}[\log p_{\theta,V}(o,Y\given X)-\log q(Y)]+\KL(\mathbb Q\Vert\boldsymbol\Pi)$
(\MTeq{eq:elbo-identity}) holds under (H1)--(H4), with equality exactly at the posterior as a measure;
equal coordinate marginals are not enough (\MT{thm:elbo-exact-identity}). The extended form exists because
the classical split can be undefined where the bound is not: a Cauchy trial law against a Gaussian
posterior, or an infinite-entropy countable law that is its own posterior
(\MT{prop:elbo-integrability-witnesses}). In the program's Gaussian models (H4) holds, and these witnesses
are heavy-tailed pathologies.

\subsection{Correlated recognition and total correlation}
\label{comp:ch05-tc}

For any law $Q$ of blocks and any reference laws $\varrho_j$,
\[
\KL\Bigl(Q\Bigm\Vert\bigotimes_j\varrho_j\Bigr)=\operatorname{TC}(Q)+\sum_j\KL(Q_j\Vert\varrho_j),
\qquad
\operatorname{TC}(Q)=\KL\Bigl(Q\Bigm\Vert\bigotimes_jQ_j\Bigr)
\]
in $[0,+\infty]$ (\MT{thm:elbo-total-correlation-chain}). The total correlation vanishes exactly for
product laws and is $+\infty$ for two perfectly correlated continuous copies. As a standard orientation, a
Gaussian with scalar blocks and covariance $\Sigma$ has $\operatorname{TC}=\tfrac12\log(\prod_j\Sigma_{jj}/\det\Sigma)$.
If each agent reports its own entropy and the population adds them, it computes a pseudo-ELBO with
\[
\widetilde\Lelbo=\Lelbo_V(\mathbb Q_{V,o,X};X,o)+\operatorname{TC}(\mathbb Q_{V,o,X})
\]
(\MT{prop:elbo-total-correlation-signs}): the overstatement is exactly the total correlation, and the result
need not bound the evidence. Since local marginals do not determine the coupling once two are non-Dirac
(\MT{prop:prob-marginals-do-not-determine-joint}), an exact bound needs the coupling.

\subsection{Restriction, and what ELBO values cannot compare}
\label{comp:ch05-restriction}

For any family $\mathcal Q'$,
$\sup_{\mathcal Q'}\Lelbo_V^{\mathrm{ext}}=\log p_{\theta,V}(o\given X)-\inf_{\mathcal Q'}\KL(Q\Vert\boldsymbol\Pi)$
(\MT{prop:restrict-principle}): the best bound in an ansatz is the evidence minus a reverse information
projection, larger families do at least as well, and attainment is not required. Excluding the posterior
does not make the optimized gap positive, and product families are not convex, so strict convexity gives no
uniqueness on them (\MT{rem:restrict-attainment}).

Two negative results fix what ELBO values mean. Across unrelated models the difference of optimized bounds
is an evidence difference minus a tightness difference, and binary witnesses show its sign orders neither
models nor recognition families (\MT{prop:restrict-unrelated-joint-elbo}): comparing variational free
energies of two Hamiltonians computed with different ansatzes says nothing about their true free energies.
And if the joint moves with the trial law, each $Q$ has its own evidence and the family designates no model
(\MT{prop:gen-no-distinguished-target}); the witness of \MT{prop:gen-moving-target-witness} shows such
dependence need not be inconsistent. The typing prohibition (\MT{req:gen-typing-prohibition}) therefore rests
on the absence of a selected target, not on a derived contradiction. Its physical content: a
self-consistent-field Hamiltonian $H[Q]$ is not one variational problem, and self-consistency may appear only
as stationarity of one fixed functional.

\subsection{Observations as records}
\label{comp:ch05-records}

The record hypothesis (\MT{hyp:local-interaction-kernels}) gives every record a density $e^{-E_{a,o}}$
against one dominating measure. A pair energy $E(y_i,y_j)\geq0$ enters, for example, as a binary success
record observed as a success with probability $e^{-E}$. This is how a Boltzmann factor on a cyclic graph
becomes a likelihood in a normalized directed joint; it is a modeling commitment that coupling energies are
data. The environment theorem is the randomization lemma \citep{Kallenberg2021}: every observation kernel is
the output $O=F(Y,U)$ of a node with private uniform randomness $U$, and conversely
(\MT{thm:obs-agent-interaction-equivalence}). Given $o$, however, the posterior of $(Y,U)$ sits on the level set
$\{F(y,u)=o\}$, a null curve in the Gaussian example, so every smooth trial law on $(Y,U)$ has infinite cost, and
a finite problem integrates $U$ out or restricts recognition to the fiber $\{F(y,u)=o\}$, where the cost is that of the
induced law on $Y$ (\MT{rem:obs-environment-fiber}). The information a record carries is
$I(Y;O)=\int\KL(P(dY\given o)\Vert P^Y)P^O(do)$, and a statistic $S$ of the record preserves the posterior
exactly when $I(Y;O\given S)=0$ (\MT{prop:obs-conditioning-information}), standard facts \citep{CoverThomas2006}
stated in $[0,+\infty]$.

\subsection{Local free energy, exact block updates, and the potential identity}
\label{comp:ch05-local}

The local conditional free energy of a trial block law is
\[
\Fenergy^{\mathrm{ext}}_{B,o}(r_B;y_{-B})=-\log Z_B(y_{-B})+\KL\bigl(r_B(\cdot\given y_{-B})\Vert\boldsymbol\Pi_{\theta,V,o,X,B}(\cdot\given y_{-B})\bigr)
\]
(\MTeq{eq:obs-local-vfe}): the block's log partition function with the boundary frozen plus its recognition
error. It contains every record incident to the block and no other. For $B=V$ it is the collective free
energy; for a block touched by no record, $Z_B=1$. The underlying fact is the chain rule
\[
\KL\bigl(Q\Vert\boldsymbol\Pi\bigr)=\KL\bigl(Q_{-B}\Vert\boldsymbol\Pi_{-B}\bigr)+\E_{Q_{-B}}\KL\bigl(r_B(\cdot\given Y_{-B})\Vert\boldsymbol\Pi_{B}(\cdot\given Y_{-B})\bigr)
\]
(\MTeq{eq:obs-posterior-block-kl-chain}) for $Q=Q_{-B}r_B$; informally, with $Q_{-B}$ fixed only the second
term moves. For two trial laws with the same outside marginal and finite divergence,
\[
\Fenergy^{\mathrm{ext}}_o(Q')-\Fenergy^{\mathrm{ext}}_o(Q)=\E_{Q_{-B}}\bigl[\Fenergy^{\mathrm{ext}}_{B,o}(r'_B;Y_{-B})-\Fenergy^{\mathrm{ext}}_{B,o}(r_B;Y_{-B})\bigr]
\]
(\MT{thm:obs-local-global-potential}). The fixed outside marginal is load-bearing: when two correlated blocks
change at once, each moves the other's outside marginal. Zeroing the conditional term is the exact block
E-coordinate: among laws with outside marginal $R$, the divergence is minimized, uniquely when finite, by $R$
times the exact local posterior (\MT{prop:elbo-exact-e-coordinate}). This is a heat-bath update in law, not a
mean-field update, and a damped version is also monotone, so sequential sweeps do not raise the free energy.
For singletons and product profiles with finite divergence and $\E_{Q_{-i}}\lvert\log Z_{\{i\}}\rvert<\infty$,
the outside-averaged local free energies make the agents an exact potential game in the sense of
\citet{MondererShapley1996} (\MT{cor:obs-potential-game}), with no equilibrium or convergence consequence drawn.

\subsection{Accounting, and links as random variables}
\label{comp:ch05-ledger}

With a product baseline $\bigotimes_i\varrho_i$,
\[
\Fenergy^{\mathrm{ext}}_o(Q)=\operatorname{TC}(Q)+\sum_i\KL(Q_i\Vert\varrho_i)+\E_Q\sum_{a}E_{a,o}(Y_{\partial a}),
\qquad
\sum_{i\in V}H_{\{i\},o}(y_i;y_{-i})=\sum_{a}\lvert\partial a\rvert E_{a,o}(y_{\partial a})
\]
(\MT{cor:obs-correlated-ledger}). In the first identity each record appears once and mean field is the case
$\operatorname{TC}=0$, not a premise. The second says that single-agent local energies count a record of scope
$k$ exactly $k$ times: local free energies are coordinate potentials, not shares, the double counting that
region-based approximations correct with counting numbers \citep{Yedidia2005}. A single-count additive split
exists by peeling agents off in a chosen order, each term conditioned on the agents before it.

If link values $\Theta$ are latent, with prior $Z_{\mathscr S}^{-1}e^{-\mathscr S}\nu^{\mathscr T}$ for a
probability reference $\nu^{\mathscr T}$, and recognition $R(d\Theta)Q_\Theta(dy)$,
\[
\Fenergy=\E_R\Fenergy_{\mathrm{matter}}(Q_\Theta;o,\Theta)+\E_R\mathscr S+\KL(R\Vert\nu^{\mathscr T})+\log Z_{\mathscr S}
\]
(\MT{prop:elbo-random-link-vfe}): the matter free energy at frozen links, averaged, plus the expected link
action, the link entropy cost, and the link normalizer. Structural links contribute only the matter term. The
form is structurally that of lattice gauge theory on the agent graph with links integrated against
an action \citep{Creutz1983}. The abstract random-link identity is realized by the Gaussian-generator
base and metric loop action of \cref{comp:lie-links}. It needs a normalizable link prior; for a
noncompact group such as $\GL(K,\R)$, Haar measure is not a probability reference.

\subsection{How model parameters move}
\label{comp:ch05-parameters}

A parameter step is accepted when it does not lower the bound at fixed recognition
(\MT{def:elbo-generalized-m-acceptance}); the test needs no evidence value, only
$\E_{\mathbb Q}\log(p_{\theta^{\mathrm{new}}}/p_{\theta^{\mathrm{old}}})$ (\MT{cor:elbo-acceptance-evidence-free}). If
recognition equals the old posterior,
\[
\log p_{\theta^{\mathrm{new}},V}(o\given X)\geq\Lelbo^{\mathrm{ext}}_{V,\theta^{\mathrm{new}}}(\mathbb Q)\geq\Lelbo^{\mathrm{ext}}_{V,\theta^{\mathrm{old}}}(\mathbb Q)=\log p_{\theta^{\mathrm{old}},V}(o\given X)
\]
(\MT{prop:elbo-evidence-monotonicity}), the EM monotonicity of \citet{Dempster1977} in the free-energy reading
of \citet{Neal1998}. Block E-coordinates keep only the bound nondecreasing; the transfer to the evidence needs the
full refresh. A finite natural-gradient step is a proposal: on a quadratic model, from a point displaced along the
largest-curvature direction of the preconditioned spectrum, with $d_{\max}$ the corresponding eigenvalue, a step
longer than $2/d_{\max}$ lowers the bound (\MT{prop:elbo-finite-step-nonmonotonicity}), so a line search or trust
region is needed whatever the preconditioner.

\section{What is proved, and how solid it is}
\label{comp:ch05-status}

\begin{scorecard}
\MT{thm:elbo-extended-gap}, \MT{thm:elbo-exact-identity} & One bound for every $Q$; gap $\KL(Q\Vert\boldsymbol\Pi)$; equality exactly at the posterior; classical form under (H1)--(H4) & \SOLID & Standard in substance (Gibbs variational principle); equal marginals do not give equality\\
\MT{prop:elbo-subspace-support-singular} & Laws charging posterior-null sets have bound $-\infty$ & \SOLID & Closed support alone is not enough\\
\MT{prop:elbo-total-correlation-signs} & Summing local entropies overstates the bound by exactly the total correlation & \SOLID & Standard chain rule\\
\MT{prop:restrict-principle}, \MT{prop:restrict-unrelated-joint-elbo} & Best bound in a family is evidence minus a reverse projection; ELBO values of unrelated models order nothing & \SOLID & The second is an established negative result\\
\MT{req:gen-typing-prohibition} & Generative factors may not read the trial law or the posterior & \CHOICE & Justified by \MT{prop:gen-no-distinguished-target}, which is SOLID\\
\MT{hyp:local-interaction-kernels} & Every datum, interactions included, is a dominated record on a scope & \CHOICE & Coupling energies are observed data\\
\MT{thm:obs-agent-interaction-equivalence} & Observation kernels are environment nodes $F(Y,U)$ & \SOLID & Standard randomization lemma; not for inference over $U$\\
\MT{prop:elbo-exact-e-coordinate} & Exact and damped block updates do not raise the free energy & \SOLID & Sequential only\\
\MT{thm:obs-local-global-potential} & Fixed outside marginal: collective change equals averaged local change & \SOLID & Chain rule; finite divergences\\
\MT{cor:obs-potential-game} & Agents form an exact potential game & \SOLID & Singletons, product profiles only; no equilibrium claim\\
\MT{cor:obs-correlated-ledger} & Single-count ledger; singleton energies count record $a$ exactly $\lvert\partial a\rvert$ times & \SOLID & Product baseline (used in (i)--(ii))\\
\MT{prop:elbo-random-link-vfe} & Latent links add action, link entropy, normalizer & \SOLID & Normalizable link prior needed\\
\MT{prop:elbo-evidence-monotonicity}, \MT{prop:elbo-finite-step-nonmonotonicity} & Evidence monotone under full refresh; long natural-gradient steps can lower the bound & \SOLID & Standard; values only, not parameters\\
\MT{open:elbo-alternating-convergence}, \MT{open:restrict-intrinsic-criterion} & Parameter convergence; a criterion inside one bound selecting a recognition structure & \OPEN & See \cref{comp:ch05-open}\\
\end{scorecard}

The chapter has no TOY results. Its small examples (binary moving-target and model-comparison witnesses,
Cauchy and countable integrability witnesses, the uniform diagonal) are counterexamples, SOLID at full
generality as refutations.

\begin{hygiene}
Conditions that matter. Absolute continuity (H3) separates the finite from the singular branch, and it fails
for point-mass recognition of continuous latents and for smooth recognition of an environment's randomness,
both of which the program meets. Coupling energies must be negative log densities of records with a common
dominating measure. The typing prohibition rules out self-consistent-field constructions inside the joint. The
potential identity needs the outside marginal fixed and finite divergences; the game needs product profiles.
Evidence monotonicity needs a full exact refresh, which no recognition family excluding the posterior can
perform, so with restricted families only bound monotonicity survives. Latent links need a normalizable prior,
which the Haar measure of a noncompact group is not.

Hygiene. Standard Borel spaces, regular conditional versions (H1), finite positive evidence (H2), joint
measurability, almost-everywhere qualifications on local posteriors, and log-integrability (H4) hold in the
Gaussian and finite-state models of the program.
\end{hygiene}

\section{Open ends}
\label{comp:ch05-open}

Convergence of the alternating scheme (\MT{open:elbo-alternating-convergence}) is open for models with finitely
many latent states, parameters in an open cone, $C^2$ joint densities, and an exact M-step with a unique
maximizer. The classical conditions \citep{Wu1983} bear on isolated stationary points, which this class does not
supply: the evidence may be constant along a curve of parameters. It is settled by checkable conditions,
including compact upper level sets of the evidence, met by a model of parameter dimension at least two with a
non-discrete stationary set, under which every orbit converges to one stationary point; or by a model with compact
upper level sets whose orbit has more than one limit point. For block refreshes the question is weaker still,
because a null step is accepted.

An intrinsic criterion (\MT{open:restrict-intrinsic-criterion}) is open: nothing internal to one evidence bound at
a fixed model is known to select a nondegenerate recognition structure, such as a partition of agents. Gap
minimization selects the largest family, replacing the joint gives no common evidence, and evidence-preserving
channels order bounds without selecting a partition. It is settled by a functional on channels or partitions,
relative to one joint and not monotone toward an endpoint, with a nondegenerate optimizer, or by a proof of
endpoint degeneracy for a stated class. For the program, the free energy alone does not say which agents belong
together; formation (\cref{comp:ch13}) is posed through declared families and dynamics instead.

\section{The same evidence bound with latent links}
\label{comp:ch05-latent-links}

When the link values are inferred, their law belongs inside the population recognition law.
Writing that law as $Q(d\mathring\omega,dy)=R(d\mathring\omega)Q_{\mathring\omega}(dy)$, the
collective free energy has the conditional decomposition
(\MT{prop:link-latent-normalization}),
\[
\Fenergy_V[Q]
=\KL(R\Vert\mathbb P^{\mathscr T}_\theta)
+\E_R\Fenergy_V[Q_{\mathring\omega};\Theta(\mathring\omega),o].
\]
The second term is the held-link free energy averaged over the generator recognition law; the
first pays for that law relative to the generator prior. For the Gaussian base and loop action,
it is $\KL(R\Vert\nu^{\mathscr T})+\E_R\mathscr S+\log Z_{\mathscr S}$ when those terms are
finite. Omitting its entropy or its normalizer changes the evidence bound
(\MT{prop:elbo-random-link-vfe}). This is a \SOLID{} conditional identity under the stated
evidence and density hypotheses, not an independent link penalty appended to an agent objective.
\Cref{comp:lie-links} explains the prior and highlights the field equation derived from this bound.

\begin{verdictbox}{\DELIVERED}
The technical chapter says it gives one extended evidence bound for every recognition law, a local--global
potential identity under which the local objectives form an exact potential game, observations as interaction
records, and the rules for model updates; it declines existence, tightness, convergence, and meta-agent claims.
The proved statements deliver this. The bound is the Gibbs variational principle in the extended reals, the
potential identity is the relative-entropy chain rule, the environment theorem is the randomization lemma, and
evidence monotonicity is EM monotonicity. None of these is new mathematics; each is a known identity applied with
care, and the contribution is typing discipline: one fixed joint, correlated recognition with its
total-correlation cost, a value on the singular branch, and explicit negative results on what ELBO values can
compare. One headline sentence is broader than its statement: the potential game holds for single agents with
product recognition laws, while the block identity behind it covers correlated conditionals one block at a time.
There is no parameter convergence, no criterion that selects a partition, and no statement about simultaneous
updates.
\end{verdictbox}

\chapter{Attention and Selective Communication}
\label{comp:ch06}

The attention constructions below can be read conditional on the link generators. The joint
recognition dynamics also updates the generator law (\MT{prop:link-recognition-flow}), and a recognized source contributes its
attention weight to the link current (\MT{cor:link-delta-rule}). A conditional attention identity
does not justify replacing a random exponential link by the exponential of its mean
(\MT{prop:link-conditional-transfer}; \cref{comp:lie-links}).

\section{The question and the answer}
\label{comp:ch06-question}

The technical chapter, \mtc{ch:selective-communication}, asks which agent a receiver listens to and what it
hears. For the first, attention is a posterior over a latent pointer. Each receiving row $i$ carries a
categorical label $J_i$ naming one incoming edge copy, with a structural prior row $\pi_i$ whose support is its
source set; the union of supports is the interaction graph, and a mask is a zero prior weight. The row's record
has likelihood proportional to $\exp[-D_{ij}(y)/\tau_i]$ when $J_i=j$, with $D_{ij}$ a comparison energy and
$\tau_i$ a temperature, so Bayes' formula gives the softmax
\[
\beta^P_{ij}(y)=\frac{\pi_{ij}e^{-D_{ij}(y)/\tau_i}}{\sum_k\pi_{ik}e^{-D_{ik}(y)/\tau_i}}
\]
(\MT{prop:obs-attention-elbo}). A recognition row held constant in the latent state, as a transformer computes
one row per pass, is optimally a different softmax, over averaged energies, at an exact expected relative-entropy
cost (\MT{prop:obs-attention-constant-row-gap}). For the second, communication is a composed Markov kernel,
encoder then transmission then decoder \citep{Shannon1948}, covariant under frame changes and separate
recoordinatizations of each message space (\MT{cor:obs-selective-kernel-gauge}). On the linear tier the composite
$\mathsf L_e=\mathsf L^{\rm dec}_e\mathsf L^{\rm tr}_e\mathsf L^{\rm enc}_e$ transforms like a link but may be
singular (\MT{prop:geo-selective-covariance}, \MT{prop:geo-selective-loops}). For a Gaussian comparison record the
exact receiver update has the transformer shape, but its value map is fixed by the query map, the key map, the
precision, and the receiver covariance, its residual is precision-weighted, and its weights use a predictive
covariance (\MT{prop:obs-selective-receiver-update}).

The row calculus feeds \cref{comp:ch08}, the event law is what \cref{comp:ch11} coarse-grains, and comparison
records supply endpoint energies to \cref{comp:ch10}. The scope paragraph declines three claims: that transformer
attention as a whole is derived (feed-forward blocks and layer normalization have no counterpart), that declared
per-head rows follow from a label model, and that an interface, its rank, or its coding establishes capacity,
reconstruction of the sender's belief, or shared meaning. Open are multi-head label models, edge creation and
deletion, and formation with evolving interfaces.

\paragraph{Is attention derived or declared?}
The softmax is derived; the model yielding it is declared. Given a pointer with a prior and a record exponential
in a comparison energy, the softmax is Bayes' formula, exactly as mixture responsibilities are \citep{Bishop2006};
attention as a latent alignment variable is known in machine learning \citep{Deng2018LatentAlignment}. Not derived:
why the joint should contain a pointer, why its record is exponential in $D_{ij}/\tau_i$, or which comparison
energy to use. The chapter shows the form is inhabited by normalized records (a binary success record, a Gaussian
comparison record). For this program one more caveat matters: the softmax over divergences between \emph{current
beliefs} is the generative posterior row with the law points held fixed, not a recognition optimum of a finite
bound, because point-mass recognition of a law point has infinite cost against any non-atomic law-point prior; and
the normalized law-point prior that would make such divergences evidence-bound terms at all is open
(\MT{cor:obs-selective-softmax}, \MT{prop:prob-law-point-point-mass}, \MT{open:prob-law-point-prior}).

\paragraph{Why so many attention objects?}
The prior $\pi$, the generative row $\beta^P(y)$, the recognition row $\beta^Q$, the receiver occupancy $\alpha$,
and the event law $\eta=\alpha\beta$ are laws of different things; $\beta^Q$ and $\beta^P$ differ by an exact Jensen
gap, and coarse-graining acts on $\eta$, not on rows. The table in \cref{comp:ch06-objects} sets them side by side.

\paragraph{What is the temperature?}
A parameter of the record's likelihood, not a heat-bath temperature: it sets how sharply the row concentrates, it
weights the prior relative entropy in the row functional, and in the Gaussian comparison record only
precision over temperature enters (\cref{comp:ch06-jensen}).

\paragraph{The query map, the key map, and the matrix link?}
The chapter fixes their types, transformation laws, and linear-algebraic limits, shows how the independent link
enters the key map, and, in one model, ties the value map to the query and key maps. It selects, learns, or derives
none of them (\cref{comp:ch06-sqk}).

\section{The picture}
\label{comp:ch06-picture}

A receiver has a hidden pointer; its attention record is evidence about which source the pointer names, and
$D_{ij}$ measures how badly source $j$ explains that record. The posterior over the pointer is a Gibbs distribution
over sources with degeneracies $\pi_{ij}$, energies $D_{ij}$, and temperature $\tau_i$.

\begin{analogy}{exact}
Canonical ensemble over sources, and mixture responsibilities. The row functional
$\sum_j\beta_{ij}D_{ij}+\tau_i\KL(\beta_i\Vert\pi_i)$ is energy minus temperature times entropy relative to the
degeneracies; its unique minimizer is the Gibbs row, and its minimum is $\tau_i$ times the row free energy
$-\log\sum_j\pi_{ij}e^{-D_{ij}/\tau_i}$ (\MTeq{eq:obs-selective-row-functional}). The same formula gives mixture
responsibilities \citep{Bishop2006}. Its natural-gradient flow on the simplex is the replicator equation
(\MT{cor:obs-attention-replicator}), in the sense of \citet{Amari1998}.
\end{analogy}

\begin{analogy}{exact}
Kalman update with unknown measurement origin. In the Gaussian comparison-record model, the receiver's mean given
the pointer is a Kalman update with gain $\Sigma_iS_i^{Q\top}(\Sigma^{\rm pred}_i)^{-1}$ (part (c) of
\MT{prop:obs-selective-receiver-update}), the pointer posterior weights sources by their innovation likelihoods,
and the posterior mean averages the per-source updates: the structure of data association in tracking.
\end{analogy}

\begin{analogy}{structural}
Mean-field Ising free energy. Integrating the pointer out leaves a log-sum-exp potential; in the witness of
\MT{prop:obs-attention-reduced-hessian} it is $(\omega+1)x^2/(2\tau)-\log\cosh(\delta x/\tau)$, the Curie--Weiss
form, whose curvature at the symmetric point changes sign when $\delta^2>\tau(\omega+1)$. With one latent variable
and no thermodynamic limit, this is a bifurcation of posterior modes, not a phase transition.
\end{analogy}

\begin{analogy}{structural}
General linear link fields. $\mathsf L_e$ transforms as $\mathsf L_e'=\mathsf g_i\mathsf L_e\mathsf g_j^{-1}$, like a
hopping term's link, without being a group element, as in linear lattice gauge theory
\citep{Wetterich2014LinearLattice}, whose compact-group and continuum conclusions the technical chapter says do not
follow here. The trace of a loop product of communication maps is invariant, like a Wilson loop, but is not the
holonomy of the independent links (\MT{prop:geo-selective-loops}).
\end{analogy}

\section{Objects and notation}
\label{comp:ch06-objects}

\begin{decoder}
Row $i$ or $(i,h)$; source set $\mathcal N_i$ & A slot selecting one source (a head row shares agent $i$'s state); the support of its prior & Neighborhood; mask\\
$J_i$; $\pi_{ij}$ & Latent pointer; its structural prior, fixed here & Hidden assignment; degeneracy\\
$D_{ij}(y)$; $\tau_i$; $\mathrm c_i$ & Recognition-free comparison energy; temperature; source-independent record factor & Energy; temperature; prefactor\\
$\beta^P_{ij}(y)$; $\beta^Q_{ij}$ & Posterior row at latent state $y$; recognition row, constant in $y$ & Gibbs occupation at fixed $y$; mean-field occupation\\
$\alpha^x_i$; $\eta^x_{ij}=\alpha^x_i\beta^x_{ij}$ & Receiver occupancy; law of a (receiver, source) event; mark $x$ is $P$ (generative at $y$), $\Pi$ (posterior-averaged), or $Q$ (recognition) & None; joint bond occupation\\
$\mathsf E^v_e$, $\mathsf N^v_e$, $\mathsf D^v_e$; $K^v_e$ & Encoder, transmission, decoder kernels; composed value kernel & Shannon channel\\
$\mathsf L^{\rm enc}_e$, $\mathsf L^{\rm tr}_e$, $\mathsf L^{\rm dec}_e$, $\mathsf L_e$ & Linear factors and their composite $E_j\to E_i$ & General linear link\\
$\mathsf M^{\rm val}$, $\mathsf M^{\rm out}$, $U_i$ & Value extraction, output placement, agent frame & $W_V$, $W_O$\\
$F_e=\mathsf L_e\Theta_e^{-1}$; $\mathsf L_p$ & Receiver processing after the link; loop product of communication maps & Adjoint field; Wilson-loop-like, not holonomy\\
$H^{\rm cmp}_e$, $S^Q_e$, $S^K_e$; $d_e$ & Comparison space, query and key maps; KL between their images of the two beliefs & $W_Q$, $W_K$ as measurements\\
$V_e$; $W_e$, $\widetilde W_i=W_i/\tau_i$ & Quadratic disagreement; comparison precision; tempered precision & Elastic energy; stiffness\\
Comparison record at $o_i=0$ & Pseudo-observation clamping agreement of query and chosen key & Soft constraint as data\\
$\Sigma^+_i$, $\Sigma^{\rm pred}_i$, $\mathsf L^{\rm val}_e$; $R_e$ & Posterior and predictive covariances, value map; $\Sigma_i^{-1}$-minimal right inverse of $S^Q_i$ & Kalman covariances, gain times observation map; metric pseudo-inverse\\
\end{decoder}

Glyphs collide. $\tau_i$ is neither time nor the Langevin temperature $\mathrm T$ of the formation chapters. The
posterior $\boldsymbol\Pi$, the role mark $\Pi$, and the source prior $\pi_{ij}$ are different objects; $V$ is the
agent set and $V_e$ an energy; $E_i$ is a sample fiber, $E_{a,o}$ a record energy, $\mathsf E^v_e$ an encoder, $\E$
an expectation; $L_o$, $\mathsf L_e$, and $\Lelbo$ are a likelihood, a communication map, and the ELBO. The
\emph{link} is always the invertible group-valued edge variable $\Theta_e$ of \MT{def:geo-graph-links}; $\mathsf L_e$,
which the technical chapter calls factored matrix communication, is not a link in that sense. The value
\emph{kernel} $K^v_e$ is a Markov kernel; the value \emph{map} $\mathsf L^{\rm val}_e$ is a matrix. An \emph{event} is a
sampled (receiver, source) edge. ``Derived'', ``declared'', ``absent'' in the transformer table are standings, not
proof statuses.

\begingroup\small
\begin{longtable}{@{}>{\raggedright\arraybackslash}p{0.08\textwidth}>{\raggedright\arraybackslash}p{0.21\textwidth}>{\raggedright\arraybackslash}p{0.14\textwidth}>{\raggedright\arraybackslash}p{0.47\textwidth}@{}}
\toprule
\textbf{Object} & \textbf{Law of} & \textbf{Depends on} & \textbf{Why it differs from the others}\\
\midrule
$\pi_i$ & The pointer before any record & Structure & Its support is the interaction graph; it is the prior in Bayes' formula and never reads recognition.\\
$\beta^P_i(y)$ & The pointer given state and record & $y$ & Exact posterior at a state; its average over $y$ is not a softmax of averaged energies.\\
$\beta^Q_i$ & The pointer under recognition & Averaged energies & Differs from $\beta^P$ by the Jensen gap; equals $\beta^P$ at the mean state when energy differences are affine.\\
$\alpha^x_i$ & Which receiver an event belongs to & A declared law & Not supplied by the per-receiver model; a structural selection law or the marginal of a single joint edge label.\\
$\eta^x_{ij}$ & The (receiver, source) event & $\alpha$, $\beta$ & What coarse-graining pushes forward; rows are recovered by disintegration where $\alpha>0$, and a row alone is not recursive state.\\
\bottomrule
\end{longtable}
\endgroup

\section{The formulas, explained}
\label{comp:ch06-formulas}

\subsection{The softmax, the recognition row, and the temperature}
\label{comp:ch06-jensen}

The declared model (\MT{def:obs-attention-label-model}) adds independent pointers with priors $\pi_i$ and gives row
$i$ the record density $L_i(o_i\given y,J_i=j)=\mathrm c_i(o_i,y)\exp[-D_{ij}(y)/\tau_i]$
(\MTeq{eq:obs-attention-likelihood}). Under exclusivity (no other factor reads a pointer) the posterior factors by
rows into the softmax $\beta^P_{ij}(y)$, strictly positive exactly on the source set (\MT{prop:obs-attention-elbo}).
With nonnegative energies, an included self edge with $D_{ii}=0$, and equal priors, no source outweighs the self
edge: any energy that vanishes on self-comparison, such as a divergence of a belief from itself, makes
self-attention maximal under a uniform prior. One label per row also fixes occupancy: every receiver listens to
exactly one source, and attending to none needs a null label or another joint label model.

Restricting recognition to rows independent of $y$, row $i$ contributes
\begin{keyequation}{Attention as recognition of a source label}
\[
\begin{aligned}
\Fenergy^{\rm att}_i(\beta^Q_i)&=\KL(\beta^Q_i\Vert\pi_i)+\frac1{\tau_i}\sum_j\beta^Q_{ij}\E_{\mathbb Q_Y}D_{ij},\\
\beta^{Q\star}_{ij}&\propto\pi_{ij}e^{-\E_{\mathbb Q_Y}D_{ij}/\tau_i}.
\end{aligned}
\]
\end{keyequation}
(\MTeq{eq:obs-attention-vfe}), a Gibbs row over averaged energies. With $R^Q_i$ the row free energy at averaged
energies and $R^P_i(y)$ that at $y$, the price of the constant row is
\[
R^Q_i-\E_{\mathbb Q_Y}R^P_i(Y)=\E_{\mathbb Q_Y}\KL\bigl(\beta^{Q\star}_i\Vert\beta^P_i(Y)\bigr)\geq0
\]
(\MT{prop:obs-attention-constant-row-gap}), zero exactly when every difference $D_{ij}-D_{ik}$ is almost surely
constant. The reason, informally, is that the row free energy is concave in the energies. This is the mean-field
decoupling of pointer and state, with an exact gap. If energy differences are affine in a statistic of the state,
$\beta^{Q\star}$ is the posterior row at the recognition mean (\MTeq{eq:obs-attention-affine-row}). For transported
quadratic energies $\tfrac12r_{ij}^\top W_ir_{ij}$, $r_{ij}=y_i-\rho_b(\Theta^b_{ij})y_j$, with uncorrelated receiver and
source, it is the row at the mean with source $j$'s logit lowered by
$\Tr(W_i\rho_b(\Theta^b_{ij})\operatorname{Cov}(Y_j)\rho_b(\Theta^b_{ij})^\top)/(2\tau_i)$: uncertain sources are
penalized. Computing the softmax at the mean state, as a transformer does, is exact in the affine case, and for transported
quadratic energies with uncertain sources it is not. Without
the restriction the optimal pointer law is the product of posterior rows, and correlated pointer laws pay a
total-correlation term.

The temperature scales the record's exponent: as $\tau_i\to0$ the row concentrates on minimal energy (hard
attention), and as $\tau_i\to\infty$ it returns to the prior; in the row functional it weights the prior relative
entropy, the canonical-ensemble role. In the Gaussian comparison record the noise covariance is $\tau_iW_i^{-1}$, so
$\tau_i$ is a noise scale and only $\widetilde W_i=W_i/\tau_i$ enters: temperature and precision are not separately
identifiable there. Nothing fixes $\tau_i$; the transformer's $1/\sqrt{\dim H^{\rm cmp}}$ is the declared choice
$\varsigma^2\tau_i=\sqrt{\dim H^{\rm cmp}_i}$ for key covariance $\varsigma^2I$ (\MT{tab:obs-transformer-correspondence}).

\subsection{Receiver occupancy and the event law}
\label{comp:ch06-events}

A row is a conditional law of source given receiver, so coarse-graining needs a receiver law; the event law
$\eta^x_{ij}=\alpha^x_i\beta^x_{ij}$ supplies it (\MT{def:obs-attention-event-law}). The per-receiver model does not
determine $\alpha$, which enters either as a structural selection law fixed before observation or as the receiver
marginal of a single joint edge label in a different model (\MT{hyp:rg-single-edge-attention}); the two give different
coarse rows (\mtc{prop:rg-attention-per-receiver-collapse} has masses $1/2$ against $3/4$). With uniform rows and no
self edges, blocks of size $n$ among $N$ agents have $\beta^{\rm c}_{II}=(n-1)/(N-1)$ and
$\beta^{\rm c}_{IJ}=n/(N-1)$ (\MT{cor:obs-attention-self-edge}): a coarse block listens to itself even when no fine agent
does.

\subsection{The reduced potential, its curvature, and row flows}
\label{comp:ch06-curvature}

Summing the pointer out, or optimizing a constant row, leaves $R_i(x)=-\log\sum_j\pi_{ij}e^{-E_{ij}(x)/\tau_i}$, with
\[
\nabla R_i=\frac1{\tau_i}\sum_j\beta_{ij}\nabla E_{ij},
\qquad
\nabla^2R_i=\frac1{\tau_i}\sum_j\beta_{ij}\nabla^2E_{ij}-\frac1{\tau_i^2}\operatorname{Cov}_{\beta_i}(\nabla E_i)
\]
(\MT{prop:obs-attention-reduced-hessian}). The gradient is the frozen-weight gradient (the weights' own dependence
drops out, as in the Hellmann--Feynman theorem); the Hessian is frozen-weight curvature minus the variance of the
sources' forces, the fluctuation term of any $-\log Z$. The correction scales as $\tau_i^{-2}$ against $\tau_i^{-1}$, so
low temperature favors loss of convexity. In the witness (scalar latent with baseline $\mathcal N(0,\tau/\omega)$,
sources $\pm1$, record $\mathcal N(x-j\delta,\tau)$) the curvature at the origin is $(\omega+1-\delta^2/\tau)/\tau$; at
$(\omega,\delta,\tau)=(1,2,1)$ the frozen value is $2$ and the reduced value $-2$, though prior and energies are convex.

At fixed state recognition the row follows the replicator equation with cost
$\mathsf c_{ij}=\log(\beta^Q_{ij}/\pi_{ij})+\E_{\mathbb Q_Y}D_{ij}/\tau_i$ and dissipation
$-\mathsf r_i\operatorname{Var}_{\beta^Q_i}(\mathsf c_{ij})\leq0$ (\MT{cor:obs-attention-replicator}); a moving state
adds a driving term of no definite sign. Each weight is its initial value times a finite exponential, so no support
changes at finite time: edges are neither created nor destroyed (\MT{open:obs-edge-creation}).

\subsection{Communication kernels, factored maps, and loops}
\label{comp:ch06-communication}

An interface on $e:j\to i$ is a triple of kernels $\mathsf E^v_e$, $\mathsf N^v_e$, $\mathsf D^v_e$ composed into the
value kernel $K^v_e$ (\MT{def:obs-selective-interface}); message recoordinatizations cancel, and
$K^{v\prime}_e(s',\cdot)=(\mathsf g_i)_\#K^v_e(\mathsf g_j^{-1}s',\cdot)$ (\MT{cor:obs-selective-kernel-gauge}). On the
linear tier (\MT{def:geo-selective-matrix-link}) the composite obeys $\mathsf L_e'=\mathsf g_i\mathsf L_e\mathsf g_j^{-1}$,
message bases do not enlarge the structure group, rank is at most the smallest factor rank, and a distinguished
sending subspace needs a metric (\MT{prop:geo-selective-covariance}). The shared-head case
$\mathsf L_e=U_i\mathsf M^{\rm out}\mathsf M^{\rm val}U_j^{-1}$ is a transformer's value and output maps in agent frames;
a sum of heads factoring through spaces $H_h$ has rank at most $\sum_h\dim H_h$.

Given an invertible link, $\mathsf L_e=F_e\Theta_e$ with $F_e'=\mathsf g_iF_e\mathsf g_i^{-1}$; this exists for every choice
of link, so communication does not recover the link (\MT{prop:geo-selective-loops}). The reverse copy is specified
separately and need not be an inverse, an adjoint, or of equal rank. A loop product $\mathsf L_p$ is conjugated at the
base vertex even when singular, so rank, trace, and characteristic polynomial are invariant, but it records processing
and loss of distinguishability, not holonomy: two sites with trivial holonomy already give round trips that keep one
direction and discard the other, and parallel $\GL(2)$ copies give nontrivial holonomy with arbitrary singular
communication. The shared-head family contains no $\Theta_e$, so for a loop of length $r$ its product is
$U_{i_0}(\mathsf M^{\rm out}\mathsf M^{\rm val})^rU_{i_0}^{-1}$, which depends on the loop only through its length: shared
value and output maps are blind to the links around a loop.

\subsection{Comparison maps and the dot product}
\label{comp:ch06-comparison}

Comparison is separate from content: query and key maps $S^Q_e:E_i\to H^{\rm cmp}_e$, $S^K_e:E_j\to H^{\rm cmp}_e$ and the
role divergence $d_e=\KL((S^Q_e)_\#q_i\Vert(S^K_e)_\#q_j)$ (\MT{def:obs-selective-comparison}), a cospan that maps nothing
back into the receiver. Reference maps $S^Q_{\rm ref}U_i^{-1}$, $S^K_{\rm ref}U_j^{-1}$ reproduce projected query--key
comparison; with independent links the key can be $S^K_{\rm ref}U_i^{-1}\Theta_e$, which enters only through the per-edge
twist $U_i^{-1}\Theta_eU_j$ (\MT{prop:geo-edge-twist-invariance}) and reduces to the reference form at trivial twist. The
divergence is an energy of law points, admissible only with law points latent; then the softmax over role divergences is
the posterior row at a latent configuration, and the recognition optimum reduces to it for a point-mass recognition
law, whose free energy is infinite when the baseline gives the point zero mass (\MT{cor:obs-selective-softmax}, \MT{prop:prob-law-point-point-mass}). For Gaussian law points
with a common projected key covariance $\Sigma^K$,
\[
d_{ij}=d_i^\circ-(\mu^Q_i)^\top(\Sigma^K)^{-1}\mu^K_j+\tfrac12(\mu^K_j)^\top(\Sigma^K)^{-1}\mu^K_j
\]
(\MTeq{eq:obs-selective-dot-product}): a dot-product logit plus a key-norm bias. The derived logit is minus a squared
query--key distance; it equals the transformer's dot product only for constant whitened key norms or a compensating
prior. The natural gradient of a quadratic disagreement $V_e$ fixes one receiving map
$\mathsf G_i^{-1}S^{Q\top}_eW_eS^K_e$ (\MT{prop:obs-selective-response}), so a declared attention insertion descends that
energy only when its value map has this form. When content also depends on the pointer, its conditional relative
entropy changes the optimal routing (\MT{cor:obs-selective-content-kl}); exact inference routes on content too, the
transformer on query and key only.

\subsection{The exact receiver update}
\label{comp:ch06-receiver}

The comparison record $o_i\given k_i,J_i=j,y_{-i}\sim\mathcal N(S^Q_ik_i-S^K_ek_j,\tau_iW_i^{-1})$ is observed at $o_i=0$
(\MTeq{eq:obs-selective-attention-record}): a pseudo-observation clamping agreement of the receiver's query with the
chosen key, with density $\mathrm c_i\exp[-V_e/\tau_i]$ at zero. Let the receiver's conditional law given all else be
$\mathcal N(\mu_i,\Sigma_i)$, let the copies share the query map and precision, and let this record be the only factor
reading the pointer. With $\Sigma^+_i=(\Sigma_i^{-1}+S^{Q\top}_i\widetilde W_iS^Q_i)^{-1}$,
$\Sigma^{\rm pred}_i=\widetilde W_i^{-1}+S^Q_i\Sigma_iS^{Q\top}_i$, and $r_{ij}=S^Q_i\mu_i-S^K_ek_j$,
\[
w_{ij}\propto\pi_{ij}e^{-\frac12r_{ij}^\top(\Sigma^{\rm pred}_i)^{-1}r_{ij}},
\qquad
\E[k_i\given y_{-i},o]=\Sigma^+_i\Sigma_i^{-1}\mu_i+\sum_jw_{ij}\mathsf L^{\rm val}_ek_j,
\qquad
\mathsf L^{\rm val}_e=\Sigma^+_iS^{Q\top}_i\widetilde W_iS^K_e
\]
(\MT{prop:obs-selective-receiver-update}). The weights integrate the receiver state out, so they use the record noise
plus the receiver's own projected uncertainty; $\Sigma^+_i\Sigma_i^{-1}\mu_i=(I-\Sigma^+_iS^{Q\top}_i\widetilde W_iS^Q_i)\mu_i$
is the receiver's mean shrunk by the record's information, in place of an identity skip connection; and
$\mathsf L^{\rm val}_e$ is a Kalman gain times the key map. In the scalar example with sources at $0$ and $2$,
$w_{i2}=e^{-1}/(1+e^{-1})$ while the posterior row at the mean gives $e^{-2}/(1+e^{-2})$, and the value map is $1/2$. For a
full-row-rank query map, as the record noise vanishes, $\mathsf L^{\rm val}_e\to R_eS^K_e$, the key map followed by the metric right inverse of the
query map (\SG{prop:geo-selective-lift}). The update is exact for one head, one receiver, and fixed source states: a
heat-bath step on one coordinate, not joint inference over all agents.

\subsection{The transformer correspondence and heads}
\label{comp:ch06-transformer}

The technical correspondence table (\MT{tab:obs-transformer-correspondence}) gives each row a standing; the last column is
the companion's reading.

\begingroup\small
\begin{longtable}{@{}>{\raggedright\arraybackslash}p{0.16\textwidth}>{\raggedright\arraybackslash}p{0.26\textwidth}>{\raggedright\arraybackslash}p{0.09\textwidth}>{\raggedright\arraybackslash}p{0.41\textwidth}@{}}
\toprule
\textbf{Transformer} & \textbf{Theory counterpart} & \textbf{MT} & \textbf{Companion reading}\\
\midrule\endhead
\bottomrule\endfoot
Token position & Agent; embedding as a latent coordinate & declared & A naming choice.\\
Query, key projections & Reference maps; role divergence; dot product plus key-norm bias & derived & Logit form derived for Gaussian law points with common key covariance; the maps are declared.\\
One key map for all edges & Links enter only through the twist; shared map is trivial twist & derived & Correct; nontrivial twist and holonomy have no counterpart, and nothing is said about their effect.\\
Softmax row & Posterior row of the pointer; at a latent configuration over law-point divergences & derived & From a declared exponential likelihood; over current beliefs it is a posterior at fixed law points, never a finite-cost recognition optimum against a non-atomic law-point prior.\\
Mask & Zero prior weight & derived & Immediate.\\
Temperature, $1/\sqrt d$ & $\tau_i$; scaled dot product one choice & declared & Free; not identifiable from precision in the Gaussian record.\\
Value, output, residual & Mean operator declared; comparison-record update fixes value map and residual weight & derived & What is derived is a different, tied value map and a shrinkage residual; the transformer's free value and output maps are declared.\\
Multiple heads & Declared per-head rows & absent & Per-head softmaxes are exact at a fixed receiver state and are not the integrated label posterior once the shared state is integrated out.\\
Feed-forward, layer norm & None & absent & None.\\
\end{longtable}
\endgroup

The table has no row for stacking layers or for training. On heads: the declared mean operator
(\MTeq{eq:obs-selective-attention-residual}) is multi-head attention in operator form with declared rows. Two heads with
their own pointers and comparison records have factorized pointer posteriors at fixed receiver state, but not after the
state is integrated out; in the scalar instance with sources at $0$ and $a$ the odds ratio between agreeing and
disagreeing label pairs is $e^{a^2/3}$ (\MT{prop:obs-selective-head-coupling}). The general Gaussian criterion is the log
cross ratio $-(\mathsf a_j-\mathsf a_{j'})^\top\Lambda^{\rm rec}_{12}(\mathsf b_l-\mathsf b_{l'})$
(\MT{rem:obs-head-coupling-criterion}); correlated projected queries alone do not force dependence.

\subsection{What survives an interface}
\label{comp:ch06-information}

Information about the source cannot grow along encoder, channel, and decoder given receiver side information
(\MT{prop:obs-selective-information-chain}). Agreement on communicated coordinates leaves a private conditional divergence,
$\KL(q\Vert p)=\KL(q_M\Vert p_M)+\int\KL(q(\cdot\given M=z)\Vert p(\cdot\given M=z))q_M(dz)$
(\MT{prop:obs-selective-private-kl}), so agents agreeing on everything communicated may disagree by any amount. For surjective linear
measurements $M_j$, $M_i$, $M$ that retain some coordinates (interfaces), a linear map descends to the retained coordinates
exactly when $\mathsf L_e\ker M_j\subseteq\ker M_i$, a linear flow exactly when its generator
preserves $\ker M$, a moving measurement exactly when $\dot M_t+M_t\mathcal A_t$ vanishes on $\ker M_t$, and a low-rank interface
is not by itself a Markov blanket (\MT{prop:geo-selective-descent}). The flow condition is the linear analogue of exact
lumpability \citep{KemenySnell1976}.

\subsection{What the chapter establishes about the query map, the key map, and the matrix link}
\label{comp:ch06-sqk}

Established: the types of $S^Q_e$ and $S^K_e$ (\MT{def:obs-selective-comparison}); the transformation law of $S^Q_e$
under frame change and comparison recoordinatization, stated in Supplement G (\SG{prop:geo-selective-lift}), with the law of
$S^K_e$ following the same pattern by the reference-map form and the kernel-level law \MT{cor:obs-selective-kernel-gauge};
the types and transformation laws of $\mathsf L_e$, with message coordinates internal (\MT{prop:geo-selective-covariance});
the entry of the independent link into attention through the key map
$S^K_{\rm ref}U_i^{-1}\Theta_e$, only via its per-edge twist, so that a shared key map is the trivial-twist case; the separation
of the matrix link from the gauge link, $\mathsf L_e=F_e\Theta_e$ with $\Theta_e$ unrecoverable, loop products that are not
holonomy, and holonomy independent of communication rank; the tie between value and comparison, the receiving map
$\mathsf G_i^{-1}S^{Q\top}_eW_eS^K_e$ of a quadratic disagreement and the exact value map
$\Sigma^+_iS^{Q\top}_i\widetilde W_iS^K_e$ of the Gaussian comparison-record model; and linear limits (rank bounds, descent
conditions, rank not being capacity).

Not established: any rule that selects or learns $S^Q_e$, $S^K_e$, or $\mathsf L_e$; the technical chapter states that
nothing there selects the content maps. \MT{def:obs-selective-interface-learning} lists what a learned or random interface
must carry (every varying normalizer in its score, or a normalized prior and declared dynamics with a regular metric) and
constructs no learning rule. A declared content map $\mathsf L_e$ and the comparison-induced value map are two mechanisms the
chapter does not unify; Supplement G gives $\mathsf L_e$ its own linear Gaussian record (\SG{prop:obs-selective-gaussian-record}).
Evolving interfaces are \MT{open:obs-selective-formation}. The linear statements are on real vector fibers and use no
compactness, so they cover $\GL(K,\R)$ links in the defining representation; complex fibers are not treated separately.

\section{What is proved, and how solid it is}
\label{comp:ch06-status}

\begin{scorecard}
\MT{def:obs-attention-label-model} & Attention is a latent pointer with a structural prior and an exponential record & \CHOICE & Exclusivity; one label per row fixes occupancy\\
\MT{prop:obs-attention-elbo} & Posterior row is the softmax; constant-row optimum is a softmax over averaged energies & \SOLID & Standard in substance (Bayes)\\
\MT{prop:obs-attention-constant-row-gap} & Exact Jensen-gap cost of a constant row & \SOLID & Zero only for constant energy differences\\
\MT{def:obs-attention-event-law}, \MT{hyp:rg-single-edge-attention} & Occupancy and event law are declared & \CHOICE & The two choices give different coarse rows\\
\MT{cor:obs-attention-self-edge}, \MT{cor:obs-attention-replicator} & With uniform rows and occupancy, a fine no-self convention is not preserved under blocking; row flow is replicator & \SOLID & Flow never changes the support\\
\MT{prop:obs-attention-reduced-hessian} & Reduced curvature is frozen curvature minus force variance; convexity can be lost & \SOLID & Witness one-dimensional (\cref{comp:ch06-toy})\\
``Seed of symmetry breaking'' (MT summary) & The curvature correction drives later instability & \VAGUE & No statement behind it\\
\MT{def:obs-selective-interface}, \MT{def:geo-selective-matrix-link}, \MT{def:obs-selective-comparison} & Kernels, factored maps, query and key maps, mean operator & \CHOICE & All declared; nothing selects them\\
\MT{cor:obs-selective-kernel-gauge}, \MT{prop:geo-selective-covariance}, \MT{prop:geo-selective-loops} & Covariance; rank bounds; $\mathsf L_e=F_e\Theta_e$; loops are not holonomy & \SOLID & Elementary; small examples are refutations\\
\MT{prop:obs-selective-response}, \MT{cor:obs-selective-softmax} & Receiving map of a disagreement; softmax over role divergences; dot product plus key-norm bias & \SOLID & Common key covariance; point-mass recognition singular\\
\MT{prop:obs-selective-receiver-update} & Tied value map, shrinkage residual, predictive weights & \SOLID & Gaussian comparison record; one head; fixed sources\\
\MT{prop:obs-selective-head-coupling}, \MT{rem:obs-head-coupling-criterion} & Shared receiver state couples head pointers; exact criterion & \SOLID & Refutes per-head factorization\\
\MT{prop:obs-selective-information-chain}, \MT{prop:obs-selective-private-kl}, \MT{prop:geo-selective-descent} & Data processing; private divergence; descent conditions & \SOLID & Standard in substance\\
\MT{open:obs-edge-creation}, \MT{open:obs-multihead-attention}, \MT{open:obs-selective-formation}, \MT{open:prob-law-point-prior} & Edge kinetics; head recursion; moving interfaces; law-point prior & \OPEN & See \cref{comp:ch06-open}\\
\end{scorecard}

\begin{hygiene}
Conditions that matter. Label exclusivity makes the posterior factor by rows; the exponential record form produces the
softmax; one label per row forces a declared $\alpha$; the constant-row restriction carries an exact cost. Energies must
not read recognition, so divergences between beliefs need latent law points and hence the open law-point prior. The
receiver update needs a Gaussian conditional receiver law, a shared query map and precision, and fixed sources; the dot
product needs a common key covariance; the link factorization needs invertible links; the interaction graph is fixed.

Hygiene. Finite energies with finite recognition means, measurability, standard Borel spaces, $C^2$ energies for the
curvature formula, and versions chosen off null sets hold in the program's Gaussian and finite models.
\end{hygiene}

\section{Toy versus general}
\label{comp:ch06-toy}

The small examples (the one-dimensional curvature witness, the scalar receiver and two-head instances, the two-site loop,
the parallel $\GL(2)$ copies) each refute a general implication, and as refutations they are SOLID. The mechanism-level
content, that integrating out a pointer subtracts a force variance from the curvature and that a shared receiver state
couples heads, is proved in general form. The threshold $\delta^2=\tau(\omega+1)$ and the odds ratio $e^{a^2/3}$ belong to
their instances, and nothing here shows that attention-induced loss of convexity yields symmetry breaking in networks of
hundreds of agents with $\GL(K)$ links, or how strong head coupling is at scale; those consequences are conjectural until
proved there.

\section{Open ends}
\label{comp:ch06-open}

Edge creation and deletion (\MT{open:obs-edge-creation}): the graph is the prior's support and the replicator flow
preserves supports. The problem fixes a threshold rule (delete a copy whose recognition weight falls below a threshold,
create one whose averaged comparison energy falls below a threshold, renormalize) and asks for a normalized joint with the
edge set as a structural coordinate under which the rule's master equation does not increase the collective free energy, or
a proof that no joint of that class works for any such rule. In plain terms: is there a free energy that edge kinetics
descend?

Label-level multi-head attention (\MT{open:obs-multihead-attention}): heads sharing a receiver state have coupled pointers,
the event law has no head index, and under exact coarse-graining the rows a parent must read through one coordinate multiply
with depth unless a head-merging channel, with its own conditional relative entropy, is declared. The problem asks for an
exact recursion with bounded head count in the coupled class, or a block whose posterior mean no comparison-record model with
fewer heads than rows reproduces.

Formation with evolving interfaces (\MT{open:obs-selective-formation}): for Gaussian agents whose query and key maps are
fixed-rank interfaces learned slowly at mobility $\epsilon_{\mathrm M}$, with noisy mean dynamics at temperature $\mathrm T$
and pointers redrawn from posterior rows, does a finite-dimensional Markov parent reproduce the path law of the communicated
coordinates with vanishing path relative entropy on horizons of order $1/\epsilon_{\mathrm M}$, uniformly in small
$\mathrm T$? Known are the full-dimension alternative and the pointwise descent criterion for a supplied interface path.

The law-point prior (\MT{open:prob-law-point-prior}) belongs to the probability chapter but bears on this one: without it,
divergences between beliefs are separately declared potentials rather than evidence-bound terms. Even with it, a softmax
computed at fixed law points is a generative posterior quantity or a function of parameters, not a recognition optimum,
because point-mass recognition of a non-atomic law point has infinite cost (\MT{prop:prob-law-point-point-mass}).

\begin{verdictbox}{\DELIVERED}
The technical chapter says that attention enters the joint as a latent source label whose generative posterior is the
softmax row, that the optimal constant recognition row is a different softmax at an exact cost, that communication is a
normalized encoder--channel--decoder kernel, and that the exact receiver update of a comparison record fixes the value map
by the query and key maps; it declines to claim that transformer attention as a whole, per-head rows, or channel capacity is
derived. The proved statements deliver exactly this. The substance is standard: the softmax is Bayes' formula for a declared
pointer, the receiver update is Gaussian conditioning with a Kalman gain, the row flow is the replicator equation. Specific
to the theory are where the independent link enters (the key map, through its twist), the separation of communication maps
from links and of loop products from holonomy, and the finding that exact inference ties the value map to the query and key
maps, shrinks the residual by precision, and weights by predictive likelihood. One table standing is generous: the
``derived'' value-and-residual row derives a tied substitute, not the transformer's free value and output maps. The chapter
does not select or learn the query, key, or communication maps, and gives no label-level multi-head model, no edge kinetics,
and no formation with moving interfaces; the softmax over current beliefs is a posterior at fixed law points rather than a
recognition optimum, and the law-point prior that would make belief divergences evidence-bound terms is open.
\end{verdictbox}

\chapter{Fiber Information Geometry and Inference Histories}
\label{comp:ch07}

The same configuration geometry includes the Gaussian recognition law of latent link generators.
Its Fisher natural-gradient flow is first order and decreases the collective free energy
(\MT{prop:link-recognition-flow}; \cref{comp:lie-links}). The episode-history construction adds
inertia to a most probable history across episodes, a different optimization problem.

\section{The question and the answer}
\label{comp:ch07-question}

\Cref{comp:ch05} left each agent at $c_*$ as a pair of points on law fibers, a belief law $q_i^b$
and a model law $q_i^m$, with one collective free energy for the population. \MT{ch:relational-inference}
adds what a dynamics and a comparison of scales need: a geometry on the law fibers, and a rule that
turns a free energy into a path with a length. It asks how the collective free energy selects an
inference history and what duration the Fisher geometry assigns to it, and, for meta-agents, what
happens to both when a block of agents is replaced by a coarser description.

The answer has five parts. The Fisher metric and the cubic Amari--Chentsov tensor of each law fiber
are gauge invariant, hence frame-independent at $c_*$ (\MT{prop:pb-statistical-tensor-descent}), but
the Fisher metric is chosen: none of the axioms that single it out is imposed. A Markov kernel that
does not depend on the parameter replaces each score by its conditional expectation given the coarse
record and loses Fisher information by exactly the expected conditional covariance of the score
(\MT{thm:cg-fisher-contraction}), the data-processing inequality for Fisher information with its
defect written out. A free energy selects a path only after a metric or mobility is chosen; the
steepest-descent ray then defines an oriented orbit whose Fisher arc length, the information clock,
is reparameterization invariant (\MT{thm:hist-fisher-clock-invariance}). The image of a history under
one parameter-independent kernel, such as the law of a meta-agent when that law is such an image,
accumulates no more duration (\MT{thm:hist-record-clock-contraction}), while a coarse orbit
recomputed from its own free energy is comparable with the fine one only through an oriented
semiconjugacy, whose time change cancels from every duration (\MT{thm:hist-oriented-semiconjugacy},
\MT{thm:hist-duration-relation}). Finally, a three-agent parity witness shows that the collective
free energy and its Fisher metric are not functions of single-agent laws
(\MT{prop:hist-parity-hyperedge-vfe}, \MT{cor:hist-parity-fisher-residual}), and two finite theorems
show that naturality under local kernels forces the product law and that, once correlated refinements
are admitted, no reference-free selector exists (\MT{thm:hist-local-markov-product-uniqueness},
\MT{thm:hist-correlated-refinement-nogo}).

The chapter consumes the collective free energy and the local--global identity of \cref{comp:ch05}.
\Cref{comp:ch08} uses its mobility and clock results, \cref{comp:ch09} the Fisher-loss theorem, and
\cref{comp:ch11,comp:ch12} the coarse configuration map and record clocks. Most of the mathematics is
standard (Fisher information under Markov kernels, arc length, gradient flows with a mobility,
{\L}ojasiewicz finite length); the chapter's contribution is the typing at $c_*$ and a counterexample
at each boundary where a standard statement stops applying to the program's objects.

\section{The picture}
\label{comp:ch07-picture}

The Fisher metric measures distinguishability: to second order the relative entropy between $p$ and
$p+dp$ is $\tfrac12g^F(dp,dp)$. For an exponential family $p_\lambda\propto e^{\lambda\cdot A(y)}$
the score in direction $\lambda_a$ is the centered observable $A_a-\langle A_a\rangle$ and the Fisher
metric is $\operatorname{Cov}(A)$, the static susceptibility matrix up to a factor of temperature.
Most of the chapter consists of statements about fluctuations of scores.

\begin{analogy}{exact}
Fisher loss under coarse observation is the law of total covariance. In an exponential family the
fine information is $\operatorname{Cov}(A)$, the coarse one $\operatorname{Cov}(\E[A\given\bar Y])$,
and the loss $\E\operatorname{Cov}(A\given\bar Y)$: the fluctuation of the conjugate observable inside
a coarse cell, invisible to a probe that sees only the record $\bar Y$. For a stochastic Kadanoff block
rule applied to spins in a field, the Fisher information of the coupling $\beta h$ seen through the
block spins is the variance of the conditional mean magnetization. \MT{thm:cg-fisher-contraction} holds for every
regular family, and it requires a blocking rule that does not depend on the parameter.
\end{analogy}

A free energy alone does not define a dynamics: in overdamped physics
$\dot x=-\mathsf M\nabla\Fenergy$, the potential fixes the equilibria and the mobility the paths. The
natural gradient \citep{amari1998natural} takes a chosen metric as inverse mobility, and no
fluctuation--dissipation relation ties that mobility to anything in the model; it is a modeling
choice, and an anisotropic one can change the orbit (\cref{comp:ch07-formulas-orbits}).

\begin{analogy}{structural}
The Fisher duration of a path of laws is the thermodynamic length of \citet{Crooks2007}; for Gibbs
laws and paths in natural parameters the mathematical object is the same. What does not transfer is
its physical use, bounds on dissipation for driving at finite speed, which needs a protocol and a
physical time. The chapter supplies neither and declines to call its duration time
(\cref{comp:ch07-formulas-clock}).
\end{analogy}

\begin{analogy}{structural}
A loop holonomy $\exp(t\xi)$ near the identity acting on an agent's law costs relative entropy
$\tfrac12t^2\mathfrak k_q(\xi,\xi)$, where $\mathfrak k_q(\xi,\eta)=g^F_q(\xi_{\mathcal B}(q),
\eta_{\mathcal B}(q))$ vanishes exactly on the isotropy algebra of the law
(\MT{cor:pb-orbit-fisher-form}). This has the algebraic form of the gauge-boson mass matrix
$M^2_{ab}\propto\langle T_av,T_bv\rangle$, which vanishes on the unbroken generators: both pull a
positive form back along an orbit map. The correspondence is one of form; no mass, link dynamics, or
symmetry breaking is claimed.
\end{analogy}

\begin{analogy}{exact}
At its center the parity witness (\MT{def:hist-three-agent-parity-lift}) is the Gibbs law of six
Ising spins $\sigma=(-1)^y$ with a single six-spin coupling, $\tfrac1{64}(1+\bar\varkappa\prod\sigma)
\propto e^{J\prod\sigma}$ with $\tanh J=\bar\varkappa$. Every correlation of fewer than six spins
vanishes, so a mean-field description from single-spin magnetizations reproduces all marginals and
misses the coupling; the free energy and Fisher metric do not (\cref{comp:ch07-formulas-lift}). The
identification is exact at the center only.
\end{analogy}

\section{Objects and notation}
\label{comp:ch07-objects}

\begin{decoder}
Law fiber $\mathcal B_x$, $x\in\{b,m\}$ & Manifold of laws an agent's belief or model can take at $c_*$ & Statistical manifold of one site\\
Score $\mathbf s_Z$ & Derivative of $\log p$ along $Z$; mean zero & Centered conjugate observable\\
Fisher metric $g^F$ & Covariance of scores & Susceptibility matrix\\
Amari--Chentsov tensor $T^F$ & Third moment of scores & Third cumulant; not a torsion\\
$\mathfrak k_q$, $\mathfrak g_q$ & Cost of moving a law along a group direction; directions that do not move it & Mass matrix; unbroken generators\\
Markov coarse morphism $\mathscr K$, kernel $K$ & Coarse observation by one fixed, frame-covariant stochastic rule & Stochastic block-spin map\\
Vertical Fisher defect $\Delta_F^{\mathscr K}$ & Fisher information lost by $\mathscr K$ on the fibers at $c_*$ & Within-block fluctuation\\
Configuration manifold $\mathcal Q_i$; metric $\mathsf G$, weights $w_b,w_m$ & Belief and model laws of an agent together; chosen metric on them & Configuration space; inverse mobility\\
Exact lift $\iota_i$, extraction $\operatorname{ext}_i$ & Supplied rule building a joint recognition law from a configuration, and its left inverse & Closure relation\\
$\mathsf G_{\mathrm J}$, $\mathsf G_w$, $\mathfrak K$, $\mathfrak C$ & Joint Fisher metric; weighted marginal metrics; genuinely joint residual; signed cross term & Higher-order versus pairwise correlations\\
Natural selector & Rule from marginals to a joint law commuting with every admitted kernel & Closure; mean field when forced\\
$\varkappa$, $\bar\varkappa=\varkappa/64$ & Parity amplitude; six-bit correlation at the center & Six-spin coupling\\
History; natural-gradient history & Oriented curve modulo reparameterization; orbit of $-\operatorname{grad}^{\mathsf G}\Fenergy$ & Trajectory without time label; overdamped path\\
Fisher duration $\mathfrak t$ & Fisher arc length from a chosen origin (the information clock) & Thermodynamic length\\
Orthogonal clock $\mathfrak t^\perp$ & Unit-rate clock whose level sets are orthogonal to descent & Synchronous time slicing\\
$\mathsf F_\ell$; oriented semiconjugacy $(\mathsf F_\ell,\mathsf a_\ell)$ & Fine-to-meta-agent configuration map; it carries the fine flow onto the coarse flow up to an increasing time change & Exact dynamical coarse-graining\\
Depth $\ell$ & Number of coarse-graining steps & Blocking steps; not time\\
\end{decoder}

``Natural'' has two unrelated senses in this chapter: Amari's metric-raised gradient, and a selector
that commutes with a class of maps in the category-theoretic sense. ``Metric'' is a chosen Riemannian
metric on laws or configurations, never a spacetime metric. ``History'' is an orbit with its
parameterization forgotten; ``duration'' and ``clock'' mean Fisher arc length, not time.
``Vertical'' means along the fibers over the single base point, so no connection enters. A ``lift''
reconstructs a joint law from marginal data; it is not the horizontal lift of a connection. Glyphs
collide across adjacent chapters: $\operatorname{par}$ is the parity character here and the parent map
in \cref{comp:ch08}, and $\mathsf a$ is a rate factor here and a proposal weight there. $\mathfrak K$
(residual tensor), $\mathfrak k_q$ (orbit form), $\mathscr K$ (coarse morphism), and $K$ (its kernel)
are four different objects, and ``semiconjugacy'' is the dynamical-systems term.

\section{The formulas, explained}
\label{comp:ch07-formulas}

\subsection{Fiber tensors, and why the metric is a choice}
\label{comp:ch07-formulas-tensors}

\begin{align*}
g^F_{x,p}(Z_1,Z_2)&=\E_p[\mathbf s_{Z_1}\mathbf s_{Z_2}],
&
T^F_{x,p}(Z_1,Z_2,Z_3)&=\E_p[\mathbf s_{Z_1}\mathbf s_{Z_2}\mathbf s_{Z_3}].
\end{align*}
The metric is the covariance of the score and the cubic tensor its third moment, the second and third
cumulants of conjugate observables in an exponential family. Both are preserved by the gauge action
because that action is a pushforward by a map of samples (\MT{prop:pb-statistical-tensor-descent}).
Invariance of the metric alone would not suffice: a M\"obius isometry of the normal family's
hyperbolic Fisher metric preserves $g^F$ and not $T^F$, so a symmetry of the metric that does not
relabel samples is not a statistical symmetry. The Fisher metric itself is not derived.
\MT{prop:ig-fisher-characterizations} collects the uniqueness theorems of \v{C}encov type
\citep{Cencov1982}, and the scope paragraph states that none is imposed; on configurations the metric
carries channel weights $w_b,w_m>0$ that nothing fixes (\MT{def:hist-configuration-manifold}), and a
rate ratio enters the same way in \cref{comp:ch08}. Every orbit and duration below inherits these
choices.

\subsection{Relative entropy to second order, and the cost of a holonomy}
\label{comp:ch07-formulas-orbit}

\[
\KL\bigl(q\Vert\widehat\rho(\exp t\xi)q\bigr)=\frac{t^2}{2}\mathfrak k_q(\xi,\xi)+O(t^3),
\qquad
\mathfrak k_q(\xi,\eta)=g^F_q\bigl(\xi_{\mathcal B}(q),\eta_{\mathcal B}(q)\bigr).
\]
This is \MT{cor:pb-orbit-fisher-form}, resting on the standard fact \citep{AmariNagaoka2000} that
the Hessian of relative entropy is the Fisher metric (\MT{prop:pb-kl-divergence-jets}). The
velocity $\xi_{\mathcal B}(q)$ of the law under the subgroup generated by $\xi$ has squared Fisher
length $\mathfrak k_q(\xi,\xi)$; the null space is the isotropy algebra $\mathfrak g_q$, on which the
left side vanishes for every $t$. For a Gaussian $\mathcal N(\mu,\Sigma)$ with $\GL(K)$ acting
linearly on samples the velocity is $(\xi\mu,\xi\Sigma+\Sigma\xi^\top)$, so $\mathfrak g_q$ consists of
the $\xi$ with $\xi\mu=0$ and $\xi\Sigma+\Sigma\xi^\top=0$; for a centered isotropic law it is all of
$\mathfrak{so}(K)$, so a rotation holonomy is invisible to that agent while a squeeze is always
detected. In the running three-agent loop this is the second-order form of law coherence up to
holonomy stabilizers (\MT{def:geo-law-coherence}).

\subsection{Coarse observation: the Fisher-loss theorem and its limits}
\label{comp:ch07-formulas-contraction}

For a regular family $\{P_\lambda\}$ of laws of $Y$ and a Markov kernel $K$ producing a record
$\bar Y$, the same kernel for every $\lambda$, \MT{thm:cg-fisher-contraction} states
\[
\bar{\mathbf s}_{\lambda_0}(\bar Y)=\E\bigl[\mathbf s_{\lambda_0}(Y)\given\bar Y\bigr],
\qquad
\mathsf I_{\mathsf Y}(\lambda_0)-\mathsf I_{\bar{\mathsf Y}}(\lambda_0)
=\E\operatorname{Cov}\bigl(\mathbf s_{\lambda_0}(Y)\given\bar Y\bigr)\succeq0 .
\]
The coarse score is the best estimate of the fine score from the record, and the second identity is
the law of total covariance for a mean-zero score. Equality in a direction $z$ holds exactly when
$z^\top\mathbf s_{\lambda_0}$ is a function of $\bar Y$. A bijective $K$ loses nothing; a $K$ that
ignores its input loses everything; discarding an independent component loses exactly its Fisher
information. No common dominating measure is needed \citep{Pollard2013}. This is the Fisher
counterpart of the data-processing inequality \citep{CoverThomas2006}, standard in substance.
Equality at one parameter is not recovery: \MT{rem:cg-fisher-equality-not-recovery} discards a
Bernoulli bit whose score vanishes at $\lambda=0$ but which depends on $\lambda$ at second order, and
no fixed kernel reconstructs the fine family near $\lambda=0$.

The one substantive hypothesis is parameter independence of the kernel, and the manuscript shows that
natural coarse maps violate it. A parameter-dependent channel adds a term to the record score
(\MTeq{eq:hist-parameter-dependent-channel-score}), and then the record can move while the fine law
stands still: a constant input fed to a channel emitting a bit with success probability
$(1+e^{-t})^{-1}$ has zero fine speed and record speed $(2\cosh(t/2))^{-1}$
(\MT{thm:hist-record-clock-contraction}). Gaussian Galerkin aggregation, substituting $y=Sc$ into an
energy with precision $\Lambda$, is the pushforward by $P_\Lambda=(S^\top\Lambda S)^{-1}S^\top
\Lambda$ (\SG{prop:ig-galerkin-pushforward}). On a fiber of fixed precision this is one kernel, the
theorem applies, and the loss vanishes exactly on directions in $\operatorname{range}(S)$. Across
precisions one kernel exists exactly when $\operatorname{range}(\Lambda S)$ does not change, and the
supplement exhibits a change of one self term inside a cluster that changes it, so along histories
with such moves the theorem does not apply. No parameter-independent kernel sends two normal laws to their joint mean-and-variance
barycenter (\SG{prop:hist-normal-barycenter-kernel}), and fitted approximations or
generative--recognition mismatches qualify only when a kernel is exhibited
(\MT{hyp:pb-markov-morphism}). A fixed stochastic blocking rule can only lose distinguishability; a
state-adapted rule is outside the theorem, and outside it a record can register motion the fine law
does not have.

On fibers the statement becomes the vertical defect $\Delta^{\mathscr K}_F=g^F-(T\mathscr K)^*\bar g^F
\succeq0$ with $\Delta^{\mathscr K}_F(Z,Z)=\E\operatorname{Var}(\mathbf s_Z\given\bar Y)$
(\MT{thm:pb-pullback-fisher-defect}), and defects compose as a cocycle,
$\Delta_F^{\mathscr K_{12}\circ\mathscr K_{01}}=\Delta_F^{\mathscr K_{01}}+(T\mathscr K_{01})^*
\Delta_F^{\mathscr K_{12}}$: losses along a chain add, each pulled back to the finest level
(\MT{thm:pb-fisher-defect-cocycle}). The cubic defect has no sign, so coarse-graining gives no
monotonicity for $T^F$.

\subsection{Configurations, lifts, and why marginals do not suffice}
\label{comp:ch07-formulas-lift}

An agent's configuration $Q_i=(q_i^b,q_i^m)$ has squared speed $w_bg^F_b(\dot q_i^b,\dot q_i^b)+
w_mg^F_m(\dot q_i^m,\dot q_i^m)$ (\MT{def:hist-configuration-manifold}); on the population the metric
is either the block sum, which treats agents as independent, or the pullback of a joint Fisher metric
along an exact lift (\MT{hyp:hist-exact-vfe-lift}). The lift is a smooth map from agent $i$'s
configurations to conditional recognition kernels of a block given the rest, with an extraction map
that recovers the configuration; it carries every correlation the configuration does not determine.
With the outside marginal fixed, the lifted objective differs from the collective free energy by a
constant (\MT{thm:obs-local-global-potential}), so both have one natural-gradient ray for one metric.
The substance is that the lift must be supplied and that different lifts give different geometry.
\MT{prop:hist-general-score-residual} quantifies this. With joint score $\mathbf s$ and embedded
marginal scores $\mathbf s_i=\E[\mathbf s\given Y_i]$,
\[
\mathsf G_{\mathrm J}-\sum_iw_i\mathsf G_i=\mathfrak K-\mathfrak C+\sum_i(1-w_i)\mathsf G_i,
\qquad
\mathfrak K=\E[\mathbf r\mathbf r^\top],
\quad
\mathbf r=\mathbf s-\sum_i\mathbf s_i,
\quad
\mathfrak C=\sum_{i\neq j}\E[\mathbf s_i\mathbf s_j^\top].
\]
The residual score $\mathbf r$ is the part of the joint score that is not a sum of single-block
scores, so $\mathfrak K\succeq0$ measures genuinely collective information; $\mathfrak C$ collects
overlaps between blocks' marginal scores and has no sign; the last term is weight mismatch. For
independent blocks $\mathfrak K=\mathfrak C=0$ and unit weights are forced. For dependent blocks no
inequality holds either way (a correlated bivariate normal gives a difference of both signs), and one
extraction can have right inverses with different pullback metrics ($\lambda_1\mapsto(\lambda_1,0)$
and $\lambda_1\mapsto(\lambda_1,\lambda_1)$ induce $d\lambda_1^2$ and $2d\lambda_1^2$ on
$\mathcal N((\lambda_1,\lambda_2),\mathrm I_2)$).

The parity witness isolates a pure higher-order dependence. Six bits, one per agent
$i\in\{1,2,3\}$ and channel $x\in\{b,m\}$, with success probabilities $\lambda_{i,x}$, get the lift
\[
p_\varkappa(y;\lambda)
=\prod_{(i,x)}\lambda_{i,x}^{y_{i,x}}(1-\lambda_{i,x})^{1-y_{i,x}}
+\varkappa(-1)^{\sum y_{i,x}}\prod_{(i,x)}\lambda_{i,x}(1-\lambda_{i,x}),
\qquad|\varkappa|<1,
\]
a product law plus a parity term whose amplitude factor keeps the law positive
(\MT{def:hist-three-agent-parity-lift}). Summing out any bit kills the parity term, so every proper
marginal is the product marginal (\MT{thm:hist-three-agent-parity-lift}). The lift is a smooth right
inverse of the one-bit marginal map, equivariant under complementing both bits of an agent and under
agent permutations, while a single-bit complement flips the sign of $\varkappa$. At the center, with
$\bar\varkappa=\varkappa/64$, \MT{cor:hist-parity-fisher-residual} gives
\[
\mathsf G_{\mathrm J}=\frac{4}{1-\bar\varkappa^2}\mathrm I_6,
\qquad
\mathsf G_{\mathrm J}-\sum_{(i,x)}\mathsf G_{i,x}=\frac{4\bar\varkappa^2}{1-\bar\varkappa^2}\mathrm I_6 .
\]
Pairs of bits are independent, so $\mathfrak C=0$ and the excess is all $\mathfrak K$; weighted
marginal metrics match only with the $\varkappa$-dependent weight $1/(1-\bar\varkappa^2)$, and since
the one-bit families do not depend on $\varkappa$, no rule that reads them returns every joint
metric. For the free energy an engineered hyperedge record
$K_{\mathrm{hyp}}(1\given y)=\tfrac12(1+\bar\varkappa(-1)^{\sum y})$ with uniform prior has the center
law as exact posterior (\MT{def:hist-parity-hyperedge}); the posterior attains $\Fenergy_1=\log2$,
while the uniform law, the product of the same one-bit marginals, has
$\Fenergy_1(\operatorname{Unif})-\log2=-\tfrac12\log(1-\bar\varkappa^2)$
(\MT{prop:hist-parity-hyperedge-vfe}). The two laws share all one-bit data and the evidence and differ
in free energy. With $\varkappa$ promoted to a coordinate the joint rank is seven and the marginal
rank six, with $\partial_\varkappa$ invisible (\MT{cor:hist-promoted-parity-rank}); combined with the
finite-family clause of \MT{thm:hist-retained-fisher-quotient}, a history moving only $\varkappa$ has
positive joint Fisher speed while every single-agent clock stands still. The numbers are small by
construction ($|\bar\varkappa|<1/64$ gives a Fisher excess below $10^{-3}$ and a free-energy excess of
at most about $1.22\times10^{-4}$ nats), so the witness is qualitative. Its message for meta-agents is
that a law assembled from constituents' single-agent laws cannot see such a dependence, while the
collective free energy and joint Fisher metric can.

Could a canonical rule supply the dependence? A selector maps one-coordinate marginals to a joint law
with those marginals, and is natural on a class of kernels if it commutes with them
(\MT{def:hist-finite-marginal-selection}). Naturality under all tensor products of single-coordinate
kernels forces the product law, because kernels from a one-point set prepare each coordinate
independently (\MT{thm:hist-local-markov-product-uniqueness}); a larger class admits a natural
selector exactly when every added kernel preserves products. Admitting two kernels that split a fair
bit into correlated bits, with $\mathrm c=\tfrac13$ and $\mathrm c=\tfrac12$, leaves no natural
selector, since both give fair marginals with different joints
(\MT{thm:hist-correlated-refinement-nogo}). Both results are elementary; their force lies in the
admitted class. What replaces an absolute selector appears only in an unnumbered paragraph: relative
to a positive reference law, a relative-entropy minimizer over a convex moment-constrained family is
unique when it exists, and the reference conditional completion of a coarse marginal follows from the
chain rule (\MT{lem:prob-kl-chain-rule}); both take the reference as extra data. This is maximum
entropy relative to a prior measure \citep{jaynes2003probability}, standard. In physics terms, without
a Hamiltonian the only closure that commutes with independent local relabelings is mean field, and a
correlated closure needs a reference measure. For meta-agents the dependence among constituents is
supplied by the generative model or a declared lift, never manufactured from configurations, in line
with the standing refusal of a canonical coupling selector (\MT{sec:intro-conventions}).

\subsection{Natural-gradient histories: the orbit depends on the mobility}
\label{comp:ch07-formulas-orbits}

\[
\mathsf G\bigl(\operatorname{grad}^{\mathsf G}\Fenergy,Z\bigr)=d\Fenergy[Z],
\qquad
\dot Q(t)=-\mathsf a(t)\operatorname{grad}^{\mathsf G}\Fenergy\bigl(Q(t)\bigr),
\quad\mathsf a(t)>0 .
\]
The gradient raises the index of $d\Fenergy$ with the metric, $\mathsf G^{-1}\nabla\Fenergy$ in
coordinates, and the arbitrary speed $\mathsf a$ makes the pair an oriented orbit rather than a
trajectory (\MT{def:hist-oriented-history}). A scalar mobility only reparameterizes orbits; an
anisotropic one can bend them: for $\Fenergy=\tfrac12(z_1^2+z_2^2)$ ordinary descent runs along
straight lines, descent with mobility $\operatorname{diag}(1,2)$ along the parabolas
$z_2=z_2^0(z_1/z_1^0)^2$ (\MT{prop:hist-scalar-mobility-orbit}), and descent orbits are not
geodesics. For the program this means that the critical points are metric independent but which one
an orbit from a given start reaches is a property of the orbit, and nothing in the chapter shows it is
mobility independent. Every dynamics the program runs carries a chosen mobility: the four-cycle
kinetics of \cref{comp:ch08,comp:ch13} use identity mobility with supplied rates, which is the Fisher
natural gradient for the unit-variance belief means but not for the model means, whose Fisher rate
would be $\tfrac12$ (\MT{prop:obs-slow-model-link-dichotomy}).

The remaining statements keep the objects well defined. Noncritical orbits are embedded, with the
free energy a strictly decreasing coordinate (\MT{prop:hist-vfe-orbit-embedded}). A semidefinite
metric cannot be inverted silently (\MT{prop:hist-semidefinite-gradient-obstruction}); a quotient must
come first (\MT{thm:hist-retained-fisher-quotient}). For softmax logits the pulled-back Fisher tensor
$\operatorname{diag}(p)-pp^\top$ has null direction $\mathbf 1$, and the Moore--Penrose solution picks
the gauge $\sum_jz_j=0$. For a free energy invariant under a group action the gradient is orthogonal
to the orbits, so histories never move along gauge directions, and for a free, proper, isometric
action they descend to the quotient with the same duration
(\MT{prop:hist-invariant-history-horizontal}). Explicit Euler steps below
$2/\mathrm c_{\mathsf G}$ descend and stay within $O(\Delta t)$ of the orbit on compact intervals
without lying on it (\MT{rem:hist-discrete-histories}).

\subsection{Fisher duration, and why it is not called time}
\label{comp:ch07-formulas-clock}

\begin{align*}
\dot{\mathfrak t}&=\sqrt{\mathsf G(\dot Q,\dot Q)},
&
\mathfrak t_{Q,t_0}(t)&=\int_{t_0}^{t}\dot{\mathfrak t}(t')dt',
\\
\frac{d\mathfrak t}{dt}&=\mathsf a\lVert\operatorname{grad}^{\mathsf G}\Fenergy\rVert_{\mathsf G},
&
\frac{d\Fenergy}{d\mathfrak t}&=-\lVert\operatorname{grad}^{\mathsf G}\Fenergy\rVert_{\mathsf G}.
\end{align*}
The clock is arc length from a chosen origin (\MT{def:hist-fisher-clock}). Along a natural-gradient
history the clock rate is the speed factor times the gradient norm, and the free-energy decrease per
unit clock is the gradient norm alone, a property of the orbit. The clock is invariant under
orientation-preserving reparameterization and stops on critical or null segments
(\MT{thm:hist-fisher-clock-invariance}). For a block-diagonal metric the collective clock satisfies
$\max_i\mathfrak t_i\leq(\sum_i\mathfrak t_i^2)^{1/2}\leq\mathfrak t\leq\sum_i\mathfrak t_i$. In the
unit-variance location family the path $\mu(t)=e^{-t}$ has duration $1$ over infinite parameter time,
and a realized update is at least as long as the Fisher distance between its endpoints
(\MT{prop:hist-realized-length-versus-distance}). Total duration is finite for a gradient history with
an accumulation point at which the objective is real analytic, by the {\L}ojasiewicz argument of
\citet{absil2005convergence}, and infinite for a smooth example that spirals toward a circle
(\MT{prop:hist-analytic-finite-duration}).

A clock with level sets orthogonal to the flow is stronger. With $\mathsf n=\lVert
\operatorname{grad}^{\mathsf G}\Fenergy\rVert_{\mathsf G}$ and $\bm\alpha=-d\Fenergy/\mathsf n$, an
orthogonal clock solves $d\mathfrak t^\perp=\bm\alpha$; since $d\bm\alpha=\mathsf n^{-2}d\mathsf n
\wedge d\Fenergy$, it exists exactly when $\mathsf n$ is locally a function of $\Fenergy$ and the
periods of $\bm\alpha$ vanish (\MT{thm:hist-global-clock-exactness}). For $\Fenergy=z_1z_2$ on the
first quadrant none exists on any open set, although a unit-speed clock built from the cross-section
$z_1z_2=1$ exists on the whole quadrant (\MT{prop:hist-cross-section-clock}).

\begin{analogy}{structural}
With the level sets of $\Fenergy$ labeled by $\tau=-\Fenergy$, $\bm\alpha=\mathsf n^{-1}d\tau$: the
arc length between neighboring slices is a lapse $\mathsf n^{-1}$ times $d\tau$. A time function whose
differential is the unit normal one-form exists exactly when the lapse is constant on each slice, the
same Frobenius-type condition as for a synchronous slicing in general relativity, here in Riemannian
signature and with no causal content.
\end{analogy}

The manuscript refuses to call the duration physical time, and the refusal is correct. Its standing
scope (\MT{sec:intro-conventions}) says an identification would need a named target clock, a
calibration map with an uncertainty model, causal and compositional consistency, and discriminating
evidence; none is supplied. The chapter's own results add reasons: the duration depends on the chosen
metric and origin, stops at critical points, is not additive over agents, generally has no orthogonal
(simultaneity-like) version, and runs slower for a meta-agent than for its constituents by an unknown
ratio. \MT{prop:hist-coordinate-independence} separates coarse-graining depth $\ell$ (a discrete
index), the orbit parameter $t$ (a choice), and the duration $\mathfrak t^{(\ell)}$ (arc length for a
chosen metric and origin), and notes that even a strictly increasing duration need not be a regular
coordinate ($z=t^2$ gives $\mathfrak t=t^2$). ``Clock'' here is a name for arc length.

\subsection{Record clocks, semiconjugacy, and depth}
\label{comp:ch07-formulas-semiconjugacy}

Applied at every instant of a history $P_t$ observed through one fixed kernel, the Fisher-loss theorem
gives, for the record history $\bar P_t=P_tK$,
\[
\dot{\mathfrak t}_P^2-\dot{\mathfrak t}_{\bar P}^2
=\E_t\operatorname{Var}_t\bigl(\mathbf s_t(Y)\given\bar Y\bigr)\geq0,
\qquad
\operatorname{len}(\bar P)\leq\operatorname{len}(P)
\]
(\MT{thm:hist-record-clock-contraction}), with equality on an interval of finite length exactly when
the fine score is a function of the record at almost every instant. This covers a meta-agent exactly
when its law is the image of the same fine history under one fixed kernel, which natural Gaussian
meta-agent maps can fail (\SG{prop:hist-normal-barycenter-kernel}). A coarse configuration map acting
channelwise through Markov morphisms with equal weights gives the configuration-level version
(\MT{thm:hist-pointwise-contraction-lift}).

A meta-agent that runs its own natural-gradient dynamics on its own free energy has a recomputed
orbit, not an image. The chapter compares the two through
\[
T_Q\mathsf F_\ell\mathscr V_\ell(Q)=\mathsf a_\ell(Q)\mathscr V_{\ell+1}(\mathsf F_\ell Q),
\qquad\mathsf a_\ell>0
\]
(\MT{def:hist-oriented-semiconjugacy}): the pushed fine velocity is a positive multiple of the coarse
velocity. Then fine orbits map onto coarse orbits with time change $\psi_Q(t)=\int_0^t\mathsf a_\ell$
(\MT{thm:hist-oriented-semiconjugacy}), and the coarse duration of the image is the intrinsic length
of the coarse orbit arc, independent of $\mathsf a_\ell$ (\MT{thm:hist-duration-relation}). A constant
map satisfies the relation trivially, so a nonconstant shared history is a separate requirement
(\MT{prop:hist-noncollapse}); with $\mathsf F_\ell=\arctan$ on $\R$ both flows are complete and only
part of the coarse orbit is reached; and the coarse duration is no larger than the fine one exactly
when $\lVert T\mathsf F_\ell\mathscr V_\ell\rVert\leq\lVert\mathscr V_\ell\rVert$ along the orbit
(\MT{thm:hist-duration-criterion}), which fiberwise Fisher contraction does not give by itself,
because configuration weights are separate data and a recomputed orbit is not an image.

\MT{prop:hist-natural-gradient-sufficiency} says where the relation holds. Equal objectives do not
suffice, since a change of metric breaks it. It holds for surjective submersions that are
functionally compatible, $\Fenergy_\ell=\chi_\ell\circ\Fenergy_{\ell+1}\circ\mathsf F_\ell$ with
$\chi_\ell'>0$, and horizontally conformal with dilation $\mathsf d_\ell$, with
$\mathsf a_\ell=\chi_\ell'\mathsf d_\ell^2$; quotients by free, proper, isometric symmetries of the
objective have $\mathsf a_\ell\equiv1$. Compatibility forces the fine objective to be constant along
the fibers of the coarse map, so the route misses maps that discard directions on which the fine
objective depends: when the coarse objective is the fiberwise infimum with unique fiberwise
minimizers, compatibility with $\chi_\ell=\operatorname{id}$ holds on no open set, and in the
manuscript's example it holds for no $\chi_\ell$. The chapter declines to assert a semiconjugacy for such maps and routes the comparison through
the reduced slow flow of \cref{comp:ch08}, a conjugacy on the graph of conditional minimizers
(\MT{prop:obs-rate-separated-flow}); whether the actual flow follows that graph is
\MT{open:obs-slow-manifold}.

\section{What is proved, and how solid it is}
\label{comp:ch07-status}

\begin{scorecard}
\MT{prop:pb-statistical-tensor-descent}, \MT{cor:pb-orbit-fisher-form} & Fisher and Amari--Chentsov tensors are frame independent at $c_*$; a small group move costs $\tfrac12t^2\mathfrak k_q(\xi,\xi)$, null on the isotropy algebra. & \SOLID & Standard in substance; needs an action by sample maps.\\
\MT{def:hist-configuration-manifold} & The configuration metric, channel weights, and the Fisher metric itself are chosen. & \CHOICE & Uniqueness theorems exist and are not imposed (\MT{prop:ig-fisher-characterizations}).\\
\MT{thm:cg-fisher-contraction} & Coarse score is the conditional expectation of the fine score; Fisher loss is the expected conditional covariance. & \SOLID & Parameter-independent kernel; equality is local sufficiency, not recovery; defects form a cocycle (\MT{thm:pb-fisher-defect-cocycle}).\\
\MT{thm:hist-retained-fisher-quotient} & Null space of a pulled-back Fisher tensor; positive quotient bundle on constant-rank strata. & \SOLID & Local; a global quotient needs a Hausdorff leaf space (open).\\
\MT{hyp:hist-exact-vfe-lift} & Dependence between agents enters through a supplied lift. & \CHOICE & Different lifts give different metrics (\MT{prop:hist-general-score-residual}, SOLID).\\
\MT{thm:hist-local-markov-product-uniqueness}, \MT{thm:hist-correlated-refinement-nogo} & Local naturality forces the product law; two correlated split kernels leave no natural selector. & \SOLID & Negative, elementary, relative to the admitted class; the reference-relative replacement is an unnumbered remark.\\
\MT{thm:hist-three-agent-parity-lift} & A positive smooth symmetric correlated right inverse of the one-bit marginal map exists. & \TOY & Three agents, binary channels, finite symmetry group, $|\bar\varkappa|<1/64$.\\
\MT{prop:hist-parity-hyperedge-vfe}, \MT{cor:hist-parity-fisher-residual} & Free energy and joint Fisher metric are not functions of single-agent marginals (and evidence). & \SOLID & A refutation by a toy counterexample; effect second order in $\bar\varkappa$.\\
\MT{prop:hist-scalar-mobility-orbit} & Scalar mobility keeps orbits; anisotropic mobility, including a rate ratio, can change them. & \SOLID & Orbit selection depends on the chosen metric.\\
\MT{thm:hist-fisher-clock-invariance}, \MT{thm:hist-global-clock-exactness} & Duration is reparameterization invariant; an orthogonal clock exists iff the gradient norm is locally a function of the free energy. & \SOLID & Standard in substance; metric and origin chosen; orthogonal clocks fail generically.\\
\MT{thm:hist-record-clock-contraction} & A same-path image under one fixed kernel runs no faster, by the exact conditional-variance defect. & \SOLID & Covers a meta-agent only if its law is such an image; fails for parameter-dependent channels.\\
\MT{thm:hist-oriented-semiconjugacy}, \MT{thm:hist-duration-relation} & If the coarse map intertwines descent fields up to a positive rate, orbits correspond and the rate cancels from durations. & \COND & Hypothesis: the semiconjugacy, not shown for any information-discarding meta-agent map.\\
\MT{prop:hist-natural-gradient-sufficiency} & Semiconjugacy holds for compatible conformal submersions and symmetry quotients; for a fiberwise-infimum coarse objective, compatibility holds on no open set. & \SOLID & The positive part excludes the maps formation needs.\\
\end{scorecard}

\begin{hygiene}
Conditions that can fail in the program's cases: parameter independence of the coarse kernel along
the whole history (fails for Galerkin aggregation when $\operatorname{range}(\Lambda S)$ moves, for
joint barycenters, for parameter-dependent channels); the chosen metric, weights, and rate ratio,
which orbits and durations inherit; the supplied lift; the semiconjugacy, a hypothesis for every
information-discarding map; a Hausdorff leaf space for global quotients; an accumulation point for
finite total duration; and regularity (\MT{hyp:pb-regular-models}), which excludes singular strata,
so that coarse families of mixture type, singular where components coincide, need checking.

Hygiene a physicist can take for granted: standard Borel spaces, differentiability in quadratic mean,
integrability of scores to third order, domination for differentiating divergences, $C^1$ metrics,
$C^2$ objectives, local Lipschitz coarse fields, and real analyticity for finite duration, which
Gaussian relative entropies and log-sum-exp attention terms satisfy in natural coordinates.
\end{hygiene}

\section{Toy versus general}
\label{comp:ch07-toy}

Most small examples in the chapter are planar refutations with general logical force; they show
which implications fail, not what large networks do. The parity witness has a general negative part, that no rule reading single-agent
marginals returns the free energy or joint Fisher metric for every law, and a TOY positive part, the
existence of a smooth symmetric correlated lift for three agents with binary channels. Its
mechanism-level content, conjectural at large $N$, is that dependences of higher order than any
retained marginal enter the collective free energy and joint Fisher metric through $\mathfrak K$ while
staying invisible to marginal-based coarse descriptions and single-agent clocks; how large
$\mathfrak K$ is at $N\geq512$ in the program's generative models, and whether such dependences are
generic there, is not addressed. The constants $\tfrac{65536}{16383}$ and $\tfrac{4}{16383}$, the bound
$|\bar\varkappa|<1/64$, the ranks seven and six, and the $57$-dimensional fiber are small-$N$ specifics.
The witness's hypotheses (finite alphabets, a discrete symmetry group) exclude continuous law fibers
with connected groups such as $\GL(K,\R)$, $\GL(K,\mathbb C)$, or $\mathrm{SU}(N)$, and whether
correlated lifts equivariant under a connected group exist there is open
(\MT{open:hist-parity-lift-scope}). Everything else (tensor descent, Fisher loss, record clocks, clock
invariance, the orthogonal-clock criterion, the semiconjugacy calculus) is general and independent of
$N$; its hypotheses admit nonabelian and noncompact groups acting by sample maps, such as linear
$\GL(K)$ on Gaussian samples.

\section{Open ends}
\label{comp:ch07-open}

\MT{open:hist-parity-lift-scope} asks, for a stated regular family and a stated class of morphisms
containing non-product-preserving kernels, for a smooth right inverse of the marginal map natural under
that class, and for correlated lifts equivariant under a connected structure group on continuous
fibers; a construction or no-go theorem settles each. \MT{open:hist-retained-fisher-quotient} asks for a
checkable criterion under which the null foliation of a pulled-back Fisher tensor has a Hausdorff leaf
space, and for block data on the quotient fixed by a typing of agents and channels, since the law and
its Fisher tensor alone select no agent splitting (\MT{prop:hist-law-fisher-agentization-nogo}).
\MT{open:hist-meta-clock-ratio} asks for the ratio of an $N$-constituent meta-agent's duration to its
constituents' collective duration, as a function of $N$ and of the correlation of constituent
velocities, in the coherent and incoherent regimes. The physicist's expectation, that coherent motion
survives coarse-graining while incoherent motion is suppressed with $N$, is exactly what is unproved,
and the problem warns that a divergence computation does not settle it because expected arc-length
ratios are not square roots of expected divergence ratios. Finally, for information-discarding maps
the comparison depends on \MT{open:obs-slow-manifold} (\cref{comp:ch08-open}); a semiconjugacy for
such maps is declined without an open problem of its own.

\begin{verdictbox}{\NARROWER}
\textbf{Claimed.} The opening says that a chosen metric turns the collective free energy into an
oriented orbit with a reparameterization-invariant information clock; that a parameter-independent
Markov channel loses Fisher information by an exact conditional-covariance defect, so that a
same-path Markov image such as the law of a meta-agent accumulates no more duration; that a recomputed
coarse orbit compares with the fine one only through an oriented semiconjugacy whose rate cancels from
durations; and that a parity witness shows the free energy and its Fisher tensor are not functions of
marginal laws.

\textbf{Delivered.} Each statement is proved at its stated generality. Much is standard mathematics in
new typing (Fisher data processing, arc length, gradient flows under a mobility, {\L}ojasiewicz),
with a counterexample at each boundary. The parity witness and the selection theorems deliver their
negative claims exactly.

\textbf{Gap.} ``Such as the law of a meta-agent'' outruns the theorems: the record clock covers a
meta-agent only when its law is a parameter-independent Markov image of the fine history, which
Galerkin aggregation across precisions and joint barycenters are not. For meta-agents that run their
own dynamics on information-discarding maps, which formation needs, the chapter's sufficient route to
a semiconjugacy does not apply (for a fiberwise-infimum coarse objective it holds on no open set),
none is asserted, and the comparison passes to a reduced flow whose tracking is open. The meta-agent clock ratio is open. Every orbit and duration rests on a chosen
metric and, for collectives, a supplied lift; neither the orbit nor the clock is a property of the
free energy alone.
\end{verdictbox}

\chapter{Two-Timescale Dynamics}
\label{comp:ch08}

For latent links, generator recognition means can be assigned to the slow block, but that
assignment needs a rate declaration. On the translation-instance invariant leaf the unmodified
Fisher flow gives belief means, model means, and link means a common rate
(\MT{rem:link-rates}). The dynamical-link chapter therefore does not by itself
establish the separation or tracking assumptions explained below (\cref{comp:lie-links}).

\section{The question and the answer}
\label{comp:ch08-question}

\MT{ch:two-timescales} asks what the collective free energy supplies for separating fast inference of
beliefs from slower evolution of models, priors, and parents. A physicist expects adiabatic
elimination: a small rate ratio $\varepsilon$, a reduced slow dynamics with the fast variables at
their conditional optimum, and a theorem that the full dynamics follows it as $\varepsilon\to0$. The
chapter proves the first two in general. The third holds only on held matching sheets of a four-agent
toy, where the chapter verifies the hypotheses of Tikhonov's theorem; elsewhere it is open.

The results are these. Any positive-definite block mobility gives descent, and in a product
recognition family an agent descending its own local free energy moves along a block of collective
descent (\MT{prop:obs-local-block-descent}). The envelope identity gives the reduced slow flow as the
natural gradient of the free energy minimized over the fast block (\MT{lem:obs-envelope-slow-flow}).
The rate-separated flow agrees with the reduced flow on the graph of conditional minimizers, where
projection to the slow coordinates is a conjugacy (\MT{prop:obs-rate-separated-flow}); this is the
route \cref{comp:ch07} uses for information-discarding maps. On the four-cycle the model channel lags
the belief channel on an interval of order $1/\varepsilon_{\mathrm m}$ exactly when the links are slow
too (\MT{prop:obs-slow-model-link-dichotomy}). In a finite discrete setting, parent-sensitive sweeps
descend (\MT{thm:rg-memory-inference-descent}), the parent is Markov under instantaneous
within-parent refresh (\MT{thm:rg-memory-prepared-parent}), and it has memory under partial refresh
(\MT{rem:rg-memory-partial-refresh}).

The free energy chooses none of the split, the rate ratio, or the mobilities, and a rate ratio can
change the orbit (\MT{prop:hist-scalar-mobility-orbit}). The scope paragraph declines tracking beyond
the held sheets, tracking by sampled dynamics with source jumps, persistence of a belief-coherent,
model-incoherent pair under such sampling at the formation temperature, convergence of the alternating
scheme, a uniform rate for the parent sweep, and closed dynamics for a static parent datum. The chapter
uses \cref{comp:ch05,comp:ch06,comp:ch07} and feeds \cref{comp:ch11,comp:ch13}. The finite-depth
formation family realizes one scalar slow parameter, with dynamics
\MTeq{eq:rg-formation-own-law-dynamics}, not slow per-agent model laws, so the belief--model split is
instantiated only on the four-cycle, with supplied rates.

\section{The picture}
\label{comp:ch08-picture}

The picture is Born--Oppenheimer: fast electrons relax at fixed nuclear positions, and the nuclei move
on the resulting surface. Beliefs play the electrons; models and generative parameters play the
nuclei.

\begin{analogy}{structural}
$F_{\mathrm{red}}(s)=\min_fF(f,s)$ plays the Born--Oppenheimer surface \citep{Born1927}, and the
envelope identity $dF_{\mathrm{red}}=\partial_sF(f^\star(s),s)$ is the Hellmann--Feynman statement for
a variational ground state: the slow force is the explicit partial derivative at the conditional
optimum, because the implicit variation of the fast variables drops out at a stationary point. What
does not transfer: the adiabatic theorem that justifies Born--Oppenheimer and its nonadiabatic
corrections; the chapter proves only the zeroth-order reduced dynamics. $F_{\mathrm{red}}$ is a
minimum over a recognition family, not a partition function; this is zero-temperature elimination,
not integrating out.
\end{analogy}

\begin{analogy}{structural}
The rate-separated flow is a singularly perturbed gradient system, and Tikhonov's theorem
\citep{Tikhonov1952} is the standard justification of adiabatic elimination: if the frozen fast
subsystem has a uniformly stable equilibrium and the fast start lies in its basin, the slow variables
converge to the reduced flow as $\varepsilon\to0$; Fenichel's theory \citep{fenichel1979geometric}
adds a nearby invariant slow manifold. The chapter states these hypotheses in
\MT{open:obs-slow-manifold} and verifies them only on held four-cycle sheets; sampled dynamics with
source jumps lie outside such a theorem.
\end{analogy}

Two features have no textbook counterpart. When the kinetic metric couples beliefs and models, as
non-product recognition makes it do, ``make the belief block fast'' is ambiguous: slowing the model
rows of the velocity and rescaling the metric by a congruence give different slow sets, and only the
second keeps a gradient structure (\MT{rem:obs-rate-separation-witnesses}). On the four-cycle the
picture is lattice gauge theory on a ring: each residual is a neighbor difference minus a link, and
both channels are compared through the same link, so a fast link cancels a slow channel's mismatch
and the fast beliefs follow it. Model disagreement measured through a link is link relative and
survives only if the link is slow too.

\begin{analogy}{structural}
The refresh $\mathsf D=\mathsf C\mathsf J$ is the discrete Nakajima--Zwanzig projector onto
conditional equilibrium \citep{Nakajima1958,Zwanzig1960}. Projected dynamics has a memory term unless
the unresolved part is reset at every step; a partial reset leaves memory, and the exact discrete
recurrence is \MT{thm:rg-projection-memory} (\cref{comp:ch11}). The prepared-parent theorem is a
lumpability statement \citep{KemenySnell1976} in intertwining form, valid from conditional-equilibrium
starts.
\end{analogy}

\section{Objects and notation}
\label{comp:ch08-objects}

\begin{decoder}
Fast block $f$; slow block $s=(s_{\mathrm m},\theta)$ & Declared fast coordinates (beliefs, source rows) and slow ones (model recognition coordinates, generative parameters) & Electrons and nuclei\\
$f^\star(s)$, $\mathcal S^\star$ & Conditional minimizer over the fast block; its graph & Slaved variables; adiabatic surface\\
$F_{\mathrm{red}}(s)$ & Free energy minimized over the fast block & Born--Oppenheimer surface\\
$\mathsf G_{\mathrm f}$, $\mathsf G_{\mathrm s}$, $\varepsilon$, $\mathsf G_\varepsilon$ & Block metrics; rate ratio; $\varepsilon\mathsf G_{\mathrm f}\oplus\mathsf G_{\mathrm s}$ & Mobilities; singular-perturbation parameter\\
Rate-separated, reduced flow & Natural gradient for $\mathsf G_\varepsilon$; natural gradient of $F_{\mathrm{red}}$ & Singularly perturbed system; adiabatic dynamics\\
Row-rated field, congruence rescaling & Two ways to make a block fast when the metric couples blocks & Choice of coordinates for elimination\\
Accepted M-coordinate & Parameter step that does not lower the evidence bound (\MT{def:elbo-generalized-m-acceptance}) & Generalized M step\\
Four-cycle: $\bar k$, $\bar m$, $\check\omega$; $r^b_i$, $r^m_i$ & Belief means, model means, link translations on a ring; neighbor difference minus link & Matter and link fields; covariant difference\\
$S$, $S^{(0)}$, $\mathrm T$ & Attention assignment; one held perfect matching; temperature & Quenched bonds; temperature\\
$\varepsilon_{\mathrm m}$, $\varepsilon_{\check\omega}$ & Supplied rates of model means and links & Mobilities\\
$\xi$, $z=\operatorname{par}(\xi)$ & Fine state and its parent (complete decorated block label) & Microstate, macrostate\\
$\mathsf M$, $\mathsf C$, $\mathsf J$, $\mathsf D=\mathsf C\mathsf J$, $\mathsf T=\mathsf M\mathsf D$ & Metropolis move; projection to parent; posterior lift given parent; refresh; one sweep & Update; lumping; disaggregation; projector\\
$\overline{\mathsf T}=\mathsf J\mathsf M\mathsf C$; prepared start $\nu\mathsf J$ & Parent transition kernel; fine law in conditional equilibrium within parents & Lumped chain; local equilibrium\\
$\mathsf D_\omega$; $D^{(1)}_N$ & Refresh with probability $\omega<1$; cost of a first-order Markov parent & Finite separation; memory\\
\end{decoder}

The chapter's ``slow manifold'' is Tikhonov's and Fenichel's attracting set; the graph
$\mathcal S^\star$ is different and at finite $\varepsilon$ generally not invariant, so the chapter's
exactness statements hold on it, not near it. ``Rate'' is a chosen mobility factor, ``declared'' means
chosen, and ``exact route'' means the conjugacy on $\mathcal S^\star$, not a tracking theorem. Kernels
act on row vectors from the right, so $\mathsf T=\mathsf M\mathsf D$ means move, then refresh. The
sweep $\mathsf T$ and temperature $\mathrm T$ are different glyphs; $\omega$ is the refresh fraction,
and the four-cycle links, written $\omega$ in \cref{comp:ch13}, are $\check\omega$ here;
$\operatorname{par}$ is the parent map (the parity character in \cref{comp:ch07}); the rate ratio
$\varepsilon$ is unrelated to the formation regime parameter.

\section{The formulas, explained}
\label{comp:ch08-formulas}

\subsection{Descent with any mobility, and when it is local}
\label{comp:ch08-formulas-descent}

\[
\frac{d}{dt}F=-\sum_{i\in V'}(\nabla_{\lambda_i}F)^\top\mathsf G_i^{-1}(\nabla_{\lambda_i}F)\leq0 .
\]
Each moving block $i\in V'$ descends the collective free energy $F(\lambda)$ with its own
positive-definite metric, rates absorbed (\MT{prop:obs-local-block-descent}); the chain rule gives the
sign. The content is locality. In a product family $Q_\lambda=\bigotimes_iQ_{i,\lambda_i}$ the partial
gradient $\nabla_{\lambda_i}F$ equals the gradient of agent $i$'s outside-averaged local free energy,
so selfish local descent is collective descent, an exact potential game
(\MT{thm:obs-local-global-potential}), and the block Fisher metrics make it the Fisher natural
gradient. Outside product families moving $\lambda_i$ changes the others' marginals, the lift of
\cref{comp:ch07} is needed, and the Fisher metric has off-diagonal blocks (for a bivariate normal
location family with correlation $\mathrm c_2$, the entry $-\mathrm c_2/(1-\mathrm c_2^2)$). ``Each
agent follows its own natural gradient'' is the collective natural gradient only when recognition
factorizes.

\subsection{Fast and slow blocks, the envelope identity, and the reduced flow}
\label{comp:ch08-formulas-envelope}

\begin{align*}
F(f,s)&=-\log p_{\theta,V}(o\given X)
+\KL\bigl(Q_{f,s_{\mathrm m}}\Vert\boldsymbol\Pi_{\theta,V,o,X}\bigr),
\qquad s=(s_{\mathrm m},\theta),
\\
F_{\mathrm{red}}(s)&=\min_fF(f,s),
\qquad
dF_{\mathrm{red}}(s)=\partial_sF\bigl(f^\star(s),s\bigr),
\qquad
\dot s=-\mathsf G_{\mathrm s}(s)^{-1}\partial_sF\bigl(f^\star(s),s\bigr).
\end{align*}
The evidence term depends only on the generative parameter $\theta$; the gap term depends on
everything. The fast block carries belief recognition coordinates and attention source rows, the slow
block model recognition coordinates and $\theta$, which moves the joint law between evaluations and
never inside one (\MT{req:gen-typing-prohibition}); the assignment to blocks is declared
(\MT{def:obs-fast-slow-blocks}). \MT{lem:obs-envelope-slow-flow} gives the second line: at the
conditional minimizer $\partial_fF=0$, so the implicit dependence of $f^\star$ drops out, and the
reduced flow is the $\mathsf G_{\mathrm s}$-natural gradient of $F_{\mathrm{red}}$. A $C^1$ branch
exists where the fast Hessian is positive definite, and on a compact fast domain with unique minimizer
the identity holds without differentiability of $f^\star$ \citep{Danskin1967}. In the limiting case
where the slow block is $\theta$ alone and the fast family contains the exact posterior, the gap
vanishes at the minimum, $F_{\mathrm{red}}(\theta)=-\log p_{\theta,V}(o\given X)$, and the reduced
flow is gradient ascent on the log evidence, a continuous-time exact EM.

\subsection{The rate-separated flow, and what is not proved about it}
\label{comp:ch08-formulas-rate}

\[
\varepsilon\dot f=-\mathsf G_{\mathrm f}^{-1}\partial_fF,
\qquad
\dot s=-\mathsf G_{\mathrm s}^{-1}\partial_sF .
\]
This is the natural gradient for $\mathsf G_\varepsilon=\varepsilon\mathsf G_{\mathrm f}\oplus
\mathsf G_{\mathrm s}$, so it descends for every $\varepsilon$, and its orbit can depend on
$\varepsilon$ (\MT{prop:obs-rate-separated-flow}). On $\mathcal S^\star$ the fast component vanishes
and the slow component equals the reduced flow; the reduced field lifted to $\mathcal S^\star$ is
tangent to it, and projection to the slow coordinates is a conjugacy with rate one, which is
\MTeq{eq:obs-reduced-semiconjugacy}; the $\mathsf G_\varepsilon$-length of the lifted orbit tends to
the slow length. The proposition does not say that a solution started off $\mathcal S^\star$ approaches
it and follows the reduced flow; that is Tikhonov's content and is not proved at this generality.

\MT{rem:obs-rate-separation-witnesses} shows why care is needed. Off $\mathcal S^\star$ the slow
velocity need not be parallel to the reduced gradient, and at finite $\varepsilon$ the graph is not
invariant. With a correlated metric the naive construction fails: for $F=[(f-s)^2+s^2]/2$, whose
conditional minimizer is $f=s$, and a metric with unit diagonal and off-diagonal
$\mathrm c_2\in(0,1)$, slowing the model row of the natural-gradient velocity gives the fast critical
set
\[
\Bigl\{f=\frac{1+2\mathrm c_2}{1+\mathrm c_2}s\Bigr\},
\]
attracting at rate $1/(1-\mathrm c_2)$ and not the conditional minimizers; at $\mathrm c_2=\tfrac12$
and $\varepsilon<7-4\sqrt3\approx0.072$ this field increases $F$ at an explicit point, so it is the
natural gradient of no metric. The congruence rescaling $\Delta_\varepsilon\mathsf G\Delta_\varepsilon$,
$\Delta_\varepsilon=\operatorname{diag}(\sqrt\varepsilon,1)$, descends for every $\varepsilon$ and
recovers the conditional minimizers at $\varepsilon=0$. With correlated metrics a two-timescale model
must state how the rate ratio enters the metric; this is a modeling requirement, since correlated
belief--model metrics are what non-product recognition produces.

\subsection{The four-cycle: slow models need slow links}
\label{comp:ch08-formulas-cycle}

On the ring of \MT{def:rg-formation-cycle-experiment}, agents have belief laws $\mathcal N(\bar k_i,1)$
and model laws $\mathcal N(\bar m_i,\tfrac12)$, edges $e_i$ from $i$ to $i+1$ carry link translations
$\check\omega_i$, and the residuals are $r^b_i=\bar k_{i+1}-\bar k_i-\check\omega_i$ and
$r^m_i=\bar m_{i+1}-\bar m_i-\check\omega_i$. Each receiver attends one neighbor; the energy $U_S$
collects attended edge costs and temperature terms; the matching $S^{(0)}$ selects $e_0$ and $e_2$.
\MT{prop:obs-slow-model-link-dichotomy} runs
\[
d\mathbf x=-\bigl(\partial_{\bar k}U_S,\varepsilon_{\mathrm m}\partial_{\bar m}U_S,
\varepsilon_{\check\omega}\partial_{\check\omega}U_S\bigr)dt
+\sqrt{2\mathrm T}\bigl(dW^{\bar k},\varepsilon_{\mathrm m}^{1/2}dW^{\bar m},
\varepsilon_{\check\omega}^{1/2}dW^{\check\omega}\bigr),
\]
with the assignment redrawn at a rate-one clock. The rates are supplied; a Fisher natural gradient
would move model means at rate $\tfrac12$.

Part (i) holds the matching. The mean then follows the noiseless rate-separated flow with
$\varepsilon=\varepsilon_{\mathrm m}$, fast block $\bar k$, slow block $(\bar m,\check\omega)$. On each
selected edge, in the difference coordinates $(\bar k^-_i,\bar m^-_i,\check\omega_i)$, the Hessian block
is
\[
\begin{pmatrix}4+3\mathrm T&-2\mathrm T&-2\sqrt2\\-2\mathrm T&4+3\mathrm T&-2\sqrt2\\
-2\sqrt2&-2\sqrt2&4+\mathrm T\end{pmatrix},
\]
with fast curvature $4+3\mathrm T\geq4$ uniformly in temperature, so with equal model and link rates
the Tikhonov hypotheses hold on that block with constant $4$. On the pair means $\bar k^+_i$, however,
the fast curvature is only $3\mathrm T$: at the small temperatures of the formation regime that
direction relaxes at a rate of order $\mathrm T$, so on the full sheet the constant is of order
$\mathrm T$ and not uniform as $\mathrm T\to0$.

Part (ii) sets $\mathrm T=0$ and starts with every residual at $-1$. Each selected-edge block has a zero
mode, on which both residuals vanish, and rates
$\varsigma_\pm=2(1+\varepsilon_{\mathrm m}+\varepsilon_{\check\omega})\pm2\sqrt{(1-\varepsilon_{\mathrm m})^2
+\varepsilon_{\check\omega}^2}$. A fast link ($\varepsilon_{\check\omega}=1$) gives
$|r^m_i(t)|\leq e^{-(4-2\sqrt2)t}$ for every $\varepsilon_{\mathrm m}$, removing the model disagreement
on the belief clock even with frozen model means; fast model means give
$r^m_i=-e^{-(4+4\varepsilon_{\check\omega})t}$. With slow means and equally slow links
($\varepsilon_{\check\omega}=\varepsilon_{\mathrm m}\leq\tfrac12$),
\[
|r^b_i(t)|\leq(1+\varepsilon_{\mathrm m})e^{-4t}+\varepsilon_{\mathrm m},
\qquad
\bigl(1-\tfrac32\varepsilon_{\mathrm m}\bigr)e^{-6\varepsilon_{\mathrm m}t}-\tfrac32\varepsilon_{\mathrm m}
\leq|r^m_i(t)|\leq e^{-(6\varepsilon_{\mathrm m}-2\varepsilon_{\mathrm m}^2)t}+\tfrac32\varepsilon_{\mathrm m},
\]
so for $\varepsilon_{\mathrm m}\leq1/25$ each matched pair has $|r^b_i|\leq2\varepsilon_{\mathrm m}$ and
$|r^m_i|\geq\tfrac12$ from $\tfrac14\log(4/\varepsilon_{\mathrm m})$ to $1/(12\varepsilon_{\mathrm m})$:
beliefs agree while models disagree, for a time of order $1/\varepsilon_{\mathrm m}$.

Part (iii) returns to the sampled process along any assignment path: a model residual of size at least
one drops below $\tfrac12$ before time $t$, on the event that residuals and configuration stay bounded,
with probability at most $2\exp[-1/(192\mathrm T\bar\varepsilon t)]$,
$\bar\varepsilon=\max(\varepsilon_{\mathrm m},\varepsilon_{\check\omega})$, for $t$ up to order
$1/\bar\varepsilon$. It is a bound on a joint event, not an unconditional persistence probability. The
mechanism-level content, conjectural beyond the four-cycle, is that a link equilibrating on the belief
clock removes any disagreement measured through it; the rates $4-2\sqrt2$ and $6\varepsilon_{\mathrm
m}$, the window, and the constant $192$ are four-cycle specifics.

\subsection{Parents: the prepared parent and the refresh hypothesis}
\label{comp:ch08-formulas-parent}

The last section is finite (\MT{def:rg-memory-parent-sweep}): fine states $\xi$, a parent map
$\operatorname{par}$, a positive prior and likelihoods with posterior $p$, and unobserved records with
laws $\psi_\xi$. The kernels are
\begin{align*}
&\mathsf M_{\xi\xi'}=\mathsf a_{\xi\xi'}\min\Bigl(1,\frac{p(\xi')}{p(\xi)}\Bigr)\quad(\xi\neq\xi'),
\qquad
\mathsf C_{\xi z}=\mathbf 1_{\operatorname{par}(\xi)=z},
\qquad
\mathsf J_{z\xi}=\frac{p(\xi)\mathbf 1_{\operatorname{par}(\xi)=z}}{p_Z(z)},
\\
&\mathsf D=\mathsf C\mathsf J,
\qquad
\mathsf T=\mathsf M\mathsf D,
\qquad
\Fenergy(q\otimes\psi';o)=\KL(q\Vert p)+\E_q\KL(\psi'_\xi\Vert\psi_\xi)-\log\mathcal Z_o .
\end{align*}
$\mathsf M$ is a Metropolis move with symmetric proposals; $\mathsf C$ reads the parent; $\mathsf J$
draws a fine state from the posterior given the parent; $\mathsf D$ refreshes within the parent;
$\mathsf T$ is one sweep. The free energy shows that matching record conditionals is a specialization
of the joint objective. \MT{thm:rg-memory-inference-descent} shows that the kernels preserve $p$, that
the refresh is the exact block update given the parent, that any $p$-stationary kernel lowers
$\KL(q\Vert p)$ by an explicit reverse-conditional divergence, that the sweep lowers it strictly unless
$q$ is in conditional equilibrium on each connected component, and that convergence is geometric on
each component, with no uniform rate.

The prepared-parent theorem (\MT{thm:rg-memory-prepared-parent}) reads
\[
\mathsf J\mathsf T=\overline{\mathsf T}\mathsf J,
\qquad
(\nu\mathsf J)\mathsf T^n=\nu\overline{\mathsf T}^n\mathsf J,
\qquad
\overline{\mathsf T}=\mathsf J\mathsf M\mathsf C .
\]
From a fine law $\nu\mathsf J$ in conditional equilibrium within parents, one sweep returns a fine law
in conditional equilibrium with parent law $\nu\overline{\mathsf T}$. The parent history is Markov with
the $p_Z$-reversible kernel $\overline{\mathsf T}$, and from any start after the first refresh, since
after a refresh the fine state's law given the whole parent history depends only on the current parent.
This is standard in substance: a chain resampled from block-conditional equilibrium at every step is a
function of its block chain. The manuscript's parallel with the continuous reduced conjugacy (fast
variable at its conditional optimum given the slow one) is structural; no theorem links the two.

Everything rests on \MT{hyp:rg-memory-instantaneous-refresh}, exact within-parent resampling at every
sweep, which the manuscript calls adiabatic elimination with an infinitely shorter timescale. The
partial refresh $\mathsf D_\omega=(1-\omega)I+\omega\mathsf C\mathsf J$, $\omega<1$, keeps
normalization, reversibility, and the parent, and loses the Markov property
(\MT{rem:rg-memory-partial-refresh}). With $\operatorname{par}(1)=\operatorname{par}(2)=z_1$,
$\operatorname{par}(3)=z_2$, uniform $p$, and one proposal between $2$ and $3$,
\[
\mathbb P(Z_2=z_2\given Z_1=z_1,Z_0=z_1)=\tfrac16+\tfrac{\omega}{12},
\qquad
\mathbb P(Z_2=z_2\given Z_1=z_1,Z_0=z_2)=\tfrac12-\tfrac{\omega}{4}.
\]
State $2$ is the border of $z_1$ with $z_2$. A parent just arrived from $z_2$ has its fine state at the
border with probability $1-\omega/2$, a parent already in $z_1$ with probability $\tfrac13+\omega/6$;
the parent's past says where in the fiber the fine state sits, and that predicts the next move. The
dependence $(1-\omega)/3$ vanishes only at $\omega=1$. The parent path law is still retained exactly
as a function of the fine chain (\MT{sec:rg-memory-hidden-history}), and the cost of forcing a
first-order Markov parent is the conditional mutual information $D^{(1)}_N$ of
\MT{thm:rg-memory-projection}, positive here at $N=3$.

\section{What is proved, and how solid it is}
\label{comp:ch08-status}

\begin{scorecard}
\MT{def:obs-fast-slow-blocks} & The fast--slow split, rate ratio, and block mobilities are declared. & \CHOICE & The free energy selects none; a rate ratio can change the orbit.\\
\MT{prop:obs-local-block-descent} & Any positive-definite block mobility descends; in product families local descent is a block of collective descent. & \SOLID & Standard in substance; correlated recognition breaks block diagonality.\\
\MT{lem:obs-envelope-slow-flow} & Envelope identity; the reduced slow flow is the natural gradient of the reduced functional. & \SOLID & Standard (envelope theorem, Danskin); needs a $C^1$ branch of minimizers.\\
\MT{prop:obs-rate-separated-flow} & On the graph of conditional minimizers the rate-separated flow has the reduced slow component; projection is a conjugacy. & \SOLID & Kinematic: says nothing about approaching the graph.\\
\MT{rem:obs-rate-separation-witnesses} & Off the graph no semiconjugacy; the graph is not invariant; correlated metrics move the slow set and can break descent; congruence repairs both. & \SOLID & Negative witnesses in two and three dimensions.\\
\MT{prop:obs-slow-model-link-dichotomy} & On the four-cycle a belief-agreeing, model-disagreeing interval of order $1/\varepsilon_{\mathrm m}$ exists exactly when model means and links are both slow; held sheets meet the Tikhonov hypotheses. & \TOY & Four agents, scalar Gaussian means, translation links, supplied rates; (ii) at zero temperature on a held matching.\\
\MT{open:obs-slow-manifold} & Tracking of the reduced flow by the rate-separated or sampled flow. & \OPEN & Posed for the sampled four-cycle class; learning outside.\\
\MT{thm:rg-memory-inference-descent} & Kernels preserve the posterior; refresh is the exact block update; sweeps descend and converge geometrically on components. & \SOLID & Standard in substance; no uniform rate.\\
\MT{hyp:rg-memory-instantaneous-refresh} & Every sweep resamples the fine state exactly given the parent. & \CHOICE & Idealized infinite separation; substantive.\\
\MT{thm:rg-memory-prepared-parent} & From a prepared start the parent history is Markov with a reversible lumped kernel. & \COND & Hypothesis: instantaneous refresh. Standard in substance.\\
\MT{rem:rg-memory-partial-refresh} & Under partial refresh the parent is not Markov; the past dependence is $(1-\omega)/3$. & \SOLID & Refutation by a three-state witness; constant specific to it.\\
\end{scorecard}

\begin{hygiene}
Conditions that can fail in the program's cases: the declared split, rate ratio, and mobilities; a
$C^1$ branch of minimizers with fast curvature uniformly bounded below, which fails after source rows
are optimized even when every source energy and frozen-weight curvature is positive
(\MT{prop:obs-attention-reduced-hessian}) and on the full four-cycle sheet, where the constant is
$3\mathrm T$; a block-diagonal metric with $f$-independent slow block, which correlated recognition
breaks; a product family for local descent; instantaneous refresh for the Markov parent; and, on the
four-cycle, a held matching in (i) and (ii), zero temperature in (ii), bounded residuals and
configuration in (iii), and no learning.

Hygiene: $C^1$ and $C^2$ regularity, solutions staying in the open domain, integrable log normalizers,
finite state spaces with positive prior and likelihoods, and proposal row sums at most $\tfrac12$,
which secures positive holding.
\end{hygiene}

\section{Toy versus general}
\label{comp:ch08-toy}

General and independent of $N$: block descent, the envelope identity and reduced flow, the kinematic
agreement on $\mathcal S^\star$, the parent sweep's descent, and the Markov parent under instantaneous
refresh. The negative witnesses (a slow velocity off the graph, the correlated-metric slow set, the
three-state partial refresh) are refutations with general force; their constants $7-4\sqrt3$ and
$(1-\omega)/3$ are specific. The four-cycle dichotomy is TOY: four agents, scalar Gaussian means with
fixed variances, translation links ($\R$ acting on means, an abelian noncompact group), supplied rates,
a held matching in (i) and (ii), zero temperature in (ii). These exclude the target of $512$ or more
agents with nonabelian links such as $\GL(K,\R)$ or $\mathrm{SU}(N)$. The mechanism, that a link fast
on the belief clock absorbs disagreement measured through it so that model incoherence with belief
coherence needs slow links as well as slow models, is plausible on larger graphs and unproved there;
the rates, the window, the bound $\varepsilon_{\mathrm m}\leq1/25$, and the persistence constant are
four-cycle specifics. The only place the open tracking problem is closed, the held sheets, is itself a
toy, as the chapter says.

\section{Open ends}
\label{comp:ch08-open}

\MT{open:obs-slow-manifold} asks when Tikhonov's and Fenichel's theorems apply to the rate-separated
flow at $c_*$. Its three hypotheses (a $C^2$ objective with $C^1$ block metrics and an
$f$-independent slow metric; a $C^1$ branch of minimizers with fast curvature at least $\mathrm c_1$
relative to the fast metric on a compact set; a reduced orbit inside that set and a fast start in the
basin) are verified only on held sheets. Source-row optimization can destroy the curvature bound,
sampled dynamics with source jumps lie outside a deterministic theorem, and no explicit bound in
$\varepsilon$ on the slow orbit or its Fisher duration exists. The problem is posed in one class, the
sampled four-cycle with slow model means and links at equal rates and no learning: settled positively
by a bound on the distance between the sampled slow coordinates and the reduced flow of the assignment
path over a bounded slow-time interval, explicit in $\varepsilon_{\mathrm m}$, $\mathrm T$, and the
number of assignment changes, tending to zero, with probability tending to one when changes per unit
slow time are bounded; negatively by rates and temperatures tending to zero along which that distance
stays bounded below with probability bounded below. Learning and correlated belief--model metrics are
outside the posed class; for the latter the chapter establishes only that the rate declaration must
name its slow set. \MT{open:elbo-alternating-convergence} asks for convergence of the discrete
alternating scheme, and \MT{open:rg-pointwise-parent-dynamics} asks how a parent's dynamics follow
from the network it summarizes, which the prepared parent answers only under instantaneous refresh.
The uniform rate of the parent sweep and persistence under sampling with jumps at the formation
temperature are declined without open problems.

\begin{verdictbox}{\NARROWER}
\textbf{Claimed.} The opening promises descent for every positive-definite block mobility, local in
product families; a reduced slow flow conjugate to its lift on the graph of the conditional optimum,
the exact route for information-discarding maps; a Markov parent under instantaneous refresh, failing
under finite refresh; and, on the four-cycle, a model channel lagging the belief channel exactly when
links are slow too.

\textbf{Delivered.} All of it, as stated. The general results are the kinematics of an adiabatic
reduction (envelope theorem, reduced natural gradient, agreement on the graph of minimizers) and its
discrete, infinitely separated case (lumpability), all standard in substance, plus clean witnesses
that correlated metrics move the slow set and finite separation creates memory. The four-cycle
dichotomy is a complete toy result with explicit rates and a pathwise probability bound.

\textbf{Gap.} The dynamical content of a two-timescale theory, that the actual flow at small
$\varepsilon$ follows the reduced one, holds only on the four-cycle's held sheets, where the chapter
verifies Tikhonov's hypotheses; it is not established on the full sheet (fast curvature of order the
temperature), with source jumps, for correlated metrics, or with learning. The ``exact route'' is exact
on the graph of minimizers and unproved as dynamics. The Markov parent rests on a named idealization,
and parent dynamics under finite refresh is open. The split, rate ratio, and mobilities are declared,
and the formation theorem does not use the belief--model split. The scope paragraph declines each of
these; the title promises more than the theorems deliver.
\end{verdictbox}

\chapter{Coarse Channels}
\label{comp:ch09}

\section{The question and the answer}
\label{comp:ch09-question}

\MT{ch:coarsegraining} opens Part II of the main manuscript. It asks what a coarse description of the
agent network at the base point $c_*$ keeps of three laws: the normalized generative joint
$\mathbb P(do,dy)$ of \mtc{ch:generative}, its selected posterior $\boldsymbol\Pi_o$, and a recognition
law $\mathbb Q_o$. The answer is one identity. Take a coarse channel $C:\mathsf Y\rightsquigarrow\mathsf Z$,
a Markov kernel that reads the latent coordinates but neither the observation nor the recognition law,
and push all three laws through it. The evidence $p(o)$ is then the same at both scales, the pushed
posterior is a posterior of the pushed joint, and the free energy splits exactly
(\MT{thm:rg-exact-coarse-vfe}):
\[
\Fenergy_{\mathbb P}(\mathbb Q_o)=\Fenergy_{\mathbb P^c}(\mathbb Q_oC)+\Delta(o),
\qquad
\Delta(o)\in[0,+\infty],
\]
where $\Delta(o)$ is what recognition gets wrong, relative to the posterior, inside the coarse cells the
channel no longer resolves. Applied to a finite block $I$ of agents and a parent state space, the same
identity produces a complete normalized parent triple, the pointwise parent datum
(\MT{thm:rg-pointwise-parent-datum}), on which \cref{comp:ch10,comp:ch11,comp:ch12,comp:ch13} build.
Around this core the chapter collects kernel tools (composition, data processing with its equality
case, Bayes recovery), proves that three operations physicists all call coarse-graining obey different
theorems, treats symmetry, defines scale diagrams, and ends with labeled mixtures whose retained label
carries Fisher information of order $N^2$ at a tie.

With the latent links of \cref{comp:lie-links}, the fine state includes the generators:
the same theorem applies to a channel on $(\mathring\omega,y)$. A tree transport that reads links
belongs to that joint channel. Keeping or discarding generators changes the conditional remainder,
and the exact identity supplies no Gaussian closure after they are marginalized
(\MT{prop:link-conditional-transfer}).

The verdict in advance: the core identity is the relative-entropy chain rule applied twice to the joint
law of fine and coarse variables. It is standard in substance; the contribution is careful typing and a
clean statement of what the identity does not give. Exact means an identity with an explicit
nonnegative remainder. It does not mean closure of the coarse law in a model family, a selected
channel, or parent dynamics, and the chapter claims none of these.

\section{The picture}
\label{comp:ch09-picture}

A coarse channel is a noisy detector attached to the latent coordinates: given a fine configuration $y$
it emits a coarse value $z$ with probability $C(y,dz)$, and a deterministic statistic such as a block
sum is the noiseless case. Pushing a law through it is what a stochastic block-spin rule does to a Gibbs
measure.

\begin{analogy}{exact}
\textbf{One step of Kadanoff blocking.} A block-spin weight $T(\sigma'\given\sigma)\geq0$ with
$\sum_{\sigma'}T(\sigma'\given\sigma)=1$ maps $e^{-H}$ to
$e^{-H'(\sigma')}=\sum_\sigma T(\sigma'\given\sigma)e^{-H(\sigma)}$, and this trace condition leaves the
partition function invariant \citep{kadanoff1966scaling,Cardy1996}. A coarse channel is such a weight on
the agents' latent coordinates, and evidence preservation (\MT{thm:cg-evidence-preserving-channel}) is
the trace condition with $p(o)$ in the role of $Z$. The correspondence is exact for one step, and
everything Kadanoff and Wilson \citep{Wilson1974} add to that step is absent. There is no lattice and no
length to rescale: the agents sit at one base point and blocks are vertex sets of a graph. The coarse
law is not projected onto a finite set of couplings (an effective action first appears in
\cref{comp:ch11}). No single rule is iterated, so there is no flow on a parameter space, no fixed point,
and no exponent. The standing scope declines all of these (\MT{sec:intro-conventions}); the asymptotic
apparatus is in Supplement B.
\end{analogy}

The split itself is simple. Recognition and posterior are two laws on the fine space; the channel
coarsens both; their relative entropy splits into the relative entropy of the coarse laws plus the
relative entropy of the two conditional laws of the fine configuration given each coarse value, averaged
over coarse recognition. The common term $-\log p(o)$ carries this over to free energies. The coarse free
energy is never the larger: the coarse description looks better only because it no longer sees the
errors it discarded, which the manuscript calls bookkeeping, not learning.

A two-spin example (the companion's arithmetic) shows which errors survive. Let the posterior be the
Ising pair $\boldsymbol\Pi(s_1,s_2)\propto e^{Js_1s_2}$ and recognition the uniform product law, a
mean-field guess, so that $\KL(\mathbb Q\Vert\boldsymbol\Pi)=\log\cosh J$. A channel keeping the bond
$s_1s_2$ finds both laws uniform inside each bond cell: $\Delta=0$ and the coarse gap is the whole
$\log\cosh J$. A channel keeping only $s_1$ finds both coarse laws equal to $(1/2,1/2)$: the coarse gap
is zero and $\Delta=\log\cosh J$, so the one-spin description is perfectly calibrated and the entire
error of that guess sits in the discarded cell.

\begin{analogy}{exact}
\textbf{Constrained (Landau) free energy.} On finite spaces and for a deterministic statistic $c$, the
coarse free energy of a coarse law $q'$ is the minimum fine free energy over fine laws with image $q'$
(\MT{prop:cg-stochastic-infimum}, part (i)); for $q'=\delta_z$ it is
$-\log[p(o)\boldsymbol\Pi_o(c=z)]$, the constrained free energy at fixed order parameter. For a noisy
channel this fails (part (ii)): the coarse free energy of $q'$ can lie strictly below every fine free
energy compatible with $q'$.
\end{analogy}

Physicists call three different moves coarse-graining: integrating out (pushing a law through a
kernel), restricting the configurations that enter an energy (fields constant on blocks), and
compressing an operator onto a subspace (a multigrid coarse operator). \MT{sec:cg-three-operations}
proves that they obey different theorems, so a proof about one transfers nothing to another.

\begin{analogy}{structural}
\textbf{Feshbach projection.} The Galerkin coarse propagator reproduces the fine dynamics on kept
signals exactly when the operator never maps discarded directions into kept ones
(\MT{prop:cg-message-passing-residual}), which in Feshbach language is the vanishing of the coupling
that would generate a self-energy. The block-operator algebra is the same; energy-dependent
self-energies are not treated.
\end{analogy}

At the end of the chapter the retained coordinate is a label $\gamma$ from a finite menu whose log weight
grows like $N$. The weights are then a softmax of extensive arguments, like the weights of coexisting
phases with free energies $Nf_\gamma$: a strict winner hides the label, a tie makes it a very sensitive
function of the parameter, and a window of width $1/N$ interpolates.

\begin{analogy}{structural}
\textbf{Finite-size rounding of a first-order transition.} \citet{BorgsKotecky1990} prove for lattice
models with finitely many coexisting phases that the transition is rounded over a window of width of
order inverse volume; at coexistence the variance of the total conjugate extensive quantity is then of
order volume squared. The retained-label theorems (\MT{sec:cg-retained-label-asymptotics}) are the same
softmax calculus with no lattice and no Gibbs measure; extensive growth of the log weights is assumed,
not derived from an agent model. The manuscript names this correspondence itself.
\end{analogy}

\section{Objects and notation}
\label{comp:ch09-objects}

\begin{decoder}
coarse channel $C:\mathsf Y\rightsquigarrow\mathsf Z$ & Markov kernel on latent coordinates, stochastic by default, reading neither $o$ nor $\mathbb Q_o$; composites act on the right & block-spin weight $T(\sigma'\given\sigma)$\\
selected posterior version $\boldsymbol\Pi_o$ & a regular conditional law of $y$ given $o$, fixed at every admitted $o$ & conditional Gibbs state\\
extended free energy $\Fenergy_{\mathbb P}(\mathbb Q_o)$ & $-\log p(o)+\KL(\mathbb Q_o\Vert\boldsymbol\Pi_o)$ for every law; $+\infty$ off absolute continuity & Gibbs--Bogoliubov functional\\
bridge lift $\widehat{\mathbb Q}_o(dy,dz)$ & joint law $\mathbb Q_o(dy)C(y,dz)$ of fine and coarse values & joint law of spins and block spins\\
discarded information $\Delta(o)$ & expected relative entropy of recognition to posterior given the coarse value & information lost inside blocks\\
Bayes reverse kernel $R_Q$ & reverse channel by Bayes' rule under $Q$; recovering a pair means restoring both laws & classical Petz map; sufficiency\\
parent state space $\mathsf Z_A=\mathsf B_A\times\mathsf M_A\times\boldsymbol\Xi_A\times\mathsf H_A$ & parent belief samples, model presentations, evaluator inputs, retained marks of the candidate parent $A$ of block $I$ & block variables plus bookkeeping\\
evaluator $\operatorname{ev}_A$ & map from a parent model value to a normalized kernel generating $(b_A,o,h_A)$ & coupling-to-Hamiltonian map\\
energy precomposition & substitute a restricted configuration into an energy, $\bar E(z)=E(\iota z)$ & constrained configurations\\
Galerkin data $S,M,L,H_c$ & prolongation, reference form, operator, coarse propagator & multigrid coarse operator\\
scale diagram & functor from resolutions to kernels; composites along routes agree & semigroup property of blocking\\
finite labeled family & mixture whose component label stays in the sample & multiphase ensemble with phase label\\
log mass $\lambda_\gamma$, weight $W_\gamma$ & log weight of label $\gamma$; softmax of the $\lambda$ & $-\beta F_\gamma$; phase weight\\
extensive size $N$, density $\mathsf s_\gamma$ & index along which log masses grow linearly; $\lim\lambda_\gamma/N$ & volume; free-energy density\\
strict winner, tie, window $\mathsf w$ & one maximal density; several; $\theta=\theta_0+\mathsf w/N$ & one phase; coexistence; rounding window\\
natural selector & data-dependent partition rule compatible with relabeling and gauge & choice of blocking scheme\\
\end{decoder}

Exact means an identity with a remainder, never a truncation-free flow; scale, resolution, and depth
are indices with no length attached. $C$ is always the channel, never a covariance. $\Delta$ is the
discarded information here and a matrix (the cut excess) in \cref{comp:ch10}, and $W$ is a label weight
here and an edge form there. $N$ is an extensive size that the manuscript allows, but does not require,
to be the number of agents. In the house style, declared marks an input the theory does not select (the
block $I$, the label $A$, the kernel $C_A$), which is substantive; selected version and admitted
observation (finite positive evidence) are null-set bookkeeping, and vertical over the identity of
$c_*$ is automatic at one base point.

\section{The formulas, explained}
\label{comp:ch09-formulas}

\subsection{Channels, contraction, and recovery}

A channel sends $P$ to $(PK)(B)=\int K(x,B)P(dx)$, and $P(KL)=(PK)L$ (\MT{prop:cg-markov-category}):
stochastic matrices on standard Borel spaces. Every channel contracts relative entropy; when
$\KL(P\Vert Q)<\infty$, equality holds exactly when $dP/dQ$ depends on $x$ only through the output
(\MT{thm:cg-dpi-equality}), and then the Bayes reverse kernel of $Q$ also recovers $P$
(\MT{cor:cg-pairwise-bayes-recovery}). This is the information-loss theorem of
\citet{Kullback1951,Csiszar1967} and the commutative case of \citet{Petz1988}. Finiteness cannot be
dropped: a three-point example has $\KL(PK\Vert QK)=\KL(P\Vert Q)=+\infty$ and no common recovery
(\MT{cor:cg-dpi-infinite-equality-warning}). On the singular branch of the theory, where recognition is
not absolutely continuous, identities hold as $+\infty=+\infty$ and say nothing.

\subsection{The exact coarse free energy}

The central display (\MTeq{eq:rg-vfe-chain-rule}) is
\begin{keyequation}{Coarse free energy and discarded information}
\[
\begin{aligned}
\Fenergy_{\mathbb P}(\mathbb Q_o)
&=\Fenergy_{\mathbb P^c}(\mathbb Q^c_o)+\Delta(o),\\
\Delta(o)&=\int_{\mathsf Z}\KL\bigl(\widehat{\mathbb Q}_o(\cdot\given z)\Vert\widehat{\boldsymbol\Pi}_o(\cdot\given z)\bigr)\mathbb Q^c_o(dz).
\end{aligned}
\]
\end{keyequation}
The first term on the right is computed entirely at the coarse scale, from the pushed joint (evidence
$p(o)$, posterior $\boldsymbol\Pi_oC$) and the pushed recognition law $\mathbb Q^c_o=\mathbb Q_oC$. The
remainder compares, for each coarse value $z$, what recognition and posterior say about the fine
configuration given $z$, weighted by how often recognition produces $z$. If $C$ is the identity,
$\Delta=0$. If $C$ sends everything to one point, the coarse free energy is $-\log p(o)$ and
$\Delta=\KL(\mathbb Q_o\Vert\boldsymbol\Pi_o)$: the whole gap is discarded. For a deterministic
statistic, $\Delta=0$ exactly when recognition has the posterior's shape inside every fiber
$c^{-1}(z)$. For Gaussian laws and a linear projection $(z,w)\mapsto z$, $\Delta$ is the expected
relative entropy of the two conditional Gaussians of $w$ given $z$.

Is this more than the chain rule and data processing? In mathematical substance, no: the proof
disintegrates the bridge lift over $y$, where the common conditional $C(y,\cdot)$ drops out, and over
$z$, where the coarse term and $\Delta$ appear (\MT{lem:prob-kl-chain-rule}). The theorem adds typing
that later chapters use: the pushed posterior is a coarse posterior version at every declared
observation, including exceptional ones, which matters because the theory evaluates at the one observed
record; the identity holds in the extended reals without absolute continuity; with finite fine relative
entropy, $\Delta(o)=0$ exactly when one kernel recovers both recognition and posterior; and recovery of
a whole family needs the hypothesis for every member at once (\MT{rem:cg-support-boundaries} handles the
support boundaries). For the ELBO (\MTeq{eq:cg-elbo-monotone}),
\[
\Lelbo^{\mathrm{ext}}_{\mathbb P^c}(\mathbb Q_oC)
=\log p(o)-\KL(\mathbb Q_oC\Vert\boldsymbol\Pi_oC)
\geq\log p(o)-\KL(\mathbb Q_o\Vert\boldsymbol\Pi_o)
=\Lelbo^{\mathrm{ext}}_{\mathbb P}(\mathbb Q_o),
\]
and \MT{thm:cg-evidence-preserving-channel} draws the right conclusion: the coarse ELBO is higher by
$\Delta(o)$ at unchanged evidence, which is neither model improvement nor evidence ascent, and ELBOs of
separately declared models compare tightness only at equal evidence.

On finite spaces, where \MT{prop:cg-stochastic-infimum} is stated, and for a deterministic statistic
the coarse free energy is the constrained minimum of the fine one, attained by $q'(c(y))\boldsymbol\Pi_o(y\given c(y))$. For a noisy channel it is not: with a uniform
binary posterior, the symmetric channel with flip probability $1/4$, and $\mathbb Q=(1,0)$, the only fine
law with image $(3/4,1/4)$ is $\mathbb Q$, and its free energy exceeds the coarse one by
$2\log2-\frac34\log3\approx0.56$ (\MT{prop:cg-stochastic-infimum}).

\subsection{The pointwise parent datum}

For a block $I$, a parent label $A$, and a parent kernel $C_A:\mathsf Y_I\rightsquigarrow\mathsf Z_A$,
\[
\mathbb P_A(do,dz\given X)=\int_{\mathsf Y_I}C_A(y,dz)\mathbb P_I(do,dy\given X),
\qquad
\boldsymbol\Pi_{A,o,X}=\boldsymbol\Pi_{I,o,X}C_A,
\qquad
\mathbb Q_{A,o,X}=\mathbb Q_{I,o,X}C_A,
\]
and $\Fenergy_I=\Fenergy_A+\Delta_A$ (\MT{thm:rg-pointwise-parent-datum}). The corollary takes
$\mathbb Q_{I,o,X}\ll\boldsymbol\Pi_{I,o,X}$ as an input, which is harmless because the exact coarse free
energy identity holds without it. Does a parent law always
exist? Yes, trivially: a pushforward exists for every declared block and kernel. Is the parent an agent
with its own dynamics? No, and the statement says what is missing. The block, label, and kernel are
inputs that nothing selects (\MT{open:cg-natural-selector}). The parent's model coordinate is generative
only through an evaluator; disintegrating $\mathbb P_A$ always supplies one, so the parent always has a
generative mechanism, but it is whatever $\mathbb P_A$ says, and an evaluator declared independently
(say, a Gaussian agent with stated parameters) agrees with $\mathbb P_A$ only under a separate
compatibility hypothesis that disintegration does not validate. The parent's belief and model marginals
do not determine its law unless one is a point mass and the auxiliary factors are trivial. And the
construction is static: no fine evolution induces a parent evolution
(\MT{open:rg-pointwise-parent-dynamics}, the defining open problem of the phase, \cref{comp:ch11}).
How holonomy enters the parent is the subject of \cref{comp:ch10}.

\subsection{Three operations}

Energy precomposition, $\bar E(z)=E(\iota z)$ and $\bar P(dz)=e^{-\bar E(z)}\bar\nu(dz)/\bar Z$, needs a
declared reference $\bar\nu$. On a countable space with $E=-\log p$ it is conditioning on
$X\in\iota(\mathsf Z)$, and conditioning can increase relative entropy: $(1,1,98)/100$ and
$(1,2,97)/100$ are about $0.003$ apart, their restrictions $(1/2,1/2)$ and $(1/3,2/3)$ about $0.059$
(\MT{prop:cg-hard-identification}). So no kernel produces precomposed laws, and data processing, evidence
preservation, and ELBO monotonicity do not transfer. On a continuous fiber the restricted set is
typically null, as for block-constant configurations, and the answer depends on the reference:
$\bar E(z)=z^2$ gives $\mathcal N(0,1/2)$ against Lebesgue measure and $\mathcal N(0,1/4)$ against
$e^{-z^2}dz$.

Galerkin restriction (\MT{def:cg-galerkin-data}) takes an injective prolongation $S$, a form $M\succ0$,
and an operator form $L$, with $H=M^{-1}L$, $S^+=(S^\top MS)^{-1}S^\top M$, and $H_c=S^+HS$; the form
$S^\top LS$ is the aggregation multigrid coarse operator \citep{Notay2010Aggregation}. The residual
$S^+H-H_cS^+=S^+H(I-SS^+)$ vanishes exactly when (\MTeq{eq:cg-galerkin-exactness}) $H$ maps the
$M$-orthogonal complement of $\operatorname{range}S$ into itself; for symmetric $L$ this is invariance of
$\operatorname{range}S$, and without symmetry the two conditions are independent
(\MT{prop:cg-message-passing-residual}). As an illustration outside the manuscript, for block aggregation
with $M=I$ and $L$ a graph's adjacency matrix it is the equitable-partition condition: every vertex of a
block has the same number of neighbors in each block. Restricting a recognition family is yet another
operation: it keeps the joint and optimizes the same ELBO over a subset.

\subsection{Symmetry and scale diagrams}

An equivariant channel, $K(gx,B)=K(x,g^{-1}B)$, commutes with the action on laws
(\MT{prop:cg-equivariant-channels}); at $c_*$ gauge covariance is equivariance for $\rho_b,\rho_m$, and
relabeling covariance is equivariance for permutations. The Bayes reverse kernel is defined only almost
everywhere, and whether it can be chosen equivariantly everywhere is settled for finite spaces, for
compact stabilizers by Haar averaging (\MT{prop:cg-equivariant-bayes-versions}), and for any stabilizer
when the posterior has a weakly continuous version (\MT{cor:cg-continuous-version-equivariant}).

A scale diagram (\MT{def:rg-scale-diagram}) is a functor from resolutions to kernels, with right-acting
composites $K_{\ell r}=K_{\ell k}K_{kr}$. On a chain this is automatic (\MT{prop:rg-scale-functoriality}).
It has content on a poset that is not a chain, where two routes must agree, and for a blocking rule
indexed by block factor, where it is the semigroup law $K_{b_1b_2}=K_{b_1}K_{b_2}$. Block sums satisfy it;
majority rule with ties broken by the first member does not, since $(0,0,1,0,1,0)$ goes to $1$ by pairs
then triples and to $0$ by triples then pairs. Mutually inverse kernels between finite spaces are
permutations, so no stochastic isomorphism links levels of different sizes. The Dobrushin contraction
bound is not in this chapter: it is part of \MT{thm:rg-action-lp-contraction} (\cref{comp:ch11}), and
the scale diagram here carries no contraction rate.

\subsection{Labeled families and Fisher information}

A finite labeled family (\MT{def:rg-thermo-label-family}) keeps the label in the sample,
$P_\theta(\gamma,x)=W_\gamma(\theta)p_{\gamma,\theta}(x)$ with $W_\gamma\propto e^{\lambda_\gamma}$. With
$\xi_\gamma=\nabla\lambda_\gamma$ and $\bar\xi=\sum_\gamma W_\gamma\xi_\gamma$, the score is
$\xi_\gamma-\bar\xi+\mathbf s_{p_\gamma}(x)$ and the Fisher tensor is
(\MTeq{eq:rg-thermo-finite-labeled-fisher})
\[
g^F(P)=\operatorname{Cov}_W(\xi_\gamma)+\sum_\gamma W_\gamma g^F(p_\gamma),
\]
label information plus weighted branch information (\MT{prop:rg-thermo-finite-labeled-fisher}), the law
of total covariance for the score. If $\lambda_\gamma=-\beta F_\gamma(\theta)$, then $\xi_\gamma$ is the
extensive quantity conjugate to $\theta$ in phase $\gamma$ and the first term is its between-phase
variance, the coexistence part of a susceptibility. Observed-data information is complete minus missing
information, Louis's identity \citep{Louis1982} averaged over records
(\MT{prop:rg-thermo-finite-missing-information}); a channel acting within branches never touches the
label information (\MT{prop:rg-thermo-finite-fisher-loss}); and along a word of channels the Fisher
losses and the discarded relative entropies telescope (\MT{prop:rg-thermo-account-word}), dividing by an
extensive size giving per-agent bookkeeping, not a thermodynamic density.

The asymptotics assume extensive growth, $\lambda_{\gamma,N}=O(N)$ in $C^3$ with
$\lambda_{\gamma,N}/N\to\mathsf s_\gamma$. With the branch densities $N^{-1}g^F(p_{\gamma,N})$ converging in
$C^2$, a strict winner makes the Fisher density per $N$ converge to the winning branch's, with
exponentially small errors (\MT{thm:rg-thermo-label-strict}). At a tie where the weights converge,
$W_{\gamma,N}(\theta_0)\to\bar{\mathsf p}_\gamma$, and the branch Fisher tensors are $O(N)$,
$N^{-2}g^F(P_N)(\theta_0)\to\operatorname{Cov}_{\bar{\mathsf p}}(\mathsf v_\gamma)$ with
$\mathsf v_\gamma=\nabla\mathsf s_\gamma(\theta_0)$, of rank at most
$\min(d,\lvert\operatorname{supp}\bar{\mathsf p}\rvert-1)$ (\MT{prop:rg-thermo-label-tie-scale}). In the
window $\theta=\theta_0+\mathsf w/N$, given center ratios $\mathrm c_\gamma$ of the tied weights, the
label law tends to a softmax of $\mathsf w\cdot\mathsf v_\gamma+\log\mathrm c_\gamma$ and the Fisher
tensor in $\mathsf w$ to (\MTeq{eq:rg-thermo-label-crossover-fisher})
\[
D^2_{\mathsf w}\log\sum_{\gamma\in S}\mathrm c_\gamma e^{\mathsf w\cdot\mathsf v_\gamma},
\]
the Hessian of a finite log-partition function (\MT{thm:rg-thermo-label-crossover}); the center ratios
are extra finite-$N$ data. \MT{prop:rg-thermo-label-tie-counterexamples} shows, with one
parameter-independent branch law throughout, that extensive growth alone gives none of the three
conclusions (diverging density $N/4$ at a linear tie, nonconverging weights, sub-extensive scales, a rare
tied label with diverging density, uniform convergence without convergence of second derivatives), and
\MT{rem:cg-label-observation} identifies the $N^2$ scale as information about the label: in the binary
examples the unlabeled configuration carries none.

What is $N$, and is this a thermodynamic limit? $N$ indexes a sequence of finite labeled families whose
log masses are assumed to grow linearly; the number of agents is offered as an example. The label set is
fixed while $N$ grows. The theorems are the finite-mixture calculus of phase coexistence, not a
thermodynamic limit of the agent network: there is no infinite-volume measure, and nothing derives
extensive growth or a fixed label set from an agent model (Supplement B instantiates the three regimes
for its lattice law with $N=L^2$ cells, deriving the branch inputs while still assuming a growth
condition on the raw label weights, \SB{cor:rg-thermo-native-label-limits}; that lattice carries no agent
sections). The natural label sets of an $N$-agent
network, such as its partitions, grow with $N$; the theorems apply to a fixed finite menu of coarse
alternatives whose log weights are supplied as extensive.

\section{What is proved, and how solid it is}
\label{comp:ch09-status}

\begin{scorecard}
\MT{def:cg-coarse-channel} & A coarse channel is a kernel on latent coordinates reading neither recognition nor observation, applied in common to the three laws. & \CHOICE & Fitted coarse joints, different channels per law, and observation coarsening are outside it.\\
\MT{prop:cg-markov-category,prop:cg-equivariant-channels} & Kernels compose and push laws functorially; equivariant kernels commute with the action. & \SOLID & Standard (Kleisli category of the Giry monad \citep{Giry1982}).\\
\MT{thm:cg-dpi-equality,cor:cg-pairwise-bayes-recovery,cor:cg-dpi-infinite-equality-warning} & Data processing with its equality case; pairwise Bayes recovery; nothing follows from $+\infty=+\infty$. & \SOLID & Standard in substance (Kullback, Csisz\'ar, commutative Petz). Needs finite divergence.\\
\MT{thm:rg-exact-coarse-vfe} & Evidence and posterior version preserved; $\Fenergy=\Fenergy^c+\Delta$ in the extended reals; $\Delta=0$ iff the discarded conditionals agree, and with finite divergence iff one kernel recovers both. & \SOLID & Standard in substance (chain rule); typing new. One common channel, observation untouched.\\
\MT{thm:cg-evidence-preserving-channel} & Coarse ELBO exceeds fine ELBO by $\Delta$ at equal evidence. & \SOLID & The rise is discarded information, not model improvement.\\
\MT{prop:cg-stochastic-infimum} & Constrained-minimum reading holds for deterministic statistics and fails for noisy channels. & \SOLID & Part (ii) is a two-state counterexample.\\
\MT{thm:rg-pointwise-parent-datum} & Every declared block and parent kernel give a normalized parent triple with the chain rule; evaluator by disintegration; marginals do not fix the law. & \SOLID & Static; block, label, kernel are inputs; a declared evaluator needs a compatibility hypothesis.\\
\MT{prop:cg-hard-identification} & Precomposition is conditioning, can raise relative entropy, is reference dependent on null sets. & \SOLID & Negative: pushforward theorems do not transfer.\\
\MT{prop:cg-message-passing-residual} & Galerkin propagation is exact iff the discarded complement is invariant. & \SOLID & Standard linear algebra.\\
\MT{prop:cg-equivariant-bayes-versions,cor:cg-continuous-version-equivariant} & Equivariant Bayes reverse kernels for finite spaces, compact stabilizers, continuous posteriors. & \SOLID & General noncompact case open.\\
\MT{prop:rg-scale-functoriality} & Functoriality empty on chains, contentful on other posets and rule semigroups; tie-broken majority fails. & \SOLID & A resolution index; no rescaling, no flow.\\
\MT{prop:rg-thermo-finite-labeled-fisher,prop:rg-thermo-finite-missing-information,prop:rg-thermo-account-word} & Label-plus-branch Fisher decomposition, missing information, telescopes. & \SOLID & Standard in substance (total covariance; Louis).\\
\MT{thm:rg-thermo-label-strict,prop:rg-thermo-label-tie-scale,thm:rg-thermo-label-crossover} & Strict winner hides the label; a tie gives $N^2$ information; the $1/N$ window gives a log-sum-exp Hessian. & \COND & Extensive growth, branch convergence, center ratios supplied, not derived; label set fixed as $N$ grows.\\
\MT{prop:rg-thermo-label-tie-counterexamples} & Each conclusion fails without its hypotheses. & \SOLID & Negative.\\
\MT{open:cg-equivariant-versions,open:cg-natural-selector} & Equivariant versions for noncompact stabilizers; a natural partition selector. & \OPEN & See \cref{comp:ch09-open}.\\
\end{scorecard}

\begin{hygiene}
Conditions that matter. The channel must be one kernel for all three laws, blind to the observation and
to recognition; different channels per law, fitted coarse joints, and observation coarsening (even
invertible) change the evidence event or the comparison target and are outside the theorem. The
difference form and the recovery characterization need finite fine relative entropy, which fails on the
singular branch. The constrained-minimum reading needs a deterministic statistic and is stated for
finite spaces. Fisher contraction needs a parameter-independent channel, since a parameter-dependent
one can create information (\MT{thm:hist-record-clock-contraction}). The label asymptotics need
extensive growth, convergence of the tied weights, $C^2$ convergence of branch densities in the strict
and window cases (only $O(N)$ branch information at the tie), center ratios in the window, and a fixed
label set. Equivariant reverse kernels
need a compact stabilizer or a continuous posterior. Hygiene: standard Borel spaces, measurable kernels,
the choice of versions on null sets, fixed-support strata, and $C^3$ log masses.
\end{hygiene}

The chapter has no toy content: its small examples (three-point infinite equality, the two-state noisy
channel, three-point conditioning, six-agent majority, binary tie families) are counterexamples, which
establish negative statements at full generality.

\section{Open ends}
\label{comp:ch09-open}

\MT{open:cg-equivariant-versions} asks for everywhere-defined equivariant Bayes reverse kernels under a
noncompact stabilizer when the posterior has no continuous version, jointly measurable across orbits;
it is settled by an equivariant measurable selection theorem or a counterexample in the stated class.
An elementary remark outside the manuscript bounds its reach for the program's Gaussian models: for a
faithful representation the stabilizer of a nondegenerate Gaussian law is compact, so the Haar-averaged
version of \MT{prop:cg-equivariant-bayes-versions} applies. The open case concerns laws with noncompact
stabilizers, such as degenerate ones or nonfaithful actions.

\MT{open:cg-natural-selector} is the substantive one. It asks for a data-dependent partition, membership
law, or structural channel that is natural under relabeling (relabeling the input relabels the unlabeled
output) and reads gauge data only through invariant scores. No analyzed cost selects one, and the reason
is monotonicity. For fixed recognition the coarse free energy never increases along a composition of
channels, so minimizing it favors the channel to a point, which discards everything and returns
$-\log p(o)$; minimizing $\Delta$ favors the identity; an unpenalized sum of within-block divergences is
minimized by singletons (\MT{rem:cg-structural-vs-marginal}); the Gaussian recognition costs of
Supplement G select trivial endpoints or need an external coefficient, a price per block much like a
chemical potential; and common recovery of a whole experiment needs simultaneous hypotheses. The problem
is settled by a relabeling-natural, gauge-compatible functional that is not monotone under refinement and
has a nondegenerate optimizer, or by an impossibility theorem in a stated class. In Kadanoff--Wilson
renormalization the blocking is supplied by lattice geometry and never selected; here no geometry
supplies it. Parent dynamics, \MT{open:rg-pointwise-parent-dynamics}, is posed in
\mtc{ch:agent-network-rg} and discussed in \cref{comp:ch11}.

\begin{verdictbox}{\DELIVERED}
The chapter says that one coarse channel applied to the joint, the posterior, and a recognition law
preserves the evidence and the posterior version and splits the free energy into the coarse free energy
plus the discarded conditional information; that this specializes to a candidate parent; that
pushforward, precomposition, and Galerkin restriction are different operations; and that a retained label
carries information of order $N^2$ at a tie. Each is proved as stated. The central theorem is standard in
substance (the relative-entropy chain rule and commutative Petz sufficiency, carefully typed), and the
chapter is explicit that exactness buys no closure, no selected channel or block, and no parent dynamics.
Two wordings outside the statements promise more than the statements give. The chapter opening places
the $N^2$ information at a tie among $N$ agents, while the theorems concern a fixed finite label set with
supplied extensive log masses. The introduction (\MT{sec:intro-reading-map}) says this chapter separates
Markov channels, energy restriction, and parameter aggregation; the chapter separates pushforward,
precomposition, and Galerkin restriction, and never treats parameter aggregation.
\end{verdictbox}

\chapter[Gauge-Retained Network State and Fixed Sectors]{Gauge-Retained Network State\\and Fixed Sectors}
\label{comp:ch10}

\section{The question and the answer}
\label{comp:ch10-question}

Grouping agents into a block does not average away their gauge data: the links $\Theta^x_e$ are
independently variable group elements on the edge copies of the interaction multigraph, one family per channel
$x\in\{b,m\}$, and closed walks carry holonomy. \MT{ch:network-state} asks what a coarse description of
a block must keep of these data and which part of the block's state survives exact aggregation. The
answers come in six pieces. A rooted spanning tree turns the block's gauge data into one root-framed
datum (member coordinates carried to the root, the based holonomy, dressed boundary edges) that
transforms under a single root gauge transformation (\MT{prop:rg-rooted-datum}). The configurations
that survive exact aggregation form the kernel of the connection Laplacian, which with positive-definite
weights is isomorphic to the holonomy-fixed vectors at the root (\MT{prop:cg-kernel-holonomy}). The
members' marginal laws admit one parent law exactly when, carried to the root, they coincide in a
holonomy-invariant law (\MT{thm:cg-holonomy-kl-marginal}); compact holonomy gives a closed form, and for
Gaussian parents noncompact holonomy generically allows no parent at all
(\MT{cor:cg-holonomy-noncompact-gaussian}). Cut links compress exactly into rectangular endpoint maps,
and into a single invertible link with new self terms exactly when an explicit matrix is invertible, with
one balanced choice that has no law-level counterpart and no continuous extension to singular cuts. The
frame of a meta-agent is invisible to its members and cannot be computed from their frames by any
symmetric rule that does not depend on a reference point (equivalently, any symmetric left-equivariant
rule) (\MT{prop:cg-frame-sector}). A declared noisy record at one base point distinguishes loop
holonomies (\MT{thm:cg-loop-operational-holonomy}).

When the links are latent generators, as in \cref{comp:lie-links}, the root transports, dressed edges,
and holonomies are random marks with their joint law. The Gaussian-parent obstruction then becomes an
almost-sure statement: for $G=\GL^+(K)$ in its defining representation, a connected block with an
internal loop has infinite holonomy-conditioned Gaussian score under every generator law absolutely
continuous with respect to Lebesgue measure (\MT{cor:link-holonomy-parent-almost-sure}). This concerns
the holonomy-fixed Gaussian parent family. The exact pushed parent law of \cref{comp:ch09} still
exists, and need not belong to that family.

These results feed \cref{comp:ch11,comp:ch12,comp:ch13,comp:ch14}; Supplement G realizes them for
Gaussian families. The scope paragraph fences them: the fixed-sector rank is an admissibility criterion,
not a cost, keeps singletons admissible, and does not require beliefs to agree; the marginal-law score
takes its block and weights as inputs, selects no partition, supplies no kernel, and does not recover
the dependence structure of a joint law; the operational record recovers no link values.

\section{The picture}
\label{comp:ch10-picture}

At one base point the agent network is a lattice gauge system over a zero-dimensional base: sites are
agents, links $\Theta^x_e\in G$ live on edge copies, and a passive gauge transformation acts by
$\Theta^x_e\mapsto g_i\Theta^x_eg_j^{-1}$ for $e:j\to i$. Because each agent also carries a frame, the
links read relative to the agents' own frames (the per-edge twists) are invariant, and the complete
invariant is the tuple of twists up to one simultaneous conjugation, which contains the loop holonomies
up to conjugation (\MT{prop:geo-edge-twist-invariance}); structurally this is the composite link of a
gauge-Higgs model with a group-valued Higgs field. Blocking means choosing a root, carrying every
member's data to it along a spanning tree, and keeping the loops the tree does not use.

\begin{analogy}{exact}
\textbf{Maximal-tree gauge and Wilson loops.} Gauging every tree link to the identity, uniquely up to
one element per block, turns each non-tree link into the holonomy of the loop it closes
(\MT{prop:rg-rooted-datum}). This is the maximal-tree gauge of lattice gauge theory \citep{Creutz1977},
which the manuscript cites, and the retained holonomy homomorphism carries the block's Wilson loops and
more: the based holonomies up to one root conjugation.
\end{analogy}

\begin{analogy}{exact}
\textbf{Zero modes of the covariant lattice Laplacian.} The connection Laplacian is the covariant
kinetic energy of a matter field: each edge penalizes the mismatch between a member's vector and its
neighbor's vector carried across the link. Its zero modes are the covariantly constant configurations,
which exist with a given root value exactly when every loop returns that value to itself, as a charged
field on a ring has a constant mode only for trivial flux. The kernel is the degree-zero sheaf cohomology
of the block \citep{HansenGhrist2019}, and for orthogonal links and scalar weights the operator is the
vector diffusion of \citet{SingerWu2012} (\MT{prop:cg-kernel-holonomy}).
\end{analogy}

To give a block one parent law, carry each member's marginal law to the root. The result depends on the
walk only through a holonomy action, so an honest parent must be holonomy invariant, and for a compact
holonomy group the best invariant parent in forward relative entropy is the transported mixture averaged
over the group: a twirl, with the block's represented holonomy in the role of the symmetry. A planar
rotation loop forces a zero parent mean and an isotropic covariance; a loop that stretches volume admits
no Gaussian parent at all.

Parallel links between two blocks become one effective link only at a price. Their quadratic energy
completes to a single link, the weight-averaged twist, plus a self-energy measuring how much the twists
disagree. Twists $I$ and $-I$ with equal weights average to zero and leave the blocks coupled only
through self terms: the two transports interfere destructively. Requiring the effective link to balance
its self terms across the cut makes it unique; for orthogonal twists and scalar weights it is the
orthogonal polar factor of the mean twist.

\begin{analogy}{structural}
\textbf{Link smearing with projection to the group.} Lattice practice replaces a link by a sum of
parallel transports and projects the sum back to the group, often by polar decomposition. The balanced
link in the orthogonal case is such a projection (\MT{prop:cg-balanced-cut-presentation}, part (iii)).
The projection jumps where the sum is singular, and \MT{open:cg-cut-compression} proves that no
continuous rule avoids this.
\end{analogy}

The frame of a meta-agent is pure gauge for its members, and it cannot be computed as a symmetric,
reference-independent average of their frames.

\begin{analogy}{exact}
\textbf{The phase of the mean of the roots of unity.} The phases $1,\omega,\ldots,\omega^{n-1}$ with
$\omega=e^{2\pi i/n}$ are permuted by multiplication by $\omega$, so a symmetric rotation-covariant average
would have to be fixed by $\omega$, which is impossible. The frame no-go (\MT{prop:cg-frame-sector},
part (ii)) is this argument for any group with an element of order $n$.
\end{analogy}

Finally, the operational record reads, with Gaussian noise, the displacement $(\mathrm{Hol}-I)z$ of a
probe agent's state carried around a loop, as an interferometer reads the holonomy seen by a test
particle. The analogy is structural, and the record kernel is declared to have this form rather than
derived from the agents' dynamics.

\section{Objects and notation}
\label{comp:ch10-objects}

\begin{decoder}
levelwise fiber state $\mathfrak X_\ell$ & all fiber data at depth $\ell$: group, frame torsor, associated fibers, multigraph, links, holonomy groups, cross maps & configuration space of a lattice gauge and matter theory\\
scale intertwiner & homomorphism $\kappa_\ell:G_\ell\to G_{\ell+1}$ with compatible fiber maps & embedding of gauge groups under blocking\\
cross map $\Phi,\widetilde\Phi$ & declared map between belief and model bundles & none\\
tree transport $\Theta^x_{I\leftarrow i}$ & link product along the tree path from $i$ to the root & parallel transporter\\
based holonomy group $\mathrm{Hol}^x_I(r)$; $\mathrm{Hol}^x(p_e)$ & all loop products at the root; the loop closed by one non-tree edge & Wilson-loop data; independent loop\\
block-dressed edge $\widehat\Theta^x_e$ & boundary link with tree transports attached at both ends & link between block roots\\
holonomy-stabilized law & law invariant under every holonomy element & invariant (singlet) state\\
averaging maps $\mathsf C_x,\mathsf S_x$ & weighted mean of transported features and its right inverse & gauge-covariant block averaging\\
connection Laplacian $L_I$; fixed fiber $\ker L_I$, rank $f_I$ & covariant quadratic mismatch energy; its zero modes and their number & covariant lattice Laplacian; zero modes\\
full, partial, no fixed sector & $f_I=K$, $0<f_I<K$, $f_I=0$ & flat, partly flat, fully twisted\\
marginal-law score $\mathfrak D^x_I$ & least weighted forward relative entropy from transported marginals to an invariant parent & none\\
Haar-symmetrized mixture $\bar M$ & mixture of transported marginals averaged over the holonomy group & twirled state\\
exact common recovery & one fixed reverse kernel restores every law of an experiment & sufficient statistic\\
cut blocks $\Lambda_{II},\Lambda_{IJ},\Lambda_{JJ}$; excess $\Delta$ & quadratic coupling of two block roots; self-energy left after completing the square & frustration of parallel links\\
anchored, no-new-anchor, balanced & one link and weight plus new self terms; without self terms; with self terms matched across the link & effective link plus mass terms\\
endpoint maps $S^Q_e,S^K_e$ & factorized edge energy through an edge space & sheaf restriction maps\\
probe $i_0$, record $r$ & agent whose state is carried around a loop; noisy readout & test particle; interferometer\\
\end{decoder}

Holonomy always means the product of independent link variables around a closed walk of the agent
multigraph, neither the holonomy of a base connection nor a topological property of the bundle. Fixed
means holonomy-fixed; sector means a subspace or its rank, not a superselection sector. $\Delta$ is now
a positive semidefinite matrix, and $W_e$ are positive semidefinite forms on the fiber. Represented means
seen through $\rho$; a nonfaithful representation can hide principal holonomy. Balanced names one
symmetry condition, not an optimum. $K$ is the fiber dimension, and $\GL(K)$ is real unless stated.

\section{The formulas, explained}
\label{comp:ch10-formulas}

\subsection{Depth maps and the rooted datum}

A fiber state (\MT{def:rg-geometric-state}) collects the depth-$\ell$ group, frame torsor, associated
fibers, multigraph, links, holonomy groups, and cross maps. A map to the next depth is a homomorphism
$\kappa_\ell$, a compatible torsor map (unique once one frame is sent to one frame), and fiber maps
$q_{\ell,x}$; the induced map of associated fibers exists exactly when
$q_{\ell,x}\circ\widehat\rho_{\ell,x}(g)=\widehat\rho_{\ell+1,x}(\kappa_\ell(g))\circ q_{\ell,x}$
(\MT{prop:rg-scale-intertwiner}), and defect pairs have gauge-invariant equality loci
(\MT{prop:rg-defect-covariance}). This is standard associated-bundle functoriality.

For a block $I$ with root $r_I$ and a spanning tree, the tree transport, the holonomy of a non-tree edge
$e:j\to i$, and the dressed edge of a boundary edge from block $J$ are
\[
\Theta^x_{I\leftarrow i}=\Theta^x_{p_{I,i}},
\qquad
\mathrm{Hol}^x(p_e)=\Theta^x_{I\leftarrow i}\Theta^x_e(\Theta^x_{I\leftarrow j})^{-1},
\qquad
\widehat\Theta^x_e=\Theta^x_{I\leftarrow i}\Theta^x_e(\Theta^x_{J\leftarrow j})^{-1},
\]
and a gauge transformation sends them to $g_{r_I}\Theta^x_{I\leftarrow i}g_i^{-1}$,
$g_{r_I}\mathrm{Hol}^x(p_e)g_{r_I}^{-1}$, and $g_{r_I}\widehat\Theta^x_eg_{r_J}^{-1}$
(\MT{prop:rg-rooted-datum}): after transport only the roots' gauge remains. The holonomy homomorphism
does not depend on the tree; a tree change multiplies transports on the left by holonomy elements. A
tree block has trivial holonomy and pure-gauge internal data; for abelian $G$ each holonomy is gauge
invariant; graph-flat links (\MT{prop:geo-trivializing-criterion}) have trivial holonomy. The manuscript
keeps the raw root-framed datum rather than a gauge quotient, since quotienting each holonomy separately
loses its orientation relative to the root features and boundary legs, and an orbit space would need a
quotient theorem it does not supply. A parent carries holonomy either by full groupoid covariance of all
fine laws, versions, kernel, and evaluator, or by placing the rooted datum in its retained marks
$\mathsf H_A$ with no invariance claim (\MT{prop:rg-parent-holonomy-alternatives}); neither selects the
block.

The natural linear coarse feature is the weighted mean of transported features, with its right
inverse,
\[
(\mathsf C_xz)_I=\sum_{i\in I}w_{Ii}\rho^{\mathbb W}_x(\Theta^x_{I\leftarrow i})z_i,
\qquad
(\mathsf S_x\hat z)_i=\rho^{\mathbb W}_x(\Theta^x_{I\leftarrow i})^{-1}\hat z_I.
\]
It is gauge covariant but tree dependent (\MT{prop:rg-linear-cross-scale}): for two vertices joined by edge copies with
$\mathrm{SO}(2)$ links $I$ and $R_\vartheta$, equal weights, and equal features $e_1$, the two trees give
averaged features of norms $1$ and $\lvert\cos(\vartheta/2)\rvert$, both invariant under root gauge
transformations. Averaging is tree independent on holonomy-fixed features and, for trivial represented
holonomy, on all features. The tree is a gauge
fixing for the retained datum but part of the model for the averaged feature, as gauge-covariant smearing
of matter fields depends on the paths unless the field is covariantly constant.

\subsection{The holonomy-conditioned marginal law}

Carry each marginal to the root, $\bar q^x_i=\widehat\rho_x(\Theta^x_{p_i})q^x_i$, and score the block by
(\MTeq{eq:cg-holonomy-kl-score})
\[
\mathfrak D^x_I=\inf_{Q\in\mathscr Q^x_{I,\mathrm{fix}}(r)}\sum_{i\in I}w_i\KL(\bar q^x_i\Vert Q),
\qquad
\mathscr Q^x_{I,\mathrm{fix}}(r)=\bigl\{Q\in\mathscr M^x_r:\widehat\rho_x(\mathsf k)Q=Q\text{ for all }\mathsf k\in\mathrm{Hol}^x_I(r)\bigr\},
\]
with $\inf\varnothing=+\infty$ and $\mathscr M^x_r$ a declared class of candidate parents.
\MT{thm:cg-holonomy-kl-marginal} proves that the score depends on neither walks, root, nor frames, and
that an attained zero means all transported marginals equal one invariant law, which then defines a
holonomy-stabilized parallel family of member laws; an unattained zero means only approximability. For
compact represented holonomy with Haar probability $d\mathsf k$, put $\bar P=\sum_iw_i\bar q^x_i$ and
$\bar M(B)=\int\bar P(\mathsf k^{-1}B)d\mathsf k$; every invariant $Q$ then satisfies
\[
\sum_{i\in I}w_i\KL(\bar q^x_i\Vert Q)=\sum_{i\in I}w_i\KL(\bar q^x_i\Vert\bar M)+\KL(\bar M\Vert Q),
\]
so the minimizer is $\bar M$ for unrestricted parents and its moment-matched Gaussian for Gaussian ones
(\MT{prop:cg-holonomy-compact-closed-form}). That the forward-relative-entropy barycenter of laws is
their mixture is standard; the symmetrization and the gauge typing are the chapter's.

On the running three-agent loop with $\mathrm{SO}(2)$ links, loop rotation $\phi$ with $\sin\phi\neq0$,
and transported marginals $\mathcal N(m_i,\sigma^2I_2)$, the only invariant Gaussians are
$\mathcal N(0,sI_2)$ and (\MT{rem:cg-holonomy-score-loop})
\[
\mathfrak D^b_I=\log\Bigl(1+\frac{1}{2\sigma^2}\sum_iw_i\lvert m_i\rvert^2\Bigr),
\qquad
s_\star=\sigma^2+\tfrac12\sum_iw_i\lvert m_i\rvert^2:
\]
the parent forgets direction and absorbs the members' means into its variance. Frame-aligned links give
$\tfrac12\log\det(I_2+\operatorname{Cov}_w(m)/\sigma^2)$, never larger, and a link of determinant $2$
gives $+\infty$. The last case is general (\MT{cor:cg-holonomy-noncompact-gaussian}): with Gaussian
parents and linear holonomy on $\R^K$, an invariant parent exists iff the represented holonomy group is
conjugate into $\mathrm O(K)$; one loop holonomy with an eigenvalue off the unit circle, or a
nondiagonalizable one, gives $+\infty$ for all marginals; and for a block containing a loop with
independent $\GL(K)$ links, $\lvert\det\rvert=1$ on that loop is Lebesgue-null, so the score is $+\infty$
almost everywhere and distinguishes configurations only after the links are restricted to a group
conjugate into $\mathrm O(K)$ or a parent family with invariant members is declared.

\subsection{The obstruction under a law of generators}

For the latent links, the determinant test is especially transparent
(\MT{cor:link-holonomy-parent-almost-sure}):
\begin{keyequation}{The determinant test for a fixed Gaussian parent}
\[
\begin{aligned}
\det\mathrm{Hol}^x(p)
&=\exp\Bigl(\sum_a\Tr\mathring\omega^x_{e_a}\Bigr),\\
\mathfrak D^x_I=+\infty
&\quad\text{if }\sum_a\Tr\mathring\omega^x_{e_a}\ne0.
\end{aligned}
\]
\end{keyequation}
Here $G=\GL^+(K)$ acts on $\R^K$ by the defining representation, and $p$ lies in a connected
block whose internal multigraph is not a forest. A loop copy, two parallel copies, or a cycle through
three or more agents supplies a closed walk with nonzero net traversal of a reversal pair. For such a
walk the trace sum is a nonzero linear functional of the generators, so its zero set is a hyperplane.
Every generator law absolutely continuous with respect to Lebesgue measure gives that hyperplane
probability zero. This includes the Gaussian-base generator prior with positive-definite covariances
and every recognition law absolutely continuous with respect to it. The score is therefore infinite
almost surely, and its expectation is infinite. This is a general result at every admitted finite
network size, rather than an obstruction inferred from a small example.

The hypotheses decide the conclusion. A singular generator law can charge the hyperplane: the point
mass at zero gives trivial holonomy and a finite score for nondegenerate Gaussian marginals. Removing
the trace direction changes the situation. For $\mathrm{SL}(2,\R)$ with cycle rank one, the corollary
exhibits a nonempty open set of generator families whose holonomy group is conjugate into
$\mathrm{SO}(2)$; any everywhere-positive generator density gives this set positive probability.
Thus noncompactness alone is not the asserted obstruction. Nor does the empty fixed Gaussian family
prevent normalized coarse laws, parents that retain the generators, or other declared parent families.
Conditionally on the generators, the held-link results apply; mixing those conditionals can lose
Gaussianity (\MT{prop:link-conditional-transfer}; see \cref{comp:gaussian-latent-links}).

Two fences. A zero score is not common recovery: it reads only marginals, and joint laws with equal
marginals cannot both be recovered through a channel that reads only marginals
(\MT{rem:cg-marginal-not-joint}). A cross map commuting with the links carries invariant belief parents
to invariant model parents (\MT{cor:cg-holonomy-cross-morphism}) without identifying the bundles.

\subsection{Structural fixed sectors}

The connection Laplacian of a block is (\MTeq{eq:cg-connection-laplacian-energy})
\[
z^\top L_Iz=\sum_{e:j\to i}(z_i-\Theta_ez_j)^\top W_e(z_i-\Theta_ez_j),
\]
with represented links $\Theta_e$ and edge forms $W_e\succeq0$. How does its kernel relate to
$\operatorname{Fix}(\mathrm{Hol}_I)$? With positive-definite weights, evaluation at the root is an
isomorphism from $\ker L_I$ onto the fixed space of the represented holonomy on each component, so $f_I$
is its dimension, and $f_I=K$ exactly when the links are graph flat on the block
(\MT{prop:cg-kernel-holonomy}). With merely semidefinite weights the kernel is set by the edgewise
constraints $W_e^{1/2}(z_i-\Theta_ez_j)=0$, and holonomy can hide in weight-null directions: a triangle
with loop product $2I$, which fixes nothing, still has a two-dimensional kernel.

Aggregation onto fixed sections is Galerkin compression: with $\iota_{\mathcal P}$ block diagonal and
spanning each $\ker L_I$, the coarse form $\iota_{\mathcal P}^\top\Lambda\iota_{\mathcal P}$ sees no internal
edges and is positive definite exactly when no fixed section is a fine zero mode
(\MT{prop:cg-fixed-section-compression}); nested sections compose (\MT{prop:cg-nested-sections-compose}).
This structural rank is independent of the marginal-law score (\MT{rem:cg-structural-vs-marginal}): a
two-vertex tree has trivial holonomy and arbitrarily distant transported marginals, while a triangle
with nontrivial $\mathrm{SO}(3)$ holonomy has $f_I=1<3$ and zero divergence for isotropic Gaussians.
Pairwise transported divergence also depends on the walk unless the laws are stabilized: with
$\mathsf k=\operatorname{diag}(1,-1,-1)$ and two copies of $\mathcal N(e_2,I_3)$, two walks give $0$ and $2$.

\subsection{Cut links}

For blocks $I,J$ in root coordinates with dressed edges $\widehat\Theta_e$ and weights $W_e\succ0$ on
the cut copies, put $\Lambda_{II}=\sum_eW_e$, $\Lambda_{IJ}=\sum_eW_e\widehat\Theta_e$, and
$\Lambda_{JJ}=\sum_e\widehat\Theta_e^\top W_e\widehat\Theta_e$. Completing the square
(\MT{lem:cg-cut-completed-square}, \MTeq{eq:cg-cut-excess}),
\[
\mathcal E_{\mathrm{cut}}=(z_I-\bar\Theta z_J)^\top\Lambda_{II}(z_I-\bar\Theta z_J)+z_J^\top\Delta z_J,
\qquad
\Delta=\sum_e(\widehat\Theta_e-\bar\Theta)^\top W_e(\widehat\Theta_e-\bar\Theta)\succeq0,
\]
with $\bar\Theta=\Lambda_{II}^{-1}\Lambda_{IJ}$ the weight-averaged twist. The excess is the weighted
spread of the twists, zero exactly when all dressed edges agree. $\bar\Theta$ transforms like a dressed
edge but can be singular ($I$ and $-I$ average to $0$): a mean of group elements need not be one. One
invertible link with new self terms at both ends (an anchored presentation) exists exactly when
$\det\Lambda_{IJ}\neq0$, positive for a $\GL^+$ link; the admissible weights form the Loewner interval
$\Lambda_{IJ}\Lambda_{JJ}^{-1}\Lambda_{IJ}^\top\preceq W\preceq\Lambda_{II}$, and the presentation is
unique exactly when $\Delta=0$, the only case without new self terms. This is Schur-complement algebra.

Requiring the self term at $J$, carried along the link, to equal the one at $I$, equivalently
$\Theta^\top\Lambda_{II}\Theta=\Lambda_{JJ}$, selects one presentation
(\MT{prop:cg-balanced-cut-presentation}): its weight is the matrix geometric mean \citep{Bhatia2007} of the
interval endpoints,
\[
W_{\mathrm{bal}}=W_{\mathrm{lo}}^{1/2}\bigl(W_{\mathrm{lo}}^{-1/2}W_{\mathrm{hi}}W_{\mathrm{lo}}^{-1/2}\bigr)^{1/2}W_{\mathrm{lo}}^{1/2},
\qquad
\Theta_{\mathrm{bal}}=W_{\mathrm{bal}}^{-1}\Lambda_{IJ},
\]
gauge covariant, the same read from either end, and for scalar weights and orthogonal links the
polar factor of $\bar\Theta$. Two negative results bound it. For Gaussian laws, with
$\sum_ew_e\KL(q_I\Vert(\widehat\Theta_e)_\#q_J)$ in place of the quadratic energy, a cut with excess has
no presentation as one divergence to an image of $q_J$ plus self terms reading one law each whose mean
part is the balanced link, for any law class at $J$ closed under covariance rescaling
(\MT{prop:cg-no-balanced-law-presentation}); the self term at $I$ would have to carry
$\Lambda_{II}-W_{\mathrm{bal}}$, which depends on the covariance at $J$. For twists $I$ and $2I$ the
excess is $\frac{w_1w_2}{2(4w_1+w_2)}\Sigma_J^{-1}$. And \MT{open:cg-cut-compression}, despite its label,
is an established no-go: two identity-weighted $\mathrm{SO}(2)$ twists rotating by $\pm\arccos t$ give
$\Lambda_{IJ}=2tI_2$ and $\Theta_{\mathrm{bal}}=\operatorname{sign}(t)I_2$, which jumps at $t=0$.

What survives exactly is the rectangular closure (\MT{thm:cg-rectangular-endpoint-closure}): factor
$W_e=(S^Q_e)^\top S^Q_e$, put $S^K_e=S^Q_e\Theta_e$, and the edge energy
$\lVert S^Q_ez_i-S^K_ez_j\rVert^2$ compresses exactly under block injections, associatively, with parallel
copies merged into one edge space of dimension
$\operatorname{rank}\sum_e\delta_e^\top\delta_e\leq\dim\mathbb V_I+\dim\mathbb V_J$. This is the
cellular-sheaf Laplacian category of \citet{HansenGhrist2019}, and the price of exactness is that the
coarse edge is a pair of rectangular maps, not an element of the structure group.

\subsection{The frame sector}

What each member sees of a coarse datum $(u_I,q_I)$ is the coordinate of $[u_I,q_I]$ in its own frame,
$\mathcal N(U_i\bar m,U_i\bar SU_i^\top)$ for the Gaussian fiber of $\GL^+(K)$ (\MT{prop:cg-frame-sector},
part (i)). Data indistinguishable to every member form one free orbit
$\{(u_I\cdot g,\widehat\rho_x(g)^{-1}q_I)\}$, of dimension $K^2$ for $\GL^+(K)$: the parent frame is pure
gauge for the members. Part (ii): if $G$ contains $v\neq e$ with $v^n=e$, no $F:G^n\to G$ is both left
equivariant, $F(gg_1,\ldots,gg_n)=gF(g_1,\ldots,g_n)$, and symmetric, because left multiplication by $v$
permutes $(e,v,\ldots,v^{n-1})$. $\GL^+(K)$ with $K\geq2$ has such elements for every $n\geq2$; on $\R^d$
the arithmetic mean works. A rule $u_I=\sigma\cdot F(a_1,\ldots,a_n)$ is independent of the reference
$\sigma$ exactly when $F$ is left equivariant, so no symmetric rule that is independent of the reference
point computes a meta-agent's frame from its members'; a symmetric rule that fixes a reference point
exists. The theory does not need one, since the parent frame is invisible; the rooted datum uses
a root, a choice that is equivariant but not permutation symmetric. The no-go concerns maps on all of
$G^n$ and does not exclude rules near consensus: as a standard fact outside the manuscript, Riemannian
centers of mass for a left-invariant metric \citep{Karcher1977} give symmetric equivariant averages of
sufficiently concentrated configurations. The obstruction is global.

\subsection{The operational record}

Fix a probe $i_0$, a closed walk at it, and labels $\omega$ with prior masses $\psi_\omega$, each
prescribing link values on the walk and hence a represented loop holonomy
$\mathrm{Hol}^{\mathrm{pr}}_\omega$ on the probe's belief and model coordinates. The record joint is
\[
P^{\mathrm{rec}}(\omega,dy,do,dr)=\psi_\omega\mathbb P(dy,do)
\mathcal N\bigl(r;(\mathrm{Hol}^{\mathrm{pr}}_\omega-I)z_{i_0}(y),\Sigma_{\mathrm r}\bigr)dr.
\]
\MT{thm:cg-loop-operational-holonomy} proves that it is normalized and that given $\omega$ and $o$ the
record is $(\mathrm{Hol}^{\mathrm{pr}}_\omega-I)Z$ plus independent noise, so two labels share a record
law exactly when their images of $Z$ agree in law; that a trivial and an invertible $\mathrm{Hol}-I$ give
different laws unless the probe's posterior state is the point mass at zero, with
$r^\top\Sigma_{\mathrm r}^{-1}r$ having mean $\dim\mathbb V_{\mathrm{pr}}$ in the first case and strictly more
in the second; and that the posterior over $\omega$ and the law of that statistic are gauge invariant.

Can an experiment at one base point detect network-loop holonomy? At the level of laws, yes, if the
experiment is modeled by a record kernel reading the loop transport of the probe's state: the theorem
shows such a record is a legitimate, gauge-covariant part of the joint whose law separates holonomies
unless the probe is pinned at zero. It does not show that any record kernel of the program's interaction
model has this form. The latent law is independent of $\omega$, so here the links do not influence the
agents' states. Holonomy in the kernel of $\rho$, or holonomy fixing the support of the probe's state, is
invisible. No link value is recovered, and no sample-size or power statement is made.

\section{What is proved, and how solid it is}
\label{comp:ch10-status}

\begin{scorecard}
\MT{prop:rg-scale-intertwiner,prop:rg-defect-covariance} & Depth maps exist iff fiber maps intertwine; defect equality loci are gauge invariant. & \SOLID & Standard associated-bundle functoriality.\\
\MT{def:rg-rooted-tree-data} & A block keeps the raw root-framed datum, not a gauge quotient. & \CHOICE & Stated after \MT{prop:rg-rooted-datum}; a design decision, not a necessity theorem.\\
\MT{prop:rg-rooted-datum,prop:rg-parent-holonomy-alternatives} & Gauge laws of transports, holonomy, dressed edges; tree links gauge to the identity; two ways a parent carries holonomy. & \SOLID & Standard in substance (maximal-tree gauge \citep{Creutz1977}); neither alternative selects the block.\\
\MT{prop:rg-linear-cross-scale} & Linear averaging of transported features is covariant but tree dependent. & \SOLID & Tree independent on holonomy-fixed features and, for trivial represented holonomy, on all features.\\
\MT{thm:cg-holonomy-kl-marginal} & Marginal-law score is walk, root, and frame independent; attained zero iff one invariant law. & \SOLID & Marginals only; no kernel, no joint recovery.\\
\MT{prop:cg-holonomy-compact-closed-form} & Compact holonomy: minimizer is the Haar-symmetrized mixture or its moment-matched Gaussian. & \COND & Compact represented holonomy: holds for $\mathrm{SU}(N)$ links, fails for $\GL(K)$ links.\\
\MT{cor:cg-holonomy-noncompact-gaussian} & Gaussian parents exist iff holonomy is conjugate into $\mathrm O(K)$; generic $\GL(K)$ loops give infinite score. & \SOLID & Negative for the noncompact target.\\
\MT{cor:link-holonomy-parent-almost-sure} & Latent $\GL^+(K)$ generators give infinite fixed-Gaussian score almost surely on a connected nonforest block. & \SOLID & Defining representation and a Lebesgue-absolutely-continuous law; not a no-go for pushed parents.\\
\MT{rem:cg-holonomy-score-loop} & Three-agent $\mathrm{SO}(2)$ loop: explicit score. & \TOY & Illustrates the compact closed form.\\
\MT{prop:cg-kernel-holonomy,prop:cg-fixed-section-compression,prop:cg-nested-sections-compose} & Fixed fiber is $\ker L_I$, the holonomy-fixed space for positive-definite weights; compression and nesting exact. & \SOLID & Standard in substance (sheaf cohomology). Semidefinite weights can hide holonomy.\\
\MT{rem:cg-structural-vs-marginal} & Structural rank and marginal-law occupation are independent; singletons minimize unpenalized divergence sums. & \SOLID & Witnesses for negative statements.\\
\MT{lem:cg-cut-completed-square,prop:cg-balanced-cut-presentation} & Completed square and excess; one-link presentations iff $\det\Lambda_{IJ}\neq0$, unique iff no excess; unique balanced one. & \SOLID & Schur complements and geometric means; balance is a chosen condition.\\
\MT{prop:cg-no-balanced-law-presentation,open:cg-cut-compression} & No balanced law-level presentation with excess; no continuous extension to singular cuts. & \SOLID & Negative; the second despite its label prefix.\\
\MT{thm:cg-rectangular-endpoint-closure} & Exact compression of cuts into rectangular endpoint maps. & \SOLID & The coarse edge is not a group element.\\
\MT{prop:cg-frame-sector} & Coarse frame is a free orbit invisible to members; no symmetric left-equivariant $G^n\to G$ under order-$n$ torsion. & \SOLID & Excludes global rules only.\\
\MT{thm:cg-loop-operational-holonomy} & A declared noisy record of the loop displacement distinguishes holonomies unless the probe is pinned at zero. & \SOLID & Record declared; latent law ignores the label; no link values or power.\\
\end{scorecard}

\begin{hygiene}
Conditions that matter. Compactness of the represented holonomy decides whether the marginal-law parent
has a closed form and, for Gaussian parents, whether it exists; it fails for generic $\GL(K)$ links. The
parent family matters: the Gaussian family forces conjugacy into $\mathrm O(K)$. Positive-definite versus
semidefinite edge weights decides whether the fixed fiber is the holonomy-fixed space. A nonfaithful
representation hides holonomy from the fixed sector and the record. One-link cut compression needs
$\det\Lambda_{IJ}\neq0$. Composing averaging maps across depths needs nested trees and multiplied
weights. The record theorem needs its declared kernel and a label-independent latent law; the covariance
branch of the parent alternatives needs covariance of every law, version, kernel, and evaluator.
Attainment of the infimum matters too: without it a zero marginal-law score gives only approximability
by invariant parents, not one invariant law. Hygiene: bimeasurable actions, metrizable compact groups,
Polish fibers, and finite second moments.
For the latent almost-sure obstruction, full Lebesgue absolute continuity of the generator law,
$\GL^+(K)$ in its defining representation, and a nonforest internal multigraph are substantive.
A generator law confined to a lower-dimensional subalgebra or to a point is outside that conclusion.
\end{hygiene}

\section{Toy versus general, and the noncompact target}
\label{comp:ch10-toy}

The one result proved on a small instance and not as a counterexample is the three-agent
$\mathrm{SO}(2)$ loop of \MT{rem:cg-holonomy-score-loop}. Its mechanism, that nontrivial compact holonomy
forces the parent mean into the fixed space and inflates the parent covariance, is general by the compact
closed form; its formulas are special to plane rotations. The chapter's other small configurations are
witnesses for general negative statements.

For the program's target of 512 or more agents with links in $\GL(K,\R)$, $\GL(K,\mathbb C)$, or
$\mathrm{SU}(N)$, the statements sort by whether they use compactness. The rooted datum, tree dependence,
the kernel-holonomy isomorphism, fixed-section compression, the cut calculus, the frame sector, and the
operational record use only group structure and linear algebra, and hold for all three groups with
complex links realified. The torsion hypothesis of the frame no-go holds for all three
($e^{2\pi i/n}I\in\GL(K,\mathbb C)$ and $\operatorname{diag}(e^{2\pi i/n},e^{-2\pi i/n},1,\ldots,1)\in
\mathrm{SU}(N)$ for $N\geq2$), a direct check of the stated hypothesis. The marginal-law parent is where
compactness decides. For $\mathrm{SU}(N)$ links the holonomy lies in a compact group, the closed form
applies (a relatively compact holonomy group has the fixed laws of its closure), and Gaussian
constituents give a finite score. For $\GL(K,\R)$ links with Gaussian parents the score is infinite for
almost every configuration on every block containing a loop; under the default complete interaction
graph every block of three or more agents contains a triangle, so for generic independent links a finite
Gaussian score is available only on blocks whose internal multigraph is a forest, unless links or parent
family are restricted. For $\GL(K,\mathbb C)$ the conjugacy criterion applies configuration by
configuration, but the almost-everywhere statement is proved for real $\GL(K)$ only. Whether link
dynamics drive holonomies into compact conjugacy classes is not addressed. By an elementary remark
outside the manuscript, in the defining representation (realified for $\mathrm{SU}(N)$) a generic real
matrix and a generic element of $\mathrm{SU}(N)$ on $\R^{2N}$ have no eigenvalue one, so with
positive-definite weights a block containing a loop with generic links has $f_I=0$ and is retained only
through the rooted datum. This is specific to that representation: \MT{def:cg-connection-laplacian}
allows any linear representation, and one with a trivial summand, such as the adjoint representation of
$\mathrm{SU}(N)$, keeps $f_I>0$ at every link configuration. The balanced cut link lies in $\GL(\mathbb V)$,
orthogonal in the scalar-weight orthogonal case; no statement keeps it in a smaller group containing the
twists.

In the latent exponential model the real-group almost-sure statement is now explicit for
$\GL^+(K)$, including loop and parallel copies, under any Lebesgue-absolutely-continuous generator
law (\MT{cor:link-holonomy-parent-almost-sure}). At every positive width of a nondegenerate Gaussian
base prior, concentrating near zero does not remove this null-set obstruction. The exactly held
identity links have admissible Gaussian parents, so the limit is singular for this particular score;
the held-link free-energy limit of \MT{thm:link-held-limit} concerns a different quantity and does
not assert continuity of the holonomy-conditioned score.

\section{Open ends}
\label{comp:ch10-open}

The chapter poses no open problem of its own; its one open-problem label names an established negative
proposition. It feeds the open questions of \cref{comp:ch09}: the natural selector
(\MT{open:cg-natural-selector}), since the fixed-sector rank keeps singletons admissible and the
marginal-law score takes the block as input; equivariant reverse kernels beyond compact stabilizers
(\MT{open:cg-equivariant-versions}); and parent dynamics (\MT{open:rg-pointwise-parent-dynamics}). Two
gaps are named without being posed: an orbit space of rooted data would need a quotient theorem, and no
statement connects the declared record to the record kernels of the interaction model.

\begin{verdictbox}{\DELIVERED}
The chapter says it determines what a coarse description of a block retains of its gauge data and which
part of its state survives exact aggregation, and its opening lists the rooted datum, the kernel-holonomy
fixed sector, the holonomy-conditioned parent law, the cut compressions with their balanced choice and two
obstructions, the frame no-go, and a record sensitive to loop holonomy. Each is proved at the stated
generality. Much is standard in substance: maximal-tree gauge fixing, the kernel of a sheaf Laplacian, the
mixture as forward-relative-entropy barycenter, Schur complements, the matrix geometric mean. The solid
new content is mostly negative: no law-level balanced presentation, no continuous extension of balanced
links, no symmetric reference-independent frame average, no Gaussian parent under noncompact holonomy. Three limits of reach
remain. What a block must retain is answered by a choice (the raw root-framed datum), not a necessity
theorem. The record distinguishes holonomies because it is declared to read them. For the $\GL(K,\R)$
target, the holonomy-conditioned Gaussian parent exists, for almost every link configuration, only on
blocks without internal loops under the stated absolute-continuity assumptions; for latent $\GL^+(K)$
generators the obstruction holds almost surely under every Lebesgue-absolutely-continuous law. Its
positive fixed-Gaussian parent results reach that target only under a restriction of the links or of
the parent family. Exact pushed parent laws and the retained random root-framed datum remain
available throughout.
\end{verdictbox}

\chapter{Exact Agent-Network Coarse-Graining}
\label{comp:ch11}

\section{The question and the answer}
\label{comp:ch11-question}

The technical chapter (\MT{ch:agent-network-rg}) applies one coarse channel, a Markov kernel from
fine to coarse configurations (\MT{def:cg-coarse-channel}), to the interacting agent network at the
base point $c_*$ and asks what the step does to the network's energy function and couplings, and
what must be carried along so that nothing is lost. The coarse effective action is minus the
logarithm of a conditional partition function: the fine Boltzmann weight averaged over the fine
configurations compatible with each coarse configuration (\MT{thm:rg-effective-action}). Written in
interaction coordinates, one component for every set of coupled agents, the step is an exact
nonlinear map from fine couplings to coarse couplings, and the derivative of any composite of such
steps along an exact trajectory has norm at most $3^{\lvert V_n\rvert}-1$, a constant fixed by the number of target agents and
independent of the number of steps (\MT{prop:rg-interaction-native-growth-bound}). What must be
kept is collected in a table (\MT{tab:rg-retained-state}): induced multi-agent couplings, the joint
law of fine and coarse variables, the joint law of attention events rather than attention rows, the
gauge marks of each block, and the path law with its memory.

Has this resolved the renormalization of the agent network at one base point? Not in the sense a
physicist gives the word, and the chapter does not say that it has. Exact and proved are the
effective action as minus the log density of the pushed evidence measure relative to the pushed
reference, for every channel; the
exact step on couplings, together with an identity for the residual of any declared truncation; the
finite-step multipliers (the linearized exact flow) and their depth-independent bound; and the
coarse path law with an exact memory recurrence, filter, and finite-memory cost. Absent are four
things a physicist would expect. There is no closed truncation: apart from sectors produced by
independence, symmetry, or couplings already carried by the coarse variables, no proper subspace of
couplings is shown to be mapped into itself, and no truncation residual is estimated
(\MT{open:rg-invariant-retained-subspace}). There is no flow on one fixed set of couplings: across a
step that reduces the number of configurations the coupling spaces have different dimensions and
admit no identification (\MT{prop:rg-interaction-finite-network-obstruction}), and the main
manuscript constructs no rescaling, fixed point, or beta function (\MT{sec:intro-conventions}).
There is no thermodynamic limit: every statement concerns a finite network, and infinite networks
are deferred to Supplement B, where the limit questions are open (\SB{open:rg-two-index-limits}).
And there is no parent dynamics: the parent exists at every instant as an exact pushforward
(\MT{thm:rg-pointwise-parent-datum}), but whether its evolution closes on itself is
\MT{open:rg-pointwise-parent-dynamics}, the defining open problem of the current phase.

The chapter uses the exact coarse free energy of \cref{comp:ch09} (\MT{thm:rg-exact-coarse-vfe}) and
the root and holonomy marks of \cref{comp:ch10}, and it feeds \cref{comp:ch12} and the formation
witnesses of \cref{comp:ch13}, whose parents are exact marked parents in the present sense.

\section{The picture}
\label{comp:ch11-picture}

Take $N$ agents at $c_*$ with states $y_i$ (belief samples, model presentations, labels) and a joint
law equal to a reference probability times a record likelihood. Read the reference $\varrho$ as the a
priori measure, $L_o=e^{-H_o}$ as a Boltzmann weight at unit temperature, $\mathsf m_o=e^{-H_o}\varrho$
as the Gibbs weight, and its mass $M_o$ as the partition function; in Bayesian terms $M_o$ is the
evidence and the normalized $\bar{\mathsf m}_o$ the posterior. A Kadanoff block map is a
deterministic channel sending each block to one block variable; a noisy block map is a stochastic
channel; both are admitted, and random or soft blocks enter through the structural channel of
\cref{comp:ch12}. The exact coarse theory is the pushforward of every law. The chapter says what
that pushforward looks like as an effective Hamiltonian with couplings, and what a reduced
description cannot discard.

\begin{analogy}{exact}
The conditional-partition formula is the position-space renormalization formula with a weight
function, in the block-spin tradition of \citet{kadanoff1966scaling} reviewed by
\citet{wilson1974renormalization}: $e^{-H'(s')}=\sum_sP(s'\given s)e^{-H(s)}$ up to normalization.
The a priori measure is explicit, so fiber-size (entropy) terms such as those of a majority rule are
carried by the pushed reference. Lengths and fields are not rescaled and nothing is truncated; the
identification is exact at the level of the formula and brings no fixed-point theory with it.
\end{analogy}

\begin{analogy}{exact}
The two-agent example of \MT{rem:rg-two-agent-example} is decimation: summing over spin $2$ of a
pair with fields $\xi_1,\xi_2$ and coupling $\xi_{12}$ gives the coarse field
$\xi_1'=\xi_1+\operatorname{artanh}(\tanh\xi_2\tanh\xi_{12})$, so the pair coupling survives only as a
field and only through the polarization $\tanh\xi_2$ of the eliminated spin. With three neighbors the
same mechanism produces a three-spin coupling (\MT{prop:rg-hyperedge-generation}), as decimation of
the two-dimensional Ising model produces four-spin couplings. This is why exact coarse-graining
leaves the pairwise world.
\end{analogy}

\section{Objects and notation}
\label{comp:ch11-notation}

\begin{decoder}
depth $\ell$, $V_\ell$ & Coarse-graining level and its finite agent set & Blocking step; not time\\
$C_\ell$, $C$ & Coarse channel & Block-spin weight function $P(s'\given s)$\\
$\varrho$, $\mathsf m_o$, $M_o$ & Reference probability; evidence measure $e^{-H_o}\varrho$; its mass
& A priori measure; Gibbs weight; partition function\\
$H_o$, $H^c_o$ & Action (minus log density relative to the reference); coarse action & Hamiltonian
at unit temperature; block Hamiltonian\\
$C^\dagger$, $\mathsf U$ & Reverse (Bayes) kernel of fine given coarse; its conditional expectation &
Conditional Gibbs measure at fixed block spins\\
$\mathcal A(\psi)$, $\mathsf U^\psi$ & Exact change of the coarse action under a bounded change $\psi$
of the fine action; tilted reverse expectation & Change of the effective Hamiltonian; perturbed
conditional ensemble\\
$\lVert\cdot\rVert_{\mathrm{osc}}$, $\delta$ & Half the range of a function; largest total-variation
distance between reverse conditionals & Sup norm of an energy modulo constants; Dobrushin
coefficient\\
$\nu_\ell$, $\mathsf P_{\ell,B}$ & Product reference law; projector onto the pure $B$-body part &
Noninteracting a priori measure; projection onto $k$-body terms\\
$\mathcal G_\ell$, $\mathsf g=(\mathsf g_B)$ & One bounded component per nonempty set of agents,
normed by $\sum_B\lVert\mathsf g_B\rVert_\infty$ & All $k$-body couplings of a finite system\\
$\operatorname{asm}$, $\operatorname{hd}$ & Couplings to action; action to couplings & Writing and
reading off a Hamiltonian\\
$\mathcal T_\ell$ & Exact interaction step & Exact coupling recursion, untruncated, unrescaled\\
$D_\ell$, $D_{n\leftarrow\ell}$ & Derivative of one step on an exact orbit; ordered product &
Linearized RG matrix along a trajectory\\
$\mathsf P^{\mathrm{ret}}_\ell$, $r^{\mathcal G}$ & Declared retained projection; its residual &
Truncation of coupling space; truncation error\\
hyperedge & Component on three or more agents & Multi-spin coupling\\
posterior bridge, bridge kernels & Joint posterior law of fine and coarse variables; its pair
conditionals & Joint law of spins and block spins\\
$\eta$, $\alpha$, $\beta$ & Event law over edge copies; receiver occupancy; attention row & Who
attends to whom; its receiver marginal; the conditional\\
$D^{(k)}_N$, $\mathrm h^{\mathrm{ex}}_k$ & Cost of a memory-$k$ path law; excess conditional entropy
& Truncation cost of a $k$-th order Markov model; excess entropy\\
\end{decoder}

Some glyphs mislead. $M_o$ is an evidence, not a magnetization, and $\mathcal G_\ell$ is a coupling
space, not a group. Three objects share a $D$: the multipliers $D_\ell$, the edge divergences $D_e$,
and the memory defects $D^{(k)}_N$. $H$ is an action early in the chapter and an entropy in the memory
theorems. The dagger in $C^\dagger$ is Bayes reversal (the $L^2$ adjoint of the lift), not complex
conjugation. Here $z$ is the coarse sample of one arrow, whereas in \cref{comp:ch12} $Z_\ell$ is the
coordinate block at depth $\ell$. ``Exact'' means without approximation, not solvable in closed form;
``closure'' means that the kept state suffices for the next exact step; a ``multiplier'' is a
derivative and a ``cocycle'' an ordered product of derivatives along an orbit, not a cohomology
class. The main manuscript avoids the words relevant, fixed point, and beta function by design.

\section{The formulas, explained}
\label{comp:ch11-formulas}

\subsection{The effective action is a conditional partition function}

For one arrow $C:\mathsf Y\rightsquigarrow\mathsf Z$, with $C^\dagger_\varrho(z,dy)$ the law of the
fine state given the coarse state under the reference joint law $\varrho(dy)C(y,dz)$,
\[
\exp[-H^c_o(z)]=\int\exp[-H_o(y)]C^\dagger_\varrho(z,dy),
\qquad
M_o=\int\exp[-H^c_o(z)]\varrho^c(dz),
\qquad\varrho^c=\varrho C
\]
(\MT{thm:rg-effective-action}). The coarse Boltzmann weight at $z$ is the a priori average of the
fine weight over the configurations the channel can send to $z$, and the partition function is the
same at both scales. For the identity channel $H^c_o=H_o$; for a channel to a point the action
becomes the number $-\log M_o$; for a deterministic block map on a finite space with uniform
reference, $e^{-H^c_o(z)}$ is the average of $e^{-H_o}$ over the fiber of $z$, the Kadanoff block
Hamiltonian with the entropy $\log\lvert\chi^{-1}(z)\rvert$ moved into the reference. The theorem holds for any channel between standard Borel spaces
and is standard in substance (the density of a pushforward is a conditional expectation). It says
nothing about the form of $H^c_o$: on a finite network it always exists but may couple every coarse
agent to every other, while in infinite volume the analogous object can fail to exist as a
quasilocal Hamiltonian \citep{VanEnterFernandezSokal1993}, which no finite-network statement reaches.

\subsection{Perturbing the fine action: cumulants and contraction}

Perturb the fine evidence measure by a bounded $\psi$. With $C^\dagger$ the reverse kernel of the
normalized joint law,
\[
\mathcal A(\psi)(z)=-\log\int e^{-\psi(y)}C^\dagger(z,dy),
\qquad
D\mathcal A(\psi)[h]=\mathsf U^\psi h,
\qquad
D^2\mathcal A(\psi)[h,k]=-\operatorname{Cov}_{C^\dagger_\psi(z,\cdot)}(h,k)
\]
(\MT{thm:rg-bounded-action-calculus}). $\mathcal A$ is the conditional cumulant generating function of
$\psi$, $\langle\psi\rangle_z-\tfrac12\langle\psi;\psi\rangle_z+\cdots$, the cumulant expansion of any
integrating-out calculation. It is real analytic on all bounded perturbations, order preserving,
passes constants through, and never expands the sup norm
(\MTeq{eq:rg-nonlinear-action-sup-contraction}). Its linear part $\mathsf U$ is a conditional
expectation: it contracts every $L^p$ norm, its $L^2$ loss is the expected conditional variance, and
it contracts oscillation by the Dobrushin coefficient $\delta$ of the reverse kernel
\citep{Dobrushin1956} (\MT{thm:rg-action-lp-contraction}). For Kadanoff blocks the last bound is
empty: under a deterministic block map with a nondegenerate coarse law the reverse conditionals at
distinct coarse states live on disjoint fibers, so $\delta=1$ and the bound says only that nothing
expands; strict contraction needs a noisy channel. Every norm here is a norm of the whole action. A
uniform perturbation of size $h$ per agent has sup norm of order $Nh$, so these statements control
the total action, not couplings per agent.

\subsection{Interaction coordinates}

Under a product reference $\nu_\ell=\bigotimes_i\nu_{\ell i}$ equivalent to the normalized law
(\MT{def:rg-network-scale}), let $\E_{\ell,B}$ average out the agents outside $B$. The Hoeffding
projectors and the assembly of an action from couplings are
\[
\mathsf P_{\ell,B}=\sum_{B'\subseteq B}(-1)^{\lvert B\rvert-\lvert B'\rvert}\E_{\ell,B'},
\qquad
\psi_{\mathsf g}=\sum_{\varnothing\ne B\subseteq V_\ell}\mathsf g_B
\]
(\MT{def:rg-hoeffding-coordinates}). $\mathsf P_{\ell,B}f$ is the part of $f$ that depends on all of
$y_B$ jointly and has zero average in each single agent of $B$: the cluster, or ANOVA, expansion of the
Hamiltonian into $k$-body terms. For $\pm1$ spins with uniform reference the $B$-component is
$J_B\prod_{i\in B}y_i$. Gibbs-measure theory also uses the vacuum normalization relative to a reference
configuration \citep{Georgii1988}; the chapter has it too (\MTeq{eq:rg-mobius-potentials}) but uses only
the Hoeffding components in the step. \MT{thm:rg-hoeffding-action-isomorphism} proves that couplings
and actions modulo constants are isomorphic, with $\lVert\operatorname{asm}_\ell\rVert\le1$ and
$\lVert\operatorname{hd}_\ell\rVert\le3^{\lvert V_\ell\rvert}-1$, sharp: for independent biased $\pm1$
coordinates the product $\prod_iY_i$, of oscillation one, has components whose norms sum to nearly
$3^n-1$. The size of an action and the total size of its couplings are equivalent only up to a factor
exponential in the number of agents. The decomposition is gauge covariant when the frame change acts
agent by agent (\MT{cor:rg-hoeffding-gauge-covariance}), which is the case for the program's vertex
gauge; preserving the product reference alone is not enough.

Two premises limit the scope, and both matter for the program. The bounded tier
(\MT{rem:rg-bounded-tier-scope}) admits finite and categorical agents and continuous couplings of
compact group-valued states, so $\mathrm{SU}(N)$-valued states are inside; it excludes every polynomial
coupling on $\R^K$, since $y_i^\top\mathsf Ny_j$ is unbounded. The Gaussian realization with $\GL(K,\R)$
or $\GL(K,\mathbb C)$ links is therefore outside this machinery; its exact marginalization (Kron
reduction, a Schur complement) is in \SG{ch:gaussian-coarsegraining}. And product equivalence at both
depths, automatic for strictly positive laws on finite products, is not preserved by every channel: a
channel cloning one binary agent into two yields a diagonal law that no product law matches
(\MT{prop:rg-product-equivalence-not-preserved}).

\subsection{The exact step and the native growth bound}

The exact interaction step is
$\mathcal T_\ell=\operatorname{hd}_{\ell+1}\overline{\mathcal A}_\ell\operatorname{asm}_\ell$
(\MT{def:rg-interaction-step}): assemble the couplings into an action, integrate out exactly modulo
constants, read off the coarse couplings. The couplings perturb a base law whose own image is carried
separately, so $\mathcal T_\ell(0)=0$ for every model and channel (\MT{lem:rg-interaction-step}). Along an exact orbit the multipliers are
$D_\ell=\operatorname{hd}_{\ell+1}\overline{\mathsf U^{\psi_{\mathsf g_\ell}}_\ell}\operatorname{asm}_\ell$,
every interior pair $\operatorname{asm}_k\operatorname{hd}_k$ cancels, and
\[
D_{n\leftarrow\ell}
=\operatorname{hd}_n\overline{\mathsf U_{n-1}^{\psi_{\mathsf g_{n-1}}}}\cdots
\overline{\mathsf U_\ell^{\psi_{\mathsf g_\ell}}}\operatorname{asm}_\ell,
\qquad
\lVert D_{n\leftarrow\ell}\rVert_{\mathcal G_\ell\to\mathcal G_n}\le3^{\lvert V_n\rvert}-1
\]
(\MT{prop:rg-interaction-native-growth-bound}). Every tilted conditional expectation has oscillation
norm at most one, assembly has norm at most one, and only the final extraction can exceed one.

What is bounded is the operator norm of the linearized composite step from all depth-$\ell$ couplings
to all depth-$n$ couplings, each measured by the sum of the sup norms of its components: a
perturbation of total size one in the fine couplings changes the total size of the depth-$n$ couplings
by at most $3^{\lvert V_n\rvert}-1$ to first order, however many steps intervene and however large the
fine network was. It is not a bound per coupling or per agent. As a number it is useful only near the
top of a blocking: with pair blocks $512$ agents admit nine steps, and the constant is $2$, $8$, $80$,
and $6560$ for targets of one, two, four, and eight blocks, but about $10^{122}$ for the first step's
target of $256$ agents. The real content is independent of $N$: in the action norm the linearized
exact flow is a product of contractions, and all apparent amplification is one change of norm at the
target. Supplement B turns this into the statement that the exact step has no relevant direction under
tempered comparisons (\SB{cor:rg-interaction-no-relevant-mode}); relevance requires an extensive
normalization declared from outside.

\begin{analogy}{structural}
In Wilson's renormalization a uniform field $h\sum_is_i$ is relevant because the per-site field grows
as $b^{y_h}h$ under blocking by $b$ in $d$ dimensions, while the whole-action size $hN$ becomes
$b^{y_h-d}hN$, which does not grow since $y_h\le d$. Relevance is a statement about couplings per
site. The native bound is the exact finite-network form of the other half: with no rescaling and no
per-site normalization, nothing grows.
\end{analogy}

\subsection{Residuals, invariant sectors, and the finite-network obstruction}

A declared truncation is a gauge-intertwining idempotent $\mathsf P^{\mathrm{ret}}$ on couplings, for
example onto single-agent and pair components (\MT{def:rg-retained-projection}), with residual
$r^{\mathcal G}_{\ell+1}=(I-\mathsf P^{\mathrm{ret}}_{\ell+1})\mathcal T_\ell(\mathsf g_\ell)$.
\MT{prop:rg-interaction-residual} proves that the coupling and action residuals vanish together, agree
in norm up to $3^{\lvert V_{\ell+1}\rvert}-1$, and vanish on the retained range exactly when that range
is mapped into its successor. It is an identity, not an estimate. The proved invariant pairs are
generated by structure, not dynamics (\MT{prop:rg-interaction-intra-block-invariance},
\MT{rem:rg-trivial-invariant-pairs}): independent blocks under a block-product channel generate no
inter-block couplings; a finite symmetry preserves its fixed sector, such as the even sector under a
global spin flip; and couplings already carried by the coarse variables of a deterministic channel
pass through unchanged. Anything beyond is \MT{open:rg-invariant-retained-subspace}.

For finite agent spaces the coupling space has dimension $\prod_i\lvert\operatorname{supp}\nu_{\ell
i}\rvert-1$, so $2^N-1$ for $N$ binary agents (\MT{prop:rg-interaction-finite-network-obstruction}). A
step that reduces the number of configurations reduces the dimension (six binary agents in three binary
blocks go from $63$ to $7$), so no identification of the two coupling spaces exists and no single space
carries a flow; if every step divides the vertex count by at least $b>1$, the depth is at most
$\log_b\lvert V_\ell\rvert$. This is not a small-network artifact. It holds at $N=512$ and at
$N=6\times10^{23}$ alike: a finite system can be blocked only logarithmically often, and its coupling
space shrinks with each blocking. Physics removes the obstruction by passing to infinite volume and
identifying couplings of the same shape across scales by translation invariance; neither ingredient is
in the main manuscript, and Supplement B records that finite networks do not carry the asymptotic
classification (\SB{prop:rg-finite-depth-no-asymptotics}). The dimension count needs finite agent
spaces; termination holds for all.

\subsection{Hyperedges and bridge kernels}

Marginalizing a variable multiplies the factors that touch it and leaves one factor on the union of its
neighbors, so iterated elimination produces factors on elimination cliques
(\MT{prop:rg-hyperedge-generation}); this is variable elimination with fill-in. Eliminating the center
$y_0$ of a star of three $\pm1$ neighbors, with field $\xi_0$ and couplings $\xi_r$, leaves a
coefficient $2\operatorname{sech}^2(\xi_0)\tanh(\xi_0)\xi_1\xi_2\xi_3+O(\xi^5)$ on $y_1y_2y_3$: nonzero
for small nonzero couplings when the eliminated spin carries a field, and identically zero at
$\xi_0=0$, where spin-flip symmetry kills odd terms. Pairwise closure fails in general; ``generic'' is
shown within the star family, not over all networks, although the mechanism is that of decimation.

The posterior bridge is the joint posterior law of fine and coarse variables, and the bridge kernels
are its pair conditionals between a fine agent and a meta-agent, Bayes adjoints of each other
(\MT{def:rg-posterior-bridge}, \MT{prop:rg-bridge-kernels}). The meta-to-agent kernel does not
reconstruct the fine posterior: when the likelihood couples blocks, a fine agent's conditional law
depends on neighboring block variables too, just as the configuration inside a block is not fixed by
its own block spin. For a deterministic block map, lifting and conditioning returns a coarse observable
unchanged; for a randomized channel it does not.

\subsection{Meta-attention}

An attention row is the conditional law of a source given a receiver, and a conditional law cannot be
coarse-grained without the law of what it conditions on, as lumping a transition matrix needs
occupation weights. The recursive state is the joint event law $\eta$ over oriented edge copies:
\[
\eta^c_{IJ}(z)=\E_{\mathbb B_o}\Bigl[\sum_{e:j\to i,\ i\in I,\ j\in J}\eta_e(Y)\Bigm|Z=z\Bigr],
\qquad
\alpha^c_I=\sum_J\eta^c_{IJ},
\qquad
\beta^c_{IJ}=\eta^c_{IJ}/\alpha^c_I
\]
(\MT{prop:rg-meta-attention-pushforward}). The next depth pushes $\eta^c$, not $\beta^c$, and events
inside a block become meta-self events, so a meta-agent attends to itself whenever its members attend
to each other.

Where the receiver weights come from is a modeling choice with consequences. In the single-edge model
(\MT{hyp:rg-single-edge-attention}) each record carries one joint receiver--source label, and the
receiver law within a block is evidence weighted,
$\alpha_i'\propto a_{i\mid I}\widetilde{\mathrm c}_i\mathcal Z_i$, where
$\mathcal Z_i=\sum_j\pi_{ij}e^{-D_{ij}/\tau_i}$ is the partition function of receiver $i$'s row and
$\widetilde{\mathrm c}_i$ the likelihood ratio of its record under the event; the block row is the
$\alpha'$-weighted average of member rows (\MT{prop:rg-attention-single-edge-weights}). In the
per-receiver model of \cref{comp:ch06}, every receiver has its own label, and a block row needs a
structural selector that draws one receiver with fixed weights independently of everything; its
posterior equals its prior, and the block row collapses to the prior-weighted average
(\MT{prop:rg-attention-per-receiver-collapse}). The two agree when $\widetilde{\mathrm c}_i\mathcal Z_i$
is constant on the block and otherwise differ ($3/4$ against $1/2$ in the chapter's two-receiver
example); the collapse is the loss of evidence weighting. Keeping all receivers' labels is exact, and
collapsing them to one row is a further channel whose cost enters \MTeq{eq:rg-vfe-chain-rule}.

\begin{analogy}{exact}
Grouping the sources of one row gives the Gibbs row of the coarse prior $\pi^c_J=\sum_{j\in J}\pi_j$ and
coarse energy $D^c_J=-\tau\log[(\pi^c_J)^{-1}\sum_{j\in J}\pi_je^{-D_j/\tau}]$
(\MT{rem:rg-attention-log-sum-exp}): the free energy of a macrostate whose microstates have
degeneracies $\pi_j$. In the single-edge model, receivers are merged with weights proportional to their
own partition functions, as subsystems are when their ensembles are merged.
\end{analogy}

The event-weighted mean of transported link representations between blocks maps the root frame of $J$
to that of $I$ covariantly, but it need not lie in the group (the mean of rotations by $\pm\pi/2$ is
zero) or be invertible, and it is only a first moment of the joint law of marks and features that
exact closure keeps (\MT{prop:rg-attention-hom-moment}). It is the blocked link of lattice gauge
theory, which leaves the group and is usually projected back; here nothing is projected.

\subsection{Path laws and memory}

Applying the channel at every time of a path is again a coarse channel, so the free-energy chain rule
holds for path laws; the ordering index is an inference index, not physical time
(\MT{sec:rg-path-space}). The coarse path law is exact but generally not Markov.
\MT{thm:rg-strong-lumpability} is the exact condition of \citet[Thm.~6.3.2]{KemenySnell1976}, on
standard Borel spaces, for Markov coarse dynamics from every initial law: equal jump probabilities into
each coarse set from all fine states with the same coarse value. \MT{rem:rg-lumpability-witnesses}
shows a chain whose coarse process is Markov from one start without lumpability, and a deterministic
three-cycle with labels $(0,0,1)$ whose stationary coarse process has memory. For a linear recursion
$\xi_{n+1}=\mathsf T\xi_n$ with resolved coordinate $w_n=\mathsf C\xi_n$, lift $\mathsf J$
($\mathsf C\mathsf J=I$), and $\mathsf Q=I-\mathsf J\mathsf C$,
\[
w_{n+1}=\mathsf C\mathsf T\mathsf Jw_n
+\sum_{k=0}^{n-1}\mathsf C\mathsf T\mathsf Q(\mathsf Q\mathsf T\mathsf Q)^{n-1-k}\mathsf Q\mathsf T\mathsf Jw_k
+\mathsf C\mathsf T\mathsf Q(\mathsf Q\mathsf T\mathsf Q)^n\mathsf Q\xi_0
\]
exactly (\MT{thm:rg-projection-memory}): a Markovian term, memory kernels, and the unresolved initial
state. This is the discrete, linear form of the projection-operator method of \citet{Nakajima1958} and
\citet{Zwanzig1960}, the last term playing the role of the noise.

For a hidden finite chain observed through labels, the exact filter is the forward algorithm, and its
terminal posterior determines every prediction but may take infinitely many values
(\MT{thm:rg-memory-filter}), so a finite-state Markov parent does not follow even in principle. The
memory-$k$ approximation whose rows are the true conditionals given the last $k$ symbols is the unique
forward-KL optimum, with cost
\[
D^{(k)}_N=\KL\bigl(\mathbb P^{\mathrm{path}}_N\Vert\widehat{\mathbb P}^{(k)}_N\bigr)
=\sum_{t=0}^{N-1}I\bigl(Z_t;Z_{0:t-k-1}\given Z_{\max(0,t-k):t-1}\bigr)
\]
(\MT{thm:rg-memory-projection}): what the discarded past still says about the present given the kept
past. The reverse divergence is not controlled (the memory-one projection of the three-cycle gives the
impossible word $000$ probability $1/6$), and using the projection as a prior changes a recognition
law's free energy by an unsigned amount. For a stationary start,
\MT{thm:rg-memory-duration} gives $D^{(k)}_N/N\to\mathrm h^{\mathrm{ex}}_k$ with
$\sum_j\mathrm h^{\mathrm{ex}}_j\le H(Y_0)$ and $\mathrm h^{\mathrm{ex}}_k\le H(Y_0)/(k+1)$, where $N$ is
now the path length: the excess entropy \citep{CrutchfieldFeldman2003} is bounded by the entropy of the
hidden state. For the program this budget is extensive, up to $512\log2\approx355$ nats for $512$
binary agents, and the theorem is stated for finite hidden spaces only. The prepared parent of
\cref{comp:ch08} is Markov when the fine state inside each parent is refreshed to its exact conditional
law at every step (\MT{thm:rg-memory-prepared-parent}), which is adiabatic elimination by hypothesis
(\MT{hyp:rg-memory-instantaneous-refresh}); with partial refresh the parent history has memory
(\MT{rem:rg-memory-partial-refresh}), whose first-order cost is $D^{(1)}_N$.

\subsection{The retained-state table and what ``exact'' means}

\MT{tab:rg-retained-state} lists each retained object with its exact operation, the statement that
derives it, and its admission hypotheses, and the chapter says it is a catalog, not a completeness
theorem. Its first row keeps the complete normalized generative, posterior, and recognition laws, and
its last the path law with its memory. Closure is therefore achieved by keeping everything; the other rows
name parts of those objects that a smaller description must not discard, several with witnesses that
they cannot be dropped in general (hyperedges, receiver laws, memory, the marked kernel rather than its
mean, and in \cref{comp:ch12} the graph datum). ``Exact'' means exact as $\sum_ye^{-H(y)}$ is exact, a
lossless change of description, not exact as one-dimensional Ising decimation is, a finite set of
couplings with a closed recursion; no reduced retained set is shown to close. Blocks and channel are
inputs, and must be: at the uniform law on seven outcomes, no rule that sees only the law and its
Fisher metric and treats all outcomes alike can select three two-dimensional blocks
(\MT{prop:hist-law-fisher-agentization-nogo}).

\section{What is proved, and how solid it is}
\label{comp:ch11-scorecard}

\begin{scorecard}
\MT{thm:rg-effective-action} & Coarse action is minus the log of a conditional partition function;
evidence preserved & \SOLID & Any channel; standard in substance\\
\MT{thm:rg-bounded-action-calculus}, \MT{thm:rg-action-lp-contraction} & Analytic, nonexpansive action
increment; cumulant derivatives; $L^p$ and Dobrushin contraction & \SOLID & Whole-action norms;
$\delta=1$ for deterministic block maps\\
\MT{thm:rg-hoeffding-action-isomorphism}, \MT{cor:rg-hoeffding-gauge-covariance} & Couplings and
actions isomorphic; extraction constant $3^{\lvert V\rvert}-1$, sharp; gauge covariant & \SOLID & ANOVA
standard; sharp constant specific\\
\MT{rem:rg-bounded-tier-scope}, \MT{prop:rg-product-equivalence-not-preserved} & Polynomial couplings on
$\R^K$ excluded; product equivalence can fail & \SOLID & Negative; Gaussian $\GL(K)$ couplings are out
of scope\\
\MT{lem:rg-interaction-step}, \MT{prop:rg-interaction-rn-gauge-covariance} & Exact step fixes the origin,
has the tilted derivative, is gauge covariant & \SOLID & Bounded tier; covariance of a.e.\ classes\\
\MT{prop:rg-interaction-native-growth-bound} & Multipliers telescope; norm at most
$3^{\lvert V_n\rvert}-1$ at every depth & \SOLID & Bounded tier; vacuous as a number for large targets\\
\MT{prop:rg-interaction-finite-network-obstruction} & No identification of coupling spaces across a
reducing step; blocking terminates & \SOLID & Negative; every finite network, not only small ones\\
\MT{prop:rg-interaction-residual}, \MT{rem:rg-trivial-invariant-pairs} & Residual identity; independence,
symmetry, and pullback sectors are invariant & \SOLID & No residual estimate; trivial pairs only\\
\MT{rem:rg-two-agent-example} & Decimating one binary agent & \TOY & Two agents; textbook mechanism\\
\MT{prop:rg-hyperedge-generation} & Pairwise closure fails (star), not universally (parity) & \SOLID &
Negative by counterexample; ``generic'' within the star family\\
\MT{prop:rg-bridge-kernels} & Bridge kernels adjoint; no reconstruction of the fine posterior & \SOLID &
Standard in substance\\
\MT{prop:rg-meta-attention-pushforward}, \MT{prop:rg-attention-per-receiver-collapse} & Push event laws;
evidence-weighted versus prior-weighted receivers & \SOLID & Inside two declared label models\\
\MT{thm:rg-strong-lumpability}, \MT{thm:rg-projection-memory}, \MT{thm:rg-memory-filter} & Lumpability;
discrete Nakajima--Zwanzig; exact filter & \SOLID & Standard in substance\\
\MT{thm:rg-memory-projection}, \MT{thm:rg-memory-duration} & Memory-$k$ cost; excess entropy at most
hidden entropy & \SOLID & Finite spaces; stationary start; budget extensive in $N$\\
\MT{tab:rg-retained-state}, read as closure & Exact closure keeps the listed objects & \VAGUE & Rows
proved or declared; sufficiency trivial; no minimality claim\\
\MT{prop:hist-law-fisher-agentization-nogo} & Law and Fisher metric cannot select blocks & \SOLID &
Negative\\
\MT{def:rg-network-scale}, \MT{def:rg-retained-projection} & Channel, blocks, reference, retained
projection are inputs & \CHOICE & Declared, not selected\\
\MT{open:rg-invariant-retained-subspace}, \MT{open:rg-pointwise-parent-dynamics} & Closed truncation;
closed parent dynamics & \OPEN & See \cref{comp:ch11-open}\\
\end{scorecard}

The only toy positive is the two-agent decimation; the other small instances are counterexamples, hence
solid negative results. The mechanisms they display (fielded eliminated agents generate multi-agent
couplings, hidden coordinates produce memory, evidence weighting differs from prior weighting)
plausibly survive at large $N$ but are proved only where stated. The general results hold at every
finite $N$; their only $N$-dependence is in $3^{\lvert V_n\rvert}-1$ and $H(Y_0)$, the quantities a
large-$N$ theory would have to replace with intensive ones.

\begin{hygiene}
The conditions that can fail in the program's models are these. The bounded tier admits finite,
categorical, and compact group-valued states ($\mathrm{SU}(N)$ included) and excludes the quadratic
couplings of $\R^K$-valued agents, so the coupling results say nothing about the Gaussian $\GL(K,\R)$ or
$\GL(K,\mathbb C)$ realization. Product equivalence at both depths fails for channels that impose hard
constraints among coarse agents. The channel and blocks are declared and must not read the recognition
law. Everything is finite-network, with no limit taken. The memory results need finite hidden spaces
and, for the budget, a stationary start. The rest is hygiene (standard Borel spaces, measurable reverse
kernels, almost-everywhere versions, integrability), with one item touching noncompact groups:
pointwise equivariant conditional versions are settled only for finite and compact stabilizers and
weakly continuous posteriors (\MT{open:cg-equivariant-versions}), so for $\GL(K)$ covariance holds for
almost-everywhere classes.
\end{hygiene}

\section{Open ends}
\label{comp:ch11-open}

\MT{open:rg-invariant-retained-subspace} asks for a closed truncation. It is posed for finite agent
spaces with full-support product references, at least two blocks, and deterministic block-product
channels whose block maps read every member of their block (majority or block-sum maps, not
decimation), for product base laws and for all base laws. It is settled by a proper invariant pair,
invariant for every base law in an open set, whose range is not inside the sum of the pullback sector
and any grouping sector, or by a proof that generically every proper invariant pair lies in such a
sum. The exclusion of decimation sets aside the classical exactly closed recursion of the
one-dimensional Ising chain \citep{goldenfeld1992lectures}.

\MT{open:rg-pointwise-parent-dynamics} is the defining problem. The exact parent exists at every
instant as the image of the fine law, so parents do evolve, but only as images of the fine evolution;
the question is whether that evolution is a dynamics of its own. In a differentiable case the
obstruction is the semiconjugacy defect
\[
\delta_t=\partial_t\mathsf F_t+D\mathsf F_t\mathscr V_t-\overline{\mathscr V}_t\circ\mathsf F_t
\]
between the fine field pushed through the moving coarse map and a candidate parent field; for Markov
evolutions the exact condition is semigroup intertwining. Existing statements answer the unrestricted
question both ways (lumpable kernels lump, a parent carrying the full filter reproduces every path law,
the three-cycle admits no first-order parent), so it is posed for a class: switching diffusions at
$c_*$ of the four-cycle type of \cref{comp:ch13} (finitely many source assignments, a quadratic energy
on each assignment's sheet, overdamped diffusion, a rate-one clock redrawing the assignment), with a
parent that keeps the assignment and a projection of the configuration but hides coordinates that the
receiver rows of a non-negligible assignment read. Candidate parents are time-homogeneous Markov
processes of the retained variables, possibly with a regular finite-dimensional filter
(\MT{def:rg-regular-finite-filter}), started from the exact image. It asks whether every member admits
such a parent whose path law agrees in total variation as the regime parameter $\epsilon\to0$, on fixed
windows up to the first expected switch and on whole histories with diverging numbers of switches; a
member for which every such parent stays a bounded distance away settles it negatively. In physics terms: does adiabatic elimination with finitely many auxiliary variables close
the retained dynamics when hidden coordinates drive the switching? The one positive instance is the
four-cycle's inherited parent on one episode window (\MT{prop:rg-formation-endogenous-cycle}), a toy
case. The limit is $\epsilon\to0$ at any finite size, not $N\to\infty$.

Three gaps are not posed as open problems of the main manuscript: no truncation residual is estimated;
the unbounded tier on $\R^K$ exists only in the Gaussian case of Supplement G; and a thermodynamic limit
of the exact step belongs to Supplement B, where fixed-depth and diagonal limits and universality are
open (\SB{open:rg-two-index-limits}).

\begin{verdictbox}{\NARROWER}
The chapter says it determines what one exact step does to the agent network at $c_*$ and which state
must be kept so that the step closes, with an effective action, residual, finite-step multipliers, and
native growth bound; the introduction calls this exact coarse-graining of the agent network. The proofs
deliver those objects: the effective action for every channel, its analytic nonexpansive calculus, an
exact step on couplings with an exact residual identity, a depth-independent derivative bound, exact
memory identities and costs, and witnesses that hyperedges, receiver laws, marked kernels, and memory
cannot be dropped in general. The delivery is narrower than the wording in four respects. Closure is
achieved by keeping the complete laws; no smaller closed set and no truncation with a bounded error are
proved. The coupling machinery holds only in the bounded tier, which excludes the program's Gaussian
$\GL(K)$ couplings. Everything is finite-network, so there is no flow on fixed couplings (proved
impossible across reducing steps) and no thermodynamic limit. And parent dynamics is open. The scope
paragraph says most of this; ``which state must be kept so that the step closes'' is where a physicist
will read more than is there. Much of the content is standard in substance; specific are the typing at
$c_*$, the sharp constant and telescoping bound, the obstruction, the attention-event bookkeeping, and
gauge covariance throughout.
\end{verdictbox}

\chapter{The Multiscale Graph VFE and Closing the Tower}
\label{comp:ch12}

\section{The question and the answer}
\label{comp:ch12-question}

The technical chapter (\MT{ch:full-graph-meta-agent-vfe}) asks how to place everything exact
coarse-graining produces (adaptive attention event laws and links, random partitions, and the influence
of parents on their children) into one variational free energy across depths, and how a tower of depths
ends. Its answer is one normalized generative law and one recognition law over the whole tower, with
free energy
\[
\Fenergy^{\mathrm{ext}}_{\mathrm{tower}}[\mathbb Q_\lambda;X_*,o]
=-\log p_\theta(o\given X_*)
+\KL\bigl(\mathbb Q_\lambda(\cdot\given o,X_*)\Vert\boldsymbol\Pi_{\theta,o,X_*}\bigr)
\]
(\MT{def:full-graph-tower-vfe}); everything else in the chapter is exact bookkeeping of this object or
is labeled as something else. For a tower built bottom-up by coarse-graining channels, both laws factor
top-down automatically, each depth conditioned on the complete next-depth state (coordinates and graph
datum) and the partition mark (\MT{lem:full-graph-adaptive-reversal}); dropping the graph datum costs a
sum of mutual informations (\MT{prop:full-graph-dropped-datum}); and the free energy splits into
observation, top-prior, partition, graph, and parent--child relative entropies
(\MT{thm:full-graph-vfe-decomposition}). Partitions may be random, from a finite catalog under two
inequivalent normalizations (\MT{prop:rg-geometry-weights}). The tower cannot close on itself: a finite
acyclic model has a parentless latent with a declared prior (\MT{prop:obs-declared-root-unavoidable}).
What remains of the idea that the agents constitute their own top lives at three loci (the inference
loop, learning of the top prior, and delayed feedback across episodes), each an evaluation of fixed
normalized joints, and delayed feedback with a declared update rule is itself the free energy of one
lifted joint (\MT{prop:full-graph-episode-lifted-joint}). Whether descent of this free energy produces a
persistent hierarchy is open (\MT{open:full-graph-hierarchy}).

Do parents act back on their children in a proved sense? In four senses of different strength. In a
declared top-down hierarchy the downward kernel reads the parent coordinates, so the parent enters every
child's prior by construction, a modeling choice; parent and child optimization are then coupled
exactly when a mixed derivative is nonzero, which the chapter proves only as a criterion in a binary
example (a kernel odds ratio different from one), declining any family-wide claim without a proof in
that family (\MT{prop:full-graph-downward-coupling}). In an adaptive tower the downward kernel is the
Bayes reversal of the bottom-up coarse-graining, and the chapter declines to read it causally. Through
the inference loop, each child's effective prior at a mean-field fixed point is formed from the parent's
factor, itself formed from all children (\MT{rem:full-graph-inference-loop}), a standard fact. Across
episodes, the parent inferred at one episode sets the children's prior at the next inside one declared
joint (\MT{prop:full-graph-episode-joint}). Not proved is a parent with its own dynamics acting back over
time: that is \MT{open:rg-pointwise-parent-dynamics} and the third result required by
\MT{open:full-graph-hierarchy}.

\section{The picture}
\label{comp:ch12-picture}

Picture depths $\ell=0,\ldots,L$: agents at depth $0$, meta-agents at depth $1$, a declared top at depth
$L$. The state at each depth includes, besides coordinates, the graph datum: who attends to whom in both
channels, the links on every edge copy, and the root and holonomy marks of each block. Going up, a
structural channel takes the complete depth-$\ell$ state and produces a partition mark $R_\ell$ and the
depth-$(\ell+1)$ state. Read downward, the same joint law is a Markov chain in decreasing depth whose
transitions are the conditional laws of finer given coarser states. The free energy of the whole stack
is one Gibbs variational functional, and its per-depth decomposition is the chain rule.

\begin{analogy}{exact}
Bayes reversal of an adaptive tower is time reversal of a Markov chain with depth as time: the chain
$Y_0\to(R_0,Y_1)\to(R_1,Y_2)\to\cdots$ read backward is Markov, and its downward transition is the
reverse block-spin kernel of \cref{comp:ch11}. No downward causation follows, as time reversal of a
diffusion does not make the future cause the past.
\end{analogy}

\begin{analogy}{exact}
The two normalizations of a partition catalog are the two ways to combine Gibbs measures on several
sectors: a mixture of separately normalized measures with weights fixed by hand (hierarchical), or one
Gibbs measure on the disjoint union, in which each sector weighs its a priori weight times its partition
function, $\mathsf a_\gamma e^{-F_\gamma}$ (global). The first fixes the probabilities of the phases; the
second lets their free energies take part.
\end{analogy}

\begin{analogy}{exact}
The coarse links of \MT{def:full-graph-coarse-links} are maximal-tree gauge fixing inside each block
\citep{Creutz1977}: tree links go to the identity, the remaining internal links carry the block's Wilson
loops as self-loops of the coarse vertex, and boundary links connect block roots. Unlike smearing of
gauge links, nothing is averaged or projected, and every parallel copy is kept.
\end{analogy}

\section{Objects and notation}
\label{comp:ch12-notation}

\begin{decoder}
$Z_\ell$ & Full coordinates at depth $\ell$ (beliefs and models) & Block variables\\
$\mathfrak N_\ell=(\eta_\ell,\Theta^{b,\ell},\Theta^{m,\ell},\mathfrak h_\ell)$ & Graph datum: event laws,
links, root and holonomy marks & Coarse lattice with gauge field and bond occupations\\
$Y_\ell=(Z_\ell,\mathfrak N_\ell)$, $\mathbf Y$ & Complete state at depth $\ell$; all depths together &
Block configuration; multiscale configuration\\
$R_\ell$, $\gamma$ & Structural mark (realized partition or membership pattern) & Random blocking
pattern\\
$C^{\mathrm{struct}}_{\ell+1\leftarrow\ell}$ & One kernel from $Y_\ell$ to $(R_\ell,Y_{\ell+1})$ applied to
all three laws & Configuration-dependent block transformation\\
adaptive tower & Tower generated bottom-up by structural channels & Iterated blocking\\
$K^\ell_\downarrow$ & Law of depth-$\ell$ coordinates given the next depth & Prolongation; reverse
block-spin kernel\\
$\Fenergy^{\mathrm{att},x}_{i,\ell}$ & Row free energy: mean edge divergence plus $\tau$ times relative
entropy to the prior & Helmholtz free energy of a receiver's source choice\\
$\mathcal J_{\mathrm{comp}}$ & Energies appended without a normalizer & Lyapunov function, not a free
energy\\
$\mathfrak G$, $\mathsf a_\gamma$, $\mathcal Z_\gamma$ & Partition catalog; label weights; branch
partition functions & Sectors, their weights and partition functions\\
$\mathfrak m_\ell$, $\mathfrak M_\ell$ & Soft membership of a vertex; joint membership of an edge's
endpoints & Fuzzy block assignment\\
declared root, $\mathbb P_L$ & Top prior supplied by no other latent & Boundary condition at the top\\
episode $n$, $\theta_n$, $\mathcal U_n$ & Learning episode, parameter, declared update & Sweep of a
learning loop\\
latent route; parameter route & One joint across episodes; separate joints linked by updates & Hidden
Markov model over episodes; learning by updates\\
lifted joint, $\widetilde o_n$ & One joint with parameters and record copies latent; latent record copy
& State augmentation in filtering\\
$\operatorname{TC}$, $\epsilon_{\mathrm e}$, $\mathrm T$ & Total correlation across episodes; aggregate
update step; descent temperature & Multi-information; learning rate; temperature\\
\end{decoder}

Here $Y_\ell$ includes the graph datum and $Z_\ell$ is the coordinate block at depth $\ell$, unlike the
coarse sample $z$ of \cref{comp:ch11}. A learning episode (index $n$) is not a coherence episode of
formation (\cref{comp:ch13}). The fraktur $\mathfrak m_\ell$ is a membership kernel, unrelated to the
evidence measure $\mathsf m_\ell$ of \cref{comp:ch11}, and $R_\ell$ is a partition mark, not a rotation.
``Top-down'' names the direction of a generative arrow and, in an adaptive tower, carries no causal
meaning. ``Closing the tower'' means ending it with a declared root. ``Participation'' is interpretive
vocabulary, not a mathematical property.

\section{The formulas, explained}
\label{comp:ch12-formulas}

\subsection{One free energy and the structural channel}

The tower free energy is the evidence identity of \cref{comp:ch05} on a latent space that contains every
depth; with $L=0$ it is the population free energy of one network. The commitment, and the reason it is
a choice rather than a theorem, is that one normalized generative law covers all depths: independently
chosen per-depth energies added together do not define a free energy, and the chapter calls such a sum a
composite potential.

At depth $\ell$ the closed scale datum holds the full coordinates and the graph datum on a finite
multigraph, the support of the source prior at depth $0$ and the multigraph induced by coarse links
deeper down (\MT{def:full-graph-scale-datum}). A structural step is one recognition-independent channel
$C^{\mathrm{struct}}_{\ell+1\leftarrow\ell}:\mathsf Y_\ell\rightsquigarrow\mathsf R_\ell\times\mathsf
Y_{\ell+1}$, applied alike to the generative marginal, the posterior, and the recognition law, whose
laws of $R_\ell$ may therefore differ (\MT{def:full-graph-structural-channel}). Kadanoff blocks are a
constant mark with a deterministic block-product map. Random blocks are a mark drawn by the channel,
possibly from the fine configuration, which makes the step adaptive. What the blocking may not read is a
law: choosing blocks by inspecting the posterior or recognition law is excluded, since that is not one
channel applied to three laws. Soft blocks are membership kernels $\mathfrak m_\ell(I\given i,R_\ell)$, a
hard partition being the deterministic case, with a joint endpoint kernel $\mathfrak M_\ell$ for
correlated assignment of an edge's two ends (\MT{def:full-graph-membership-kernel}). Under independent
draws the coarse event law is
\[
\eta^{x,\ell+1}_{IJ}(y,\gamma)
=\sum_{i\ne j}\eta^{x,\ell}_{ij}(y)\mathfrak m_\ell(I\given i,\gamma)\mathfrak m_\ell(J\given j,\gamma)
+\delta_{IJ}\sum_i\eta^{x,\ell}_{ii}(y)\mathfrak m_\ell(I\given i,\gamma),
\]
the loop term using one draw for both ends of a self-edge (\MT{prop:full-graph-event-pushforward}). The
random $\eta_{\ell+1}$ is part of the state, not its mean, and unit-weight membership in several parents
replicates mass rather than splitting it. Nested assignments compose by Chapman--Kolmogorov, and a
multi-depth channel is exact only if the state carries every mark the next kernel reads
(\MT{def:full-graph-composition}).

\subsection{Adaptive reversal and the cost of dropping the graph datum}

The standing hypothesis \MT{hyp:full-graph-depth-markov} factors the generative law from the top,
\[
\begin{aligned}
\mathbb P_\theta(do,d\mathbf Y\given X_*)
={}&\mathbb P^{\mathrm{obs}}_\theta(do\given Z_0,X_*)\mathbb P_L(dY_L\given X_*)\\
&\times\prod_{\ell=0}^{L-1}\mathbb P^\ell_R(dR_\ell\given Y_{\ell+1})
\mathbb P^\ell_{\mathfrak N}(d\mathfrak N_\ell\given Y_{\ell+1},R_\ell)
K^\ell_\downarrow(dZ_\ell\given Y_{\ell+1},R_\ell,\mathfrak N_\ell),
\end{aligned}
\]
with $X_*$ suppressed in the depth factors, and recognition likewise. Each depth reads the complete next
depth and nothing deeper: a Markov property in scale, not mean field, since dependence within and between
adjacent depths is arbitrary. For a tower generated bottom-up by structural channels it is automatic
(\MT{lem:full-graph-adaptive-reversal}): deeper variables are produced from $Y_{\ell+1}$ by kernels that
see nothing finer, and successive disintegration gives the top-down factors, which are reversal (bridge)
kernels rather than independent choices. When the parent is a function of the child and the mark, the
downward kernel is supported on compatible child configurations, the constrained sum of a Kadanoff block
read backward.

If the depth-$\ell$ factors read only $(Z_{\ell+1},R_\ell)$, the reduced product differs from the tower
law by
\[
\KL\bigl(\mathbb P^{\mathrm{tow}}\Vert\mathbb P^{\mathrm{red}}\bigr)
=\sum_{\ell=0}^{L-1}I_{\mathbb P^{\mathrm{tow}}}\bigl(Y_\ell;\mathfrak N_{\ell+1}\given Z_{\ell+1},R_\ell\bigr)
\]
(\MT{prop:full-graph-dropped-datum}): what the coarse graph datum says about the fine state beyond the
coarse coordinates and the partition. The chapter's four-vertex witness, where fine event laws with equal
coarse coordinates push to a boundary event or a meta-self event, costs $\log2$. The identity is the
chain rule; its content is that coarse attention structure is part of the coarse state.

\subsection{The hierarchical decomposition and the row free energies}

Under the depth-Markov factorization, with common domination and integrability,
\[
\begin{aligned}
\Fenergy_{\mathrm{tower}}
={}&-\E_{\mathbb Q_\lambda}\log p^{\mathrm{obs}}_\theta(o\given Z_0,X_*)
+\KL(\mathbb Q_L\Vert\mathbb P_L)\\
&+\sum_{\ell=0}^{L-1}\E_{\mathbb Q_\lambda}\bigl[
\KL(\mathbb Q^\ell_R\Vert\mathbb P^\ell_R)
+\KL(\mathbb Q^\ell_{\mathfrak N}\Vert\mathbb P^\ell_{\mathfrak N})
+\KL(\mathbb Q^\ell_Z\Vert K^\ell_\downarrow)\bigr]
\end{aligned}
\]
(\MT{thm:full-graph-vfe-decomposition}): the data misfit at the finest depth, the top's recognition law
against its declared prior, and at each depth the costs of the partition mark, the graph datum, and the
coordinates against their parent-conditioned priors, averaged under recognition. With one depth it is
misfit plus prior relative entropy. It is the measure-level chain rule, standard in substance; its value
is that every candidate hierarchy has exact accounting.

With transported divergences $D^{x,\ell}_e$, copy priors $\pi^{x,\ell}_i$, and temperatures
$\tau^x_{i,\ell}$, the row free energy
$\Fenergy^{\mathrm{att},x}_{i,\ell}=\sum_e\beta^{x,\ell}_eD^{x,\ell}_e+\tau^x_{i,\ell}\KL(\beta^{x,\ell}_i\Vert\pi^{x,\ell}_i)$
(\MT{def:full-graph-transported-divergences}) is uniquely minimized by the Gibbs row
$\beta^*_e\propto\pi_ee^{-D_e/\tau}$ with minimum $-\tau\log\mathcal Z$, infinite-energy edges included
(\MT{prop:full-graph-gibbs-rows}): the Gibbs variational principle for $U-\tau S$ with the prior as
degeneracy. With $\tau\ne1$ this energy-unit form is not an independently weighted sector of one
evidence bound. Appending such row energies, partition terms, cross-scale terms, and a link regularizer
to a fixed population free energy gives the composite potential of
\MT{def:full-graph-composite-potential}, a descent potential with no evidence interpretation; the same
terms from a normalized Gibbs prior inside the joint belong to the tower free energy instead.

\subsection{Random partitions and the two normalizations}

A finite catalog $\mathfrak G$ gives each label $\gamma$ a carrier, a reference $\nu_\gamma$, a raw factor
$\widetilde p_\gamma$, a weight $\mathsf a_\gamma$, and a likelihood, with branch partition function
$\mathcal Z_\gamma=\sum_y\nu_\gamma(y)\widetilde p_\gamma(y)$ (\MT{def:rg-geometry-ensemble}):
\[
\mathbb P_{\mathrm h}(\gamma,y,o)=\bar{\mathsf a}_\gamma p_\gamma(y)L_\gamma(o\given y),
\qquad
\mathbb P_{\mathrm f}(\gamma,y,o)=\frac{\mathsf a_\gamma\nu_\gamma(y)\widetilde p_\gamma(y)L_\gamma(o\given y)}{\sum_\delta\mathsf a_\delta\mathcal Z_\delta}.
\]
Both are normalized with the same within-branch laws, but the label weight is $\bar{\mathsf a}_\gamma$
hierarchically and proportional to $\mathsf a_\gamma\mathcal Z_\gamma$ globally; the joints coincide
exactly when all $\mathcal Z_\gamma$ agree, rescaling one branch's raw factor changes only the global
joint, and the two free energies differ only in normalizers (\MT{prop:rg-geometry-weights}). A catalog
creates no edge. The partition term of \MT{def:full-graph-partition-factor} enters the free energy only
when its Gibbs factor is part of the generative law, and no canonical partition selector follows from the
free energy (\MT{open:cg-natural-selector}).

\subsection{Coarse links}

Under a hard partition with one root and spanning tree per block, every fine edge copy survives between
the blocks of its endpoints, with the dressed link
$\Theta^{x,\ell+1}_e=\Theta^{x,\ell}_{I\leftarrow i}\Theta^{x,\ell}_e(\Theta^{x,\ell}_{J\leftarrow j})^{-1}$
(\MT{def:full-graph-coarse-links}). By \MT{prop:full-graph-coarse-links}, no link content is lost: internal
tree copies have identity links, internal non-tree copies carry the based holonomy of their fundamental
loops, closed-walk products reproduce the block holonomy, and coarse divergences and rows are gauge
invariant. A loop divergence vanishes on a non-tree copy exactly when the parent law is fixed by that
holonomy. Once holonomy is nontrivial the trees are part of the coarse channel, not a gauge choice,
although the divergences do not depend on them for holonomy-stabilized parent laws. Overlapping
membership is excluded, since an edge would need a root transport into every block containing an
endpoint.

\subsection{How a parent enters its children's prior}

In a declared top-down hierarchy whose downward kernel factors over parents and reads the parent
coordinates, the cross-scale term of the free energy is
\[
\Fenergy^\ell_{\mathrm{cross}}
=\E_{\mathbb Q_\lambda}\KL\bigl(\mathbb Q^\ell_Z(\cdot\given Y_{\ell+1},R_\ell,\mathfrak N_\ell,o,X_*)\Vert
K^\ell_\downarrow(\cdot\given Y_{\ell+1},R_\ell,\mathfrak N_\ell,X_*)\bigr)
\]
(\MT{prop:full-graph-downward-coupling}), and the parent can also couple to child marks through the
partition and graph factors. A kernel that ignores the parent contributes no parent--child
cross-derivative through this factor, and normalization alone implies no coupling. When the kernel reads
the parent and the recognition parameterization transmits the dependence, the mixed derivative can be
nonzero, so optimization is coupled both ways although the generative arrow points down. For a binary
parent and child with product recognition $q_b(1)=u$, $q_y(1)=v$, the mixed derivative is
$-\log[K_\downarrow(1\given1)K_\downarrow(0\given0)/(K_\downarrow(0\given1)K_\downarrow(1\given0))]$,
zero exactly at unit odds ratio. That criterion in a two-variable example is the whole of the proved
bidirectional coupling.

\subsection{Closing the tower and the three loci}

In a finite directed acyclic model where every record has a latent parent, a topological order starts at
a parentless latent whose factor depends only on parameters and structural data
(\MT{prop:obs-declared-root-unavoidable}). The proof is one line of graph theory and the result an
established negative: within one finite acyclic joint the top cannot be constituted by the agents it
organizes. Same-time reciprocal structure among latents is admissible only as a globally normalized
undirected model with its partition function (\MT{sec:gen-undirected}, \MT{prop:gen-gibbs-counterexample});
an energy appended without its normalizer is not an evidence bound.

The first locus is the inference loop (\MT{rem:full-graph-inference-loop}): mean-field coordinate updates
of parent and child read each other \citep[Sec.~10.1.1]{Bishop2006}, and at a fixed point a child's
effective prior is formed from the population's beliefs through a latent inside the joint. This is
mean-field self-consistency, as when each spin of a mean-field magnet sees a field set by its neighbors,
whose fields include its own; it is a schedule against one fixed joint, not participation in time. The
second is learning the top prior (\MT{prop:full-graph-aggregate-prior}): with independent episodes
sharing a top-prior parameter, the optimum over an unrestricted family is the aggregate recognition law,
the variational form of empirical Bayes \citep{Robbins1956} and the aggregate-posterior prior of
\citet{HoffmanJohnson2016} and \citet{TomczakWelling2018}.

The third is delayed feedback. In the latent route (\MT{def:full-graph-episode-joint}) one joint spans
episodes, and the parent read out of episode $n$ parameterizes episode $n+1$, so
\[
\mathbb P^{\mathrm{ep}}\bigl(\mathbf Y^{(n+1)}\in\cdot\given o_{0:n}\bigr)
=\int\mathbb P^{(n+1\leftarrow n)}(\cdot\given z,X)
\mathbb P^{\mathrm{ep}}\bigl(Z^{(n)}_{\mathrm{par}}\in dz\given o_{0:n}\bigr)
\]
(\MT{prop:full-graph-episode-joint}): the children's next prior is the filtered parent pushed through the
transition. The joint is acyclic, so it still needs a declared root at episode $0$. In the parameter
route each episode has its own joint and the parameter is updated by a declared map of the selected
recognition coupling. \MT{prop:full-graph-episode-lifted-joint} shows that such a route is one joint:
make the parameters and latent record copies $\widetilde o_n$ latent states with Dirac transitions and
evaluate the update at the copy, not at the recognition law being scored. The lifted joint's evidence is
the product of the episode evidences, its posterior the product of episode posteriors, and
\[
\Fenergy^{\mathrm{lift}}[\mathbb Q]
=\sum_{n=0}^{n_1}\Bigl[-\log p_{\theta_n,V}(o_n\given X)
+\KL\bigl(\mathbb Q^{y}_n\Vert\boldsymbol\Pi_{\theta_n,V,o_n,X}\bigr)\Bigr]
+\operatorname{TC}(\mathbb Q^{y}),
\]
the sum of the episode bounds plus a cross-episode total correlation that vanishes at product laws. In
physics terms this is state augmentation: a learning loop becomes filtering once the parameters are
latent. The point is typing: no generative factor reads the recognition law being scored, so the typing
prohibition of \cref{comp:ch04} (\MT{req:gen-typing-prohibition}) holds.

\section{What is proved, and how solid it is}
\label{comp:ch12-scorecard}

\begin{scorecard}
\MT{def:full-graph-tower-vfe} & One generative and one recognition law over the whole tower & \CHOICE &
The modeling commitment; the rest is accounting\\
\MT{def:full-graph-structural-channel}, \MT{def:full-graph-membership-kernel} & One structural channel
for three laws; blocking reads states, not laws; soft membership & \CHOICE & Covers Kadanoff, random,
adaptive, and soft blocks\\
\MT{lem:full-graph-adaptive-reversal} & Adaptive towers factor top-down; downward factors are reversal
kernels & \SOLID & Standard in substance\\
\MT{prop:full-graph-dropped-datum} & Dropping the graph datum costs a sum of mutual informations &
\SOLID & Chain rule; $\log2$ witness illustrative\\
\MT{thm:full-graph-vfe-decomposition} & Five-term decomposition of the tower free energy & \SOLID &
Standard in substance; depth-Markov form\\
\MT{prop:full-graph-gibbs-rows} & Gibbs rows minimize row free energies & \SOLID & Gibbs variational
principle; $\tau\ne1$ not a sector of one bound\\
\MT{def:full-graph-composite-potential} & Appended energies give a descent potential & \CHOICE & No
evidence interpretation, as stated\\
\MT{prop:rg-geometry-weights} & The two normalizations differ by branch partition functions & \SOLID &
Finite catalog and carriers\\
\MT{prop:full-graph-event-pushforward}, \MT{def:full-graph-composition} & Event pushforward under soft
membership; composition & \SOLID & Product form needs independent draws\\
\MT{prop:full-graph-coarse-links} & Coarse links keep all holonomy; trees belong to the channel & \SOLID
& Hard partitions only\\
\MT{prop:full-graph-downward-coupling}, general part & No parent dependence, no cross-derivative;
normalization gives no coupling & \SOLID & Declared top-down hierarchy\\
\MT{prop:full-graph-downward-coupling}, criterion & Bidirectional coupling iff odds ratio $\ne1$ & \TOY &
Binary parent and child; no family theorem\\
\MT{prop:obs-declared-root-unavoidable} & A finite acyclic tower has a declared root & \SOLID & Negative;
every record needs a latent parent\\
\MT{rem:full-graph-inference-loop} & Parent and child updates read each other & \SOLID & Standard\\
\MT{prop:full-graph-aggregate-prior} & Optimal top prior is the aggregate recognition law & \SOLID &
Empirical Bayes in substance\\
\MT{prop:full-graph-episode-joint}, \MT{prop:full-graph-episode-lifted-joint} & Delayed feedback is the
free energy of one joint & \SOLID & Declared measurable selection and update\\
\MT{sec:obs-participatory}, read as participation & The loci realize what a cycle was reaching for &
\VAGUE & Interpretation; each identity proved\\
\MT{open:full-graph-hierarchy}, \MT{open:full-graph-episode-accounting} & Persistent hierarchy; feedback
by labels alone & \OPEN & See \cref{comp:ch12-open}\\
\end{scorecard}

Toy content is confined to the binary odds-ratio criterion and the four-vertex $\log2$ witness; only the
first stands in for a general claim, which the chapter explicitly does not make.

\begin{hygiene}
Four conditions matter. The depth-Markov form is automatic for adaptive towers but must be built into a
declared hierarchy. The structural channel must not read a law, which excludes posterior-dependent
blocking. Coarse links are constructed for hard partitions only. And no generative factor may read the
recognition law being scored, which is why the lifted joint evaluates its update at a latent record
copy. A fifth concerns reading, not truth: energy-unit row free energies with $\tau\ne1$ and the
composite potential are not evidence bounds. The rest is hygiene: standard Borel spaces and regular
conditional versions, common domination and integrability, extended-real conventions, joint measurability
of selector and update, and discrete records for the displayed episode decomposition.
\end{hygiene}

\section{Open ends}
\label{comp:ch12-open}

\MT{open:full-graph-hierarchy} asks for the hierarchy the program is after: a normalized partition law
and parent--child kernels at $c_*$ such that descent of the tower free energy produces persistent nested
blocks, with pushforwards and link lifts that compose, a controlled closure residual, and non-flat
holonomy kept rather than forced flat. What holds is the exact accounting, the fast--slow typing of
\cref{comp:ch08}, nested formation under a supplied hierarchy of interaction strengths
(\MT{thm:finite-depth-adaptive-formation}), and one endogenous level on the four-cycle
(\MT{prop:rg-formation-endogenous-cycle}). Unproved is that descent produces a second level that
depends on the first. The problem is posed for models with a partition catalog inside both laws,
recognition coordinates descending by reversible diffusion at temperature $\mathrm T$ with rate-one mark
redraws, and $\mathrm T$-independent data, as $\mathrm T\to0$. A positive answer needs three results: a
two-level lineage formed with probability bounded below whose second level is recursive in the
manuscript's strict sense (it depends on the first by a margin over counterfactual handoffs and is not a
supplied-scale sequence) and is a coherence episode (held by a restoring interaction against thermal
noise, on horizons short compared with its stationary recurrence or with probability exceeding its
stationary value by a margin); an intrinsic partition selector; and a separated slow manifold on which
coarser variables feed back consistently into finer kernels. A negative answer is a proof that no second
level produced by descent in this class meets these conditions. The recursion and coherence rules are
decoded in \cref{comp:ch13}. The limit is zero temperature on networks of any finite size, not
$N\to\infty$.

\MT{open:full-graph-episode-accounting} asks whether delayed feedback can be carried by the labels alone.
In a label-sufficient latent route the episode model is that of a parameter route with the aggregate
update
$\theta_{n+1}=(1-\epsilon_{\mathrm e})\theta_n+\epsilon_{\mathrm e}\boldsymbol\Pi^{\mathrm{lab}}_{\theta_n,V,o_n,X}$,
and the parent read-out sees only the label and a preference coordinate. Does some such route's
filtered label law follow the aggregate update to $o(\epsilon_{\mathrm e})$ per episode, uniformly over
histories and horizons of order $1/\epsilon_{\mathrm e}$, or does every route miss by order
$\epsilon_{\mathrm e}$? The records are restricted either to include an uninformative value while moving
the label posterior by an amount bounded below on average, or to have generic likelihood vectors (no
nonzero polynomial vanishes on their cone). In physics terms: can a stochastic-approximation learning
rule be represented, to first order in its step, as filtering in a fixed latent model whose memory is a
label and one preference coordinate? No natural partition selector is known either
(\MT{open:cg-natural-selector}), and parent dynamics (\cref{comp:ch11-open}) is what would make the
static and episodic senses of parents acting back dynamical.

\begin{verdictbox}{\DELIVERED}
The chapter says it places adaptive event laws and links, random partitions, and the influence of parents
on their children inside one variational free energy, factors that free energy for adaptive towers, and
closes the tower with a declared root, with participation at three loci. It does exactly that: one Gibbs
variational functional for the tower, automatic top-down factorization, an identity for the five-term
decomposition, an exact cost for dropping the graph datum, the two partition normalizations
distinguished exactly, coarse links that keep all holonomy, a proof that a declared root is unavoidable,
and a proof that delayed feedback with a declared update is the free energy of one joint. Two remarks keep
the verdict honest. What is delivered is accounting and typing, largely standard in substance (chain
rule, Markov-chain reversal, Gibbs variational principle, mean-field coordinate ascent, empirical Bayes,
state augmentation); nothing is proved about what descent of this free energy does, and the hierarchy it
would have to produce is \MT{open:full-graph-hierarchy}. And ``the influence of parents on their
children'' is delivered as a typed possibility: in a declared hierarchy the parent enters the children's
prior by construction, bidirectional coupling is proved only as a binary criterion, and in adaptive
towers the downward kernel is a Bayes reversal the chapter declines to read causally. ``Closing the
tower'' means closing it by declaration; the chapter proves that a finite acyclic tower cannot close
itself.
\end{verdictbox}

\chapter{Formation Witnesses}
\label{comp:ch13}

This chapter of the companion mirrors \MT{ch:finite-depth-formation}, with the proofs in
\MT{app:finite-depth-formation-proofs}. It is where the program's central word, formation, receives
theorems, and where the theorems most need to be read against their hypotheses: the headline result
forms collectives out of a hierarchy the model is given, not one it finds. The technical chapter's scope
paragraph declines each stronger reading; the opening paragraph of the introduction summarizes the
chapter more strongly than its theorems support.

\section{The question and the answer}
\label{comp:ch13-question}

Earlier chapters establish that an exact parent exists: coarse-graining the agent network by a Markov
channel gives a normalized parent law that preserves the evidence and splits the free energy exactly
(\MT{thm:rg-exact-coarse-vfe}; see \cref{comp:ch11}). They also establish that this parent is in
general not an autonomous Markov process (\MT{thm:rg-projection-memory}), and how a parent's own
dynamics is fixed by the network it summarizes is the defining open problem of the phase
(\MT{open:rg-pointwise-parent-dynamics}). The formation chapter asks something more concrete. In one
fully normalized experiment at the base point, can coherent groups of agents form one after another
while each parent is started from the complete law of the level below, keeps every loop coordinate and
every attention label, and learns its own copy of a model parameter from its own law? Three statements
answer, and they must be read together.

\MT{thm:finite-depth-adaptive-formation} says yes for a chain of \(2^{L-1}\) triangles at every finite
depth \(L\ge2\): the number of coherent components passes through
\(3\cdot2^{L-1}\to2^{L-1}\to\cdots\to1\) with probability tending to one, and every error made by
replacing a level with its parent vanishes, separately for the free energy, the learning signal, the
attention history, and the continuous path. The succession is put in by hand. The interaction
strengths are supplied as \(\mathrm J_\ell=\epsilon^{-\mathrm Q_\ell}\) with a separation doubly
exponential in the depth, the initial state is a prepared displacement of a common mean, the learned
parameter is pinned by a calibration record, and the attention record is saturated. The chapter says so:
formation in this family is hierarchical relaxation of the prepared mean, and attention and learning
are perturbations the proof controls, not causes (\MT{rem:rg-formation-saturation}). What the theorem
establishes is that exact marked parents can be carried through such a relaxation without their errors
accumulating, at every finite depth, with constants that depend on the depth in an uncontrolled way.

\MT{prop:rg-formation-endogenous-cycle} removes the supplied hierarchy. On a homogeneous cycle of four
agents, where an edge carries energy only while one of its endpoints attends to it, a pair forms by a
feedback between attention and belief relaxation, with probability at least
\(3e^{-1/4}/512\approx4.6\times10^{-3}\) uniformly in the small parameter, and an exact parent follows
the episode for one unit of time. The pairing is seeded by the random initial attention assignment, the
law of the process stays symmetric, there is one level only, and nothing is learned.
\MT{prop:rg-gated-depth-one} shows where this mechanism stops: in the whole gated family on any finite
graph, the coherent blocks of the stationary law are, up to an error vanishing with the temperature, the
components of a minimum edge cover, a forest of stars, and every further merge costs a factor of the
temperature.

The chapter therefore proves that exact parents can track a supplied hierarchy, that one endogenous
level can form on a small graph, and that in the family containing that mechanism the stationary law
stops at one level. It does not prove that collectives form through the agents' coupled inference in
more than one level, the program aim of \MT{sec:intro-program-aim}; that is
\MT{open:rg-formation-general}, and whether learning or loop holonomy can decide which components form
is \MT{open:rg-formation-adaptive-family}. The seeded contrast formation from every Gaussian initial
law, the pairing-selection theorem, and the parent-activation continuation results belong to the
interpretation chapter of the technical manuscript and are treated in \cref{comp:ch14}.

\section{The picture}
\label{comp:ch13-picture}

\paragraph{The triangle chain: springs of very different stiffness.}
Give each agent a belief mean \(\bar k_i\) and a model mean \(\bar m_i\), and each oriented edge
\(e:j\to i\) a translation \(\omega_e\), the link, which shifts the sender's frame into the receiver's.
The mismatch on an edge is \(r_e=\bar k_i-\bar k_j-\omega_e\), and the energy is a sum of squared
mismatches weighted by spring constants. Inside each triangle the springs, and the spring tying each
agent's belief to its own model, have strength \(\mathrm J_1\). The bridge between triangles \(j\) and
\(j+1\) has strength \(\mathrm J_{\ell(j)}\) with \(\ell(j)=2+v_2(j)\), where \(2^{v_2(j)}\) is the largest
power of two dividing \(j\); with four triangles (\(L=3\), twelve agents) the three bridges have levels
two, three, and two. The bridges thus sit on the levels of a binary tree, like the couplings of Dyson's
hierarchical model \citep{Dyson1969}, a resemblance that is structural only. The strengths are
separated by enormous factors, \(\mathrm J_1\gg\mathrm J_2\gg\cdots\gg\mathrm J_L=1\).

The system starts displaced, every mismatch of order \(\mathrm A\), with thermal fluctuations of size at
most \(\sqrt{\mathrm T}=\epsilon^3\mathrm A\), and relaxes by overdamped Langevin dynamics. The stiffest springs
relax first, on times of order \(1/\mathrm J_1\) up to a logarithm, and the triangles become internally
coherent; on times of order \(1/\mathrm J_2\) the level-two bridges relax and triangles pair up; and so on,
until on times of order one the chain is a single component. In the twelve-agent example the counts are
\(12\to4\to2\to1\). An edge is coherent when its conditional root-mean-square mismatch is below
\(\eta_L\mathrm A\), a small fixed fraction of the initial displacement, so formation here means that a
group's internal relative displacement has decayed below a resolution set by the preparation. The mean
dynamics is linear and, relative to the displacement, deterministic up to corrections of order
\(\epsilon^3\), and the grouping is fixed entirely by which springs are strong. A parent is the law with
the already relaxed distortions integrated out; it keeps the loops, the gauge directions, and the
complete attention assignment.

\begin{analogy}{exact}
Linear response after a field is switched off, with adiabatic elimination of the fast modes. The
reference preparation of \MT{def:rg-formation-regime}, a Gaussian with mean \(\mathrm A\mathbf v_0\) and
covariance \(\mathrm TH^{-1}\) on the non-gauge coordinates, is exactly the Gibbs state of
\(\tfrac12\mathbf x^\top H\mathbf x-\mathbf h^\top\mathbf x\) in the static field
\(\mathbf h=\mathrm AH\mathbf v_0\), and the chain's dynamics is the Ornstein--Uhlenbeck relaxation after
the field is switched off at time zero. Integrating out relaxed modes is exact Gaussian integration, with
the Schur complement \(H_\ell\) as retained stiffness. The identification is a theorem for the Gaussian
core, where the record is off and the parameter frozen and the mean follows
\(\bar{\mathbf x}(t)=e^{-H_{\ell'}t}\bar{\mathbf x}(0)\) (\MT{prop:rg-formation-gaussian-core}); in the full
theorem the record and learning add perturbations that the proof bounds.
\end{analogy}

\paragraph{Why the regime is so extreme.}
The reason, informally, is that two requirements pull against each other. Coherence must be many thermal
widths wide, so \(\mathrm J_\ell\mathrm A^2/\mathrm T\) must diverge. But a parent integrates out fast modes
while slower modes are still moving, and the fast modes lag behind the position the slow modes dictate.
The parent's error is that lag measured in thermal units, and at low temperature a small lag is many
thermal widths. In the Gaussian core the loss is exactly of order
\((\mathrm A^2/\mathrm T)(\mathrm J_{\ell+1}^4/\mathrm J_\ell^3)\) (\MTeq{eq:rg-formation-core-window}); keeping
it small while \(\mathrm A^2/\mathrm T\) is large forces each level's stiffness to beat the next by a large
power, and iterating over \(L\) levels gives a separation doubly exponential in \(L\). Separately, the
regime makes the network's stored energy, of order \(\mathrm A^2\mathrm J_1\), negligible against the
unit-scale calibration term. That is what makes the amplitude tiny, and it is what makes learning blind
to the network.

\paragraph{The four-cycle: attention feedback and a frustrated loop.}
Four agents sit on a ring, each receiver attending one of its two neighbors. The prepared state puts
every mean at zero and a unit translation on every link, so every edge has mismatch \(-1\) and the links
carry a circulation of four. If the assignment is a perfect matching, only the two matched edges carry
energy: they contract, while the unattended edges absorb the circulation and stretch. Each receiver's
partner edge soon becomes much cheaper than its exterior edge, and at \(\mathrm T=\epsilon^8\) the attention
posterior then keeps the partner. Attention caused the contraction and the contraction justifies the
attention. Nothing in the energies prefers one pair; the initial assignment, drawn uniformly, decides, and
the law stays invariant under the symmetries of the ring at every time.

\begin{analogy}{structural}
A frustrated plaquette, with symmetry broken along single trajectories. The flux around the loop forbids
all four bonds from being satisfied at once, so the strain must sit somewhere, and attention decides
where. The broken symmetry is of the finite-volume kind: each sample path picks a pair and holds it over
the episode while the ensemble law stays symmetric, as a finite Ising magnet below its ordering
temperature has a symmetric Gibbs state although single trajectories sit in one well. No thermodynamic
limit is taken; law-level selection is the subject of \cref{comp:ch14}.
\end{analogy}

\paragraph{Why the gated mechanism stops at one level.}
On a fixed attention assignment, a sheet, the energy stiffens exactly the attended residuals; every other
direction, including any rigid translation of one block relative to another, is held only by the reference
prior, with stiffness of order \(\mathrm T\) in energy units. All sheets share the minimum energy zero, so a
sheet's stationary weight is its phase-space volume, of order \(\mathrm T^{\mathrm k_S/2}\) with \(\mathrm k_S\)
the number of stiff directions. The law prefers the sheets with the fewest stiff directions, and each
attended edge joining two blocks adds two of them and costs a factor \(\mathrm T\). Blocks float relative to
each other and nothing pulls them together.

\begin{analogy}{structural}
Order by disorder. Among degenerate ground states, harmonic fluctuations favor those with the most soft
modes, because each stiff mode contributes a factor proportional to \(\sqrt{\mathrm T}\) to the partition
function; this is how fluctuations lift a ground-state degeneracy in frustrated magnets. Here the
degenerate states are the attention sheets and the selection favors fragmentation. The counting is the
same; the degenerate set is a finite family of label sectors rather than a continuous manifold.
\end{analogy}

\section{Objects and notation}
\label{comp:ch13-notation}

The technical chapter uses a local notation table (\MT{tab:rg-formation-notation}). The decoder covers
what a first reading needs.

\begin{decoder}
Configuration \(\mathbf x=(\bar k,\bar m,\omega)\) & Belief means, model means, link translations; agent laws are Gaussian with fixed variance, so only means move & Lattice fields: two matter fields on sites, links on bonds \\
Link \(\omega_e\), group \((\R,+)\) & Independent edge variable shifting the sender's frame into the receiver's & Noncompact abelian gauge link \\
Residual \(r_e\), cost \(\tfrac12r_e^2\) & Edge mismatch; the cost is a relative entropy between unit Gaussians & Covariant difference; bond energy \\
Loop \(h_q=\mathfrak h_q/\sqrt3\) & Normalized total translation \(\mathfrak h_q\) around a triangle & Plaquette flux \\
\(\mathcal U,\mathcal H,\mathcal K\) & Normal (gauge-invariant, flux-free) distortions; loops; pure-gauge directions & Massive modes; flux sector; gauge orbit \\
Stiffness \(H\); strengths \(\mathrm J_\ell=\epsilon^{-\mathrm Q_\ell}\) & Quadratic form of the energy; supplied spring constants, one per level & Dynamical matrix; coupling hierarchy \\
\(\mathrm A\), \(\mathbf v_0\), margin \(\mathrm m_L\) & Size and direction of the initial displacement; unrelaxed bridges stay at least \(\mathrm m_L\mathrm A/2\) from zero & Applied field switched off at time zero \\
Calibration record \(o_{\mathrm{cal}}=2\) & Precise observation of the learned parameter itself & External clamp \\
Learned scalar \(\vartheta\) & Global multiplier of the stiffness, \(\vartheta H\), learned by natural gradient & Slowly ramped coupling \\
Source, assignment \(S\), sheet & Which neighbor each receiver attends; one fixed assignment and its slice of state space & Label field; sector \\
Source record \(E_i\), \(\phi_S\), \(B_0\) & Bounded observation of the attended edge's mismatch, total effect at most \(B_0=10^{-4}\) & Weak bounded potential \\
\(\boldsymbol\Pi\), \(q^\star\), \(q_t\) & Physical posterior; best law of the restricted family; evolving recognition law & Gibbs state; variational optimum; time-dependent density \\
Marked channel \(C_{\ell\leftarrow\ell'}\), parent & Keep retained coordinates and every label, integrate out the rest & Exact decimation \\
Schur complement \(H_\ell\) & Stiffness of retained coordinates after minimizing over discarded ones & Tree-level effective stiffness \\
Discarded loss \(\Delta_{\ell':\ell}\) & Mean relative entropy of the discarded coordinates' actual law to their equilibrium law, given the retained ones & Adiabatic-elimination error \\
Window \(I_\ell\), activation, \(k_L\) & Interval on which level \(\ell\) is coherent; start of the next parent; how many relaxation times to wait & Timing of decimations \\
Coherent edge, component & Conditional root-mean-square residual at most \(\eta_L\mathrm A\); cluster of such edges & Cluster at a resolution \\
\end{decoder}

Some traps. \(\mathrm T\) is the temperature, not the frame comparison \(T_{ij}\). \(\mathrm J_\ell\) is a spring
constant while calligraphic \(\mathcal J\) is the set of assignments, and a source is the neighbor a
receiver attends, not a source term in a field theory. \(H\) is a stiffness matrix, not a Hamiltonian or a
holonomy. The root is the set of gauge directions, not the root of a tree. A parent is a marginal law,
not an agent: the scope paragraph declines that parents are coarse agents with their own belief and model
fibers. Learning is the deterministic drift of one scalar. Coherence is measured against the preparation
amplitude \(\mathrm A\), not the thermal width. Complete, exact, marked, and retained are emphasis on
objects that keep all their information; declared marks a modeling choice, as in the declared parent
kinetics. NOT-CLAIMED marks a declined claim, and declining is not refuting.

\section{The formulas, explained}
\label{comp:ch13-formulas}

\subsection{The carrier}
\label{comp:ch13-carrier}

The chain of \MT{def:rg-formation-carrier} has \(N=3\cdot2^{L-1}\) agents and configuration space
\(\R^{10\cdot2^{L-1}-1}\). By \MT{lem:rg-formation-carrier},
\[
\mathbf x^\top H\mathbf x=\mathrm J_1\sum_i\mathfrak d_i(\mathbf x)^2+\sum_e\mathrm J_{\ell(e)}r_e(\mathbf x)^2,
\qquad
|r_q|^2=|O_\triangle r_q|^2+h_q^2 .
\]
The first sum ties each belief to its own model through the discrepancy
\(\mathfrak d_i=\bar k_i-\bar m_i\), with the strongest spring; the second weights each edge mismatch by its
level. A passive frame change \(\bar k_i\mapsto\bar k_i+g_i\), \(\bar m_i\mapsto\bar m_i+g_i\),
\(\omega_e\mapsto\omega_e+g_i-g_j\) leaves every \(r_e\) and \(\mathfrak d_i\) unchanged, the
translation-group case of the frame invariance proved in general
(\MT{cor:elbo-independent-frame-invariance}). The second identity splits a triangle's squared mismatches
into two gauge-invariant contrasts and the flux \(h_q^2\): the vertex means telescope around the triangle,
so the mismatches sum to \(-\mathfrak h_q\) whatever the means are, and a triangle with flux cannot be
fully satisfied. Space splits orthogonally into \(\mathcal U\oplus\mathcal H\oplus\mathcal K\), and \(H\)
vanishes exactly on \(\mathcal K\).

\begin{analogy}{exact}
The quadratic form is the action of a Gaussian, noncompact lattice abelian gauge-Higgs model in
St\"uckelberg form, in the lattice formulation of \citet{WilsonConfinement1974} but with no plaquette
term: flux costs energy only through the matter coupling, and the model channel is a second matter field
tied on site to the first. The group is \((\R,+)\), abelian and one-dimensional, so nothing about
nonabelian or compact groups can be read off.
\end{analogy}

\subsection{One normalized experiment and its three laws}
\label{comp:ch13-experiment}

\MT{def:rg-formation-experiment} makes the carrier a probability model. Each agent has a model sample
\(m_i\sim\mathcal N(\bar m_i,\sigma^2_{\mathrm{ev}})\) and, given it, a belief sample
\(k_i\sim\mathcal N(m_i,1-\sigma^2_{\mathrm{ev}})\). The recognition family fixes the coupling of the two
samples to its optimal form, at the price \(\mathfrak d_i^2/(2(1-\sigma^2_{\mathrm{ev}}))\) per agent
(\MT{prop:rg-formation-normalized-experiment}), summing to \(\tfrac12\mathbf x^\top R_{\mathrm{ev}}\mathbf x\).
The configuration prior is Gaussian with precision \(\Omega_0(\vartheta)/\mathrm T\), where
\[
\Omega_0(\vartheta)=\vartheta H+P_{\mathcal K}-\mathrm TR_{\mathrm{ev}} .
\]
The first term is the physics, a stiffness multiplied by the learned
\(\vartheta\); the second normalizes the gauge directions; the third cancels the restriction price so that
the optimizer has a temperature-independent energy. The last two are bookkeeping.

Three observations complete the joint. A root observation gives the gauge coordinates an independent
Gaussian law, a form of gauge fixing. A calibration record \(o_{\mathrm{cal}}=2\) from
\(\mathcal N(\vartheta,\mathrm T)\) observes the learned parameter itself; it is what the parameter will
learn. And each receiver has a binary record reading the edge it attends,
\[
\Pr(E_i=1\mid\mathbf x,S_i=j)=\tfrac12\exp\Bigl[-\delta\chi\Bigl(\sqrt{\mathrm J_1/\mathrm T}r_{ij}(\mathbf x)\Bigr)\Bigr],
\qquad \chi(r)=1-e^{-r^2/2},\qquad \delta=10^{-4}/N .
\]
The argument of \(\chi\) is the mismatch in units of the thermal width \(\sqrt{\mathrm T/\mathrm J_1}\), and
\(\chi\) saturates at one. Summed over receivers, the configuration-dependent part \(\phi_S\) of the record's
negative log-likelihood lies in \([0,10^{-4}]\), so attention differs from its prior by a factor of at most \(e^{\pm10^{-4}}\)
(\MTeq{eq:rg-formation-source-saturation}): attention is blind to the configuration by construction. The
attention prior is uniform except at one receiver \(i_0\), whose prior is
\(\operatorname{sig}(4(\vartheta-3/2))\); the learned parameter controls one attention weight, through the
prior only.

Three laws coexist (\MT{def:rg-formation-laws}; \MT{prop:rg-formation-laws}): the physical posterior
\(\boldsymbol\Pi\), the restricted optimizer \(q^\star\), and the evolving recognition law \(q_t\). The
scaled free energy is
\(\Fenergy_{\mathrm T}(q,\vartheta)=\mathrm T\KL(q\Vert q^\star_\vartheta)+\mathrm C_{\mathrm T}(\vartheta)\),
and the optimizer differs from the posterior by at most
\(2e^{B_0}N\mathrm T^2/((1-\sigma^2_{\mathrm{ev}})\mathrm J_1)\) in free energy
(\MTeq{eq:rg-formation-physical-optimizer-gap}), negligible here. The theorem is proved for both
\(\boldsymbol\Pi\) and \(q^\star\).

\subsection{Marked parents and the discarded loss}
\label{comp:ch13-parents}

The parent at level \(\ell\) comes from the deterministic channel
\(C_{\ell\leftarrow\ell'}((y,S),\cdot)=\delta_{(P_\ell y,S)}\), which projects onto the retained space
\(W_\ell\) (unrelaxed normal directions, all loops, all gauge directions), keeps the whole assignment, and
leaves the observations alone (\MT{def:rg-formation-parents}). The parent law is a marginal, with sheet
potential
\[
V_{\ell,S}(y,\vartheta)=\tfrac12\vartheta y^\top H_\ell y+\mathrm C_\ell^{\mathrm G}(\vartheta)+\mathrm T\psi_{\ell,S}(y,\vartheta),
\qquad 0\le\psi_{\ell,S}\le B_0,
\]
by \MT{prop:rg-formation-parent-potential}: the Schur-complement stiffness, the Gaussian determinant of
the integrated modes (their entropy, which depends on \(\vartheta\)), and the record averaged over the
discarded coordinates, still below \(10^{-4}\). The chain rule of relative entropy gives what a parent
loses,
\[
\Fenergy^{\ell'}_{\mathrm T}(q^{\ell'},\vartheta)-\Fenergy^{\ell}_{\mathrm T}\bigl((C_{\ell\leftarrow\ell'})_\#q^{\ell'},\vartheta\bigr)
=\mathrm T\Delta_{\ell':\ell}(q^{\ell'},\vartheta),
\]
where \(\Delta_{\ell':\ell}\) is the mean relative entropy of the discarded coordinates' actual conditional
law to the optimizer's conditional. It vanishes exactly when the discarded coordinates are in conditional
equilibrium, the condition for exact adiabatic elimination. These identities are standard in substance.

The parent is not a coarse agent. At the top level the retained space has dimension \(4\cdot2^{L-1}\): all
\(N\) gauge directions and all \(N/3\) loops, with every normal direction integrated out
(\MTeq{eq:rg-formation-dimensions}), alongside the full assignment with \(2^N\) values. For \(L=2\) the
carriers go \(\R^{19}\to\R^9\to\R^8\) with 64 assignments throughout, and the hidden conditionals are not
Gaussian once the record is on (\MT{rem:rg-formation-two-triangle}). The top parent is the zero-mode and
flux sector of the Gaussian theory together with the complete label field.

\subsection{Kinetics and learning}
\label{comp:ch13-kinetics}

The microscopic kinetics (\MT{def:rg-formation-kinetics}) are overdamped Langevin dynamics on each sheet
with a rate-one clock that redraws the whole assignment from the attention posterior,
\[
\mathscr L_\vartheta^0\mathsf f(\mathbf x,S)=\mathrm T\Delta_{\mathbf x}\mathsf f-\nabla_{\mathbf x}U_{\vartheta,S}\cdot\nabla_{\mathbf x}\mathsf f
+\sum_{S'}\beta_{S'}(\mathbf x,\vartheta)\bigl[\mathsf f(\mathbf x,S')-\mathsf f(\mathbf x,S)\bigr],
\]
a diffusion at temperature \(\mathrm T\) with unit mobility plus a heat-bath update of the labels. The parent
generator has the same form with the parent potential and inherited jump rows, conditional averages of
the fine rows, which are reversible and compose across levels (\MT{prop:rg-formation-rate-composition}).
The parent kinetics is declared, not derived: the time evolution of a true marginal carries memory
(\MT{thm:rg-projection-memory}; compare \citet{Zwanzig1960}), and the chapter checks its Markov closure
afterwards through conclusions (c) and (e).

Each level learns its own copy of \(\vartheta\) by natural gradient,
\[
\dot\vartheta_\ell=-\epsilon\mathsf m(\vartheta_\ell)^{-1}g_\ell(q^\ell_t,\vartheta_\ell),
\qquad
\mathsf M(\vartheta)=\frac1{\mathrm T}+\frac12\Tr\bigl[(\Omega_0(\vartheta)^{-1}H)^2\bigr]+16\pi_{i_0,+}(1-\pi_{i_0,+}),
\]
with \(\mathsf m=\mathrm T\mathsf M\) (\MT{def:rg-formation-learning};
\MT{prop:rg-formation-score-dissipation}). The metric is dominated by the calibration term, so
\(\mathsf m\approx1\), and the covector \(g_\ell\) (\MTeq{eq:rg-formation-complete-score}) is \(\vartheta-2\)
plus the stored elastic energy \(\tfrac12\E y^\top H_\ell y\), at most of order \(\mathrm A^2\mathrm J_1\),
plus entropic and attention terms of order \(\mathrm T\). Hence
\(\vartheta(t)\approx2-e^{-\epsilon t}\). Along the coupled dynamics the free energy decreases through
three nonnegative dissipation terms, the Fisher information of the diffusion, the symmetrized entropy
production of the label jumps, and the natural-gradient term \(\epsilon\mathsf m^{-1}g_\ell^2\)
(\MTeq{eq:rg-formation-dissipation}): an H-theorem, standard in substance.

\subsection{The regime}
\label{comp:ch13-regime}

The regime of \MT{def:rg-formation-regime} is
\[
\mathrm Q_L=0,\quad \mathrm Q_\ell=2\mathrm Q_{\ell+1}+20,\quad \mathrm J_\ell=\epsilon^{-\mathrm Q_\ell},\quad
\mathrm A^2=\epsilon^{3\mathrm Q_1+20},\quad \mathrm T=\mathrm A^2\epsilon^6,\quad t_*=\frac2\epsilon .
\]
Then \(\mathrm Q_1=20(2^{L-1}-1)\) and \(\mathrm T=\epsilon^{60\cdot2^{L-1}-34}\) (\MT{rem:rg-formation-regime}).
Since \(N=3\cdot2^{L-1}\), this is \(\mathrm T=\epsilon^{20N-34}\): the temperature exponent is linear in the
number of agents, and the stiffness range is \(\epsilon^{-20(N/3-1)}\). The last two columns of the table
are arithmetic on the class of schedules in which the sharpness problem is posed
(\cref{comp:ch13-core}): the least admissible \(\mathrm Q_1\) there, rounded, and its temperature.

\begin{center}
\small
\begin{tabular}{@{}rrrrrrr@{}}
\toprule
\(L\) & \(N\) & \(\mathrm Q_1\) & \(\mathrm A\) & \(\mathrm T\) & least \(\mathrm Q_1\), base \(4/3\) & its \(\mathrm T\)\\
\midrule
2 & 6 & 20 & \(\epsilon^{40}\) & \(\epsilon^{86}\) & 14 & \(\epsilon^{68}\)\\
3 & 12 & 60 & \(\epsilon^{100}\) & \(\epsilon^{206}\) & 25.3 & \(\epsilon^{102}\)\\
4 & 24 & 140 & \(\epsilon^{220}\) & \(\epsilon^{446}\) & 40.4 & \(\epsilon^{147}\)\\
8 & 384 & 2540 & \(\epsilon^{3820}\) & \(\epsilon^{7646}\) & 171 & \(\epsilon^{539}\)\\
9 & 768 & 5100 & \(\epsilon^{7660}\) & \(\epsilon^{15326}\) & 235 & \(\epsilon^{730}\)\\
\bottomrule
\end{tabular}
\end{center}

The ratios the proof actually uses are the stiffness range \(\epsilon^{-\mathrm Q_1}\), the noise-to-signal
ratio \(\mathrm T/\mathrm A^2=\epsilon^6\), the learning mobility \(\epsilon\), and the stored energy relative to
calibration, \(\mathrm A^2\mathrm J_1=\epsilon^{2\mathrm Q_1+20}\). The window of level \(\ell\) is
\(I_\ell=[t_\ell,2t_\ell]\) with \(t_\ell=k_L\log(1/\epsilon)/\mathrm J_\ell\); the last closes at order
\(k_L\log(1/\epsilon)\), long before the horizon \(2/\epsilon\) on which the parameter moves. The entrance
constant must exceed thresholds proportional to \(\mathrm Q_1\) (\MTeq{eq:rg-formation-common-entrance}), and
\(\epsilon_L\) is not computed. The regime is a formal asymptotic regime at each fixed depth, with no limit in
which \(N\) grows at fixed \(\epsilon\).

\subsection{The theorem, conclusion by conclusion}
\label{comp:ch13-theorem}

\MT{thm:finite-depth-adaptive-formation} assumes the experiment and kinetics at a fixed depth (H1), the
regime (H2), an initial law in the preparation neighborhood (H3), and scheduled activations (H4). The
preparation neighborhood consists of smooth transports of the Gaussian reference that move each sample by
less than \(\epsilon\sqrt{\mathrm T}\) in the energy norm (\MTeq{eq:rg-formation-preparation-transport}): a
thermalized cloud around the displaced mean, slightly deformed, possibly depending on \(\epsilon\) and
\(L\). The reader's phrase ``coherence is linear hierarchical relaxation of the prepared mean'' has a
precise meaning here. Up to perturbations the proof bounds, the mean obeys
\(\dot{\bar{\mathbf x}}=-\vartheta H\bar{\mathbf x}\) from \(\mathrm A\mathbf v_0\). In the flag coordinates of
\MTeq{eq:rg-formation-weighted-flag}, layer \(\ell\) has stiffness of order \(\mathrm J_\ell\) and couples to the
slower layers only through strengths of order \(\mathrm J_{\ell+1}\), so the solution decays level by level. On \(I_\ell\) the layers up to \(\ell\) have decayed
below \(\eta_L\mathrm A\), while the higher bridges have moved by at most of order
\(\mathrm J_{\ell+1}t_\ell\mathrm A\to0\) and keep their prepared margin.

Conclusion (a) gives unique non-exploding solutions for every process, and learned parameters in
\([1,21/10]\) that stay within \(\mathrm C_L\epsilon^{2\mathrm Q_1+20}\) of the calibration curve
\(2-e^{-\epsilon t}\) (\MTeq{eq:rg-formation-parameter-bound}): learning is relaxation to the calibrated value,
the same at every level because no level can see the network.

Conclusion (b) is the formation statement. Except on an event of probability at most
\(\mathrm C_L(1+t_*)\exp(-\mathrm c_L/\epsilon^6)\), the fine process and its parents show the counts
\(3\cdot2^{L-1}\to2^{L-1}\to\cdots\to1\) (\MTeq{eq:rg-formation-agent-counts}): on \(I_\ell\) every edge of
level at most \(\ell\) is coherent, every belief agrees with its model to within \(\eta_L\mathrm A\), and every
higher bridge stays a fixed multiple of \(\mathrm A\) from zero; the final component persists to the horizon,
and the loops stay random and retained. The exponent \(\epsilon^{-6}=\mathrm A^2/\mathrm T\) is the
signal-to-noise ratio. By \MT{rem:rg-formation-saturation} the coherence proof uses the record only through
\(0\le\phi_S\le B_0\) and the parameter only through \(1\le\vartheta\le21/10\), so it gives the same
components with the record off and the parameter frozen at one. Attention and learning play no role in (b).

Conclusion (c) bounds the free energy discarded at each activation, \(\Delta_{\ell-1:\ell}\le\mathrm C_L\epsilon^{14}\)
early and \(\mathrm C_L\epsilon^2\) at the last level, and the difference of consecutive learning signals by
\(\mathrm C_L\mathrm T\epsilon^7\) and \(\mathrm C_L\mathrm T\epsilon\) (\MTeq{eq:rg-formation-conditional-score-bounds}):
the parents are faithful in free energy and in learning gradient. The exponent 14 is
\(\mathrm A^2\mathrm J_{\ell+1}^2/(\mathrm T\mathrm J_\ell)=\epsilon^{-6}\epsilon^{20}\), the general Fisher-budget
estimate; the Gaussian core shows the true early loss is much smaller.

Conclusion (d) shows the model is not static: from every activation to the horizon the distinguished
receiver's mass on source \(+\) rises by more than one half, and the first hidden conditionals change by
relative entropy more than \(1/200\) (\MTeq{eq:rg-formation-source-movement}). Both are driven by
\(\vartheta\) moving toward two, through the prior \(\operatorname{sig}(4(\vartheta-3/2))\) and the stiffness
\(\vartheta H\), and both happen after formation, since \(\vartheta\le1+\mathrm C_L\epsilon\log(1/\epsilon)\)
until the last window closes. This is an engineered stress test showing that the lineage tracks kernels
that really move, not a mechanism of formation.

Conclusion (e) says that after each activation the retained coordinates of the fine process and of the
parent stay within \(\mathrm C_L\mathrm A\epsilon^2\) of one deterministic path except with probability
\(\mathrm C_L(1+t_*)e^{-\mathrm c_L/\epsilon^2}\) (\MTeq{eq:rg-formation-common-path-concentration}), and that
their attention histories can be coupled to disagree with probability at most \(\mathrm C_L\epsilon^6\)
(\MTeq{eq:rg-formation-two-path-bounds}). The resolution \(\mathrm A\epsilon^2=\sqrt{\mathrm T}/\epsilon\) lies
between the thermal width and the displacement. This small-noise statement is weaker than closeness of
path laws, which the chapter declines along with exact lumpability: the parent reproduces the fine
trajectory at a resolution coarser than the fluctuations it carries.

\MT{rem:rg-formation-saturation} fixes the reading of the whole theorem. The record is saturated, since at a
residual of order \(\mathrm A\) its argument is of order \(\epsilon^{-(3+\mathrm Q_1/2)}\), so it cannot tell a
coherent edge from an incoherent one at the coherence resolution; learning is calibration; and the loops
have root-mean-square size of order \(\sqrt{\mathrm T/\mathrm J_1}\), so holonomy is thermal noise. The
chapter's own summary says the reductions succeed in a regime ``in which nothing adaptive decides the
outcome''.

\subsection{The Gaussian core and whether the doubling is needed}
\label{comp:ch13-core}

\MT{prop:rg-formation-gaussian-core} computes the loss with the record off and the parameter frozen. Part (a)
gives \(\Delta=\tfrac1{2\mathrm T}\mathfrak v^\top H_{zz}^{-1}\mathfrak v\), with \(\mathfrak v\) the force on the
discarded coordinates: the lag of the fast modes behind their slaved position, weighted by compliance, in
thermal units. Part (b) gives, uniformly on each early window,
\[
\Delta_{\ell-1:\ell}(t)=\frac{\mathrm A^2}{\mathrm T}\frac{\mathrm J_{\ell+1}^4}{\mathrm J_\ell^3}
\bigl[\mathfrak y_\ell(\mathbf v_0)+o(1)\bigr],
\qquad \mathfrak y_\ell(\mathbf v_0)>0 .
\]
Informally, the next level moves at rate \(\mathrm J_{\ell+1}\) and pulls through couplings of order
\(\mathrm J_{\ell+1}\), the current level answers with compliance \(1/\mathrm J_\ell\), so the lag force is of
order \(\mathrm A\mathrm J_{\ell+1}^2/\mathrm J_\ell\) and its compliance-weighted square gives the display;
the general estimate behind (c) is larger by \((\mathrm J_\ell/\mathrm J_{\ell+1})^2\). Part (c) turns this into
a criterion: with \(\mathrm A^2/\mathrm T=\epsilon^{-\mathrm Q_{\mathrm T}}\) and a large enough entrance constant,
the core loss vanishes on every early window exactly when \(3\mathrm Q_\ell-4\mathrm Q_{\ell+1}>\mathrm Q_{\mathrm T}\) for every \(\ell<L\), and is
\(O(\epsilon^{\mathrm Q_\Delta})\) exactly when
\(3\mathrm Q_\ell-4\mathrm Q_{\ell+1}\ge\mathrm Q_{\mathrm T}+\mathrm Q_\Delta\), forcing
\(\mathrm Q_1\ge(\mathrm Q_{\mathrm T}+\mathrm Q_\Delta)[(4/3)^{L-1}-1]\).

Is the doubling required? Not shown, and the chapter says so. It enters the proof through the full Fisher
budget of \MT{lem:formation-conditional-entropy} and the history estimate of
\MT{lem:formation-history-coupling}. What is proved is that in the Gaussian core, for this preparation and
schedule, a separation growing like \((4/3)^L\) in the exponent is necessary. Whether the full family works
under the base-\(4/3\) schedules \(3\mathrm Q_\ell\ge4\mathrm Q_{\ell+1}+20\), \(\mathrm Q_\ell\ge14\), with
depth-independent constants, is \MT{open:rg-formation-regime-sharpness}. Either way the separation diverges
with the number of agents: since \(L-1=\log_2(N/3)\), \((4/3)^{L-1}=(N/3)^{\log_2(4/3)}\approx(N/3)^{0.415}\), so
the exponent of the stiffness range grows linearly in \(N\) under the doubling and like a power of \(N\) under
the best schedule the core allows. That is the quantitative sense in which this family has no large-\(N\)
limit.

\subsection{The four-cycle witness}
\label{comp:ch13-cycle}

\MT{def:rg-formation-cycle-experiment} puts four agents on a ring with edge cost and record
\[
\mathcal D_{e_i}(\mathbf x)=\tfrac12\bigl[(r^b_i)^2+(r^m_i)^2\bigr],
\qquad
\Pr(E_i=1\mid\mathbf x,S)=\tfrac12\exp\bigl[-\mathcal D_{e(i,S_i)}(\mathbf x)/\mathrm T\bigr],
\]
the cost adding belief- and model-channel relative entropies and the record reading only the attended
edge, with weight \(1/\mathrm T\). An attended edge is stiff against the thermal scale; an unattended edge
carries no energy at all. This is the gating. The prior is uniform, nothing is learned, every coordinate
has unit mobility, \(\mathrm T=\epsilon^8\) with \(\epsilon\le1/100\), the formation time is
\(t_\epsilon=2\log(1/\epsilon)\), and the horizon is \(t_\epsilon+1\). The preparation is an equal mixture of
narrow Gaussians around \(\pm\mathbf x^\circ\), where \(\mathbf x^\circ\) has all means zero and all links one,
with a uniform assignment, so it names no pair. On the path that holds a matching from \(\mathbf x^\circ\),
\(r_{\mathrm{int}}=-e^{-(8+\mathrm T)t}\) and \(r_{\mathrm{ext}}=e^{-\mathrm Tt}[-\tfrac32+\tfrac12e^{-8t}]\)
(\MTeq{eq:rg-formation-cycle-center}): matched mismatches decay at rate eight, exterior ones go to about
\(-3/2\) and then decay at the thermal rate. The circulation, equal to minus the sum of the belief mismatches,
drops from four to about three and then relaxes only at rate \(\mathrm T\); it is what stretches the
unattended edges.

Part (a) of the proposition gives the episode: with probability at least \(p_0=3e^{-1/4}/512\) the assignment
is one perfect matching throughout; at time zero each matched pair has belief relative entropy above
\(1/4\), so the contrast is entered during the episode; after \(t_\epsilon\) the matched relative entropies are
at most \(\epsilon^4\); and after time \(1/4\) every exterior belief relative entropy exceeds \(1/2\). The pair
is held with restoring rates \(8+\mathrm T\) and \(4+5\mathrm T\), and the kinetics started from the
stationary law realize the event with probability at most \(4e^{-1/(4\mathrm T)}<p_0/2\). Part (b) is the
parent channel, discarding the three fast normal directions of the attended edges and keeping the
assignment, the loop, and one fast coordinate. Part (c) bounds the discarded conditional relative entropy
by \(\bar\eta_\epsilon\) (\MTeq{eq:rg-formation-cycle-constants}). Part (d) puts the parent activated at
\(t_\epsilon\) within total variation \(2p_\epsilon/(1-p_\epsilon)+4\sqrt{8/3}e^{-1/\mathrm T}\) of the native
retained path law conditioned on the same prefix, assignment path included. The constants are explicit at every \(\epsilon\le1/100\); at
\(\epsilon=1/100\) they give a path-law error of about \(1.5\) percent and a discarded loss of about \(0.045\)
nats, both vanishing as \(\epsilon\to0\). Unlike the chain, the comparison here is in total variation of path
laws, because on a held matching the hidden coordinates are exactly independent Ornstein--Uhlenbeck modes.

\begin{center}
\small
\begin{tabular}{@{}>{\raggedright\arraybackslash}p{0.26\textwidth}>{\raggedright\arraybackslash}p{0.54\textwidth}>{\raggedright\arraybackslash}p{0.14\textwidth}@{}}
\toprule
Ingredient & Role & Origin \\
\midrule
Attention-contraction feedback & Attended edges contract, contracted edges win attention, the matching persists & endogenous \\
Energy scale & All edges have the same cost structure & homogeneous \\
Which pair forms & Initial assignment, drawn uniformly; the law stays symmetric & random seed \\
Circulated preparation & A flux around the ring that survives on the episode's time scale and stretches the unattended edges & supplied \\
Source gating & Energy only while attended, with weight \(1/\mathrm T\) & supplied \\
Temperature \(\epsilon^8\) & Makes attention an arg-min of edge costs & supplied \\
Euclidean kinetic metric & Beliefs, models, and links all move at rate one; models are not slow & supplied \\
Duration-one episode; no learning & The pair is proved held for one time unit; parameters fixed & supplied \\
\bottomrule
\end{tabular}
\end{center}

One observation places the witness. The four-cycle belongs to the gated family, so by
\MT{prop:rg-gated-depth-one} its stationary law already sits on the two perfect-matching sheets, up to a
fraction of order \(\mathrm T\) of its mass, with the matched pairs as its coherent blocks. The paired block
partition the episode reaches is the typical partition of equilibrium; what distinguishes the episode from
equilibrium coherence is its entrance clause (the pairs are incoherent at time zero) and the seeded choice
of matching. Physically, the witness is relaxation from
a prepared symmetric incoherent start into one of two paired equilibrium sectors, the choice made by a
sampled label and held by attention feedback. That is a genuine coupled-inference mechanism at one level,
not a collective the equilibrium law would not already contain.

\subsection{The gated family and depth one}
\label{comp:ch13-gated}

\MT{def:rg-gated-translation-family} places the four-cycle's laws, records, and kinetics on any finite
simple graph without isolated vertices, with a product prior independent of \(\mathrm T\). On a sheet the
energy is \(V_S=\tfrac12\mathbf x^\top H_S\mathbf x\), where \(H_S\) sums over receivers the residual pair of
the attended edge, and both laws are Gaussian on each sheet with precision \(H_S/\mathrm T+\Omega\)
(\MTeq{eq:rg-gated-sheet-law}). By \MT{prop:rg-gated-depth-one}, (a) every sheet has a stiff band of order
one on the \(\mathrm k_S\) attended directions and a soft band of order \(\mathrm T\) elsewhere, and the prior
enters only sheet weights, never stiffness; (b) the configuration marginal is a Gibbs law within
\(\mathrm T\sum_i\log(1/\min_j\pi_{ij})\) of the energy \(V_{\min}=\sum_i\min_j\mathcal D_{e(i,j)}\), in which each
receiver pays only for its cheapest edge and which is flat under rigid translation of any block of a
covering edge set; (c) the counting is
\[
\operatorname{rank}\mathbf a_F=|F|+N-\operatorname{cc}(F),
\qquad
\mathrm k_S\ge2(N-\mathfrak n),
\qquad
q_\Omega\{S:\operatorname{cc}(\operatorname{att}(S))\le\mathfrak n-m\}\le\mathrm K_m\mathrm T^m,
\]
with equality in the middle exactly for minimum edge covers, the spanning forests of \(\mathfrak n\) stars,
and sheet weights proportional to \(\mathrm T^{\mathrm k_S/2}\) up to controlled factors
(\MTeq{eq:rg-gated-sheet-weight}); (d) at resolution \(\mathrm a\) the stationary block partition differs from
a minimum-edge-cover partition with probability at most
\[
\mathrm K_1\mathrm T+\frac{|E|\|\Omega\|}{2\mathrm s_*^2}\mathrm a^2+4|E|e^{-\mathrm a^2/(4\mathrm T)},
\]
the mass of merged sheets, the chance that an unheld edge between blocks falls inside the resolution, and
the chance that a held edge fluctuates outside it. On an even ring or an even complete graph the blocks are
the pairs of a perfect matching, and on a star graph the block is the star: the stationary law fragments
into the smallest blocks the graph allows.

Is this large-\(N\), and is it static? It holds on every finite graph, so it applies verbatim at 512 or more
agents and is not a small-network result. But it concerns \(\mathrm T\to0\) at a fixed graph with
graph-dependent constants. \(\mathrm K_1\) contains the factor \((1+\|\Omega\|/\mathrm h_*)^{N-\mathfrak n}\),
exponential in \(N\), times a ratio whose numerator runs over every merged assignment
(\MTeq{eq:rg-gated-merge-cost}), and the second term grows with the number of edges. Nothing keeps the bound
small as \(N\) grows at fixed temperature, so the limit of many agents at fixed temperature, where the number
of ways to merge competes with the cost of each merge, is not addressed. The statement is also purely
static: it has nothing on relaxation times, metastability, or how long sampled paths can hold a larger block.
Since the dynamic question is posed only outside the family (\MT{open:rg-formation-general}), a transient
second coalescence inside the gated family is neither excluded nor posed.

\subsection{The authoring rules R13 to R16 in physics terms}
\label{comp:ch13-rules}

Four authoring rules govern what the technical documents may call learning-driven selection, recursion, a
coherence episode, and a settled open problem.

R13, the Curie rule, is Curie's principle applied to learning. If model, records, parameterization, and
kinetics share a symmetry, a claim that learning selects a structure the symmetry does not fix must name
the symmetry-breaking datum it amplifies and compare against a frozen-learning run, or else declare the
selecting mechanism, for example a rule acting on sampled paths, and then the claim concerns sampled paths.
A unique, equivariant law-level dynamics from a symmetric start cannot produce an asymmetric law, the
finite-volume fact that a finite system's Gibbs state is symmetric. The rule is substantive: it is why the
four-cycle result claims only a sampled-path pairing.

R14 defines recursion by a counterfactual. One declares a handoff time, a first-stage event with its
blocks, a second-stage event whose block is a union of first-stage blocks, a statistic such as the second
stage's probability, and a class of handoff states that includes, with probability bounded below, states
from histories without first-stage formation. The second stage is recursive when it has probability
bounded below and its statistic after actual first-stage formation beats its value after a counterfactual
handoff by a margin. If it happens with probability tending to one from every handoff in the class, or its
proof uses only properties that counterfactual handoffs share, the sequence is a supplied-scale sequence:
the hierarchy is in the Hamiltonian, as in any system whose separated couplings order each on its own scale
whether or not the previous scale has ordered. The physical distinction is between a hierarchy of
couplings and a hierarchy of products, whose stage-two units are made of stage-one units and would not
form otherwise. The definition is substantive. The chapter does not run the test on the triangle chain, but
each level's coherence follows from that level's stiffness acting on the prepared displacement, so in
R14's terms the chain is a supplied-scale sequence; this is the companion's reading.

R15 defines a coherence episode. The restoring clause requires a named interaction whose rate times the
squared resolution, divided by the temperature, diverges on the sheets the kinetics select: by
equipartition a mode of stiffness \(k\) has width \(\sqrt{\mathrm T/k}\), so the block sits many thermal widths
inside its resolution, held by energy rather than luck. The horizon clause requires horizons short compared
with the time equilibrium takes to revisit the event, or a probability exceeding the equilibrium value by a
margin. Restored coherence that fails the horizon clause, for instance because equilibrium has it with
probability tending to one, is restored equilibrium coherence, not formation; coherence failing the
restoring clause is coincidental. The recurrence branch is how a physicist says a broken-symmetry state is
real, since tunneling times diverge; the margin branch admits the transients of relaxation from any
prepared, atypical start. The rule thus excludes thermal accidents and equilibrium ordering from the word
formation but not preparation-driven relaxation. The four-cycle meets both clauses, the horizon clause
through its entrance requirement. In the chain the intermediate windows meet the margin branch, but the
final single component is what the restricted optimizer looks like, every residual being thermal under it,
so by R15's own terms the chain's last stage is restored equilibrium coherence. The chapter does not
evaluate R15 for the chain; this is the companion's reading.

R16 requires each open problem to be posed in a stated class with two nontrivial settling branches, neither
met by an existing statement or a trivial construction. It is bookkeeping for well-posedness, and it is why
the open problems below are long and cannot be closed by pointing at the witnesses already proved.

\section{What is proved, and how solid it is}
\label{comp:ch13-scorecard}

\begin{scorecard}
\MT{lem:rg-formation-carrier} & Gauge-invariant residuals; triangle energy splits into contrasts and flux; orthogonal split into normal, loop, gauge directions & \SOLID & Standard in substance \\
\MT{prop:rg-formation-laws} & Normalized joint, quadratic restriction cost, free energy as relative entropy to the optimizer, gap \(O(N\mathrm T^2/\mathrm J_1)\); with \mtc{prop:rg-formation-normalized-experiment} & \SOLID & Standard in substance \\
\MT{prop:rg-formation-parent-potential} & Parents are exact marginals with Schur stiffness and bounded tilt; free energy drops by exactly \(\mathrm T\Delta\); with \mtc{prop:rg-formation-rate-composition} and \mtc{prop:rg-formation-score-dissipation}, reversible inherited rows and an H-theorem & \SOLID & Standard in substance: Gaussian integration, chain rule \\
\MT{def:rg-formation-kinetics} & Langevin plus heat-bath label jumps; declared parent kinetics; own-law learning (\mtc{def:rg-formation-learning}) & \CHOICE & The parent kinetics is a Markov closure \\
\MT{def:rg-formation-regime} & Doubling recursion, amplitude, temperature \(\epsilon^{20N-34}\), horizon, prepared mean, calibration & \CHOICE & Where the outcome is put in \\
\MT{thm:finite-depth-adaptive-formation}(b) & Component counts \(3\cdot2^{L-1}\to\cdots\to1\) on the windows & \TOY & Triangle chain, \((\R,+)\), scalar fibers, supplied hierarchy; linear relaxation; constants depend on \(L\) \\
\MT{thm:finite-depth-adaptive-formation}(a), (d) & Learning follows calibration; attention mass and hidden kernels move after formation & \TOY & Movement engineered; no effect on formation \\
\MT{thm:finite-depth-adaptive-formation}(c), (e) & Discarded loss and score error vanish; histories couple; paths concentrate & \TOY & Resolution \(\sqrt{\mathrm T}/\epsilon\), not total variation of path laws \\
\MT{rem:rg-formation-saturation} & Record saturated, learning calibration-dominated, loops thermal; coherence unchanged without record and learning & \SOLID & Negative content fixing the theorem's reading \\
\MT{prop:rg-formation-gaussian-core} & Exact loss; two-sided asymptotic; base-\(4/3\) separation criterion & \TOY & Sharp, for the Gaussian core of the chain only \\
\MT{prop:rg-formation-endogenous-cycle} & A pair forms by attention feedback with probability at least \(4.6\times10^{-3}\); exact parent over one time unit & \TOY & Four agents; gating, seed, circulation, \(\mathrm T=\epsilon^8\), no learning; symmetric law \\
\MT{prop:rg-gated-depth-one} & Stationary blocks are minimum-edge-cover components; each merge costs \(\mathrm T\) & \SOLID & No-go inside the gated family only; stationary; \(\mathrm T\to0\) at fixed graph \\
\MT{open:rg-formation-regime-sharpness} & Is the doubling needed; depth-uniform constants & \OPEN & Chain family, base-\(4/3\) schedules \\
\MT{open:rg-formation-adaptive-family} & Learning or holonomy changes, in law, which block forms & \OPEN & Four-cycle with learned rows \\
\MT{open:rg-formation-general} & A recursive (R14) second coalescence that is an episode (R15) & \OPEN & Outside the gated family \\
\end{scorecard}

\begin{hygiene}
Conditions that decide the chain's outcome, several chosen precisely to switch off the feedback the
program is interested in: the supplied strengths with the doubling, the prepared displacement with a strict
margin, the tie \(\mathrm T=\mathrm A^2\epsilon^6\), the calibration record, the saturated record, and source
palettes confined to each triangle. Conditions that exclude the program's targets: the group \((\R,+)\) with
scalar fibers, unit mobility (models as fast as beliefs), and the uncomputed \(\epsilon_L\) with
depth-dependent constants. For the four-cycle the conditions that matter are the supplied rows of the table
in \cref{comp:ch13-cycle}; for the gated family they are exclusive single-label attention, a scalar
translation group, and temperature-independent energies and prior ratios.
Hygiene: the root observation and gauge weight, the cancellation term \(-\mathrm TR_{\mathrm{ev}}\), the
preparation neighborhood of smooth transports, positivity and non-explosion of densities, the choice of
entrance constant and tolerance above their thresholds, the Gaussian reference, and \(\mathrm T\le1\).
\end{hygiene}

\section{Toy versus general}
\label{comp:ch13-toy}

None of the three results has hypotheses that include the program's target, 512 or more agents with
nonabelian, noncompact groups. For the triangle chain the mechanism-level content concerns coarse-graining,
not formation: exact marked parents can be carried through a hierarchical relaxation with separately
vanishing errors, provided the stiffness separation beats the adiabatic lag in thermal units; that this
tension appears in every construction of the type is plausible but unproved. The doubling, the exponents
20, 6, and 14, the constants \(10^{-4}\), \([1,21/10]\), \(1/200\), and \(57/100\), the sigmoid prior, and the
chain topology are specifics, and the theorem, though valid for arbitrarily many agents, is not a
large-\(N\) result. For the four-cycle the mechanism-level content is the attention-contraction feedback,
with an order-one flux deciding which edges stretch; it is local and could operate in large graphs, but
whether it forms many pairs at once, how neighboring pairs interact, and whether the law picks them are not
addressed, and the floor \(3e^{-1/4}/512\), the rates 8 and 4, and the one-unit horizon are specifics. For
the gated family the mechanism-level content, that exclusive gating puts the label into a sheet's weight and
never its stiffness, is proved on every finite graph; it is the chapter's one result that is not about a
small network, and it does not reach non-exclusive or multi-head attention, temperature-dependent
couplings, interactions that stiffen the directions between blocks, nonabelian links, dynamics, or many
agents at fixed temperature.

On the reader's own intuitions the chapter is neutral. Formation in the chain is continuous: the dynamics is
smooth relaxation, and the discrete counts come from the resolution \(\eta_L\mathrm A\) and scale separation,
not from a transition. No phase transition, condensation, or thermodynamic limit appears; every result is at
fixed finite \(N\) with \(\mathrm T\) or \(\epsilon\) tending to zero. The gated no-go shows fragmentation at fixed
\(N\) and low temperature in one family, and rules condensation neither in nor out elsewhere or as
\(N\to\infty\) at fixed temperature.

\section{Open ends}
\label{comp:ch13-open}

The scope paragraph's declined claims map onto three open problems. Order-one or nonabelian holonomy, a
model-channel record, slow per-agent models, and a record informative at the coherence scale are absent
from the chain and reappear as the ingredient list of \MT{open:rg-formation-adaptive-family}: a family in
which learning or holonomy changes, in law, which components form needs at least one of a record
standardized at the coherence scale, for example through \(\chi(r/(\eta_L\mathrm A))\); model means and links
both slow relative to beliefs, since a fast link or fast model erases the model disagreement on the belief
clock (\MT{prop:obs-slow-model-link-dichotomy}; see \cref{comp:ch08}); loops of order \(\mathrm A\) or
nonabelian links; or interacting model-channel records. The problem is posed on the four-cycle with learned
attention rows, as a total-variation difference, bounded below, between the conditional laws of the first
formed block in the learned and frozen arms, or between initial means differing only in their loop
coordinate, or a proof that these differences vanish for almost every mean. In a slower-mobility regime the
learned and frozen first-block laws coincide in the limit, so learning selects nothing there
(\MT{prop:phil-early-block-null}), and a parent activated after hidden thermalization has separately
vanishing errors on the seeded sub-event (\MT{prop:phil-early-seeded-parent}); the original coupled mobility
and the loop comparison remain open (\cref{comp:ch14}).

Depth-uniform constants and a non-doubling regime are \MT{open:rg-formation-regime-sharpness}; a positive
answer would still leave a stiffness range growing like a power of \(N\). Arbitrary families and initial laws,
mixed-level source support, permanent membership, and control as the depth tends to infinity are folded into
\MT{open:rg-formation-general}: a second endogenous coalescence, with learned parameters and from
non-matching assignments, in a source-gated model with temperature-independent data outside the gated
family, recursive in the sense of R14, held in the sense of R15, and followed by parents with separately
vanishing errors, or a proof that none exists in the class. A dynamic second coalescence inside the gated
family falls between this problem and the depth-one proposition. Intrinsic selection of the contrasts is
\MT{open:cg-natural-selector}, a persistent multilevel hierarchy formed by descent is
\MT{open:full-graph-hierarchy}, and behind all of them stands \MT{open:rg-pointwise-parent-dynamics}: the
chapter's parents are declared Markov closures checked against the fine process in two special families,
which does not settle how a parent's dynamics follows from the network it summarizes.

\begin{verdictbox}{\NARROWER}
\textbf{What the chapter says it does.} Its opening asks whether coherent components can form in succession
in one normalized experiment while every parent is activated from the complete preceding law, keeps every
loop and label, and learns its own parameter, and says two witnesses answer yes under explicit restrictions.
The opening paragraph of the introduction describes a finite-depth family that forms successive meta-agents
under free-energy descent with learning, and the program aim is that collectives form through the agents'
coupled inference. The introduction's reading map, by contrast, describes the chapter accurately.

\textbf{What the statements give.} For the chain the question as posed is answered yes
(\MT{thm:finite-depth-adaptive-formation}): exact marginal parents that keep loops and labels and learn from
their own laws track the fine process with separately vanishing errors at every finite depth. The succession
itself is linear relaxation of a supplied displacement through a supplied stiffness hierarchy; attention is
within \(10^{-4}\) of its prior, learning follows calibration and moves after formation, loops are thermal,
and the coherence proof works with attention and learning switched off (\MT{rem:rg-formation-saturation}).
The meta-agents are marginal laws, not coarse agents; the temperature is \(\epsilon^{20N-34}\) and the
constants depend on the depth without control. The four-cycle witness (\MT{prop:rg-formation-endogenous-cycle})
and the depth-one no-go (\MT{prop:rg-gated-depth-one}) are delivered as stated.

\textbf{The gap.} Nothing shows a collective forming because of coupled inference beyond one level. The one
endogenous witness is a sampled-path event on four agents with a symmetric law and no learning, whose end
state equilibrium already favors, and the opening's ``two witnesses answer yes'' fits the chain but not the
four-cycle, which has neither succession nor learning. Recursive endogenous formation
(\MT{open:rg-formation-general}), law-level selection by learning or holonomy
(\MT{open:rg-formation-adaptive-family}), and the depth-uniform control that any statement about hundreds of
agents needs (\MT{open:rg-formation-regime-sharpness}) are open. The technical chapter's scope paragraph is
candid about all of this; it is the introduction's opening summary and the program aim that promise more
than these theorems deliver.
\end{verdictbox}

\chapter{Interpretation and Its Limits}
\label{comp:ch14}

This chapter covers \MT{ch:philosophy} together with the proof appendices that carry its four-cycle
results, from \MT{sec:phil-pairing-selection-proof} to \MT{sec:phil-early-parent}. The technical
chapter is short, but it is where the program's interpretive claims meet its most delicate formation
theorems, and eight of the propositions treated here are stated, not only proved, in those appendices.

\section{The question and the answer}
\label{comp:ch14-question}

The technical chapter asks what the preceding chapters commit to, and which readings of them an
observation or a proof could count against. Its answer has two layers that should be kept apart. The
first is interpretive and proves almost nothing: it sorts readings into declared rules of reading
(definitions), commitments the mathematics uses (standing hypotheses), and unsettled questions (open
problems). Its verdicts are that the architecture is local-first rather than background-free, that the
frame-change invariants are compatible with several structural-realist positions and select none, that
environment messages create no agents, and that nothing proved before the formation chapter
distinguishes the program from an ordinary hierarchical model with learned priors, whereas formation by
the inference could. The second layer is mathematical and substantial: a line of results on the
homogeneous four-cycle of \MT{def:rg-formation-cycle-experiment} asking whether learning, rather than a
random seed or a supplied hierarchy of strengths, decides which collective forms.

The answer of the second layer fits in one paragraph. On the four-cycle, formation in the technical
sense, meaning that a block enters coherence with a probability above its stationary value, happens
from every Gaussian initial law with a probability bounded below, but that probability is supplied by
the preparation's random initial assignment and source learning plays no role in it
(\MT{prop:phil-seeded-contrast-formation}). Law-level selection does happen when learning is slow: at a
small fixed temperature, as the learning mobility tends to zero, the learned source prior picks one of
the two pairings of the cycle for almost every Gaussian initial mean, while the arm with a frozen prior
ends symmetric (\MT{conj:phil-contrast-formation}, a theorem despite its label). That selection
amplifies an asymmetry already present in the initial law, it appears only at late times, and at those
times the selected pair's coherence is a property of equilibrium, not a formation event
(\MT{prop:phil-selected-pair-persistence}); on the early horizon where formation happens, slow learning
leaves the law of the first formed block unchanged (\MT{prop:phil-early-block-null}). Early formation
without law-level selection, late law-level selection without formation: the two halves do not meet,
and making them meet is the open problem \MT{open:rg-formation-adaptive-family}. The remaining results
concern the parent. The fixed parent channel of the four-cycle can be activated with separately
vanishing conditional and dynamical errors, both after the late selection and on the early seeded
episode, provided the activation time waits for the discarded modes to thermalize; activation at the
moment the coherence criterion is first met provably fails.

In the chain of the theory this chapter follows the formation chapter (\cref{comp:ch13}), whose
finite-depth theorem forms collectives through a supplied hierarchy of interaction strengths and whose
endogenous four-cycle witness forms one pairing from a symmetric preparation without learning. The
present chapter asks whether learning can do the choosing. \Cref{comp:tab-ch14-line} maps the
statements by time horizon and observable, the organizing principle the technical text uses but never
tabulates.

\begin{sciencetable}{The four-cycle continuation line. ``Early'' is the original episode horizon, of order $\log(1/\epsilon)$ at temperature $\mathrm T=\epsilon^8$; ``late'' is the selection window $t\sim\epsilon_{\mathrm L}^{-1}\log(1/\epsilon_{\mathrm L})$ at fixed or slowly vanishing temperature.}
\label{comp:tab-ch14-line}
\small
\begin{tabularx}{\textwidth}{>{\raggedright\arraybackslash}p{0.25\textwidth}>{\raggedright\arraybackslash}p{0.17\textwidth}>{\raggedright\arraybackslash}X>{\raggedright\arraybackslash}p{0.2\textwidth}}
\toprule
\textbf{Statement} & \textbf{Horizon} & \textbf{What is compared or bounded} & \textbf{Role of learning}\\
\midrule
\MT{prop:phil-seeded-contrast-formation} & early & probability of a contrast (entrance) event & idle \\
\MT{conj:phil-contrast-formation} & late, fixed $\mathrm T$ & one-time assignment law, learned versus frozen & amplifies the initial law \\
\MT{cor:phil-joint-selection} & late, $\mathrm T\to0$ jointly & same, along implicit mobility thresholds & same \\
\MT{prop:phil-temperature-tracking} & late & distance of the law from its instantaneous equilibrium & adiabatic following only \\
\MT{prop:phil-selected-pair-persistence} & late window $\mathrm T^{-1/4}$ & path law and coherent partition & fixes the label only \\
\MT{prop:phil-early-block-null} & early & conditional law of the first formed block & none if $\epsilon_{\mathrm L}=o(\mathrm T^{1/4})$ \\
\MT{prop:phil-actual-conditioned-parent} & late activation & discarded conditional information & parent inherits rows \\
\MT{prop:phil-actual-parent-dynamics} & late window $\mathrm T^{-1/4}$ & retained path law, parent versus fine & parent learns from its own law \\
\MT{prop:phil-early-seeded-parent} & early activation & discarded information, both directions & idle for $\epsilon_{\mathrm L}\in[0,1]$ \\
\MT{prop:phil-early-first-completion} & first completion & discarded information diverges & idle \\
\MT{prop:phil-early-parent-dynamics} & windows $\ll\mathrm T^{-1}$ & retained path law, parent versus fine & idle \\
\bottomrule
\end{tabularx}
\end{sciencetable}

\section{The picture}
\label{comp:ch14-picture}

\subsection{Scaffolding, idle posits, and what survives a change of frame}

Each agent carries two law points at the base point, a belief law and a model law, which the cross maps
do not identify (\MTeq{eq:phil-fiber-measures}). The base, the principal bundle, and the measurable
spaces are coordination scaffolding: they make the agents' interfaces comparable, but no inference law
constructs them and no law is placed on them. That is the meaning of local-first and not
background-free, and a physicist will recognize lattice gauge theory, where the lattice and the group
are fixed background and only site fields and links fluctuate. The verdict describes the present
architecture; it would change if edges became structural coordinates updated by inference
(\MT{open:obs-edge-creation}) or links were recovered from the records of a target population. Because the scaffolding
is not an inferential target, its lack of a law decides neither a Neo-Kantian reading (frames as
revisable constitutive conditions of meaning) nor a noumenal one (the bundle as structure reached only
through representations).

The chapter's one new definition is a parsimony rule: a posit is idle when changing it alters the law
of no record and of no statistic of a record in the model, and the idle-wheel criterion
(\MT{def:phil-idle-wheel}) removes idle posits. It plays the role gauge redundancy plays in physics, but
it is a declared rule of reading, close in spirit to constructive empiricism without being derived from
it, and it is relative to the records a model declares. What survives it is structural data,
\[
\mathcal I_{\mathrm{struct}}
=\Bigl\{
\text{ELBO and relative-entropy values},\
[\mathrm{Hol}^b(p)]_{\rm conj},\
[\mathrm{Hol}^m(p)]_{\rm conj},\
\bigl[(\mathring\Theta^b_e,\mathring\Theta^m_e)_{e}\bigr]
\Bigr\}
\]
(\MTeq{eq:phil-invariants}): scalars invariant because the kernels are covariant
(\MT{prop:gen-product-evidence-invariance}), loop classes in both channels, invariant because a passive
re-expression conjugates a loop product, and the per-edge twists (each link expressed in its endpoints'
own frames) modulo one simultaneous conjugation, which is finer than the loop classes because it keeps
how the twists of different edges sit relative to one another. For abelian links, as on the four-cycle,
the loop class is the loop sum $\sum_i\omega_i$. On a graph-flat family the loop classes are trivial
(\MT{prop:geo-trivializing-criterion}) and give a noumenal reading nothing to do; a nontrivial loop class
is a referent inside the interaction complex, which a declared record at $c_*$ can read
(\MT{thm:cg-loop-operational-holonomy}). The technical text finds these invariants compatible with
epistemic, moderate, and ontic structural realism and gives no argument from theory change, since they
are invariants of frame change inside one theory, whereas Worrall's structure \citep{Worrall1989}
survives a change of theory.

\begin{analogy}{exact}
The loop classes $[\mathrm{Hol}^x(p)]_{\rm conj}$ are the gauge-invariant content of a closed loop in
lattice gauge theory, the data of a Wilson loop before the trace. On the interaction graph this is an
identity of structure: links are group elements on edges, a passive re-expression acts as
$\Theta_e\mapsto g_i\Theta_eg_j^{-1}$, and only conjugacy classes of loop products are invariant
(\MT{prop:geo-edge-twist-invariance}). Nothing carries over to a continuum connection: graph-link
holonomy is neither base-connection holonomy nor bundle topology, and the technical text keeps these and
gauge redundancy apart.
\end{analogy}

\subsection{An inferred link field and its physical reading}

\Cref{comp:lie-links} changes the status of the links within this architecture. They are now latent
Lie-algebra generators under a declared normalized prior, rather than necessarily held background
values. Their Gaussian recognition stationarity balances prior stress against a matter current
(\MT{prop:link-field-equation}); for a comparison record that current has the outer-product form
of a delta rule (\MT{cor:link-delta-rule}). The generator history prior couples neighboring
episodes, and its linearization near zero has the form of a massive Euclidean vector operator, with
matter contributing a positive screening mass (\MT{prop:link-linearized-equation}). The central
equations and their approximation domain are explained in \cref{comp:lie-links} and their Gaussian
values in \cref{comp:gaussian-latent-links}.

\begin{analogy}{structural}
\textbf{A Proca-type operator on the interaction graph.} The cycle operator supplies a spatial
curl-type stiffness on graph copies, the history increment supplies a second difference in the
episode index, and the generator precision and matter response supply mass terms. This is the
form of a massive Euclidean field equation, not a derivation of a spacetime Proca field. The
episode index is an inference index. The most-probable-history inertia belongs to smoothing a
joint history; the Fisher recognition flow is first order in its own flow time
(\MT{prop:link-map-inertia}, \MT{prop:link-recognition-flow}).
\end{analogy}

The scaffolding verdict is unchanged: the base, principal bundle, frames, and reference point are
structural inputs that this inference does not construct (\MT{sec:phil-base}). Inferring links
within a declared model is a mathematical construction; recovering links from a target system's
records requires the separate operational evidence of \MT{open:phil-testability}. The held-link
model is recovered in the zero-width free-energy limit under its stated hypotheses
(\MT{thm:link-held-limit}), which connects the architectures without establishing a physical
action principle, an empirical field, or the formation mechanism below.

\subsection{Participation and why it does not yet distinguish the program}

Participation, in the technical sense, is reciprocal inferential dependence through records inside one
joint law, with no live posterior feeding a same-step generative factor (\MT{req:gen-typing-prohibition}).
It occurs in the same-time inference loop (block coordinate ascent, as in any connected hierarchical
model), in parameter learning (whose optimal prior over an unrestricted family is the aggregate
recognition law, \MTeq{eq:obs-aggregate-prior}), and in delayed feedback through a declared kernel (still
the free energy of one joint, \MT{prop:full-graph-episode-lifted-joint}). The verdict is blunt: none of
these separates the program from a hierarchical model with learned priors, and the resemblance to
Wheeler's participatory proposal \citep{Wheeler1990} is neither derived from it nor evidence for it.
Formation could separate them: collectives constituted by the coupled inference, held by a restoring
interaction, and dependent on the preceding state. Hence the formation theorems in this chapter.

\subsection{The four-cycle: a ring of agents choosing whom to listen to}

Four agents sit on a ring. Each has a scalar Gaussian belief (mean $\bar k_i$, unit variance) and a
scalar Gaussian model (mean $\bar m_i$, variance $1/2$); each edge carries a translation link $\omega_i$
in the abelian group $(\R,+)$. The cost of an edge is the transported relative entropy of its residuals.
Each agent listens to exactly one neighbor, and an edge carries energy only when an endpoint listens to
it (source gating). Of the sixteen listening assignments, the two perfect matchings, which split the ring
into two mutually attending pairs, are the candidate collectives. The kinetics is Langevin diffusion at
temperature $\mathrm T$ on the energy sheet of the current assignment, interrupted by a rate-one clock at
which every agent redraws whom it listens to from its posterior. Everything is symmetric under the
rotations and reflections of the ring.

\emph{Seeded formation.} From a Gaussian configuration and a uniformly random assignment, with small
probability the assignment is already a matching and the configuration lies near a seed point where every
edge is equally mismatched; the attended residuals then relax, the unattended ones grow, and two pairs
form. The probability is that of a lucky initial condition, independent of the small parameter and
unchanged if the prior learns (\MT{prop:phil-seeded-contrast-formation}).

\emph{Selection by slow learning.} Now let the listening prior learn slowly from the population's
recognition law. Its rest point is the recognition marginal, so learning has the form of the
expectation-maximization update of mixture weights, run in continuous time; averaged over the fast
kinetics, the listening probabilities climb the logarithm of an evidence function in the Fisher metric of
the simplex. At low temperature that evidence is largest at the two pure matchings and has a saddle at
the uniform prior, unstable exactly in the direction that favors one matching over the other. An
asymmetric Gaussian initial law leaves, during the early transient, an asymmetric imprint on the rows of
the order of the mobility; the saddle amplifies it, and after a time of order
$\epsilon_{\mathrm L}^{-1}\log(1/\epsilon_{\mathrm L})$ the prior sits at one matching while the
recognition law tracks the corresponding equilibrium. The frozen arm has nothing to amplify and relaxes
to its symmetric equilibrium (\MT{conj:phil-contrast-formation}).

\begin{analogy}{structural}
The theorem has the dynamical signature of mean-field symmetry breaking. The learned four-cycle is a
nonlinear Markov process, whose generator depends on its own law through the rows, as in the
McKean--Vlasov limit of the Curie--Weiss model; the rest point of learning is a Curie--Weiss-type
self-consistency (the listening marginal computed under the rows equals the rows), and below the
theorem's temperature threshold its symmetric solution is unstable along one direction. The limits fail
to commute in the familiar way: at fixed time the rows move by at most $\epsilon_{\mathrm L}t/2$ (part (a)
of \MT{prop:phil-seeded-contrast-formation}), so $\epsilon_{\mathrm L}\to0$ first and then $t\to\infty$
gives the symmetric frozen answer, while $t\propto\epsilon_{\mathrm L}^{-1}\log(1/\epsilon_{\mathrm L})$
gives a selected matching. The imprint, of the order of the mobility, acts as a vanishing field whose
effect survives because its amplification time diverges.

It is not spontaneous symmetry breaking in the strict sense. There is no thermodynamic limit, and the
frozen kinetics has a unique invariant law (the proof uses the Harris theorem for exactly this). Nothing
is spontaneous either: the rows learn from the law, so they are deterministic and carry no fluctuations
that could choose a branch, and the branch is fixed by a definite datum, the sign of an analytic function
of the initial mean. A symmetric initial law keeps the rows uniform forever (part (c) of
\MT{prop:phil-seeded-contrast-formation}), as a mean-field magnetization started at zero stays at zero.
The accurate description is explicit symmetry breaking by the initial data, amplified by a dynamical
instability; the technical text's word is ``amplifies''.
\end{analogy}

\begin{analogy}{structural}
Why matchings win at low temperature: every listening sheet has the same minimum energy, zero, at the
same configuration, so the sixteen sheets are degenerate. They differ in how many directions they
stiffen: a matching stiffens four residual directions, a path six, an oriented cycle seven (the residuals
around a loop obey one linear relation). Each extra stiff direction costs a factor of order
$\sqrt{\mathrm T}$ in thermal phase-space volume, so the evidence peaks at the matchings. Selection among
degenerate ground states by the volume of thermal fluctuations is the mechanism of order by disorder; the
technical documents do not use that term.
\end{analogy}

\emph{Late persistence, early null, and the parent.} Once the prior sits at a matching, the selected pairs
stay coherent, but the stationary law at those rows realizes the same persistence with probability
tending to one, and the frozen arm shows persistent pairs too, split half and half between the two
pairings (\MT{prop:phil-selected-pair-persistence}). Late, learning changes which pairing, not whether
there is pairing. Early, a slow prior has barely moved, and the law of the first formed block is the same
with or without learning (\MT{prop:phil-early-block-null}). The parent of a pairing is obtained by a fixed
channel that keeps the assignment, the loop coordinate $\sum_i\omega_i$, the slow directions, and one
fast mode (the sum of the two pairs' average residuals), and discards three fast modes (the difference of
the pairs' average residuals and each pair's belief-model residual difference). The parent is exact as a
summary when the discarded modes, given the kept ones, are thermal, and exact as a dynamics when its own
evolution reproduces the kept trajectories. Both hold after the late selection and, on the early seeded
episode, at deterministic activation times after the discarded modes have relaxed to the thermal scale;
the first moment the coherence criterion is met is too early by at least a factor of two in logarithmic
time, and there the discarded information diverges (\MT{prop:phil-early-first-completion}).

\begin{analogy}{exact}
The discarded conditional relative entropy is the error term of integrating out fast variables. The
chain rule $\KL(q\Vert\mu)=\KL(C_\#q\Vert C_\#\mu)+\mathcal H_{\mathrm{opt}}(q)$
(\MTeq{eq:phil-conditioned-parent-chain}) is an identity, so the fine divergence is the parent's
divergence plus a remainder that vanishes exactly when the fast modes are in conditional equilibrium given
the slow ones: adiabatic elimination written as an identity. The dynamical statements are the structural
analogue of a valid reduced Langevin equation, proved only on stated windows. Those windows lie on a
plateau between the thermalization time of the discarded modes, $\log(1/\epsilon)=\tfrac18\log(1/\mathrm
T)$, and the slow time $\mathrm T^{-1}$ on which weakly restored directions equilibrate and the pairing
can change; a quasi-stationary state between fast relaxation and slow escape is what a physicist calls
metastability.
\end{analogy}

\section{Objects and notation}
\label{comp:ch14-notation}

\begin{decoder}
Idle posit & A posit changing the law of no record or record statistic (\MT{def:phil-idle-wheel}). & Gauge redundancy; unobservable absolute frame. \\
$[\mathrm{Hol}^x(p)]_{\rm conj}$ & Conjugacy class of the product of links around closed walk $p$ in channel $x$. & Wilson loop before the trace. \\
Frame, law, mean agreement & Three inequivalent notions of agents agreeing (\MT{prop:geo-three-agreements}). & Pure gauge; equal states up to transport; equal first moments. \\
$\mathbf x=(\bar k,\bar m,\omega)$ & Four-cycle configuration: belief means, model means, link translations; twelve coordinates. & Site fields and abelian link fields. \\
$r^b_i$, $r^m_i$, $\mathcal D_{e_i}$ & Belief and model residuals on edge $i$ and the edge cost, a transported relative entropy. & Covariant difference and its energy. \\
$S$, $\sigma_i$ & Listening assignment; $\sigma_i=+1$ when $i$ listens to $i+1$. & Discrete label field. \\
$S^{(0)}$, $S^{(1)}$, matching & The two assignments that split the ring into two mutually attending pairs. & The two dimer coverings of a plaquette. \\
$\pi_{ij}$, $p_i=\pi_{i,i+1}$ & Learned source prior (listening probabilities), one row per receiver. & Mean-field order parameter. \\
$\epsilon_{\mathrm L}$, $\mathrm T$ & Learning mobility; temperature ($\mathrm T=\epsilon^8$ in the original regime). & Slow relaxation rate; temperature. \\
$q^{\epsilon_{\mathrm L}}_t$, $q^0_t$ & Recognition law of the kinetics, learned arm and frozen arm (rows fixed at $1/2$). & Ensemble with and without feedback. \\
$q^\star$, $\mu_{\mathrm T,p}$ & Restricted optimizer $\propto\pi_Se^{-U_S/\mathrm T}$, the invariant law at fixed rows $p$. & Gibbs state at fixed couplings. \\
$\nu^M_{\mathrm T}$ & Gaussian law on the sheet of matching $M$, assignment held at $M$. & Pure phase. \\
Contrast formation & A block internally coherent and externally distinct, entered after time zero. & Nucleation of a bound pair. \\
$\mathsf j_{\mathrm T}(\mathsf m_0)$ & Transient imprint: integrated early asymmetry of the frozen law along the matching direction. & Integrated response to an initial perturbation. \\
Restored equilibrium coherence & Coherence with a restoring interaction that is typical at stationarity; not formation. & Equilibrium order. \\
$C$, $W_M$, $Z_M$ & Parent channel; kept subspace (nine dimensions); discarded subspace (three fast modes). & Blocking map; kept and integrated-out variables. \\
$\mathcal H_{\mathrm{opt}}$, $\mathcal H_{\mathrm{phys}}$ & Mean discarded conditional relative entropy against the optimizer or the physical posterior. & Error of adiabatic elimination. \\
Retained-path TV & Total variation between the parent's path law and the fine kept-coordinate path law. & Validity of a reduced dynamics. \\
$\theta$, $\mathsf t_{\mathrm c}$ & Deterministic activation time; first qualifying completion time. & Fixed measurement time; first-passage time. \\
\end{decoder}

Several words carry nonstandard senses. ``Episode'' means a coherence episode in the sense of rule R15,
not a learning episode. ``Formation'' is always a threshold event at declared resolutions (relative
entropies below $\epsilon^4$ inside, above $1/2$ across the boundary, with an entrance clause), and
``selection'' always a property of the law, never of a sample path. ``Actual'' means the law the kinetics
really produces, not a Gaussian or stationary substitute. ``Physical'' refers to the generative posterior,
``optimizing'' to the restricted optimizer. The regime parameter appears as $\epsilon$ in the appendices,
$\xi$ in the chapter text, and $\mathrm T^{1/8}$ in \MT{prop:phil-early-block-null}; the original coupled
regime sets the mobility $\epsilon_{\mathrm L}$ equal to it. The label
\texttt{conj:phil-contrast-formation} survives from when the statement was a conjecture; the printed
heading says Theorem, and it is proved.

\section{The formulas, explained}
\label{comp:ch14-formulas}

\subsection{Energy, invariant law, and learning rule}

On edge $e_i$ the residuals are $r^b_i=\bar k_{i+1}-\bar k_i-\omega_i$ and
$r^m_i=\bar m_{i+1}-\bar m_i-\omega_i$, and
\[
\mathcal D_{e_i}(\mathbf x)=\tfrac12\bigl[(r^b_i)^2+(r^m_i)^2\bigr]
=\KL\bigl(q^b_{i+1}\Vert(\Theta_{e_i})_\#q^b_i\bigr)
+\sigma^2_{\mathrm{ev}}\KL\bigl(q^m_{i+1}\Vert(\Theta_{e_i})_\#q^m_i\bigr)
\]
(\MTeq{eq:rg-formation-cycle-cost}): the relative entropy between one agent's law and its neighbor's law
carried across the link. The sheet energy $V_S=\sum_i\mathcal D_{e(i,S_i)}$ counts an edge twice if both
ends attend to it and not at all if neither does. The restricted optimizer, which is the invariant law of
the kinetics at a fixed prior, is
\[
q^\star(\mathbf x,S)\propto\pi_Se^{-U_S(\mathbf x)/\mathrm T},
\qquad
U_S=V_S+\tfrac{\mathrm T}2|\mathbf x|^2+\tfrac{\mathrm T}2\mathbf x^\top R_{\mathrm{ev}}\mathbf x.
\]
The first term is the interaction; the second comes from the unit Gaussian prior on configurations and
gives every direction a weak restoring rate of order $\mathrm T$; the third comes from the model evaluator
and penalizes belief-model discrepancy, and the physical posterior omits it. At zero temperature every
sheet has a flat zero-energy valley; at finite temperature each sheet is a Gaussian whose weight is its
phase-space volume.

The learning rule of the source prior is
\[
\dot\pi_{ij}=\epsilon_{\mathrm L}\bigl(q_t(S_i=j)-\pi_{ij}\bigr)
\]
(\MTeq{eq:phil-source-learning}): the rows relax at rate $\epsilon_{\mathrm L}$ toward the current
recognition marginal. It is the natural-gradient flow of the free energy in the prior row under the Fisher
metric of the simplex, with the recognition label marginal as rest point, as for the aggregate prior. The
right side uses the law $q_t$, not a sample, so the rows are deterministic; at $\epsilon_{\mathrm L}=0$
they stay at $1/2$, the frozen arm against which every selection claim is measured. Part (a) of
\MT{prop:phil-seeded-contrast-formation} gives well-posedness and $|\pi_{ij}(t)-1/2|\le\epsilon_{\mathrm
L}t/2$, so nothing learned can matter on horizons short compared with $1/\epsilon_{\mathrm L}$.

\subsection{Where the probability of formation comes from}

Part (b) of \MT{prop:phil-seeded-contrast-formation} bounds the probability of contrast formation below by
\[
\frac{e^{-1/4}}{16}\sum_{\pm}\mathcal N(\mathsf m_0,I_{12})
\bigl\{|\mathbf x\mp\mathbf x^\circ|<\varrho_\circ\bigr\},
\qquad \mathbf x^\circ=(0,0,\mathbf 1_4),\quad \varrho_\circ=\tfrac1{100},
\]
uniformly in the regime parameter. The factor $1/16$ is the probability that the uniform initial
assignment is a given matching; $e^{-1/4}$ is the probability that the redraw clock is silent until time
$1/4$, long enough for a cost gap to open; the Gaussian mass is the probability of starting within $1/100$
of a seed point, where all means vanish and every link translation is $\pm1$, so every edge is equally
mismatched. On this sub-event the mechanism of \MT{prop:rg-formation-endogenous-cycle} runs: attended
residuals decay at rates of order one, unattended ones grow toward $3/2$ in magnitude, and posterior
redraws keep the matching except with exponentially small probability. The bound holds equally with
frozen rows, and the technical text draws the right conclusion: a claim that learning creates a contrast
must compare the arms at the level of the law. Two calibrations: the explicit estimate
$e^{-1/4}\varrho_\circ^{12}/737280$ times a Gaussian factor is about $10^{-31}$ at $\mathsf m_0=0$, so its
only role is to be positive and $\epsilon$-independent; and the seeded sub-event is genuine formation
under R15, with a restoring interaction and a probability exceeding its value from the invariant law, at
most $4e^{-1/(4\mathrm T)}$, by a margin bounded below.

\subsection{The row flow, its saddle, and the imprint}

Averaging the learning rule over the invariant law at fixed rows gives the reduced field
\[
\mathscr V_i(p)=\bar\beta_i(p)-p_i=p_i(1-p_i)\partial_{p_i}\log\mathsf Z(p),
\qquad
\mathsf Z(p)=\sum_\sigma\pi_\sigma(p)\mathsf w_\sigma,
\qquad
\mathsf w_\sigma=\det(\mathsf A_\sigma/\mathrm T)^{-1/2}
\]
(\MTeq{eq:phil-selection-reduced-field}). Here $\bar\beta_i(p)$ is the stationary probability that
receiver $i$ listens to $i+1$, $\mathsf w_\sigma$ is the Gaussian volume of sheet $\sigma$, and
$\mathsf Z$ is the evidence of the prior. The factor $p_i(1-p_i)$ is the inverse Fisher metric of a
Bernoulli row, so the reduced flow is the replicator equation with fitness $\partial_{p_i}\log\mathsf Z$,
natural-gradient ascent of the log evidence. With the matching weight normalized to one, paths have
relative weights $\mathsf r_E,\mathsf r_N\approx\sqrt{40/7}\mathrm T$ and oriented cycles
$\mathsf r_C\approx\sqrt{5/3}\mathrm T^{3/2}$, the count of stiff directions made exact. At the uniform
prior the linearized field has one positive eigenvalue $\alpha\ge7/5$ along $\mathbf v=(1,-1,1,-1)$,
toward one matching and away from the other, and negative eigenvalues elsewhere
(\MTeq{eq:phil-selection-saddle-spectrum}); both matching vertices are sinks. The threshold
$\mathrm T_0$ is a temperature below which the sheet ratios are at most $1/8$: a sufficient condition of
the proof, not a located critical temperature, and nothing is proved above it.

The seed that pushes the rows off the saddle is the transient imprint
\[
\mathsf j_{\mathrm T}(\mathsf m_0)=\mathbf v\cdot\mathsf I(\mathsf m_0),
\qquad
\mathsf I_i(\mathsf m_0)=\int_0^\infty\bigl[q_t^0(\sigma_i=+1)-\tfrac12\bigr]dt
\]
(\MTeq{eq:phil-selection-imprint}), the time-integrated early asymmetry of the frozen law. It is an entire
function of the initial mean and changes sign under rotation of the ring by one site, so it vanishes at
every rotation-invariant mean, the seed points included. The proof shows it is not identically zero, so
its zero set is Lebesgue-null: that is the excluded set of means. After the transient the rows sit at
$1/2+\epsilon_{\mathrm L}\mathsf I(\mathsf m_0)+o(\epsilon_{\mathrm L})$; the component along
$\mathbf v$ grows like $\epsilon_{\mathrm L}e^{\alpha\epsilon_{\mathrm L}t}$ and reaches order one at
$t\approx\alpha^{-1}\epsilon_{\mathrm L}^{-1}\log(1/\epsilon_{\mathrm L})$, hence the window
$t_{\mathrm L}=\mathrm C_{\mathrm T}\epsilon_{\mathrm L}^{-1}\log(1/\epsilon_{\mathrm L})$ with
$\mathrm C_{\mathrm T}>1/\alpha$. The theorem states: there is $\mathrm T_0>0$ such that for every
$\mathrm T\in(0,\mathrm T_0]$ and Lebesgue-almost every $\mathsf m_0$ one matching $M$ satisfies
\[
\lim_{\epsilon_{\mathrm L}\to0}\ \sup_{t_{\mathrm L}\le t\le2t_{\mathrm L}}
q^{\epsilon_{\mathrm L}}_t(S\ne M)=0,
\]
while for every $\mathrm T>0$ and every $\mathsf m_0$ the frozen difference
$q^0_t(S=S^{(0)})-q^0_t(S=S^{(1)})$ tends to zero as $t\to\infty$. The order of limits is fixed
temperature, then mobility to zero, with the window scaled by the mobility; the conclusion concerns the
one-time law of the assignment on a moving window, not histories, events, or $t\to\infty$ at fixed
mobility.

\subsection{What the corollary and the tracking bound add}

\MT{cor:phil-joint-selection} lets the temperature vanish too. Along any $\mathrm T_n\downarrow0$ with
tolerances $\eta_n\to0$ there are mobility thresholds $d_n\downarrow0$, the same for almost every mean,
below which the selection error on the window (now with $\mathrm C_{\mathrm T}=2$ for all temperatures)
is at most $\eta_n$, the learned law is $\eta_n$-close in relative entropy to its instantaneous
equilibrium, and the frozen law is $\eta_n$-close to its symmetric equilibrium, so the frozen matching
probabilities tend to $1/2$. It adds a single joint path $(\mathrm T(\epsilon_{\mathrm L}),
\epsilon_{\mathrm L})\to0$ and an assignment-law contrast tending to $1/2$ in total variation. It adds no
rate: the thresholds exist by a measure-theoretic argument and have no values, they can be taken below
$\mathrm T_n^a$ for any $a>7/2$ only as a cap, and the matching may change with $n$, because nothing
controls the imprint's sign as the temperature vanishes.

\MT{prop:phil-temperature-tracking} makes the adiabatic half explicit. For $\mathrm T\le1/100$ the frozen
entropy production dominates $\mathrm c_*\mathrm T^{5/2}/[4+3\log(1/\mathrm T)+8\log(1/p_{\min})]$
times the relative entropy, and
\[
\operatorname{TV}(q_t,\mu_{p(t)})
\le e^{-\ell t/2}\sqrt{\frac{13(12+|\mathsf m_0|^2)}{4\mathrm T}}
+\frac{2\epsilon_{\mathrm L}\mathrm D(\mathrm T,\epsilon_{\mathrm L})}{\mathrm c_*\mathrm T^{7/2}},
\qquad
\ell=\frac{\mathrm c_*\mathrm T^{5/2}}{\mathrm D(\mathrm T,\epsilon_{\mathrm L})}
\]
(\MTeq{eq:phil-temperature-tracking}), with $\mathrm D$ logarithmic. The first term is memory of the
initial condition, erased at the mixing rate $\ell$; the second is the lag of a law chasing an equilibrium
that moves at a rate of order $\epsilon_{\mathrm L}/\mathrm T$, divided by the mixing rate, which is where
the exponent $7/2$ comes from. For $\epsilon_{\mathrm L}=\mathrm T^a$, $a>7/2$, tracking holds on the
window: an explicit adiabatic condition, not a selection result, since whether the equilibrium has moved
to a matching depends on the imprint, for which no lower bound as $\mathrm T\to0$ is available.

\subsection{Persistence, and why it is equilibrium}

\MT{prop:phil-selected-pair-persistence} assumes the learned law is $\delta$-close to $\nu^M_{\mathrm T}$
and the frozen law $\delta$-close to its symmetric equilibrium. Over a window $\mathsf h$ the learned path
law is then within $(\mathsf h+2)\delta$ of a stationary Ornstein--Uhlenbeck path held at $M$, the frozen
one within $\delta+16\mathrm T+31\mathrm T\mathsf h$ of an equal mixture of the two held-matching paths,
and coherence at resolution $\mathsf a$ fails with probability at most
\[
(\mathsf h+2)\delta+\mathsf b(\mathrm T,\mathsf a,\mathsf h),
\qquad
\mathsf b=8\sqrt{2/\pi}\mathsf a+(8+3400\mathsf h)\mathrm T/\mathsf a^2
\]
(\MTeq{eq:phil-selected-coherence-persistence}). The first term of $\mathsf b$ is the chance that an
exterior residual, of variance of order one, is accidentally below $\mathsf a$; the second is the chance
that thermal noise pushes an internal residual above $\mathsf a$ during the window. The choice
$\mathsf a=\mathrm T^{1/4}$, $\mathsf h=\mathrm T^{-1/4}$ puts the resolution far above the thermal scale
(restoring ratio $4\mathsf a^2/\mathrm T\to\infty$) and the window far below the hopping time
$\mathrm T^{-1}$; the error is then of order $3400\mathrm T^{1/4}$, small only far below
$\mathrm T=10^{-14}$. The last clause decides the classification: with rows near the matching vertex, the
stationary law at those frozen rows realizes the same persistence event with probability at least
$1-4(\mathsf h+2)\delta-\mathsf b$. The event is typical at equilibrium, fails the horizon clause of R15,
and is restored equilibrium coherence. The coherent-partition laws of the two arms differ by nearly $1/2$
because the frozen arm splits its equilibrium pairs evenly (\MTeq{eq:phil-coherent-partition-contrast}).

\subsection{The early first-block null}

\MT{prop:phil-early-block-null} keeps the original contrast detector on the horizon
$H_{\mathrm T}=2\mathrm T^{-1/8}$ and takes the block that completes first. Coupled with common noise,
clock, and initial sample, the arms differ only when a redraw differs, and a row displaced by
$\epsilon_{\mathrm L}t/2$ changes a redraw probability by at most $\epsilon_{\mathrm L}t$. The conditional
first-block laws differ by at most $\mathsf d_{\mathrm T}/\mathsf c(\mathsf m_0)$, with
$\mathsf d_{\mathrm T}=\min\{1,8\epsilon_{\mathrm L}\mathrm T^{-1/4}\}$ and $\mathsf c$ the seeded bound
above (at most of order $10^{-30}$), so they agree when $\epsilon_{\mathrm L}=o(\mathrm T^{1/4})$, at
extraordinarily small temperatures; at the original mobility $\mathrm T^{1/8}$ the bound is vacuous.

\subsection{Conditioning, and what a vanishing discarded term means}

The late parent starts from the full recognition law conditioned on a completed coherence prefix $E$ on
$[t_{\mathrm L},t_{\mathrm L}+1]$, pushed through the channel. \MT{prop:phil-actual-conditioned-parent}
bounds the conditioned relative entropy and the optimizing discarded term by
\[
\kappa_E=\frac{\eta+h_b(\rho)}{1-\rho},
\qquad
h_b(r)=-r\log r-(1-r)\log(1-r),
\]
and the physical discarded term by $\kappa_E+\mathrm T$, where $\eta$ bounds the unconditioned relative
entropy to equilibrium and $\rho$ the failure probability of the prefix. The binary entropy is the price
of conditioning, small only for a typical event. \MT{rem:phil-conditioned-hidden-sign} shows the price is
real: conditioning a thermal Gaussian on the sign of one discarded coordinate costs $\log2$ at every
temperature, and conditioning on the whole realized history pins the discarded coordinates and makes the
relative entropy infinite. A parent is activated from a law given a high-probability event, never from
one history. With $\eta,\delta\le\mathrm T$ and $\epsilon_{\mathrm L}\le\mathrm T^2$,
$\kappa_E=O(\mathrm T^{1/4}\log(1/\mathrm T))$.

\MT{prop:phil-actual-parent-dynamics} compares paths. The parent diffuses on the kept subspace and, at each
ring of its own clock, reinserts the discarded coordinates as fresh thermal draws, redraws the assignment
from the posterior at the reconstructed configuration, and projects back; its rows learn from its own law.
Its path law, including the assignment history and hence the loop, is within
\[
2\rho+(\mathsf h+3)\delta+\frac{4\mathsf hd}{1-4d},
\qquad d=\delta+\epsilon_{\mathrm L}\mathsf h,
\]
of the conditioned fine kept-coordinate path law (\MTeq{eq:phil-actual-parent-path-error}), which is
$O(\mathrm T^{1/4})$ for $\mathsf h=\mathrm T^{-1/4}$: the cost of conditioning, the distance of both
processes from the held-matching stationary path, and the probability that the parent leaves the matching.

\subsection{The activation clock}

On the early seeded episode, at $\mathrm T=\epsilon^8$ and any $\epsilon_{\mathrm L}\in[0,1]$,
\MT{prop:phil-early-seeded-parent} bounds the discarded information at a deterministic activation time
$\theta$ from both sides. Asymptotically,
\[
\mathcal H_{\mathrm{opt}}(q^\circ_\theta)\le\tfrac92\varrho_\circ^2e^{-8(\theta-\log(1/\epsilon))}+o(1),
\qquad
\mathcal H_{\mathrm{opt}}(q^\circ_\theta)\gtrsim\frac{\varrho_\circ^2}{7}\mathrm c_{\mathsf m}
e^{8(\log(1/\epsilon)-\theta)}-\log2
\quad(\theta\le\log(1/\epsilon)),
\]
with $q^\circ_\theta$ the law conditioned on the seeded prefix and $\mathrm c_{\mathsf m}>0$ depending
only on the initial mean. The physics is one line. The discarded discrepancy coordinates start at offsets
of order $\varrho_\circ$ and decay at rate about $4$, while their thermal width is of order
$\sqrt{\mathrm T}=\epsilon^4$; a displaced Gaussian has relative entropy of order squared offset over
variance, $\varrho_\circ^2e^{-8\theta}/\mathrm T=\varrho_\circ^2e^{-8(\theta-\log(1/\epsilon))}$. So the
discarded information is large before the thermalization time $\log(1/\epsilon)$ and small after it. The
upper bound needs $\theta-\log(1/\epsilon)\to\infty$ and $\epsilon^4\theta\to0$, the lower bound diverges
when $\log(1/\epsilon)-\theta\to\infty$, and the window $\theta=\log(1/\epsilon)+O(1)$ is undecided. The
original witness of \MT{prop:rg-formation-endogenous-cycle} activates at $2\log(1/\epsilon)$, safely on
the good side.

\MT{prop:phil-early-first-completion} treats the clock a physicist might pick first: the first time any
block meets the contrast criterion. That criterion asks for residuals below about $\epsilon^2$, far
coarser than the thermal scale $\epsilon^4$. Seeded residuals fall below $\epsilon^2$ by about
$\tfrac12\log(1/\epsilon)$, and near $\tfrac14\log(1/\epsilon)$ when a pair's initial discrepancy is small,
so the first completion occurs by $\theta_\circ=\tfrac12\log(1/\epsilon)+1$. There the discarded
information is at least of order $\mathrm T^{-5/8}$, driven by seeds whose first pair completes near
$\tfrac14\log(1/\epsilon)$ while the other matched pair, whose discrepancy the channel also discards,
still carries an imprint of order $\varrho_\circ\epsilon$. The divergence persists when the law is
conditioned on the whole contrast event. This is an established negative result: no parent of this type
activated at the first qualifying completion has a vanishing conditional error, in either arm and at any
mobility in $[0,1]$.

\MT{prop:phil-early-parent-dynamics} bounds the path error of a parent activated at a good deterministic
time by
\[
24\exp\Bigl[-\frac1{76800\mathrm T\mathsf h}\Bigr]
+4\bigl(1+\sqrt{8/3}\bigr)\Lambda_{\mathrm L}(\theta+\mathsf h)\mathsf he^{-1/\mathrm T},
\qquad
\Lambda_{\mathrm L}(t)=2e^{\epsilon_{\mathrm L}t}-1,
\]
uniformly in the mean and in $\epsilon_{\mathrm L}\in[0,1]$, with $\Lambda_{\mathrm L}$ bounding how far
learned rows can tilt a redraw. The first term is the probability that thermal diffusion over the window
moves the configuration a distance $1/20$ off its deterministic path, negligible when
$\mathrm T\mathsf h\to0$; the second is the Arrhenius probability of a redraw across a cost gap larger than
one. Windows up to $\mathrm T^{-1+\mathrm q}$, $\mathrm q\in(0,1)$, are covered, and nothing is concluded
at order $\mathrm T^{-1}$. The late proof cannot be reused here because the early activation law is at
total-variation distance essentially one from the stationary kept-coordinate Gaussian
(\MT{rem:phil-early-stationary-mismatch}), so the early proof follows the energy gap along the
conditioned path.

\section{What is proved, and how solid it is}
\label{comp:ch14-scorecard}

The first scorecard covers the interpretive layer.

\begin{scorecard}
\MT{def:phil-idle-wheel} & A posit that changes the law of no record or record statistic is removed. & \CHOICE & Declared rule of reading, not derived from constructive empiricism. \\
Local-first verdict, \MTeq{eq:phil-fiber-measures} & Agents' law points are primary; base, bundle, spaces are scaffolding with no law. & \CHOICE & Describes the present architecture; changes if edges become inferred coordinates. \\
\MTeq{eq:phil-invariants} & Evidence values, loop classes, and per-edge twists up to conjugation survive frame change. & \SOLID & Invariance proved in the cited statements; reading it as ``structure'' is interpretation. \\
Realist and Neo-Kantian readings & Compatible with the invariants; none selected. & \VAGUE & The chapter declines to select and names what support would need. \\
\MT{prop:link-field-equation}, \MT{prop:link-linearized-equation} & Latent links have a stationarity equation and a Proca-type linearized operator. & \SOLID & Declared prior, stated matter and history hypotheses; physical-field reading remains structural. \\
\MT{thm:link-held-limit} & Held-link free-energy infima and evidence are zero-width limits of latent-link quantities. & \COND & Bounded records, positivity, and the theorem's continuity and variational hypotheses. \\
\MT{hyp:phil-observational-closure} & All data reach an agent as clamped states or node messages. & \CHOICE & Free for kernels by \MT{thm:obs-agent-interaction-equivalence}; removes no conditioning. \\
\MT{hyp:phil-boundary-agents} & Boundary environment nodes are called agents. & \CHOICE & Naming only; adds no state, action, blanket, or free energy. \\
Participatory reading & Participation via inference loop, learning, delayed feedback; formation could distinguish the program. & \VAGUE & Negative half rests on proved identifications; positive half is programmatic. \\
\MT{claim:physical-law-identification} & Identify a coarse-graining object with a physical law. & \OPEN & Needs a named target system and held-out evidence; none supplied. \\
\MT{open:phil-testability} & An observation at $c_*$ that could count against the program. & \OPEN & Only model-internal separations are proved. \\
\end{scorecard}

The second covers the continuation line. Every theorem-type statement is proved only on the four-cycle,
with abelian translation links on scalar Gaussian means, so all are TOY; the two general items are marked
separately.

\begin{scorecard}
\MT{prop:phil-seeded-contrast-formation} & Contrast formation from every Gaussian initial law with probability bounded below; same with frozen rows; symmetric laws keep rows uniform. & \TOY & Probability from the matching seed; learning idle; floor about $10^{-31}$ at $\mathsf m_0=0$. \\
\MT{conj:phil-contrast-formation} & Fixed small $\mathrm T$, $\epsilon_{\mathrm L}\to0$: learned assignment law concentrates on one matching on $[t_{\mathrm L},2t_{\mathrm L}]$ for a.e.\ mean; frozen arm symmetric. & \TOY & Amplifies the initial-law asymmetry; one-time law; no rate in $\mathrm T$. \\
\MT{cor:phil-joint-selection} & Joint limit $\mathrm T,\epsilon_{\mathrm L}\to0$ along implicit thresholds; frozen matching probabilities tend to $1/2$. & \TOY & Thresholds existential; matching may vary with $\mathrm T_n$. \\
\MT{prop:phil-temperature-tracking} & Entropy-production rate of order $\mathrm T^{5/2}$; tracking for $\epsilon_{\mathrm L}=\mathrm T^a$, $a>7/2$. & \TOY & Tracking only; no power-law selection. \\
\MT{prop:phil-selected-pair-persistence} & Selected pair coherent at resolution $\mathrm T^{1/4}$ for time $\mathrm T^{-1/4}$; partition contrast near $1/2$. & \TOY & Restored equilibrium coherence; no entrance claim. \\
\MT{prop:phil-early-block-null} & For $\epsilon_{\mathrm L}=o(\mathrm T^{1/4})$ learned and frozen first-block laws agree in the limit. & \TOY & Silent at the original mobility $\mathrm T^{1/8}$. \\
\MT{prop:phil-actual-conditioned-parent} & Late activation: relative entropy and discarded information $O(\mathrm T^{1/4}\log(1/\mathrm T))$. & \TOY & Averaged over a declared prefix, not per history. \\
\MT{prop:phil-actual-parent-dynamics} & Parent path law within $O(\mathrm T^{1/4})$ of the conditioned fine kept path over $\mathrm T^{-1/4}$. & \TOY & Finite window; own-law learning; no recursion. \\
\MT{prop:phil-early-seeded-parent} & Discarded information vanishes if $\theta-\log(1/\epsilon)\to\infty$, diverges if $\log(1/\epsilon)-\theta\to\infty$. & \TOY & Seeded sub-event; window $\log(1/\epsilon)+O(1)$ open. \\
\MT{prop:phil-early-first-completion} & At first qualifying completion, discarded information at least of order $\mathrm T^{-5/8}$, also given the whole contrast event. & \TOY & Established negative result on the four-cycle. \\
\MT{prop:phil-early-parent-dynamics} & Early parent path error vanishes when $\mathrm T\mathsf h\to0$, uniformly in mean and mobility. & \TOY & Windows shorter than $\mathrm T^{-1}$. \\
\MT{lem:phil-early-discarded-information} & Discarded information under conditioning, mixing, and extra variables; a Donsker--Varadhan lower bound. & \SOLID & General; standard in substance (chain rule). \\
\MT{rem:phil-conditioned-hidden-sign} & Conditioning on a probability-$1/2$ event costs $\log2$; on a realized history, infinity. & \SOLID & General counterexample shown on a Gaussian. \\
\end{scorecard}

\begin{hygiene}
What matters is what makes the four-cycle special and what fixes the comparison. The graph has exactly two
perfect matchings; links are abelian translations on scalar means, so loop holonomy is the number
$\sum_i\omega_i$ and plays no role in any selection; edges carry energy only when attended. The temperature
is small and fixed ($\mathrm T\le\mathrm T_0$ or $1/100$ late, $\mathrm T\le10^{-16}$ for the early null,
$\mathrm T=\epsilon^8$ with $\epsilon\le1/100$ for the early parent), and the bounds are asymptotic, some
with large constants. The rows learn from the law, not from samples, which makes them deterministic and
excludes fluctuation-driven selection: a modeling commitment, not hygiene. The initial law is Gaussian
with full support, and ``almost every mean'' excludes the imprint's zero set, which contains every
rotation-invariant mean. The kinetics are a declared realization along which the free energy does not
increase, not selected by it; the parent channel is fixed by hand. The late results need
$\eta,\delta\le\mathrm T$ and $\epsilon_{\mathrm L}\le\mathrm T^2$; the early ones condition on the seeded
sub-event, not the whole contrast event. Hygiene to take on trust: standard Borel spaces, measurability
of the first-block mark, well-posedness, finite moments, density envelopes, and form-domain extensions of
functional inequalities. Observational closure costs nothing for kernels; the boundary-agent hypothesis is
a naming decision.
\end{hygiene}

\section{Toy versus general}
\label{comp:ch14-toy}

\subsection{Mechanism and small-network specifics}

Every formation statement here lives on four agents; the program's target is 512 or more agents with
nonabelian or non-compact groups such as $\GL(K,\R)$ or $\mathrm{SU}(N)$. The four-cycle grounds and
validates; it does not set direction, and none of its results transfers without proof.
\Cref{comp:tab-ch14-toy} separates plausible mechanism-level content, the companion's reading and
conjectural beyond four agents, from four-agent specifics.

\begin{sciencetable}{Mechanism-level content versus four-cycle specifics.}
\label{comp:tab-ch14-toy}
\small
\begin{tabularx}{\textwidth}{>{\raggedright\arraybackslash}p{0.16\textwidth}>{\raggedright\arraybackslash}X>{\raggedright\arraybackslash}p{0.31\textwidth}}
\toprule
\textbf{Result} & \textbf{Mechanism-level content (conjectural beyond four agents)} & \textbf{Four-cycle specifics}\\
\midrule
Seeded formation & A favorable initial condition relaxes into a coherent block; the probability is the measure of such conditions. & Seed points $\pm\mathbf x^\circ$, floor near $10^{-31}$, factors $1/16$, $e^{-1/4}$. \\
Symmetric laws & A unique equivariant law-level dynamics cannot break a symmetry of its initial data; elementary and general, stated only for the four-cycle. & The form $\mathsf m_0=(\mathrm c_1\mathbf 1_4,\mathrm c_2\mathbf 1_4,0)$. \\
Pairing selection & Slow law-level learning is natural-gradient ascent of the log evidence; low temperature favors sheets with fewest stiff directions; a transient imprint decides the branch after time $\epsilon_{\mathrm L}^{-1}\log(1/\epsilon_{\mathrm L})$. & Two matchings, one unstable direction, ratios $\sqrt{40/7}\mathrm T$, $\sqrt{5/3}\mathrm T^{3/2}$, $\alpha\ge7/5$, the imprint computation. \\
Tracking & Mobility must be small against the frozen mixing rate. & Rate $\mathrm T^{5/2}$, exponent $7/2$. \\
Persistence & A structure selected in the prior is coherent in equilibrium on windows short against the hopping time. & Resolution $\mathrm T^{1/4}$, window $\mathrm T^{-1/4}$, flux $31\mathrm T$. \\
Early null & Learning slower than formation cannot bias the first block. & Threshold $o(\mathrm T^{1/4})$, horizon $2\mathrm T^{-1/8}$. \\
Parent closure & Exact to the extent discarded modes are in conditional equilibrium; activation must follow their thermalization; a detector coarser than the thermal scale fires too early. & Time $\log(1/\epsilon)$, divergence $\mathrm T^{-5/8}$, windows $\mathrm T^{-1/4}$, $\mathrm T^{-1+\mathrm q}$. \\
\bottomrule
\end{tabularx}
\end{sciencetable}

Some hypotheses exclude the target outright. Scalar fibers mean $K=1$, and abelian links mean holonomy
cannot act, so whether holonomy selects components is not touched; it is one branch of
\MT{open:rg-formation-adaptive-family}. The stationary preference for the fewest stiff directions does
generalize within the gated family: minimum edge covers dominate the stationary law on any finite graph
(\MT{prop:rg-gated-depth-one}). But nothing is proved about which of the combinatorially many covers slow
learning selects, or whether hopping times between covers grow with $N$, which is where genuine
metastability and a large-$N$ transition could enter; the four-cycle rules none of this in or out.
Recursion is out of reach too: eligible blocks have two or three agents, so the pairs cannot coalesce
within the event class, and recursive formation in the sense of R14 is posed outside the gated family
(\MT{open:rg-formation-general}).

\subsection{What the four-cycle says about waiting out the transient}

The reader's intuition is that meta-agent formation is soft and continuous: wait out the transient and
keep the slow-to-relax mode. The four-cycle gives this a precise form at $N=4$ in three respects and
qualifies it in two.

The activation-clock results are the sharpest support. A parent on the seeded episode is exact, as a
summary and as a dynamics, when activated at a deterministic time after the discarded modes have
thermalized, and provably not at the first moment a threshold criterion is met
(\MT{prop:phil-early-seeded-parent,prop:phil-early-first-completion}). ``Wait out the transient'' is
exactly $\theta-\log(1/\epsilon)\to\infty$, and the negative result says that a first-passage clock
driven by a detector coarser than the thermal scale is the wrong clock. The waiting is bounded on the
other side too: valid activation times and windows lie below the slow time $\mathrm T^{-1}$, so the
meta-agent is a quasi-stationary description on a plateau, not a permanent object. Second, what survives
the transient late is equilibrium structure: the selected pair is coherent because the equilibrium at the
learned prior is paired. Third, the channel does keep the slow directions, relaxing at rates of order
$\mathrm T$, with the loop and the assignment, and every discarded coordinate is fast.

The first qualification is that the channel was not chosen by a slow-mode criterion: it is supplied, it
keeps one fast mode as well, and no statement derives it from relaxation rates. The second cuts deeper.
Waiting out the transient selects nothing by itself. In a symmetric system with law-level learning, which
pairing forms is a property of sample paths, and at the level of the law nothing is selected unless a
datum breaks the symmetry (part (c) of \MT{prop:phil-seeded-contrast-formation}, rule R13). Soft
continuous formation and law-level selection of a culture are different claims; on the four-cycle the
first holds and the second needs an asymmetric initial law. Whether large networks add a genuine source
of selection, such as branch choice by fluctuations under learning from finite samples, or a transition
with diverging hopping times, nothing here addresses.

\section{Open ends}
\label{comp:ch14-open}

Physical-law identification (\MT{claim:physical-law-identification}) asks for a named target system whose
data the model does not generate, a record kernel separating two coarse-graining objects that differ in a
gauge-invariant datum, and held-out evidence, robust across blocking schemes, that a baseline margin is or
is not met. Single-point testability (\MT{open:phil-testability}) asks for a test at $c_*$, fixed before
the data, whose failure would falsify nontrivial holonomy on a walk or formation of a block, or a proof
that no record statistic distinguishes the alternatives. What exists is model-internal (the holonomy
record of \MT{thm:cg-loop-operational-holonomy}, the learned-versus-frozen selection); no target system,
noise model, estimator, or baseline is supplied.

The central formation open end is \MT{open:rg-formation-adaptive-family}: at the original coupled mobility
$\epsilon_{\mathrm L}=\mathrm T^{1/8}$, a difference bounded below between the learned and frozen
conditional laws of the first formed block, or between two initial means differing only in their loop
coordinate, with parents of vanishing errors; or a proof that these differences vanish for almost every
mean. The late selection and the early null bracket this regime without entering it. Also undecided are
the activation window $\theta=\log(1/\epsilon)+O(1)$, the discarded information given the whole contrast
event at a deterministic time, windows of order $\mathrm T^{-1}$ or longer, and a power-law selection
regime, which needs a quantitative lower bound on the imprint. Beyond the four-cycle, repeated endogenous
formation (\MT{open:rg-formation-general}), dynamic closure of the pointwise parent
(\MT{open:rg-pointwise-parent-dynamics}), edge creation (\MT{open:obs-edge-creation}), and a natural
partition selector (\MT{open:cg-natural-selector}) remain open; the last two mark where the local-first
verdict would change.

\begin{verdictbox}{\NARROWER}
The technical chapter says its readings are definitions, hypotheses, and open problems, and that it proves
a proposition and a theorem that sharpen the participatory reading: formation from every Gaussian initial
law with learning idle, and law-level pairing selection by slow source learning against a symmetric
frozen arm. Every proved statement does exactly what it says, the scope paragraph is accurate, and the
interpretive verdicts (local-first, no realist position selected, no agents created by naming) correctly
describe the formalism. For the claim that matters to the program, the result is narrower than it sounds.
What is selected is one of two pairings of a four-agent ring with abelian links, in the one-time
assignment law, by amplifying an asymmetry of the initial law, at fixed small temperature with mobility
tending to zero. The selected pair's later coherence is equilibrium coherence, not formation; the early
formation that passes the formation test is seeded and unaffected by slow learning; parent closure is
proved as averaged, finite-window statements on declared prefixes, and the first-completion clock is
refuted. The chapter itself concedes that assignment selection does not establish the recursive
participatory claim. The gap is the meeting point of formation and selection,
\MT{open:rg-formation-adaptive-family}, with recursion, large networks, nonabelian holonomy, and contact
with any observed system all open.
\end{verdictbox}

\chapter{Related Work}
\label{comp:ch15}

This chapter covers \MT{ch:related-work}, a positioning chapter with one small established remark, and
adds what it does not tabulate: which of the theory's main results are standard mathematics in new
typing and which are new.

\section{The question and the answer}
\label{comp:ch15-question}

The technical chapter asks where the construction sits among existing work and what it owes. The
neighbors flow a distribution along a scale parameter (Bayesian renormalization), coarse-grain graphs
whose edges carry scalars (network renormalization, graph coarsening), constrain learned maps over a base
manifold (equivariant networks), or share, omit, or fit a group-level model (multi-agent active
inference). The present theory computes the parent law from the fine joint by an exact Markov channel at
fixed evidence (\MT{thm:rg-exact-coarse-vfe}, \MT{thm:rg-pointwise-parent-datum}), with group-valued
independent links at one point. The central comparison, with Bayesian renormalization, rests on two
statements: evidence bounds of different joints cannot be compared unless their evidences are related
(\MT{prop:gen-no-distinguished-target}), and an evidence-preserving channel applied to both the posterior
and the recognition law orders the bounds by data processing (\MT{thm:cg-evidence-preserving-channel}).
The scope paragraph states that no other construction is identified with this one and that every
qualification of another work is the technical text's reading.

The link construction adds a direct comparison with lattice gauge theory, induced actions, and
learning rules (\MT{sec:links-related-work}). The references concern particular actions,
normalizations, and updates, rather than an identification of the agent model with a physical
lattice field theory or a neural architecture. The formulas themselves are graded in
\cref{comp:lie-links}.

\section{The picture}
\label{comp:ch15-picture}

\Cref{comp:tab-ch15-map} maps the neighbors; the grade refers to the relation between the neighbor's
central object and the theory's corresponding object.

\begin{sciencetable}{Neighbors of the construction.}
\label{comp:tab-ch15-map}
\small
\begin{tabularx}{\textwidth}{>{\raggedright\arraybackslash}p{0.19\textwidth}>{\raggedright\arraybackslash}X>{\raggedright\arraybackslash}p{0.3\textwidth}>{\raggedright\arraybackslash}p{0.1\textwidth}}
\toprule
\textbf{Neighbor} & \textbf{Its object} & \textbf{Relation here} & \textbf{Grade}\\
\midrule
Bayesian renormalization & Posterior flow on parameter space along a Fisher information scale; stiff and sloppy directions. & Different construction; composition needs a common target (\SG{open:grg-bayesian-bridge}). & loose \\
Information-geometric renormalization & Channels from microscopic to observable, a metric on Hamiltonians, monotone correlation functions. & Only the classical Fisher contraction for one kernel (\MT{thm:cg-fisher-contraction}). & structural \\
Network (Laplacian) renormalization & Heat kernel $e^{-tL}/\Tr e^{-tL}$, spectral entropy, specific-heat peaks. & Spectra gauge invariant, node affinities not (\MT{rem:grg-heat-kernel-gauge}). & structural \\
Algebraic multigrid & Galerkin coarse operator $P^\top AP$ with aggregation prolongation. & Same operator, inside normalized laws (\MT{prop:cg-message-passing-residual}). & exact \\
Matrix-weighted consensus & Laplacian with matrix edge weights; consensus subspace. & The Gaussian interaction family (\SG{def:gauss-interaction-family}). & exact \\
Cellular sheaves & Sheaf Laplacian; $H^0$ as its kernel; Kron obstruction. & Represented-link energies are sheaf energies (\MT{prop:cg-kernel-holonomy}). & exact \\
Connection Laplacians & Orthogonal transports on a graph; convergence to a manifold operator. & Special case: orthogonal links, scalar weights. & exact \\
Maximal-tree gauge fixing & Tree links set to the identity; non-tree links carry loop holonomies. & The rooted datum (\MT{prop:rg-rooted-datum}). & exact \\
Lattice gauge theory & Wilson plaquette term and temporal-gauge time-like loop. & Metric-isometric loop defects match the plaquette term; history loops use identity time-like links (\MT{prop:link-loop-defect}, \MT{prop:link-history-prior}). & exact in scope \\
Induced gauge theory & Integrating out matter produces a link action. & Jointly normalized undirected matter induces an action; record extension does not (\MT{prop:link-induced-action}). & structural \\
Bayesian learning rule & Natural-gradient variational update in Gaussian natural coordinates. & Explicit Euler step of the Gaussian generator Fisher flow (\MT{prop:link-recognition-flow}). & exact \\
Equivariant networks & Weight sharing over a global group or local frames on a base manifold. & Links may be inferred as generators under priors at one point; no learned layer is thereby identified. & loose \\
Multi-agent active inference & Shared models, collective motion, group-level blankets. & Parents computed from the fine joint. & loose \\
Path-space averaging & Relative-entropy bounds for effective diffusions. & Not growing-history, switching-source comparisons. & structural \\
\bottomrule
\end{tabularx}
\end{sciencetable}

The pattern is that the linear algebra is almost entirely standard and the probabilistic and gauge typing
around it is the theory's own. Represented-link quadratic energies are sheaf Laplacian energies with
rectangular restriction maps, whose kernels are global sections; the coarse operator is that of algebraic
multigrid; loops are handled by maximal-tree gauge fixing. The theory places these operators inside
normalized joint laws with exact evidence accounting and two channels that may carry inequivalent
representations, and proves formation theorems there. Its relation to renormalization in the physicist's
sense is weaker than the vocabulary suggests: the main manuscript's renormalization is an exact
finite-step channel with no rescaling, fixed points, or flow of couplings, and the asymptotic apparatus
lives in Supplement~B.

For links, the lineage is more direct but the scope differs. The technical text relates the
metric-isometric loop defect to the Wilson plaquette normalization
\citep{WilsonConfinement1974,Creutz1983}, and the identity time-like links in the history prior to
temporal gauge \citep{KogutSusskind1975,Creutz1983}. The prior uses a Gaussian reference on the Lie
algebra in place of Haar measure; this supplies finite normalization for noncompact groups at the
cost of anchors, covariances, and a fiber metric that transform under reference changes
(\MT{prop:link-generator-prior}). A continuous class function on the general linear group cannot
separate unipotent holonomy from the identity, so a trace-only reading of every loop defect would
miss content that the metric defect detects (\MT{cor:link-no-invariant-defect}).

\begin{analogy}{exact}
Kron reduction of a Gaussian network is integrating out Gaussian sites: the Schur complement of the
precision is the precision of the marginal. The Kron obstruction of \citet{HansenGhrist2019} says that
for vector stalks an eliminated site can induce couplings no pairwise restriction maps express, much as
decimation generates couplings outside the original family. Supplement G classifies when matrix-weighted
Kron reduction stays in its family (\SG{prop:cg-kron-leaves-family}, \SG{thm:cg-congruence-diagonal-kron});
the scalar case is closed (\SG{lem:cg-scalar-kron-closure}). Galerkin closure
(\MT{thm:cg-rectangular-endpoint-closure}) avoids the obstruction by restricting the form to a subspace:
block-spin averaging rather than decimation.
\end{analogy}

\begin{analogy}{exact}
\MT{rem:grg-heat-kernel-gauge} is the lattice-gauge fact that a propagator between different sites is not
gauge invariant. Under vertex reframings $P=\operatorname{diag}(P_i)$ the generator $R^{-1}L$ is
conjugated, so spectrum, heat trace, and spectral entropies are invariant, while off-diagonal heat-kernel
blocks transform as $P_i^{-1}(\cdot)P_j$, like $\langle\phi_i\phi_j^\dagger\rangle$ without a connecting
Wilson line. A partition selector built on node affinities therefore needs invariant block statistics.
\end{analogy}

\section{Objects and notation}
\label{comp:ch15-notation}

\begin{decoder}
Information scale & In Bayesian renormalization, a scale on parameter space set by Fisher distinguishability. & A renormalization scale, by its authors' analogy. \\
Stiff and sloppy directions & Parameter directions well and poorly constrained by data. & Relevant and irrelevant directions, loosely. \\
$D^{\mathfrak O}_\theta$ & Pushed inverse Fisher metric, the diffusion tensor of the data-space flow. & Diffusion constant of a Fokker--Planck equation. \\
$\chi(t)$ & Spectral specific heat $-d\mathcal S/d\log t$; peaks mark resolution scales. & Specific heat. \\
Sheaf, stalk, restriction map & A vector space per vertex and edge, maps from vertex to edge spaces. & Matter field and parallel transport on a lattice. \\
$H^0$ & Global sections: vertex data consistent across every edge. & Covariantly constant configurations. \\
Kron reduction & Eliminating vertices by a Schur complement. & Decimation of Gaussian sites. \\
Galerkin form $P^\top AP$, aggregation & Restricting a form to a subspace of piecewise-constant block fields. & Kadanoff block-spin averaging. \\
Maximal-tree gauge & Gauge in which the links of a spanning tree are the identity. & Axial or tree gauge. \\
Evidence-preserving channel & Coarse channel that leaves the observation marginal unchanged. & Integrating out at fixed partition function. \\
Metric loop defect $\mathfrak Q_p$ & Squared displacement of the represented holonomy from identity in a declared fiber metric. & Wilson plaquette term only for metric isometries. \\
Induced action $-\log\mathsf Z_{\mathrm Y}$ & Contribution obtained by integrating matter out of a jointly normalized undirected model. & Matter-induced gauge action. \\
Generator Fisher flow & Natural-gradient evolution of Gaussian recognition means and precisions. & Variational Bayesian learning rule upon Euler discretization. \\
\end{decoder}

``Renormalization'' in the main manuscript means exact finite-step coarse-graining of the latent law;
Wilsonian words appear only when describing other authors' work. ``Channel'' means a classical Markov
kernel, not the quantum channel of \citet{beny2015information}. ``Gauge'' means passive re-expression of
frames at one point with independent links, not a connection on a base manifold, which is why the
comparison with gauge-equivariant networks \citep{cohen2019gauge} is loose.

\section{The formulas, explained}
\label{comp:ch15-formulas}

The display of Bayesian renormalization the technical text engages is the pushed inverse Fisher metric
\[
D^{\mathfrak O}_\theta=\partial\mathfrak O(\theta)g^F(\theta)^{-1}\partial\mathfrak O(\theta)^\top
\]
(\MTeq{eq:rg-bks-diffusion}), with $\mathfrak O$ the map from parameters to data. In the late-data regime
the posterior is Gaussian with covariance proportional to $(g^F)^{-1}$; the delta method pushes it to data
space with covariance proportional to $D^{\mathfrak O}_\theta$, which \citet{BermanKlingerStapleton2023}
read as the diffusion tensor of a flow on data. The technical text's reading is that this is the leading
local approximation of the pushforward, exact for affine $\mathfrak O$, and that continuing it beyond the
late-data regime is a renormalization construction, not a theorem of Bayesian asymptotics.

The comparison turns on \MT{prop:gen-no-distinguished-target}: for joints indexed by recognition laws
$Q$ and $Q'$,
\[
\Lelbo(Q';X)-\Lelbo(Q;X)
=\mathsf e(Q')-\mathsf e(Q)
-\bigl[\KL(Q'\Vert P_{\theta,Q'}(\cdot\given o,X))-\KL(Q\Vert P_{\theta,Q}(\cdot\given o,X))\bigr],
\]
with $\mathsf e$ the log evidence of each joint. A change of bound mixes a change of evidence with a change
of tightness, so the optimized bounds of two independently declared joints, such as a model at one scale
and a model whose prior is another scale's posterior, are not ordered. An evidence-preserving channel fixes
$\mathsf e$, and data processing (\MT{thm:cg-kl-dpi-extended}) then orders the divergences. In physics
terms, variational free energies of different Hamiltonians cannot be compared unless their partition
functions are related, and integrating out keeps the partition function fixed. Hence the typing
prohibition does not refute Bayesian renormalization, but inserting a posterior as the generative top
factor at another scale declares a new model, comparable only through a supplied comparison map.

For network renormalization the key object is the spectral specific heat
$\chi(t)=-d\mathcal S/d\log t$ of \citet{villegas2023laplacian} (\MTeq{eq:grg-heat-capacity}). Spectral
quantities of the reframed generator survive frame changes, but a scale diagnostic selects no partition
(\MT{open:cg-natural-selector}), and the node affinities one would use to build blocks are not invariant.
For algebraic multigrid the coarse operator $P^\top AP$ with $P$ piecewise constant on disjoint aggregates
\citep{Notay2010Aggregation} is the theory's Galerkin form. The technical documents add a criterion for
exact coarse message passing, invariance of the $M$-orthogonal complement of the coarse range under the
fine operator (\MT{prop:cg-message-passing-residual}), and the placement of the operator inside normalized
laws with gauge covariance and exact evidence accounting.

\section{What is proved, and what is standard}
\label{comp:ch15-scorecard}

The chapter proves one remark; the rest is attribution. The scorecard adds the statements it leans on and
the companion's reading of which main results are standard in substance, which is information, not
criticism: it tells the reader what a result can and cannot do.

\begin{scorecard}
\MT{rem:grg-heat-kernel-gauge} & Spectra and heat traces of the reframed generator are invariant; off-diagonal blocks are not. & \SOLID & Elementary similarity; standard in substance. \\
\MT{prop:gen-no-distinguished-target} & Differences of bounds of different joints mix evidence and tightness. & \SOLID & Algebra of the evidence identity; the point is the typing. \\
\MT{thm:cg-evidence-preserving-channel} & A recognition-independent channel preserves evidence and orders bounds. & \SOLID & Standard in substance (data processing). \\
\MT{thm:elbo-extended-gap} & One evidence bound whose gap is the relative entropy to the posterior, in the extended reals. & \SOLID & Standard in substance (Gibbs variational principle); new in its measure-level, many-agent typing. \\
\MT{thm:rg-exact-coarse-vfe} & Fine free energy equals coarse free energy plus discarded conditional information. & \SOLID & Standard in substance (chain rule). \\
\MT{thm:rg-effective-action} & Coarse likelihood as a conditional partition function; composition along a tower. & \SOLID & Standard in substance (integrating out), with care about versions. \\
\MT{thm:cg-fisher-contraction} & Fisher information contracts under a parameter-independent kernel. & \SOLID & Classical, as the technical text says. \\
\MT{prop:cg-kernel-holonomy} & The kernel of the represented-link energy is the holonomy-fixed sector. & \SOLID & Standard in substance (sheaf $H^0$, cited). \\
\MT{prop:rg-rooted-datum} & Tree links gauge to the identity; non-tree links carry loop holonomies. & \SOLID & Standard (maximal-tree gauge fixing, cited). \\
\SG{prop:gauss-interaction-energy-kernel} & Energy and kernel criterion for matrix-weighted Laplacians. & \SOLID & Standard in substance (consensus literature, cited). \\
\MT{thm:cg-loop-operational-holonomy} & A declared record's law separates loop-holonomy classes. & \SOLID & New as a construction; elementary proof. \\
\MT{prop:link-loop-defect} & Metric-isometric loop defect equals the Wilson plaquette form. & \SOLID & The identification is restricted to isometric represented holonomy. \\
\MT{prop:link-induced-action} & Undirected matter gives an induced link action; record extension leaves the prior generator marginal unchanged. & \SOLID & Normalization distinguishes the two models. \\
\MT{prop:link-recognition-flow} & Gaussian Fisher flow has the Bayesian learning rule as its explicit Euler step. & \SOLID & Stated Gaussian family and step conditions; not a rate-separation theorem. \\
\SG{open:grg-bayesian-bridge} & Relate Bayesian renormalization to Gaussian aggregation through a common monotone. & \OPEN & Needs a common state space and a transported monotone. \\
\end{scorecard}

\begin{hygiene}
The comparisons are claims about other authors' work; their accuracy depends on the cited sources, which
the companion has not rechecked, and the scope paragraph marks each qualification as the technical text's
reading. The Bayesian-renormalization contrast rests only on the evidence identity, not on how the
Gaussian pushforward is read. The sheaf and connection-Laplacian identifications hold for quadratic
energies with represented links, not for the non-Gaussian laws of the general theory, where only the
kernel and the holonomy-fixed sector are claimed.
\end{hygiene}

\section{Standard versus new}
\label{comp:ch15-standard}

Standard in substance, and valuable to know as such, are the evidence-bound identity and its gap (the
Gibbs variational principle), the local-global potential identity (coordinate ascent in mean-field
variational inference, \MT{thm:obs-local-global-potential}), attention as the posterior of a latent
source label (mixture-model responsibilities, \MT{prop:obs-attention-elbo}), the exact coarse free energy
(the chain rule), the effective action (integrating out), Fisher contraction, the kernel of
represented-link energies (sheaf cohomology in degree zero), maximal-tree gauge fixing, the Galerkin
operator of algebraic multigrid, matrix-weighted Laplacians, and the Kron obstruction for vector stalks.
For each, the technical documents supply careful typing: measure-level statements in the extended reals,
two channels with independent links, and a separation of Markov channels, energy restriction, and
parameter aggregation. The typing is the contribution, and it is where the documents' length comes from.

The new latent-link chapter makes further debts explicit (\MT{sec:links-related-work}). Its
Gaussian expectation derivatives use the identities of \citet{OpperArchambeau2009}, and its
Fisher recognition flow gives the Gaussian Bayesian learning rule of \citet{KhanRue2023} as an
Euler step. The comparison-record current has the outer-product form of the fast-weight delta
rule discussed by \citet{Schlag2021}, following \citet{Schmidhuber1992}; the correspondence is
structural because here the current is transported through the derivative of the exponential and
averaged in a latent-generator stationarity condition. Integrating jointly normalized undirected
matter gives an induced action in the sense of the mechanism used by \citet{KazakovMigdal1993}.
The action vanishes for matter attached by normalized record extension; in the Gaussian
connection-Laplacian realization it contains one half a log-determinant plus a self term
(\MT{prop:link-induced-action}). The technical text compares this determinant mechanism with the
cycle-rooted-forest formulas of \citet{Kenyon2011VectorBundleLaplacian}. These are attributions to
specific source formulas, not claims of priority for the full agent construction.

New in content are the formation results: the finite-depth theorem under a supplied hierarchy
(\MT{thm:finite-depth-adaptive-formation}), the endogenous four-cycle witness
(\MT{prop:rg-formation-endogenous-cycle}), the depth-one result for the gated family
(\MT{prop:rg-gated-depth-one}), and the selection and parent line of \cref{comp:ch14}. Their tools are
standard (the Harris theorem, logarithmic Sobolev inequalities, couplings, saddle escape) and their
statements new, within the scopes \cref{comp:ch13,comp:ch14} describe. New as constructions, with
elementary proofs, are the operational holonomy record and the use of the no-distinguished-target identity
against posterior-flow constructions. Supplement G's classification of Kron closure is set against cited
prior art for the obstruction and the scalar case; its assessment belongs to the companion's Supplement G
chapters.

\section{Open ends}
\label{comp:ch15-open}

The bridge between Bayesian renormalization and Gaussian aggregation (\SG{open:grg-bayesian-bridge}) needs
a common state space, a map between the Bayesian scale parameter and a block ratio, and a functional
proved monotone along both constructions, or a Gaussian hierarchy in a stated class with no such
functional. A combined network scheme, sketched but not posed as a numbered problem, would infer a
partition from gauge-invariant diagnostics of edge-event laws through a structural channel with a named
law and push the complete law exactly; its partition part is \MT{open:cg-natural-selector}. The
growing-history, switching-source path comparison that averaging theory does not supply is part of
\MT{open:rg-pointwise-parent-dynamics}.
The finite-dimensional group-valued closure of coarse latent-link laws is
\MT{open:link-coarse-closure}; Gaussian generators and arbitrary real logarithms cannot be
assumed through a general-linear coarse step (\cref{comp:ch12-latent-closure}).

\begin{verdictbox}{\DELIVERED}
The technical chapter says it positions the construction among existing work without identifying it with
any of them, with the Bayesian-renormalization comparison resting on the no-distinguished-target identity
and on evidence-preserving channels. It does that: each comparison is anchored to a statement of the
theory, the Bayesian-renormalization contrast follows correctly from the evidence identity, and the one
proved remark is correct and elementary. Two limits. The chapter does not say, result by result, which of
the theory's statements are standard mathematics in new typing; the companion's reading is that most of
the linear algebra and the variational identities are, and the formation results are new. The added
link comparisons identify the Wilson and temporal-gauge forms at their stated scopes, the
matter-integration mechanism, the delta-rule form, and the Gaussian Bayesian learning update;
they do not establish a physical field or learning architecture. And two
neighbors one would expect here, matrix-weighted consensus and maximal-tree gauge fixing, are treated in
Supplement G and the network-state chapter rather than in this chapter. The attributions to other
authors are the technical text's reading and have not been rechecked by the companion.
\end{verdictbox}


\part{The Supplements}
\chapter{Supplement G, Part 1: The Gaussian Realization and Its Geometry}
\label{comp:sg1}

Supplement G asks what the single-point theory becomes when every belief and model law is
multivariate Gaussian (\SG{ch:gaussian-introduction}). This chapter covers its introduction and its
first five chapters; aggregation, Kron reduction, nonflat links, the operator arrow, and reciprocal
pairs are in \cref{comp:sg2}. Three facts frame what follows. The supplement is a realization, not the
ambient theory: it asserts nothing for a non-Gaussian belief fiber beyond the exponential-family
statements of G1, and the main manuscript uses none of its results. The Gaussian lives on real fibers,
\(\R^K\) for beliefs and \(\R^{d_m}\) for model presentations with \(K\neq d_m\) allowed, and the
structure group reaches them only through its represented image \(\Gamma_b=\rho_b(G)\subseteq\GL(K)\);
some statements hold for every \(\Gamma_b\), and a named set is proved under the specialization
\(\Gamma_b=\GL^+(K)\) (\SG{hyp:gen-lg-frame-group}). And the belief precision analyzed throughout is
that of the fast belief at frozen model latents, \(p(k\given m,o,X)\), not the joint belief--model
precision.

Most of the mathematics is finite-dimensional linear algebra and Gaussian integration, much of it
textbook in substance. What is specific to the program is the typing (which objects are laws, which are
operators, which are artifacts of a chart or frame) and the analysis of when a directed linear-Gaussian
network yields the precision a physicist would write first, a matrix-weighted graph Laplacian plus local
mass terms. Only \SG{prop:exp-cumulants-fisher} is tagged STANDARD; every other statement in G1 to G5 is
ESTABLISHED with a proof, and none rests on computation. In one line each: G1 shows that boundary
points and operator fixed points carry no law; G2 gives the information form, the forest precision, and
one Schur lemma behind every later marginal; G3 shows that the Laplacian-plus-self-term shape is a
declared restriction and says when it holds, exactly in one chart and at law level on forests and pairs;
G4 prices recognition restrictions and finds that
none selects a partition; G5 separates law-level invariants from chart and frame artifacts.

\section{Exponential families and the operator layer (G1)}
\label{comp:sg1-expfamily}

\subsection{The question and the answer}

Later chapters apply algebraic maps (aggregation, rescaling, Schur complementation) to matrices in a
cone, and a matrix in the cone need not be the precision of any normalizable law. \SG{ch:expfamily} asks
what happens when a map of natural parameters leaves the natural domain \(\mathsf N\). On its boundary
nothing is left: log partition, mean parameter, Fisher metric, relative entropies, and evidence bound
diverge or are undefined (\SG{prop:exp-domain-boundary-no-law}); a map can keep every iterate inside
while its only fixed point is on the boundary (\SG{prop:exp-fixed-point-no-law}); and pinning,
quotienting, and keeping a mass coordinate are three different repairs (\SG{sec:exp-repairs}). Inside,
relative entropy is a Bregman divergence (\SG{prop:exp-kl-bregman}) and affine mean constraints have a
unique information projection (\SG{prop:exp-affine-mean-projection}).

\subsection{The picture}

\begin{analogy}{exact}
An exponential family is a Gibbs ensemble: the statistic \(\mathsf T\) lists observables, the natural
parameter \(\lambda\) their conjugate fields, \(\mathsf A(\lambda)\) is a Massieu function,
\(\nabla\mathsf A\) the thermodynamic averages, and \(\nabla^2\mathsf A=\operatorname{Cov}(\mathsf T)\)
the fluctuation--response relation, which is also the Fisher metric (\SG{prop:exp-cumulants-fisher}). For
the Gaussian, \(\mathsf T=(y,yy^\top)\) and \(\lambda=(\mathsf h,-\tfrac12\Lambda)\), a source and a
coupling matrix. The information projection is the family-restricted form of minimum relative entropy
\citep{jaynes2003probability}.
\end{analogy}

The boundary results are an infrared statement. For the scalar Gaussian with \(\mathsf T(z)=-z^2/2\),
\(\lambda\) is the stiffness and \(\mathsf N=(0,\infty)\); the map \(\lambda\mapsto\lambda/2\) keeps
every iterate normalizable, but its limit \(\lambda=0\) is the flat weight on the line
(\SGeq{eq:exp-scale-witness}), inside the operator cone \([0,\infty)\) and outside the probability
layer.

\begin{analogy}{structural}
This is the massless limit of a free field, whose zero mode has no normalizable measure. Physicists fix a
boundary value, remove the zero mode, or keep a small mass as a running coupling; G1's pinning
(\SG{def:exp-pinning}), quotienting (\SG{def:exp-quotienting}), and mass sector
(\SG{def:exp-mass-sector}) are these three repairs, and they give three different objects. A
pseudoinverse and pseudodeterminant in the full-space density formula compute the quotient normalizer,
on a smaller space with another reference measure, while the full-space integral stays infinite
(\SG{prop:exp-quotient-normalization}).
\end{analogy}

\subsection{Objects and notation}

\begin{decoder}
Canonical family \((\nu,\mathsf T,\lambda,\mathsf A,\mathsf N)\) & Laws \(\propto e^{\langle\lambda,\mathsf T\rangle}\) relative to \(\nu\); \(\mathsf N\) is where \(\mathsf A<\infty\) & Gibbs ensemble\\
Regular, minimal & \(\mathsf N\) open; no affine redundancy, so parameters determine laws & No redundant coordinates\\
Operator cone \(\mathcal W\); probability layer \(\mathcal W\cap\mathsf N\) & Parameters preserved by the maps; the part that denotes laws & Couplings with unphysical points; normalizable couplings\\
Bregman divergence; information projection & Error of the linear extrapolation of \(\mathsf A\); closest member under affine mean constraints & Availability; maximum entropy\\
\end{decoder}

Traps: the sans-serif \(\mathsf A\) is the log partition, the italic \(A_i\) of later chapters are self
terms; the bold \(\boldsymbol\mu\) is a mean parameter, for the Gaussian the pair \((\mu,M_2)\) of mean
and second moment (\SG{def:gauss-coordinate-charts}); ``operator'' means a matrix acted on by an
algebraic map, nothing quantum.

\subsection{The formulas, explained}

Inside the domain,
\[
\KL(Q_{\lambda'}\Vert Q_\lambda)=\mathsf A(\lambda)-\mathsf A(\lambda')-\langle\nabla\mathsf A(\lambda'),\lambda-\lambda'\rangle
\]
(\SG{prop:exp-kl-bregman}): the difference of log partitions minus the first-order change predicted from
\(\lambda'\). Convexity makes it nonnegative; for nearby parameters it is the Fisher quadratic, and at
finite separation it is not, a difference the supplement insists on. For a closed affine set of mean
parameters meeting the interior, the projection \(Q_{\widehat\lambda}\) is unique and every
constraint-satisfying member obeys
\[
\KL(Q_{\lambda'}\Vert Q_\lambda)=\KL(Q_{\lambda'}\Vert Q_{\widehat\lambda})+\KL(Q_{\widehat\lambda}\Vert Q_\lambda),
\]
as a squared distance splits under orthogonal projection (\SG{prop:exp-affine-mean-projection}); the
Bernoulli examples of \SG{rem:exp-projection-hypotheses} show that interior feasibility, closedness,
and convexity are each needed. The scope paragraph is a real warning: no coarse-graining statement
follows from exponential-family structure alone, since an operator arrow needs, among other things, a
declared rescaling, and a law-level arrow needs finite normalizers at every scale.

\section{The Gaussian realization at the base point (G2)}
\label{comp:sg1-realization}

\subsection{The question and the answer}

\SG{ch:generative-gaussian} asks what the directed law becomes when every kernel is linear Gaussian, and
what the precision blocks mean. The precision is at once the coupling matrix and the Fisher metric of the
mean (\SG{prop:gauss-log-normalizer-moments}); a zero block is conditional independence, not zero
covariance (\SG{prop:gauss-conditional-precision}); the marginal precision of a set of directions is a
Schur complement below the restricted precision, with equality exactly under \(\Lambda\)-orthogonality
(\SG{lem:gauss-schur}). A forest assembles into an explicit sparse precision
(\SG{prop:gauss-assembled-precision-blocks}), every kernel is frame covariant
(\SG{prop:gen-lg-kernel-covariance}), and the learning step has a closed form that fails to exist in two
named cases (\SG{prop:gen-lg-m-step}).

\subsection{The picture}

\begin{analogy}{exact}
In \(q\propto\exp(\mathsf h^\top y-\tfrac12y^\top\Lambda y)\), \(\Lambda\) is the quadratic part of a
lattice action (the inverse propagator), \(\mathsf h\) an external source, and \(\Sigma=\Lambda^{-1}\)
the propagator. A zero off-diagonal block means no direct bond while the propagator still connects the
sites through paths (\SGeq{eq:gauss-precision-not-covariance}). The Schur complement is integrating out,
exact at tree level for a Gaussian. The inequality \(\Lambda_{bb}^{-1}\preceq(\Lambda^{-1})_{bb}\)
(\SG{prop:restrict-optimum-vs-marginal}) is the Le Chatelier--Braun inequality: a block fluctuates less
with the rest clamped than with the rest free to relax, which is why a factorized recognition law, whose
block covariances are the clamped ones, underestimates every marginal fluctuation.
\end{analogy}

The directed model (\SG{def:gen-lg-forest}) is a state-space model on a tree: each nonroot belief is its
parent's belief carried by the link \(\Theta_{ij}\), multiplied by a gain \(F_i\), shifted by a bridge
term \(\mathsf B_im_i\) and an offset, plus noise of precision \(\Lambda_i\), with linear-Gaussian
observations. The product \(\Omega_i=F_i\Theta_{i\operatorname{par}(i)}\) is the transition. At frozen
model latents the belief sector is the forest
\(k_i\given k_{\operatorname{par}(i)}\sim\mathcal N(\Omega_ik_{\operatorname{par}(i)},\Lambda_i^{-1})\),
and model latents, offsets, and the model part of observations move only \(\mathsf h\): model
presentations act on beliefs as an external source, never as a coupling. This is the Born--Oppenheimer
move of the main manuscript's fast--slow split (\MT{sec:obs-fast-slow}), cleaner than usual because the
fast precision does not depend on the frozen variables at all.

\subsection{Objects and notation}

\begin{decoder}
\((\mathsf h,\Lambda)\); \(\mu=\Lambda^{-1}\mathsf h\), \(\Sigma=\Lambda^{-1}\) & Information vector, precision; mean, covariance & Source, couplings; propagator\\
Four charts & Information \((\mathsf h,\Lambda)\), natural \((\mathsf h,-\tfrac12\Lambda)\), expectation \((\mu,M_2)\), moment \((\mu,\Sigma)\) & Choices of thermodynamic variables\\
\(F_i\), \(\Theta_{ij}\), \(\Omega_i\) & Gain, represented link, transition & Transfer matrix times transporter\\
Restricted versus marginal precision & \(E^\top\Lambda E\) versus \((E^\top\Lambda^{-1}E)^{-1}\) & Clamped versus relaxed stiffness\\
M-coordinate & Maximizing expected complete log likelihood over generative parameters & Learning (EM's M-step)\\
\end{decoder}

Traps: one subscript on \(\Lambda\) is a kernel's innovation precision, two a block of the assembled
matrix (the \(W_i\) of G3 follows the same rule); symmetric-matrix coordinates scale off-diagonal entries
by \(\sqrt2\); the moment chart is not Legendre dual to the natural chart, the expectation chart is;
``pullback'' means restriction along an injective parameter map.

\subsection{The formulas, explained}

The log normalizer
\(\mathsf A(\mathsf h,\Lambda)=\tfrac12\mathsf h^\top\Lambda^{-1}\mathsf h-\tfrac12\log\det\Lambda+\tfrac n2\log2\pi\)
is a source energy mediated by the propagator, a fluctuation determinant, and a measure constant; the
Fisher information in \(\mu\) at fixed \(\Sigma\) is \(\Lambda\) (\SG{prop:gauss-log-normalizer-moments}),
so the stiffness of the energy along a direction is the distinguishability of mean shifts along it. The
conditional mean \(\E[Y_b\given Y_{-b}]=\mu_b-\Lambda_{bb}^{-1}\sum_{c\neq b}\Lambda_{bc}(y_c-\mu_c)\),
with conditional covariance \(\Lambda_{bb}^{-1}\), is the local equilibrium of block \(b\) in its
neighbors' field, and as an update the coordinate-ascent step of G4. With an agreement coupling
\(\Lambda_{bc}=-W_{bc}\preceq0\), a neighbor's deviation pulls the block toward it in the \(\Lambda_{bb}\)
inner product but not necessarily in the Euclidean one (\SG{prop:gauss-conditional-precision}, part
(iii)): agreement is relative to the
receiver's own metric, the only one available under general linear frames.

The forest precision has the blocks
\[
\Lambda_{i\operatorname{par}(i)}=-\Lambda_i\Omega_i,
\qquad
\Lambda_{ii}=\Lambda_i+\sum_{i'\in\operatorname{ch}(i)}\Omega_{i'}^\top\Lambda_{i'}\Omega_{i'},
\]
and zero between nonadjacent agents (\SG{prop:gauss-assembled-precision-blocks}): each factor puts its
noise stiffness on the child, a bond \(-\Lambda_i\Omega_i\) to the parent, and a back-reaction on the
parent. The bond need not be symmetric, which is the seed of G3. The belief--model cross block is
\(K\times d_m\), so no family of square agent blocks hosts the joint precision when \(K\neq d_m\).
Frame covariance (\SG{prop:gen-lg-kernel-covariance}) is lattice-gauge covariance of a matter action,
with gains transforming by conjugation, covariances by congruence, and observation maps absorbing the
inverse frame; it holds for every \(\Gamma_b\), and for \(\Gamma_b\subseteq\operatorname{O}(K)\) ordinary
eigenvalues are also invariant (\SG{prop:gen-lg-compact-frames}), which fails under \(\GL^+(K)\).

The learning step (\SG{prop:gen-lg-m-step}) is least squares plus residual covariance,
\(\E\log\mathcal N(t;Fw,\Sigma)\leq-\tfrac12\log\det(2\pi e\Sigma_*)\), attained when
\(\Sigma_*\succ0\); an exact regression in some direction makes the likelihood unbounded, the
variance-collapse degeneracy of maximum likelihood. The named cases are a record factor at one observed
value and a link learned inside the invertible group, whose supremum can sit at a singular map outside
the group: for \(K=1\) with independent centered \(k_i,k_j\), the optimized term
\(-\tfrac12\log(2\pi e(\E k_i^2+\mathsf t^2\E k_j^2))\) increases as \(\mathsf t\to0\) and is never
attained. Noncompactness of the link group can make learning ill posed; the scalar witness says nothing
about frequency at large \(K\). The communication formulas (\SG{prop:geo-selective-lift},
\SG{prop:obs-selective-gaussian-record}) are Gaussian conditioning and the Gaussian channel; the record's
mutual information is gauge invariant and at rank one takes every value in \((0,\infty)\).

\section{The interaction family (G3)}
\label{comp:sg1-interaction}

\subsection{The question and the answer}

The coarse-graining results of the second half act on one shape of precision, the interaction family
\(\mathcal I(V)\) (\SG{def:gauss-interaction-family}):
\[
\Lambda_{ii}=A_i+\sum_{j\neq i}W_{ij},
\qquad
\Lambda_{ij}=-W_{ij}\quad(i\neq j),
\qquad
A_i\succeq0,\quad W_{ij}=W_{ji}\succeq0.
\]
\SG{ch:gaussian} asks when the belief precision of G2 has this shape. It is a restriction, not a
consequence of linear-Gaussian directedness (\SG{prop:gauss-interaction-family-thin},
\SG{prop:gauss-global-condition-witnesses}). With
orientation-preserving frame changes, positive-determinant transitions on a forest can be gauged away and
a pair can be brought into the family exactly when \((-1)^K\det\Lambda_{12}>0\)
(\SG{prop:gauss-gain-absorption}). A reversing gain stays outside in every chart reached by
orientation-preserving frames, and a cycle with the wrong determinant sign, or for odd \(K\) the
three-agent v-structure, in every chart reached by any frames (\SG{prop:gauss-law-level-witnesses}). In the
link-trivializing chart two computable conditions decide membership
(\SG{prop:gauss-global-interaction-condition}), and there the family is adopted as a standing hypothesis
(\SG{hyp:gauss-global-interaction}).

\subsection{The picture}

A member's quadratic form is
\[
y^\top\Lambda y=\sum_{i}y_i^\top A_iy_i+\sum_{i<j}(y_i-y_j)^\top W_{ij}(y_i-y_j)
\]
(\SG{prop:gauss-interaction-energy-kernel}): \(K\)-component site variables, matrix springs
\(W_{ij}\) penalizing disagreement, and pinning or mass terms \(A_i\). Consensus configurations cost only
\(v^\top(\sum_iA_i)v\); they are the zero modes of the Laplacian part \(L\), lifted by self terms alone,
and at \(A=0\) they are the non-normalizable zero modes of G1. A singular bond leaves disagreement free
along its kernel, so \(\ker L\) can exceed the consensus subspace on a connected graph. With definite
bonds under flat comparison, one definite self term per connected component makes \(\Lambda\succ0\)
(\SG{prop:gauss-interaction-nonempty}).

\begin{analogy}{exact}
For \(K=1\), \(W_{ij}=J\) on lattice bonds and \(A_i=m^2\), a member is the massive lattice Laplacian
\(-\Delta+m^2\) of the lattice Gaussian free field; for general \(K\) it is the matrix-weighted Laplacian
of the consensus literature \citep{Trinh2018MatrixWeighted}. With nontrivial links the agreement energy
\(\sum(k_i-\Theta_{ij}k_j)^\top W(k_i-\Theta_{ij}k_j)\) is the covariant lattice Laplacian of a matter field
in a background gauge field \citep{Creutz1983}: the connection Laplacian of \citet{SingerWu2012} for
orthogonal links, a sheaf Laplacian \citep{HansenGhrist2019} for general linear ones. The program's links
are independent edge variables, as lattice links are, so this is exact at the level of the quadratic
form. G3's family is what the covariant Laplacian becomes in the pure-gauge sector: when every loop
product is the identity the links are a coboundary \(\Theta_{ij}=U_iU_j^{-1}\)
(\MT{prop:gauss-coboundary-trivialization}), and \(z_i=U_i^{-1}k_i\) turns every residual into
\(U_i(z_i-z_j)\) (\SG{hyp:gauss-flat-comparison}), giving an ordinary matrix-weighted Laplacian with
weights \(U_i^\top WU_i\). With a nontrivial Wilson loop no gauge does this; that case is excluded here
and treated in \cref{comp:sg2}.
\end{analogy}

\subsection{Objects and notation}

\begin{decoder}
\(\mathcal I(V)\), \(A_i\), \(W_{ij}\), \(L\) & Interaction family; self term; edge weight; Laplacian part & Massive matrix-weighted Laplacian; mass; bond\\
Link-trivializing chart & Coordinates \(z_i=U_i^{-1}k_i\) under flat comparison, unique up to one right factor per component & Gauge setting every link to one\\
Block reframing \(\mathbf P=\operatorname{diag}(\mathbf P_i)\) & Frame change agent by agent, \(\Lambda\mapsto\mathbf P^\top\Lambda\mathbf P\) & Site-dependent passive gauge transformation\\
Gain absorption; orientation criterion & Frames turning every transition into the identity; \((-1)^K\det\Lambda_{12}>0\) & Maximal-tree gauge; ferromagnetic bond sign (\(K=1\))\\
Chart-independent witness; v-structure & Outside the family in every admissible chart; a child with two parents & Frustrated plaquette; explaining away\\
\end{decoder}

Traps: \(\Omega_i\) is a transition, not a holonomy or frequency; \(W_i:=\Lambda_i\Omega_i\) with one
subscript is the candidate weight of the edge to the parent, \(W_{ij}\) an edge weight; ``flat'' means
trivial loop products of links at the one base point, with no base manifold in sight; an ``admissible''
chart is one reached by block reframings in the group a statement names.

\subsection{The formulas, explained}

Why the shape is not automatic is visible in one block. A receiving factor contributes the bond
\(-\Lambda_i\Omega_i\), and membership needs it to be \(-W\) with \(W\) symmetric positive
semidefinite, which nothing in a linear-Gaussian model forces. Even with links removed the gain remains:
\(\widetilde\Lambda_i=\begin{psmallmatrix}2&1\\1&3\end{psmallmatrix}\) with the quarter-turn
\(\widetilde\Omega_i=\begin{psmallmatrix}0&1\\-1&0\end{psmallmatrix}\) gives the bond
\(\begin{psmallmatrix}1&-2\\3&-1\end{psmallmatrix}\), and \(F_i=-I\) gives a symmetric bond of the wrong
sign (\SG{prop:gauss-global-condition-witnesses}). In a fixed chart every off-diagonal block must be
symmetric, \(\binom N2K(K-1)/2\) linear conditions, so the family is Lebesgue-null once \(K,N\geq2\)
(\SG{prop:gauss-interaction-family-thin}).

Thinness concerns one chart; the law-level question is whether some frame change puts the precision into
the family, and the answer is gauge-theoretic. A reframing \(k_i=\mathbf P_iz_i\) sends \(\Omega_i\) to
\(\mathbf P_i^{-1}\Omega_i\mathbf P_{\operatorname{par}(i)}\), the transformation law of a link
(\SG{prop:gauss-gain-absorption}, part (o)): gain times link behaves as one link, and no gauge
transformation can tell a gain from a link. On a tree every link configuration is pure gauge, so frames
\(\mathbf P_r=I\) at roots and \(\mathbf P_i=\Omega_i\mathbf P_{\operatorname{par}(i)}\) turn every
residual into \(z_i-z_{\operatorname{par}(i)}\), leaving pure agreement bonds
\(W=\mathbf P_i^\top\Lambda_i\mathbf P_i\) and self terms only at roots: the maximal-tree gauge applied to
transitions. The frames must lie in the frame group, which under \(\GL^+(K)\) requires
\(\det\Omega_i>0\); for a forest with invertible transitions this is also necessary (part (iii)). For a
pair the only obstruction is the sign
\[
(-1)^K\det\Lambda_{12}>0
\]
(part (ii)). Informally: a member with invertible bond has \(\Lambda_{12}=-W\), \(W\succ0\), so
\((-1)^K\det\Lambda_{12}=\det W>0\), and a reframing multiplies \(\det\Lambda_{12}\) by
\(\det\mathbf P_1\det\mathbf P_2>0\). The sign is a \(\mathbb Z_2\) gauge invariant of the edge. Admitting
orientation-reversing frames, \(\Gamma_b=\GL(K)\), removes it (part (iv)).

\begin{analogy}{exact}
For \(K=1\) this is the gauge structure of Ising spin glasses \citep{mezard1987spin}. The family is the
set of precisions with ferromagnetic couplings \(\Lambda_{ij}\leq0\) and nonnegative self terms; positive
rescalings cannot flip a bond sign (the orientation criterion); \(k_i\mapsto-k_i\) is the Mattis gauge
transformation, which makes any sign pattern on a tree ferromagnetic; and on a loop the product of bond
signs is invariant (frustration). \SG{prop:gauss-law-level-witnesses}(ii) is the \(K\times K\) version:
the sign of the product of block determinants around a cycle of \(n_\circ\) agents is invariant under
every \(\GL(K)\) reframing, and membership forces it to be \((-1)^{Kn_\circ}\). Part (iii) reads explaining
away as frustration: two parents of one child are coupled through
\(\Lambda_{jj'}=\Omega_j^\top\Lambda_i\Omega_{j'}\), and the triangle has the wrong determinant sign for
every odd \(K\). The correspondence is exact at the level of these signs and says nothing beyond them.
\end{analogy}

In the link-trivializing chart membership is decided exactly
(\SG{prop:gauss-global-interaction-condition}): each \(W_i\) must be symmetric positive semidefinite, and
each self term
\[
A_i=\Lambda_i-W_i+\sum_{i'\in\operatorname{ch}(i)}\left(W_{i'}\Lambda_{i'}^{-1}W_{i'}-W_{i'}\right)
\]
must be positive semidefinite (no \(-W_i\) at a root), with \(A_i+O_i\succeq0\) once observation
precisions are added. \(\Lambda_i-W_i\) is what a factor leaves on the child after paying for its bond,
and \(W_{i'}\Lambda_{i'}^{-1}W_{i'}-W_{i'}\) what each child's factor leaves on the parent. For a scalar
transition \(\omega\) these are \(\Lambda(1-\omega)\) and \(\Lambda\omega(\omega-1)\): a contracting gain
leaves a negative self term on the parent, an expanding one on the child, and only \(\omega=1\) is safe
factor by factor (\SG{prop:gauss-edge-local-characterization}, \SG{cor:gauss-invertible-gain-factor}).
Globally a root prior or observation can pay for a negative contribution, so factorwise membership is
strictly stronger than the hypothesis (\SG{hyp:gauss-factorwise-agreement}). A physicist can read the
condition as a ban on net tachyonic self terms.

Two consequences need emphasis. First, the adopted hypothesis is stronger than law-level membership. It
requires membership in the link-trivializing chart, and the scope paragraph says nothing selects another
chart. The prior of part (iii) of \SG{prop:gauss-global-interaction-condition}, with transition
\(\tfrac12\), is outside the family in that chart and, by gain absorption, inside it in its gain-absorbing
chart. The choice matters downstream: aggregation identifies agents by equal coordinates in the chosen
chart, so the chart decides whether agreement means link-transported equality \(k_i=\Theta_{ij}k_j\) or
equality with the prediction \(k_i=\Omega_ik_{\operatorname{par}(i)}\). That is a modeling decision, and
the supplement makes it explicitly. Second, membership is not invariant under general block reframings,
which preserve inertia but not the symmetry of off-diagonal blocks
(\SG{rem:gauss-membership-frame-dependence}); frame-independent statements need the invariants of G5.

\section{Recognition restrictions (G4)}
\label{comp:sg1-restrictions}

\subsection{The question and the answer}

Fix a Gaussian model with posterior \(p^\star=\mathcal N(\mu,\Lambda^{-1})\), \(\Lambda\succ0\). By the
restriction principle (\MT{prop:restrict-principle}), shrinking the recognition family lowers the
optimized bound by the rise in minimal relative entropy to the posterior. \SG{ch:restrictions} computes
that loss for covariance factorization (the determinant gap, \SG{thm:restrict-determinant-gap}), for mean
restrictions (\SG{prop:restrict-mean-cost}), and for identifying agents (\(+\infty\),
\SG{prop:cg-epsilon-identification-cost}), and shows that the gap selects no structure
(\SG{prop:restrict-gap-no-selection}).

\subsection{The picture}

For a Gaussian candidate,
\[
\KL(Q\Vert p^\star)=\tfrac12\left[\Tr(\Lambda\Sigma_Q)+(\mu_Q-\mu)^\top\Lambda(\mu_Q-\mu)-n-\log\det\Lambda-\log\det\Sigma_Q\right]
\]
(\SG{prop:restrict-gaussian-reverse-kl}). The mean enters only the quadratic term and the covariance only
the trace and log determinant, so mean and covariance restrictions are separate problems whose costs add
(\SG{cor:restrict-combined-costs-add}).

\begin{analogy}{exact}
The bound's gap is the slack in the Gibbs--Bogoliubov--Feynman inequality. For a Gaussian field with a
block-factorized trial law, the optimal means are exact and the optimal block covariances are the inverse
on-block stiffnesses \(\Lambda_{bb}^{-1}\) (\SG{thm:restrict-exact-block-optimum}): Weiss mean-field
theory, which drops the feedback of fluctuations through the environment, with the gap as its
free-energy error \citep{Bishop2006,Wainwright2008}. The mean-restriction cost
\(\tfrac12\mu_\perp^\top(B_\perp^\top\Lambda^{-1}B_\perp)^{-1}\mu_\perp\) is the free-energy cost of
constraining a Gaussian field's mean: half the squared displacement weighted by the inverse relaxed
susceptibility. The bare stiffness \(B_\perp^\top\Lambda B_\perp\) overstates it unless the constrained
directions are \(\Lambda\)-orthogonal to the rest (\SGeq{eq:cg-loewner}).
\end{analogy}

For two blocks the gap is their mutual information, \(-\tfrac12\sum_a\log(1-\check\rho_a^2)\) in the
canonical correlations (\SG{prop:restrict-canonical-correlation}); for more blocks it is the dual total
correlation \(h(Y)-\sum_bh(Y_b\given Y_{-b})\) of \citet{Han1978}
(\SG{cor:restrict-gap-dual-total-correlation}), built from the conditional entropies that
\(\Lambda_{bb}\) encodes, not the total correlation built from marginal ones (for three unit scalars with
pairwise correlation \(\tfrac12\), \(\tfrac12\log\tfrac{27}{16}\) against \(\tfrac12\log2\)). The gap is
observation free because a data-independent joint precision makes conditioning on \(o\) move only the
posterior mean (\SG{cor:restrict-gap-observation-free}); the mean cost is quadratic in the data
(\SG{cor:restrict-mean-cost-data}).

Identification of agents is the case that matters for coarse-graining. Let \(\operatorname{range}S\) be
the cluster-constant configurations of a hard aggregation, of dimension \(r=n_cK<n\). A recognition law
making all agents of a cluster literally equal lives on that subspace, to which the posterior, having a
Lebesgue density, gives probability zero. Absolute continuity fails (clause (H3) of
\MT{hyp:elbo-evidence-domain}), and the relative entropy is \(+\infty\). Thickening by a variance
\(\varepsilon\) in each identified direction gives
\[
\mathbb D_\varepsilon=\mathbb D_{\rm tie}+\frac{n-r}2\log\frac1\varepsilon-\frac{n-r}2
+\frac12\log\det(B_\perp^\top\Lambda^{-1}B_\perp)+\frac\varepsilon2\Tr(B_\perp^\top\Lambda B_\perp)
\]
(\SG{prop:cg-epsilon-identification-cost}), one \(\tfrac12\log(1/\varepsilon)\) per frozen continuous
coordinate, as a delta-function constraint costs entropy logarithmic in the resolution. Identification is
therefore a generative construction, the fine law conditioned on zero within-cluster disagreement
(\SG{prop:cg-aggregation-conditioning}); at finite cost a recognition law can only tie its means, the
mean-tie family with cost \(\mathbb D_{\rm tie}\) (\SG{prop:cg-mean-tie-cost}). Restricting the energy to
\(\operatorname{range}S\) and renormalizing, the ``hard trace,'' changes the evidence
(\SG{prop:restrict-hard-trace-evidence}); with a data-independent joint precision its posterior is the
image of the fine one under \(P_\Lambda=(S^\top\Lambda S)^{-1}S^\top\Lambda\), an evidence-preserving
coarse channel (\MT{thm:cg-evidence-preserving-channel}).

\subsection{Objects and notation}

\begin{decoder}
Recognition restriction; partition \(\mathcal P\) & Shrinking the candidate family; blocks with block-diagonal candidate covariance & Trial ensemble; mean-field clusters\\
Cost, gap & Loss of optimized bound, the rise of minimal relative entropy & Variational free-energy error\\
Mean tie; identification & Candidate means cluster constant; candidate supported on \(\operatorname{range}S\) & Constrained order parameter; delta-function constraint\\
Hard trace; \(P_\Lambda\) & Energy restricted to \(\operatorname{range}S\); the map carrying the fine posterior to it & Constrained partition function; Galerkin projection\\
\end{decoder}

Traps: ``restriction'' restricts the variational family, not a tensor; ``structure'' is the partition
used by the recognition law; the orthonormal basis \(B\) is unrelated to the bridge \(\mathsf B_i\);
``understates the covariance'' is said of the optimal factorized law, never of the model.

\subsection{The formulas, explained}

The determinant gap,
\[
\log p_\theta(o\given X)-\Lelbo^\star_{\mathcal P}=\frac12\left[\sum_{b\in\mathcal P}\log\det\Lambda_{bb}-\log\det\Lambda\right]\geq0,
\]
compares the product of on-block stiffness determinants with the full determinant; by Fischer's
inequality it vanishes exactly when every cross block does. When every block is a union of whole agents,
a block reframing shifts both terms by the same Jacobian, so the gap is invariant under every reframing,
whatever the frame group (\SG{prop:ig-frame-dependent-spectra-determinants}, part (iii)); a partition
that splits an agent is not.

Does any analyzed cost select a partition scale? No. The gap grows under refinement by an explicit
nonnegative amount (\SG{prop:restrict-refinement-monotonicity}) and vanishes at one block, so minimizing
it selects the unrestricted family (\SG{prop:restrict-gap-no-selection}). Non-nested partitions are
ordered by no structural feature: for
\(\Lambda=\begin{psmallmatrix}1&a&0\\a&1&b\\0&b&1\end{psmallmatrix}\) the partitions
\(\{\{1,2\},\{3\}\}\) and \(\{\{1\},\{2,3\}\}\) have equal block sizes and swap order when \(|a|\) and
\(|b|\) are exchanged (\SG{prop:restrict-nonnested-unordered}). The mean-tie cost vanishes at the finest
partition and identification is infinite for every nontrivial cluster. Each cost is extremized at an
endpoint, the Gaussian instance of \MT{open:restrict-intrinsic-criterion} and \MT{open:cg-natural-selector}.
One functional the chapter writes down is not examined in this role: the combined cost
\(\mathbb D_{\rm tie}+\mathbb D_{\rm fact}\) of \SG{cor:restrict-combined-costs-add} (mean-tie cost plus
cluster factorization gap, both displayed in \SG{prop:cg-mean-tie-cost}), with both
restrictions indexed by the same clusters, is the exact loss of one admissible family, and its terms move
in opposite directions under merging (the first cannot decrease because merging shrinks the allowed
means; the second cannot increase). The no-selection statements concern each cost separately; whether
the sum has a nondegenerate minimizer is not analyzed.

\section{Fisher geometry and gauge invariants (G5)}
\label{comp:sg1-fisher}

\subsection{The question and the answer}

\SG{ch:infogeometry} asks which geometric statements about the Gaussian family concern laws and which
concern a chart or frame. ``The Fisher mean block is the precision'' holds in the moment chart but not,
for coordinate directions, in the expectation chart (\SG{prop:ig-fisher-moment-chart},
\SG{prop:ig-fisher-expectation-chart}). Under frame changes, the generalized spectrum of the pencil
\((L,\Lambda)\) is invariant when regular (\SG{prop:grg-pencil-invariance}), while ordinary eigenvalues
and single subspace determinants are frame artifacts once the frame group contains the positive scalars
(\SG{prop:ig-frame-dependent-spectra-determinants}), and under \(\Gamma_b\supseteq\GL^+(K)\) a coupling
matrix keeps only its rank (\SG{prop:grg-residual-fiber-label}). Standard: the moment-chart metric of the
multivariate normal \citep{Skovgaard1984} with its affine isometries \citep{PineleStrapassonCosta2020},
exponential-chart duality \citep{AmariNagaoka2000}, congruence invariance and Sylvester's law
\citep{HornJohnson2013}. New: the typing of which quantities survive the represented frame group, and a
scoped statement of when a graph Laplacian is a Fisher tensor.

\subsection{The picture}

The two mean blocks differ because \(\partial/\partial\mu\) at fixed \(\Sigma\) and at fixed \(M_2\) are
different tangent vectors: the second must shrink the covariance to compensate
(\(\Sigma=M_2-\mu\mu^\top\)), and that extra motion has Fisher length, governed by the unbounded
Mahalanobis norm \(\mu^\top\Lambda\mu\). The quotient metric, a Schur complement, is exactly \(\Lambda\)
and sits below the restricted one, the clamped-versus-relaxed ordering of G2. As with \(C_p\neq C_V\), a
derivative depends on what is held fixed; an identity mixing mean and covariance sectors must name its
chart and say whether it restricts or quotients.

Under a general linear frame group there is no metric on the fiber, so the eigenvalues of one quadratic
form are chart data: rescaling all frames by \(\mathrm c>0\) multiplies every eigenvalue of \(\Lambda\) by
\(\mathrm c^2\), and a threshold on absolute eigenvalues can accept or reject any direction by choice of
\(\mathrm c\) (\SG{prop:ig-frame-dependent-spectra-determinants}, part (i)). Only relative quantities and
law functionals survive.

\begin{analogy}{exact}
Normal-mode frequencies are generalized eigenvalues of the potential form relative to the kinetic form
and do not depend on the generalized coordinates, whereas eigenvalues of the stiffness matrix alone do.
The pencil \(Lx=\mathsf d\Lambda x\) is this construction with the Laplacian part as potential and the
full precision as kinetic form; both transform by the same congruence, so the \(\mathsf d_k\) are
invariant. Each \(\mathsf d_k\in[0,1]\) is the fraction of precision along its direction that bonds supply
rather than self terms: consensus sits at \(0\), directions no self term touches at \(1\)
(\SG{prop:ig-generalized-spectrum-localization}).
\end{analogy}

\subsection{Objects and notation}

\begin{decoder}
Restriction versus quotient metric & Metric on coordinate directions versus orthogonal to a fiber & Clamped versus relaxed response\\
Block reframing \(\mathbf P\) & New-to-old frame change, \(\Lambda\mapsto\mathbf P^\top\Lambda\mathbf P\), \(\Sigma\mapsto\mathbf P^{-1}\Sigma\mathbf P^{-\top}\) & Site-dependent gauge transformation\\
Pencil \((L,\Lambda)\), regular & Pair of forms, \(Lx=\mathsf d\Lambda x\); determinant not identically zero & Normal-mode problem; no shared zero mode\\
Agent-aligned law functional; residual fiber label & Gap, mutual information, relative entropy of unions of agent blocks; congruence class of one semidefinite matrix & Gauge-invariant observable; inertia (Sylvester)\\
\end{decoder}

Traps: \(\mathsf d\) is a generalized eigenvalue, not a dimension; \(L\) is the Laplacian part, not the
bound \(\Lelbo\); a ``regular'' pencil and a ``regular'' exponential family are unrelated; the mass
\(\mathsf a\) below is a coefficient on a reference precision.

\subsection{The formulas, explained}

In the moment chart,
\[
g^F\bigl((\dot\mu_1,\dot\Sigma_1),(\dot\mu_2,\dot\Sigma_2)\bigr)=\dot\mu_1^\top\Lambda\dot\mu_2+\tfrac12\Tr(\Lambda\dot\Sigma_1\Lambda\dot\Sigma_2),
\]
and every affine map \((\mu,\Sigma)\mapsto(Q\mu+\mathrm v,Q\Sigma Q^\top)\) is an isometry; block
reframings act this way, so the Fisher geometry of the recognition family is gauge invariant for any
frame group. In the expectation chart
\(g_{[\mu\mu]}=(1+\mu^\top\Lambda\mu)\Lambda+\Lambda\mu\mu^\top\Lambda\), with eigenvalues relative to
\(\Lambda\) of \(1+\mu^\top\Lambda\mu\) (\(n-1\) times) and \(1+2\mu^\top\Lambda\mu\). The link between
Laplacian and Fisher geometry is narrow and exact:
\[
\KL\bigl(\mathcal N(\mu_i,W_{ij}^{-1})\Vert\mathcal N(\mu_j,W_{ij}^{-1})\bigr)=\tfrac12(\mu_i-\mu_j)^\top W_{ij}(\mu_i-\mu_j),
\]
summing over edges with definite weights to \(\tfrac12\mu^\top L\mu\) (\SG{prop:ig-fixed-covariance-kl}).
Outside Gaussian laws at that fixed covariance, or a frozen base law to second order, the chapter states
that a sample-space Laplacian is neither a recognition Fisher tensor nor a natural-gradient operator, a
scope limit on any use of a graph Laplacian as a preconditioner.

Pencil invariance is
\(\det(\mathbf P^\top L\mathbf P-\mathsf d\mathbf P^\top\Lambda\mathbf P)=\det(\mathbf P)^2\det(L-\mathsf d\Lambda)\).
The pencil is regular exactly when \(\ker L\cap\ker A=\{0\}\) (\SG{prop:grg-pencil-regularity}); at the
massless boundary \(A=0\) it is \((L,L)\), irregular, and says nothing about the graph until a pin,
quotient, or mass is declared. With a mass, \(\Lambda=L+\mathsf a\Lambda_{\mathrm{ref}}\) gives
\(\mathsf d=\mathsf d_{\mathrm{ref}}/(\mathsf d_{\mathrm{ref}}+\mathsf a)\)
(\SG{prop:grg-mass-pencil-transfer}), roughly the massive-propagator ratio \(p^2/(p^2+\mathsf a)\) on a
lattice; the small-eigenvalue counting exponent survives, and the limits \(\mathsf a\to0\) and
\(\mathsf d_{\mathrm{ref}}\to0\) do not commute (\SGeq{eq:grg-noncommuting-mass}). A subspace log
determinant picks up the Jacobian \(\log\det(B'^\top\mathbf P^\top\mathbf PB')\) (\SGeq{eq:ig-logdet-law}),
which orthogonal reframings kill and general linear ones do not; the gap of G4 cancels it for
agent-aligned blocks, while an agent-splitting gap moves from \(\tfrac12\log\tfrac43\) to
\(\tfrac12\log4\) in the \(N=1\), \(K=2\) witness. Under \(\Gamma_b\supseteq\GL^+(K)\) every semidefinite
matrix of rank \(r\) is congruent to \(\operatorname{diag}(I_r,0)\), so a common reframing leaves \(K+1\)
classes and no continuous modulus; under \(\operatorname{SO}(K)\) or \(\operatorname{O}(K)\) a definite
matrix keeps its \(K\) eigenvalues (\SG{prop:gen-lg-compact-frames}). With a noncompact frame group the
strength of a coupling matrix by itself is gauge; with a compact one it is physical.

One caution the wording does not make explicit. The pencil is formed after splitting \(\Lambda=L+A\) in
the chart where the family is imposed (\SG{def:ig-block-reframing}), the link-trivializing chart of
\SG{hyp:gauss-global-interaction}. Its spectrum is data of the law together with that split, hence with
its links, not of the law alone: a precision in the family in two inequivalent charts can have two
pencils. For \(N=2\), \(K=1\), \(\begin{psmallmatrix}2&-1\\-1&2\end{psmallmatrix}\) splits with \(W=1\),
\(A=(1,1)\) and has spectrum \(\{0,\tfrac23\}\), while its reframing by \(\operatorname{diag}(1,2)\),
\(\begin{psmallmatrix}2&-2\\-2&8\end{psmallmatrix}\), lies in the family with \(W=2\), \(A=(0,6)\) and
spectrum \(\{0,1\}\). This does not contradict \SG{prop:grg-pencil-invariance}, which concerns one
congruence of one pair; it limits \SG{rem:gauss-membership-frame-dependence}, which lists pencil
eigenvalues among frame-independent membership criteria, since forming the pencil presupposes a chart in
which membership already holds.

\section{What is proved, and how solid it is}
\label{comp:sg1-scorecard}

Every row is proved in the supplement or cited there to a checked source; there is no NUM statement and
no TOY positive result in this part. COND rows rest on a frame-group or graph hypothesis that fails in
cases the program cares about.

\begin{scorecard}
\multicolumn{4}{@{}l}{\textit{G1. Exponential families}}\\
\SG{prop:exp-domain-boundary-no-law} & No law, Fisher metric, or bound at a boundary point; an operator map can converge there (\SG{prop:exp-fixed-point-no-law}) & \SOLID & Negative; scalar witness suffices\\
\SG{sec:exp-repairs} & Pin, quotient, and mass sector are three distinct repairs & \CHOICE & Their two propositions are \SOLID\\
\SG{prop:exp-affine-mean-projection} & Cumulants and Fisher metric; relative entropy is Bregman; unique affine projection & \SOLID & Standard in substance; \SG{prop:exp-cumulants-fisher} is STANDARD\\
\multicolumn{4}{@{}l}{\textit{G2. Gaussian realization}}\\
\SG{hyp:gen-lg-frame-group} & Represented frame group, specialized to \(\GL^+(K)\) where named & \CHOICE & Not met by \(\operatorname{SU}(N)\) or a realified \(\GL(K,\mathbb C)\)\\
\SG{def:gen-lg-forest} & Directed linear-Gaussian forest & \CHOICE & General DAGs get one sentence\\
\SG{lem:gauss-schur} & Zero block is conditional independence; Schur complement below restricted precision & \SOLID & Standard in substance\\
\SG{prop:gen-lg-kernel-covariance} & Every kernel frame covariant for every \(\Gamma_b\); explicit forest blocks (\SG{prop:gauss-assembled-precision-blocks}) & \SOLID & Belief sector at frozen model latents\\
\SG{prop:gen-lg-m-step} & Closed-form learning step; no maximizer if the residual covariance is singular & \SOLID & Link witness is scalar: possibility, not frequency\\
\multicolumn{4}{@{}l}{\textit{G3. Interaction family}}\\
\SG{hyp:gauss-global-interaction} & The family, adopted in link-trivializing coordinates & \CHOICE & The chart fixes what agreement means\\
\SG{hyp:gauss-flat-comparison} & Belief links form a graph coboundary & \CHOICE & Automatic on forests; fails for generic links on cycles\\
\SG{prop:gauss-interaction-family-thin} & Null set in a fixed chart; not reframing invariant & \SOLID & Chart-level\\
\SG{prop:gauss-gain-absorption} & On forests, \(\GL^+(K)\) frames absorb positive-determinant transitions; pair criterion \((-1)^K\det\Lambda_{12}>0\); \(\GL(K)\) frames reach every invertible pair & \COND & Forest or pair; stated frame group; nothing sufficient for cycles or smaller groups\\
\SG{prop:gauss-law-level-witnesses} & Reversing gain, cycle determinant sign, odd-\(K\) v-structure exclude membership & \SOLID & Negative; (ii), (iii) for \(\GL(K)\); (iii) on three agents\\
\SG{prop:gauss-global-interaction-condition} & Exact membership test; prior balance sufficient, not necessary; factorwise stronger & \SOLID & Link-trivializing chart\\
\multicolumn{4}{@{}l}{\textit{G4. Recognition restrictions}}\\
\SG{thm:restrict-determinant-gap} & Exact mean, clamped covariances; loss is the gap, the dual total correlation (\SG{cor:restrict-gap-dual-total-correlation}) & \SOLID & Standard in substance\\
\SG{cor:restrict-gap-observation-free} & Gap independent of the observation; mean cost carries the data (\SG{cor:restrict-mean-cost-data}) & \SOLID & Needs data-independent joint precision\\
\SG{prop:cg-epsilon-identification-cost} & Identification costs \(+\infty\); regularized, \(\tfrac{n-r}2\log\tfrac1\varepsilon\) & \SOLID & Negative; absolute continuity\\
\SG{prop:restrict-gap-no-selection} & Gap monotone on chains, unordered off them, minimized by one block & \SOLID & Negative; combined cost not examined\\
\multicolumn{4}{@{}l}{\textit{G5. Fisher geometry and invariants}}\\
\SG{prop:ig-fisher-expectation-chart} & Mean block is \(\Lambda\) in the moment chart, larger in the expectation chart; quotient \(\Lambda\) & \SOLID & Standard in substance \citep{Skovgaard1984}\\
\SG{prop:ig-fixed-covariance-kl} & Edge relative entropy at covariance \(W_{ij}^{-1}\) is the Laplacian quadratic & \SOLID & Gaussian, fixed covariance only\\
\SG{prop:grg-pencil-invariance} & Regular pencil spectrum invariant; in \([0,1]\) & \SOLID & Invariant of law plus link-fixed split\\
\SG{prop:ig-frame-dependent-spectra-determinants} & Absolute spectra, subspace determinants frame dependent; agent gaps invariant & \SOLID & (i) needs positive scalars\\
\SG{prop:grg-residual-fiber-label} & A semidefinite coupling matrix is labeled by rank alone & \COND & Needs \(\Gamma_b\supseteq\GL^+(K)\); false for compact groups\\
\multicolumn{4}{@{}l}{\textit{Open}}\\
\SG{open:exp-mass-sector} & Mass-coordinate scaling; also \SG{open:exp-projection-commutation}, \SG{open:ig-model-space-fisher-transfer} & \OPEN & \Cref{comp:sg1-open}\\
\end{scorecard}

\begin{hygiene}
Conditions that can fail where the program lives: flat comparison (\SG{hyp:gauss-flat-comparison}),
without which the interaction hypothesis is not formulated; membership in the link-trivializing chart
(\SG{hyp:gauss-global-interaction}); the \(\GL^+(K)\) specialization (\SG{hyp:gen-lg-frame-group}),
needed by absorption and the rank-only label; frozen model latents; absolute continuity
(\MT{hyp:elbo-evidence-domain}, clause H3), which identification violates; a nondegenerate target and
regular pencil, which fail at the massless boundary \(A=0\); a data-independent joint precision; and a
nonsingular residual covariance in the learning step. Hygiene: standard Borel spaces, sigma-finite
references, measurability, regularity and minimality (automatic for the Gaussian), definite receiving
covariances (\SG{hyp:gen-lg-spd}), finite second moments, and orthonormal bases used as a computational
device.
\end{hygiene}

\section{Toy versus general, and what this part implies for \texorpdfstring{\(\GL(K)\)}{GL(K)} networks at large \texorpdfstring{\(N\)}{N}}
\label{comp:sg1-scale}

Nothing in G1 to G5 is TOY. Every theorem holds for arbitrary finite \(N\) and \(K\) (the two-agent
orientation criterion has a forest version at every \(N\)), and every small instance is a counterexample
to a universal statement, hence a solid negative result that shows a claim fails, not how often it fails
at \(N=512\). Thinness holds at every \(N\), and the cycle-sign condition
applies to any cycle with invertible blocks inside any network. Nothing here is a large-\(N\) statement
either: no thermodynamic limit, per-agent density, or scaling law.

\paragraph{The interaction family at scale.} Two conditions stand between a 512-agent Gaussian precision
with cycles and the family. Flat comparison requires every loop product of the belief links to be the
identity, automatic on a forest but violated by generic independent \(\GL(K)\) links on a graph with
many cycles; then the hypothesis behind aggregation and the operator arrow is not formulated, and the
nonflat case is in \cref{comp:sg2}. And the law-level criterion is complete only for forests and pairs;
for precisions whose graph has cycles, including those created by co-parents, only the necessary sign
condition is proved.

\paragraph{Frame groups.} For \(\GL(K,\R)\), \(\Gamma_b=\GL(K)\) contains \(\GL^+(K)\) and the positive
scalars, so every statement needing either holds. Part (iv) of \SG{prop:gauss-gain-absorption} reaches
every invertible pair, so the orientation obstruction disappears there, and the forest construction of
part (i) uses positive determinants only to keep its frames in \(\GL^+(K)\); the cycle-sign and
odd-\(K\) v-structure obstructions, stated for \(\GL(K)\) frames, remain. The gauge-invariant data of a Gaussian network under \(\GL(K,\R)\) are agent-aligned
law functionals, pencil spectra relative to the link-fixed split, cycle determinant signs, and holonomy
conjugacy classes (\MT{prop:geo-edge-twist-invariance}); absolute eigenvalues and single subspace
determinants are gauge artifacts. The supplement treats real Gaussians only. A complex group enters
through a real representation, for instance the realification on \(\R^{2K}\), whose image lies in
\(\GL^+(2K)\) and contains the positive scalars; statements proved for every \(\Gamma_b\) hold, the
rank-only label and the \(\GL^+\) sufficiency constructions do not apply as stated because their frames
need not be complex linear, and obstructions proved for larger groups still apply. For
\(\operatorname{SU}(N)\) and other compact groups, \(\Gamma_b\subseteq\operatorname{O}(K)\) after
choosing an invariant inner product: spectra are invariant, coupling matrices keep continuous eigenvalue
moduli, and no absorption result is stated.

\paragraph{Where compactness and orthogonality enter.} No theorem here assumes compactness. It is
excluded by the \(\GL^+(K)\) specialization, brings extra invariants when present, and its absence is
what the scalar link witness exploits (informally, on a compact group a continuous objective attains its
maximum, a standard fact the supplement does not state). A Euclidean structure enters as a computational
device (orthonormal bases), through orthogonal reframings that preserve spectra and subspace
determinants, through the averaging pushforward \(y\mapsto S^\dagger y\) of
\SG{cor:ig-full-rank-restriction-pushforward}, whose equality condition uses the Euclidean complement
and is chart-level under \(\GL(K)\), and through the finite part of the regularized identification cost,
whose divergent coefficient \((n-r)/2\) is frame free.

\section{Open ends}
\label{comp:sg1-open}

The supplement records three open problems here: the law of a retained mass coordinate under a scale
map, settled by an explicit induced law with a proof that its orbit stays in the natural domain or
reaches the boundary (\SG{open:exp-mass-sector}); whether the affine information projection commutes with
gauge actions or coarse parameter maps, settled by an intertwining theorem or a counterexample
(\SG{open:exp-projection-commutation}); and whether the Bayesian-renormalization information scale, built
on the Fisher metric of generative parameters, transfers to \(\Lambda\) read as the Fisher metric of a
recognition mean, settled by making the coarse mean a generative parameter with nondegenerate Fisher
metric and transferring the construction, or by a family where that metric is singular
(\SG{open:ig-model-space-fisher-transfer}). G5 declines any contraction claim for varying precision
(\SG{open:grg-galerkin-monotonicity}), and G4 is a Gaussian instance of
\MT{open:restrict-intrinsic-criterion} and \MT{open:cg-natural-selector}. Four gaps are visible but not recorded as open problems: law-level
membership for precisions whose graph has cycles; the relation of the family to a covariant Laplacian
with nontrivial holonomy; the combined mean-tie and factorization cost as a selector; and complex
Gaussians and complex or unitary frame groups beyond a real representation. These describe what the
documents do not contain; they are not proposals.

\section{Verdict}
\label{comp:sg1-verdict}

\begin{verdictbox}{\DELIVERED}
\textbf{What the part says it does.} Explicit Gaussian forms of the theory's objects: exponential-family
boundary and projection results; the information form, Schur lemma, and frame-covariant kernels of the
directed law; an exact account of when the belief precision is a matrix-weighted Laplacian plus self
terms, with that shape adopted as a hypothesis rather than derived; restriction costs, with the
statement that neither the factorization gap nor the mean cost selects a structure; and a separation of Fisher and spectral quantities into invariants and chart or
frame artifacts.

\textbf{What the statements give.} All of it, at arbitrary finite \(N\) and \(K\), with proofs. Much is
textbook in substance and presented as such. The program-specific content, the membership analysis of G3
and the invariance typing of G5, is elementary linear algebra done carefully; for \(K=1\) its sign
conditions are the Mattis-gauge and frustration conditions of spin glasses.

\textbf{The limits a reader must carry.} The interaction family is the pure-gauge sector, and the chart
in which it is imposed, not the law, decides what agreement means downstream. Law-level membership is settled on forests and pairs
and only bounded by a necessary sign condition on cycles, where a 512-agent network lives. The pencil
spectrum is data of a law with its link-fixed split, so the remark offering it as a frame-independent
membership criterion claims more than the propositions give. The combined recognition cost is not
examined as a selector. Absorption and the rank-only label need \(\GL^+(K)\) frames; under \(\GL(K,\R)\)
they hold and absorption is easier, compact groups keep more invariants and lose absorption, and complex
Gaussians are not treated. None of this is a defect in a proof, and the part states most of it itself.
\end{verdictbox}

\chapter{Supplement G, Part 2: Gaussian Coarse-Graining}
\label{comp:sg2}

This chapter covers the second half of Supplement G: score tangents and the Hermite spectrum
(\SG{ch:gaussian-hermite}), aggregation and Kron reduction (\SG{ch:gaussian-coarsegraining}), nonflat
links, relaxation, and quotient volume (\SG{ch:gaussian-nonflat}), the operator arrow
(\SG{ch:gaussian-renormalization}), reciprocal pairs, anchored stars, and loop holonomy
(\SG{ch:obstructions}), the decision metatheory (\SG{app:decision-metatheory}), and the supplement's open
problems (\SG{sec:grg-obligations}). Everything acts on the Gaussian interaction family explained with the
first half in \cref{comp:sg1}.

In short, this half of Supplement G carries out one finite step of exact coarse-graining for Gaussian agent
networks and does so completely: which step maps are exact, when the natural parameter families close under
them, where they provably do not, and which further declarations (a rescaling, a reference measure, a repair
of a singular boundary) a flow would need. It proves no flow, fixed point, or thermodynamic limit, and says
so. Most positive statements are standard mathematics in the theory's typing (Galerkin operators of algebraic
multigrid, Gaussian marginalization, Mehler's identity, block Gauss--Seidel, Tarski's decision method); the
content without a textbook counterpart among the supplement's citations is the Kron closure classification,
the bookkeeping of holonomy under coarse-graining, and a set of sharp negative results.

Three threads run through the part. Aggregation freezes each cluster into a rigid body while Kron reduction
integrates hidden agents out; both are exact, and they close on very different families. Curvature becomes
mass: a link that cannot be gauged away inside a block, around a folded loop, or across a strict cut reappears
after coarse-graining as a positive-semidefinite self term. And closure is neither selection nor a flow: a
step between spaces of different dimension has no fixed point until an identification and a rescaling are
declared.

\section{Score Tangents and the Hermite Spectrum (G6)}
\label{comp:sg2-hermite}

\subsection{The question, the answer, and the picture}

At one law there are three candidate spaces of directions of motion: the tangent space of a finite family
(for a Gaussian, mean and covariance moves), the Hilbert space \(L^2_0(\nu)\) of all centered
square-integrable scores, and the bounded perturbations \(e^{-t\varphi}\) of the main manuscript's nonlinear
action calculus (\MT{thm:rg-bounded-action-calculus}). The chapter asks how they nest at a Gaussian and what a
channel on \(b\) independent copies does to scores. It proves an exact information budget for every
parameter-independent block channel (\SG{thm:rg-score-pushforward-defect}), the exact nesting
(\SG{rem:ig-three-tangents}), which Hermite directions admit exponential tilts
(\SG{prop:ig-hermite-exponential-domain}), and the spectrum \(b^{1-k/2}\) of the linearized normalized-sum map at
the standard Gaussian (\SG{thm:rg-gaussian-hermite-spectrum}). Its point is a warning: every exact arrow
contracts Fisher information, so the eigenvalue \(b^{1/2}>1\) is created by the extensive replication factor
\(\sqrt b\), a declared normalization, and not by coarse-graining. Nothing later in Supplement G or in the main
manuscript uses these results; Supplement B builds mode exponents on the spectrum and poses nonlinear
attraction at the Gaussian block law as \SB{open:rg-hermite-nonlinear-attraction}.

\begin{analogy}{exact}
Copy a perturbed law onto \(b\) sites and push through \(\bar\xi=b^{-1/2}\sum_i\xi_i\). At \(\mathcal N(0,1)\) the
block law is again \(\mathcal N(0,1)\), and the score operator \(\mathscr L_b\) is the derivative of this central-limit
map at its fixed point. Its eigenvalues \(b^{1-k/2}\) are the scaling of cumulants under normalized summation: in
Wilsonian language the mean direction is relevant, the variance direction marginal, and every \(k\geq3\)
irrelevant. The supplement identifies its \(b=2\) case with the linearization of \citet{JonaLasinio2001} and
avoids the Wilsonian words, whose conventions belong to Supplement B.
\end{analogy}

The supplement adds the factorization \(\mathscr L_b=\mathsf U_b\mathscr I_b\). The replication lift
\(\mathscr I_bh=\sum_ih(x_i)\) has norm exactly \(\sqrt b\), since \(b\) independent sites carry \(b\) times the Fisher
information of one, and it is not a physical arrow: no parameter-independent kernel maps \(\mathcal N(x,1)\) to
\(\mathcal N(x\mathbf 1_b,I_b)\) (\SG{prop:rg-score-block-lift}). The conditional expectation \(\mathsf U_b\) only
loses information (\MT{thm:rg-action-lp-contraction}). A uniform field couples to an extensive block
magnetization; after dividing by \(\sqrt b\) the mean direction looks amplified, but blocking learned nothing new.

\subsection{Objects and formulas}

\begin{decoder}
Score \(h\); \(L^2_0(\nu)\) & Velocity \(\partial_t\log p_t\) of a path with \(L^2\)-differentiable square-root density; all centered square-integrable functions & Tangent vector; \(\delta\ln p\)\\
\(\mathscr I_b\); \(\mathscr L_b\) & Replication onto \(b\) independent sites; the block map linearized on tangents & Extensive source; linearized RG map\\
Bounded action quotient & Perturbations \(e^{-t\varphi}\), \(\varphi\) bounded, modulo constants & Bounded perturbation of a Hamiltonian\\
\end{decoder}

Traps: \(h\) is a score, not the information vector \(\mathsf h\) of later chapters; a block here is independent
copies of one scalar law, not interacting agents; the label prefixes \texttt{rg} and \texttt{grg} are historical,
and no Wilsonian flow is built anywhere in Supplement G. For a channel from \(b\) copies to an output \(z\), the
pushed score is \(\mathscr L_bh=\E[\sum_ih(x_i)\given z]\) and
\[
b\lVert h\rVert^2-\lVert\mathscr L_bh\rVert^2
=\int\operatorname{Var}\Bigl(\sum_{i=1}^bh(x_i)\Bigm| z\Bigr)\nu^{(b),c}(dz)\geq0
\]
(\SG{thm:rg-score-pushforward-defect}): Fisher information of \(b\) copies, minus what survives, equals the
expected conditional variance of the block score, the Rao--Blackwell form of information loss. Equality is
local sufficiency at one parameter, not recovery of the experiment. For the normalized sum,
\[
\mathscr L_b\widehat{\mathrm{He}}_k=b^{1-k/2}\widehat{\mathrm{He}}_k\quad(k\geq1),
\qquad
\sigma(\mathscr L_b)=\{b^{1-k/2}:k\geq1\}\cup\{0\},
\]
with \(\widehat{\mathrm{He}}_k=\mathrm{He}_k/\sqrt{k!}\), from Mehler's identity
\(\E[\mathrm{He}_k(\xi_i)\given\bar\xi]=b^{-k/2}\mathrm{He}_k(\bar\xi)\) times the replication factor \(b\)
(\SG{thm:rg-gaussian-hermite-spectrum}). The operator is compact, injective, with norm \(\sqrt b\) and dense
nonclosed range: no component is destroyed, but undoing a step multiplies degree \(k\) by \(b^{k/2-1}\), so
reconstructing fine structure is ill-posed even here.

At a nondegenerate Gaussian on \(\R^n\), the family tangent (polynomials of degree at most two, dimension
\(n+n(n+1)/2\)) and the bounded actions (dense and proper in \(L^2_0\), \SG{prop:rg-action-score-isometry}) meet
only at zero, since nonzero polynomials are unbounded (\SG{rem:ig-three-tangents}): Gaussian mean and covariance
moves lie outside the main manuscript's analytic action calculus. And a direction can be a legitimate tangent
with no exponential tilt: a Gaussian tilted by \(e^{-t\mathrm{He}_3}\), a cubic potential, is normalizable for no
\(t\neq0\), and by \(e^{-t\mathrm{He}_4}\) only for the stable sign, exactly the familiar \(\phi^3\) and wrong-sign
\(\phi^4\) instabilities (\SG{prop:ig-hermite-exponential-domain}). For the program, the spectral theorem concerns a
noninteracting block of scalar copies and says nothing about aggregating coupled \(K\)-dimensional agents.

\section{Aggregation and Kron Reduction (G7)}
\label{comp:sg2-coarse-maps}

\subsection{The question, the answer, and the picture}

Which coarse maps of the interaction family are exact, and when does the family close under them? In
link-trivialized coordinates the family is a network of matrix springs and pins,
\[
\mathsf E(z)=\frac12\sum_iz_i^\top A_iz_i+\frac12\sum_{i<j}(z_i-z_j)^\top W_{ij}(z_i-z_j),
\qquad A_i\succeq0,\quad W_{ij}=W_{ji}\succeq0
\]
(\SG{def:gauss-interaction-family}): the self term or anchor \(A_i\) pins agent \(i\) to the origin of its fiber
and \(W_{ij}\) is a matrix spring on the disagreement. Membership is declared (\SG{hyp:gauss-global-interaction});
directed models do not produce it in general, and for \(K,N\geq2\) it is a null set of precisions
(\SG{prop:gauss-interaction-family-thin}).

Hard aggregation, which identifies each cluster with one coarse coordinate, always closes: self terms and cut
weights add and internal edges drop out (\SG{prop:cg-gaussian-aggregation}); the coarse law is the fine law
conditioned on zero within-cluster disagreement (\SG{prop:cg-aggregation-conditioning}) and a
precision-dependent linear pushforward (\SG{prop:ig-galerkin-pushforward}). Kron reduction, which integrates hidden
agents out, is exact as a map of laws (\SG{prop:cg-normalized-schur-action}) but leaves the family in general
(\SG{prop:cg-kron-leaves-family}). Its closure regimes are classified in \SG{tab:cg-kron-regimes}, summarized by
the G7 rows of \cref{comp:sg2-scorecard}; the headline is that on a fixed flat graph with independent
positive-definite weights and \(K\geq2\), closure for every assignment holds exactly when every hidden component
touches at most two retained terminals (\SG{thm:cg-kron-boundary-classification}).

Aggregation makes each cluster a rigid body: internal springs are never stretched, member pins add, and cut
springs add in parallel. Kron reduction lets hidden agents move and integrates them out, so retained agents feel
effective springs softened by hidden fluctuations. One restricts the energy, the other marginalizes the law, and
the marginal precision is dominated by the restricted one (\SG{prop:restrict-marginal-vs-restricted-precision}).

\begin{analogy}{exact}
\(S^\top\Lambda S\), with \(S\) the Boolean membership matrix tensored with \(I_K\), is the Galerkin coarse operator of
aggregation algebraic multigrid \citep{Notay2010Aggregation,LivneBrandt2012LAMG}: the coarse energy is the fine
energy on rigid-block configurations. Kron reduction is integrating out Gaussian fields; for scalar weights it is
the star--mesh reduction of an electrical network \citep[Lemma~II.1]{DorflerBullo2013Kron}, and its offset
increment, led by \(\tfrac12\log\det M\), is the exact free energy of the eliminated modes.
\end{analogy}

\begin{analogy}{structural}
Aggregation is a Kadanoff block-spin map with a particular block spin: the coarse law is the law of
\(P_\Lambda y=(S^\top\Lambda S)^{-1}S^\top\Lambda y\), the rigid-block configuration closest to \(y\) in the energy norm
(\SG{prop:ig-galerkin-pushforward}(i)), not the plain block average. Since this block spin depends on the
couplings, aggregation is one parameter-independent channel only at fixed precision; no single kernel realizes
it across the interior of the interaction cone (part (v)). Conditioning on the null event of zero disagreement
also depends on the defining statistic, the Borel--Kolmogorov phenomenon behind the Jacobian of a delta-function
constraint (\SG{prop:cg-aggregation-conditioning}); the linear constraint is part of the definition.
\end{analogy}

\subsection{Objects and formulas}

\begin{decoder}
Galerkin law; \(P_\Lambda\) & Coarse law of rigid clusters; energy-weighted block value & Block spin, coupling-dependent\\
Hidden, retained, terminal, ground & Eliminated agents; kept agents; kept agents adjacent to a hidden component; a kept zero vertex carrying each pin as an edge & Interior; boundary; Dirichlet ground\\
Cone \(\mathcal W_H\) & All \(H\operatorname{diag}(d)H^\top\), \(d\geq0\), for one \(H\in\GL(K)\) & Couplings with common normal modes\\
Positive two-node presentation & Sum of positive-semidefinite terms on at most two nodes each & Nonnegative bond energies\\
\end{decoder}

Traps: ``trace'' in ``graph-exponential trace closure'' and ``hard trace'' means restricting the energy to
identified configurations, the opposite of tracing out; ``universal closure'' quantifies over weight assignments
on one graph; the pivot \(M\) is unrelated to the \(M\) of \cref{comp:sg2-operator,comp:sg2-reciprocal}.

\paragraph{Aggregation.} It gives \(A_I=\sum_{i\in I}A_i\), \(W_{IJ}=\sum_{i\in I,j\in J}W_{ij}\),
\(\mathsf h_I=\sum_{i\in I}\mathsf h_i\), and annihilates internal weights (\SG{prop:cg-gaussian-aggregation}).
Internal edges vanish because the pair energy depends only on \(z_i-z_j\): a translation-invariant pair statistic
has constant diagonal, so node parameters simply add (\SG{thm:cg-graph-exponential-closure}, which closes energies,
not normalized laws). With \(B_\perp\) an orthonormal basis of the complement of \(\operatorname{range}S\),
\[
p\bigl(dy\bigm|B_\perp^\top z=0\bigr)
=\mathcal N\bigl((S^\top\Lambda S)^{-1}S^\top\mathsf h,(S^\top\Lambda S)^{-1}\bigr),
\]
the event being \(z_i=z_j\) inside each cluster. At fixed precision the pushed Fisher metric of the mean is
\(\Lambda S\Lambda_{\mathrm c}^{-1}S^\top\Lambda\preceq\Lambda\), with equality on rigid-block directions
(\SG{prop:ig-galerkin-pushforward}(ii)): aggregation loses exactly the information in internal deformations, at
every finite depth. When the precision moves, so does the channel (\SG{open:grg-galerkin-monotonicity}).

\paragraph{Kron reduction.} With hidden block \(M=\Lambda_{\mathrm{hh}}\),
\[
\Lambda^{\mathrm{eff}}=\Lambda_{\mathrm{rr}}-\Lambda_{\mathrm{rh}}M^{-1}\Lambda_{\mathrm{hr}},\qquad
\mathfrak c^{\mathrm{eff}}=\mathfrak c+\tfrac12\log\det M-\tfrac{n_{\mathrm h}}2\log2\pi
-\tfrac12\mathsf h_{\mathrm h}^\top M^{-1}\mathsf h_{\mathrm h}
\]
(\SG{prop:cg-normalized-schur-action}): direct stiffness minus softening through hidden agents, and an offset that
adds the free energy of the eliminated modes, so raw mass and evidence are unchanged. Eliminating one agent \(k\) adds
\(W_{ik}M^{-1}W_{kj}\) between each pair of its retained neighbors. For \(K=1\) this is \(w_{ik}w_{kj}/m\geq0\), and scalar
networks always close (\SG{lem:cg-scalar-kron-closure}). For matrices it need not be symmetric:
\(W_{13}=\operatorname{diag}(1,2)\), \(W_{23}=\begin{psmallmatrix}2&1\\1&2\end{psmallmatrix}\), and unit pins
manufacture \(\tfrac1{19}\begin{psmallmatrix}9&3\\4&14\end{psmallmatrix}\), which is no spring
(\SG{prop:cg-kron-leaves-family}). For a relay with hidden pin \(Q_0\) and legs \(W,W'\) the output is symmetric exactly
when \(WQ_0^{-1}W'=W'Q_0^{-1}W\), the legs commuting after whitening by the pin; at a fixed positive-definite hidden
pin the families passing every such test are exactly the subfamilies of one cone \(\mathcal W_H\), \(HH^\top=Q_0\)
(\SG{thm:cg-independent-kron-classification}), and such a cone closes on every graph
(\SG{thm:cg-congruence-diagonal-kron}). In the frame \(H\) the network is \(K\) decoupled scalar networks.

\paragraph{The graph criterion and large \(N\).} With arbitrary independent positive-definite weights, treat each pin
as an edge to ground. A hidden component with one terminal contributes nothing; with two it is one effective matrix
spring (along a path compliances add, \(W^{-1}=\sum_jW_j^{-1}\)); with three or more it is a matrix star--mesh
transform, nonsymmetric for some assignment. So for \(K\geq2\) universal closure holds iff every hidden component
has at most two terminals, ground counted (\SG{thm:cg-kron-boundary-classification}). With a proper prior on every
hidden agent, ground is always a terminal: closure iff each hidden component touches at most one retained agent,
and otherwise the failing assignments contain an open dense set of full measure (\SG{cor:cg-kron-anchored-agents}).

This is the operative statement for integrating out agents at large \(N\). With priors on all agents, \(K\geq2\), and
independent matrix weights, integrating out a relay between two retained agents generically leaves the positive
matrix-spring family, while a hidden cluster hanging from one retained agent only stiffens its pin. Decimation,
which keeps some agents and eliminates those between them, is made of relays. The coarse law stays an exact
Gaussian; what fails is its presentation by positive pairwise springs. Degree-two hidden agents with zero pin
close in the marked nonflat class; with pins they close with a generated pin and an induced link the law does not
determine (\SG{prop:cg-nonflat-series-anchored}). Branching is worse: a proper four-agent Gaussian has a three-node
marginal that is no sum of positive terms on at most two nodes, even with new invertible links and redistributed
pins (\SG{prop:cg-kron-positive-pair-obstruction}). The action is still quadratic (M\"obius order two), so the
obstruction is positivity of the pairwise split, not a three-body term; existence is proved, not frequency.

\paragraph{Series elimination with links.} For degree-two hidden agents with zero pin
(\SG{thm:cg-nonflat-series-kron}),
\[
\Theta_{\mathrm{path}}=\Theta_1\cdots\Theta_m,\qquad
W_{\mathrm{path}}^{-1}=\sum_{i=1}^m\Theta_{\leq i-1}W_i^{-1}\Theta_{\leq i-1}^\top,\qquad
Q_{\mathrm{loop}}=(I-\mathrm{Hol})^\top W_{\mathrm{path}}(I-\mathrm{Hol}).
\]
Links multiply and compliances add after transport, resistors in series made matrix-valued and gauge covariant; a
returning path leaves a self term fixed by its holonomy. The law sees only \(Q_{\mathrm{loop}}\): \((2I,W)\) and
\((3I,W/4)\) give the same quadratic, so a folded holonomy is structural data the law does not determine.

\section{Nonflat Links, Relaxation, and Quotient Volume (G8)}
\label{comp:sg2-nonflat}

\subsection{The question, the answer, and the picture}

A twisted internal link leaves a positive-semidefinite self term \((I-\Theta_e)^\top W_e(I-\Theta_e)\) on its
cluster, which only a singular weight can hide (\SG{prop:cg-effective-support}). A Gaussian mean relaxes onto the
fixed sector exponentially, by a flow that is exactly a Fisher natural gradient (\SG{prop:cg-gaussian-fixed-relaxation}).
The Gaussian holonomy-constrained parent is the moment match (\SG{prop:cg-gaussian-forward-kl-barycenter}),
Haar-projected under compact holonomy (\SG{cor:cg-compact-holonomy-barycenter}); no parameter-independent kernel
realizes the barycenter of two normals for all parameters (\SG{prop:hist-normal-barycenter-kernel}). Hard
membership carries a Jacobian that changes raw normalizers only (\SG{prop:cg-quotient-metric}). No partition is
claimed to be selected by volume or transported divergence.

Force two agents joined by a twisted spring to be one agent. The spring wants \(z_i=\Theta_ez_j\), identification
imposes \(z_i=z_j=y\), and every \(y\) with \((I-\Theta_e)y\neq0\) stretches it: a quadratic in \(y\) alone, a mass term
for the block. The edge has become a self-loop with holonomy \(\Theta_e\). The same bookkeeping gives the folded-loop
term of \cref{comp:sg2-coarse-maps} and the generated self terms of \cref{comp:sg2-reciprocal}: frustration inside a
block becomes mass after blocking.

\begin{analogy}{structural}
For \(K=2\), \(W_e=wI\), and \(\Theta_e\) a rotation by \(\phi\), the residue is \(4w\sin^2(\phi/2)I\), the energy
\(\lvert1-e^{i\phi}\rvert^2\) of hopping back to one site with phase \(\phi\); a flux through a ring lifts the uniform
mode of a hopping Laplacian in the same way. Here the flux is the declared link and the lifted mode is the rigid
motion of the block.
\end{analogy}

\subsection{Objects and formulas}

\begin{decoder}
Weighted invisibility & \(W_e^{1/2}(I-\Theta_e)=0\): the twist acts only where the weight is blind & Twist in an uncoupled direction\\
Forward-KL barycenter & Gaussian minimizing \(\sum_iw_i\KL(q_i\Vert q)\) & Moment matching; mean-field parent\\
\(\mathsf G_{\mathcal P}\), \(\operatorname{Jac}_{\mathcal P}\) & Metric on cluster coordinates modulo translation; its volume factor & Collective-coordinate Jacobian\\
\end{decoder}

\paragraph{Internal residue.} The residue vanishes exactly when \(W_e^{1/2}(I-\Theta_e)=0\)
(\SG{prop:cg-effective-support}, from \SG{lem:gauss-internal-edge}); positive-definite weights force \(\Theta_e=I\). A
common quotient on which every edge is positive definite exists exactly when all weights share one kernel preserved
by every link, and only then is group-valued holonomy defined on the effective support. In part (v) each edge of a
triangle is invisible and the loop product is \(2I\), which fixes no vector, yet the zero-energy sections are all
constant sections: with singular weights the kernel of the connection energy, not the loop product, governs the
zero modes.

\paragraph{Relaxation.} For constant \(\mathsf G\succ0\), \(L\succeq0\), the flow \(\mathsf G\dot z=-Lz\) converges to the
\(\mathsf G\)-orthogonal projection \(\Pi_0z_0\) onto \(\ker L\),
\[
\lVert z(t)-\Pi_0z_0\rVert_{\mathsf G}\leq e^{-r_+t}\lVert z_0-\Pi_0z_0\rVert_{\mathsf G},\qquad
r_+=\lambda_{\min}^+\bigl(\mathsf G^{-1/2}L\mathsf G^{-1/2}\bigr)
\]
(\SG{prop:cg-gaussian-fixed-relaxation}): overdamped relaxation with conserved zero modes and a spectral gap, which
for \(\mathsf G=I\) is consensus dynamics. Since \(\KL(\mathcal N(z,\mathsf G^{-1})\Vert\mathcal N(w,\mathsf G^{-1}))
=\tfrac12(z-w)^\top\mathsf G(z-w)\), it is exactly the natural gradient on that family; elsewhere the Fisher metric
matches relative entropy only to second order. \(\mathsf G\) must transform with the frame; \(\mathsf G=I\) is
equivariant only under orthogonal reframings.

\paragraph{Gaussian parents.} Over nondegenerate Gaussians,
\(\sum_iw_i\KL(\mathcal N(\mu_i,\Sigma_i)\Vert\mathcal N(\mu,\Sigma))\) is minimized by
\[
\bar\mu=\sum_iw_i\mu_i,\qquad
\bar\Sigma=\sum_iw_i\bigl[\Sigma_i+(\mu_i-\bar\mu)(\mu_i-\bar\mu)^\top\bigr],\qquad
\mathcal D_I^\star=\tfrac12\Bigl[\log\det\bar\Sigma-\sum_iw_i\log\det\Sigma_i\Bigr]
\]
(\SG{prop:cg-gaussian-forward-kl-barycenter}): within-constituent plus between-mean spread, and a score equal to the
parent's entropy minus the mean constituent entropy, since \(\tfrac12\log\det(2\pi e\Sigma)\) is Gaussian entropy. Under
a holonomy group with compact closure the parent is the moment match of Haar-averaged moments
(\SG{cor:cg-compact-holonomy-barycenter}); for \(\{\pm I_2\}\) and the single law \(\mathcal N(u,I_2)\) it is
\(\mathcal N(0,I_2+uu^\top)\), so the constraint bites without disagreement. Compactness matters: a loop holonomy with
an eigenvalue off the unit circle fixes no nondegenerate Gaussian, and for independent \(\GL(K)\) links
\(\lvert\det\rvert=1\) on a loop is a null condition (\MT{cor:cg-holonomy-noncompact-gaussian}). For the program's
\(\GL(K,\R)\) links the constrained Gaussian parent is empty off a null set; a compact group such as
\(\operatorname{SU}(N)\) in its real representation meets the hypothesis. Nor is the barycenter a channel: a kernel
sending \(\mathcal N(\mu,\Sigma_0)\) to \(\mathcal N(\mathsf H\mu+\mathsf b,\Sigma_1)\) for all \(\mu\) must add noise
\(\Sigma_1-\mathsf H\Sigma_0\mathsf H^\top\) (\SG{prop:hist-normal-barycenter-kernel}), and the equal-width barycenter of two
normals needs noise \(\varsigma^2/2\), which depends on the width. Merging two agents into one of the same width is a
meta-agent map to which data-processing inequalities do not apply.

\paragraph{Quotient volume.} Modulo translation, hard membership induces
\(\mathsf G_{\mathcal P}=\operatorname{diag}(n_I)-\mathbf n\mathbf n^\top/N\), with
\(\operatorname{pdet}\mathsf G_{\mathcal P}=\lvert\mathcal P\rvert\prod_In_I/N\) and
\(\operatorname{Jac}_{\mathcal P}=(\operatorname{pdet}\mathsf G_{\mathcal P})^{K/2}\) (\SG{prop:cg-quotient-metric}): the kinetic
metric of rigid clusters with the center of mass removed (for two clusters, twice the reduced mass). The Jacobian
cancels from normalized laws and shifts raw log normalizers, so partition scores built on raw normalizers depend on
a volume convention; selection stays open (\MT{open:cg-natural-selector}).

\section{The Operator Arrow (G9)}
\label{comp:sg2-operator}

\subsection{The question, the answer, and the picture}

On operator parameters, aggregation is a positive linear semigroup action: self terms and cut weights add,
internal edges vanish, nested aggregations compose (\SG{prop:grg-aggregation-semigroup}). On the interior of the
interaction cone a step rescaled by a declared \(\zeta_\ell>0\) is a normalized Gaussian law step, realized by one
parameter-independent kernel on a fixed-precision fiber exactly when \(\zeta_\ell\leq1\)
(\SG{prop:grg-interior-law-step}). The only positive rules additive in scalar node masses are biadditive,
\(A_i=x_iA\), \(W_{ij}=x_ix_jM\), closed under every hard partition (\SG{prop:grg-positive-additive-linear}). Nothing is
asserted about infinite hierarchies, limiting rays, or iterates.

Two facts keep aggregation from being a flow. Each step changes the number of agents, so steps compose maps between
different spaces, and a fixed point needs a hierarchical identification of successive parameter spaces, which the
supplement does not supply. And the bare map multiplies stiffness by the block size: a rigid block of \(b\) agents
stretches \(b\) pins, and for equal cluster sizes
\[
b\lambda_{\min}(\Lambda)\leq\lambda_{\min}(S^\top\Lambda S),\qquad
\lambda_{\max}(S^\top\Lambda S)\leq b\lambda_{\max}(\Lambda)
\]
(\SG{prop:grg-aggregation-semigroup}(iv)). With \(b\geq2\) iterates run to infinite stiffness, so no nontrivial fixed
point can appear without a compensating rescaling. The step \(\Lambda_{\ell+1}=\zeta_\ell^{-1}S_\ell^\top\Lambda_\ell S_\ell\)
(\SG{def:grg-operator-arrow}) types that rescaling, which plays the part, structurally, of the field-rescaling step
of a Kadanoff--Wilson transformation \citep{kadanoff1966scaling,wilson1975renormalization}; with no length to rescale at
one base point, only it remains. The supplement fixes no \(\zeta_\ell\) and defers fixed rays and attraction to
Supplement B.

\begin{analogy}{exact}
The biadditive boundary member \(A=0\) is a massless Gaussian field on a finite graph: consensus configurations are
zero modes (\SG{prop:ig-generalized-spectrum-localization}), no normalized full-space law exists, and the pencil
\((L,\Lambda)\) is irregular (\SG{prop:grg-pencil-regularity}). The named repairs are the infrared regularizations of a
zero mode: keep a mass, pin a site, or quotient the zero mode. They define different problems, and a pseudoinverse
in the Lebesgue density is none of them.
\end{analogy}

\subsection{Objects and formulas}

\begin{decoder}
\(\zeta_\ell\); information-vector rule & Declared rescaling of the coarse precision; how the source is carried & Field rescaling; source rescaling\\
Interior, boundary & \(\Lambda\succ0\); a singular member such as \(A=0\) & Massive; massless\\
Declared repair & Pin, quotient, or retained mass & Infrared regulator\\
\end{decoder}

\paragraph{The step as a law.} With \(\mathsf h_{\ell+1}=S_\ell^\top\mathsf h_\ell\), the step sends
\(\mathcal N(\mu,\Lambda_\ell^{-1})\) to \(\mathcal N(\zeta_\ell P_{\Lambda_\ell}\mu,\zeta_\ell\Lambda_{\mathrm c}^{-1})\),
\(\Lambda_{\mathrm c}=S_\ell^\top\Lambda_\ell S_\ell\) (\SG{prop:grg-interior-law-step}). At \(\zeta_\ell=1\) it is the Galerkin
pushforward; for \(\zeta_\ell<1\) it adds noise, via \(y\mapsto\mathcal N(\zeta_\ell P_{\Lambda_\ell}y,\zeta_\ell(1-\zeta_\ell)
\Lambda_{\mathrm c}^{-1})\); for \(\zeta_\ell>1\) no kernel exists on that fiber, because the coarse mean family would
carry more Fisher information than the fine one. The rule \(\mathsf h_{\ell+1}=\zeta_\ell^{-1/2}S_\ell^\top\mathsf h_\ell\) makes every \(\zeta_\ell\)
a change of units. How the source is rescaled with the precision decides whether the step is a channel.

\paragraph{The biadditive family.} Under every hard partition,
\[
A_i=x_iA,\quad W_{ij}=x_ix_jM
\qquad\longmapsto\qquad
A_I=x_IA,\quad W_{IJ}=x_Ix_JM,\qquad x_I=\sum_{i\in I}x_i,
\]
and conversely a positive-semidefinite rule additive in the masses must be \(xA\) and \(xyM\), because positivity makes
each scalar projection monotone and monotone additive functions are linear (\SG{prop:grg-positive-additive-linear}).
The family is an all-to-all coupling with a rank-one scalar profile, a matrix mean-field coupling proportional to a
product of masses. What selects it is this additivity hypothesis with positivity; neither is attraction.
Translation invariance works one level earlier: difference-only pair energies make internal edges vanish and cut
weights plain sums (\SG{thm:cg-graph-exponential-closure}), which is why additivity in masses is the natural
equation. The supplement does not derive the biadditive form from translation invariance alone, and it says that
closure is not selection. Under \(\Gamma_b\supseteq\GL^+(K)\), which \(\GL(K,\R)\) satisfies, the factor \(M\) carries only
its rank (\SG{prop:grg-residual-fiber-label}); orthogonal frames would leave \(K\) eigenvalues.

\begin{analogy}{structural}
In the multiscale network renormalization of \citet{Garuccio2023Multiscale}, two blocks are linked when any pair of
members is, a nonlinear rule for probabilities, and the invariant linking probability is \(1-e^{-\delta x_ix_j}\) with
additive \(x_i\): the product of masses appears in the additive variable \(-\log(1-p_{ij})\). Gaussian weights add under
aggregation, so the equation is linear from the start and \(x_ix_jM\) is its only positive solution. The parallel is in
form only.
\end{analogy}

\section{Reciprocal Pairs, Anchored Stars, and Loop Holonomy (G10)}
\label{comp:sg2-reciprocal}

\subsection{The question, the answer, and the picture}

A tower in which each level's prior comes from the level above has a top with nothing above it; a tempting closure
lets two agents supply each other's priors (the general question is \MT{sec:full-graph-closing-tower}). For two
latents joined by edge copies in opposite directions, the kernel of the precision is the transported fixed space of
the loop holonomy (\SG{prop:obs-reciprocal-pair-kernel}), and the determinant factors through \(\det(I-\mathrm{Hol}(p))\)
(\SG{prop:obs-holonomy-determinant-factorization}). The flat unanchored fold is therefore singular for all covariances
(\SG{cor:obs-flat-fold-singular}); an anchor, or a holonomy without eigenvalue one, restores a law whose normalizer
depends on the links (\SG{prop:obs-normalizer-link-dependence}). The standard replacement, an apex with a proper prior,
has a mean-field schedule converging at the largest squared canonical correlation
(\SG{thm:obs-star-fixed-point-contraction}). A strict coarse-graining with flat clusters and bare-link cuts forces
trivial holonomy everywhere (\SG{prop:obs-no-new-anchor-exit}), a conflict the default family does not have. Only one
flat unanchored potential is singular; globally normalized cyclic Gaussians are legitimate.

\begin{analogy}{exact}
Two masses joined by springs and tied to nothing have a free center of mass and a non-normalizable Boltzmann weight.
The flat fold is exactly this: each copy charges only disagreement after transport, so transport-consistent
configurations cost nothing and the precision has a \(K\)-dimensional kernel. Mutual priors carry only relative
information; an anchor is a spring to the wall.
\end{analogy}

Holonomy is the other cure. The conditions \(y_1=\Theta_ey_2\) and \(y_2=\Theta_{e'}y_1\) agree only on the fixed space of
\(\mathrm{Hol}(p)=\Theta_{e'}\Theta_e\), so without eigenvalue one only the origin costs nothing; structurally, a flux through
a ring likewise removes the zero mode and makes the free energy depend on the flux, here through
\(\det(I-\mathrm{Hol})\). The supplement is
explicit about the price: the fiber's origin now does a prior's work, and nontrivial holonomy needs at least one link
declared independently of the frames, since frame comparisons telescope to the identity around every loop
(\MT{prop:geo-trivializing-criterion}). A declared prior is replaced by a declared link, not removed.

\subsection{Objects and formulas}

\begin{decoder}
Edge copy, reversal & One oriented interaction; parallel copies are distinct; a reversal carries the inverse link & One bond of a multigraph\\
Coboundary form & Links built from frame comparisons, telescoping around loops & Pure gauge\\
Strict no-new-anchor family & Presentations with each cut one bare link and no new self term & Coarse-graining that keeps pure gauge\\
\end{decoder}

\paragraph{Kernel, determinant, normalizer.} With \(\delta_e=(I,-\Theta_e)\), \(\delta_{e'}=(-\Theta_{e'},I)\), and
\(\Sigma_e,\Sigma_{e'}\succ0\),
\[
\Lambda=\delta_e^\top\Sigma_e^{-1}\delta_e+\delta_{e'}^\top\Sigma_{e'}^{-1}\delta_{e'},\qquad
\ker\Lambda\cong\ker(\mathrm{Hol}(p)-I),\qquad
\det\Lambda=\frac{\det(I-\mathrm{Hol}(p))^2}{\det\Sigma_e\det\Sigma_{e'}}.
\]
The kernel ignores the covariances, links enter the determinant only through the holonomy, and the link-dependent
log normalizer \(-\log\lvert\det(I-\mathrm{Hol})\rvert\) diverges on a singular locus larger than the flat point
(\(\mathrm{Hol}=\operatorname{diag}(1,2)\) is singular). Under \(\GL^+(K)\) the pair is proper for generic holonomy; under
\(\operatorname{SO}(K)\) with \(K\) odd every rotation has an axis and the unanchored pair is always singular
(\SG{cor:obs-holonomy-kernel-shrinkage}). For \(K=1\), links \(\mathsf g,\mathsf g^{-1}\), and anchor \(a_0\),
\(\det(\Lambda+a_0I)=a_0^2+a_0(\mathsf g+\mathsf g^{-1})^2\) (\SG{prop:obs-normalizer-link-dependence}): varying a declared
link changes the normalizer, so the global term cannot be treated as a per-agent constant while links are learned.
This is a change of model, not of coordinates, and it fixes no sign for the full objective.

\paragraph{The anchored star.} For an apex \(y_0\sim\mathcal N(0,A_0^{-1})\) with
\(y_i\given y_0\sim\mathcal N(\Theta_iy_0,\Sigma_i)\), precisions add, \(\Lambda_{00}=A_0+\sum_i\Theta_i^\top\Sigma_i^{-1}\Theta_i\),
and on transport-consistent configurations the energy is exactly \(v^\top A_0v\) (\SG{prop:obs-anchor-coercivity}): the
prior supplies exactly the coercivity the fold lacks. Updating constituents from the apex, then the apex, converges to
the exact mean,
\[
\lVert\mu_0^{(t+1)}-\mu_0^\star\rVert_{\Lambda_{00}}\leq\varrho_\star\lVert\mu_0^{(t)}-\mu_0^\star\rVert_{\Lambda_{00}},
\qquad
\varrho_\star=\lambda_{\max}\bigl(\Lambda_{00}^{-1/2}B_\star\Lambda_{00}^{-1/2}\bigr)<1
\]
(\SG{thm:obs-star-fixed-point-contraction}), with \(B_\star\) the part of the apex precision routed through the
constituents and \(\varrho_\star\) the largest squared canonical correlation between apex and constituents. The recursion
is block Gauss--Seidel, whose rate also governs the deterministic-scan Gibbs sampler \citep{RobertsSahu1997}. Without
observations \(\varrho_\star=1-\lambda_{\min}(\Lambda_{00}^{-1/2}A_0\Lambda_{00}^{-1/2})\); read at large \(N\) with fixed
\(A_0\), \(\Lambda_{00}\) grows and \(\varrho_\star\) approaches one, the familiar slowing of Gibbs samplers for hierarchical
models. A singular \(A_0\) without observations gives \(\varrho_\star=1\).

\paragraph{Strict flat coarse-graining against holonomy.} For a cut between trivialized clusters the completed square
(\MT{lem:cg-cut-completed-square}) is
\[
\sum_e(z_I-\Theta_ez_J)^\top W_e(z_I-\Theta_ez_J)
=(z_I-\bar\Theta z_J)^\top\Lambda_{II}(z_I-\bar\Theta z_J)+z_J^\top\Delta z_J,
\]
so the bare link required by the strict family (\MT{def:cg-no-new-anchor-pair}) exists iff the excess \(\Delta\)
vanishes, that is, iff all cut twists coincide. Nested flat clusters ending in one cluster or a bare two-cluster cut
then gauge away every link, and every closed walk has trivial holonomy (\SG{prop:obs-no-new-anchor-exit}). The default
family admits \(\Delta\), exactly the anchor the pair needed; no law-level presentation of a cut with excess is balanced
(\MT{prop:cg-no-balanced-law-presentation}), and balanced links have no continuous extension to singular cuts
(\MT{open:cg-cut-compression}, a proved no-go).

\section{The Decision Metatheory (GA)}
\label{comp:sg2-decision}

\subsection{The question, the answer, and the picture}

The Kron chapter proves specific closure theorems; the appendix asks what holds for every question of the same kind.
When fine families, coarse presentation classes, and compatibility conditions are finite semialgebraic descriptions
(finitely many polynomial equalities and inequalities, algebraic coefficients, rational matrix entries, finite graph
inventories), whether every admitted fine Gaussian has an admitted presentation of its exact marginal is decidable,
and the algorithm constructs the exact good and failure loci (\SG{thm:cg-semialgebraic-kron-classification}); the same
holds for singular laws on a declared quotient reference (\SG{thm:cg-variable-support-classification}). Three limits are
proved: no algorithm handles families given by programs that enumerate them (\SG{prop:cg-kron-enumerator-obstruction});
Borel selectors of presentations exist but continuous or gauge-equivariant ones need not
(\SG{prop:cg-kron-selector-boundary}); and feasibility across a rank change gives neither total-variation continuity
nor quadratic-mean differentiability (\SG{prop:cg-rank-change-statistical-obstruction}).

So the appendix is a completeness statement for the regime table: finite, algebraically specified closure questions
are never open in principle. The engine is quantifier elimination for the real field
(\SG{prop:cg-real-quantifier-elimination}), due to \citet{Tarski1951} and made algorithmic in
\citet{BasuPollackRoy2006}; the supplement's own work is the polynomial encoding of the exact marginal through
\(M^{-1}=\operatorname{adj}M/\det M\) and leading-minor positivity tests (\SG{lem:cg-semialgebraic-encoding}). It is also
the supplement's only treatment of singular Gaussian laws.

\begin{analogy}{structural}
A singular Gaussian needs a declared slice and measure, and the law changes unless the reference metric transforms
with the data (\SG{thm:cg-quotient-gaussian-law}), as a gauge-fixed zero mode depends on its slice and measure
\citep{FaddeevPopov1967}. The rank-change obstruction is a singular massless limit: along the precision path
\(\operatorname{diag}(1,t^2)\) a variance diverges, the Fisher information \(2/t^2\) blows up, and the laws have no limit.
\end{analogy}

\subsection{Objects and formulas}

The good and failure loci are
\[
\mathrm{Good}(x)\Longleftrightarrow\mathrm{Adm}(x)\wedge\exists y\bigl[\mathrm{Rep}(x,y)\bigr],
\qquad
\mathrm{Fail}(x)=\mathrm{Adm}(x)\wedge\neg\mathrm{Good}(x),
\]
with \(x\) the fine parameters, \(y\) a candidate presentation, and \(\mathrm{Rep}\) the polynomial statement that \(y\) is
admitted and reproduces the exact marginal, \((\det M)\Lambda_{\mathrm c}(y)=(\det M)\Lambda_{\mathrm{rr}}
-\Lambda_{\mathrm{rh}}(\operatorname{adj}M)\Lambda_{\mathrm{hr}}\), with the source likewise. Nonemptiness is decided
separately, so a vacuous closure is not mistaken for a construction, and the transcendental scalar
\(\tfrac12\log\det M\) is carried outside the formula. For singular data the covariance is the unique weighted
Moore--Penrose inverse (\SG{prop:cg-quotient-covariance-polynomial}), and unboundedly many positive pair summands repack
into at most \(n(n+1)/2\) by Carath\'eodory (\SG{prop:cg-finite-pair-presentation-bound}).

At 512 agents the appendix is a statement of principle: a fixed graph with fixed \(K\) and algebraic weights has a
definite answer, but no complexity bound useful at any practical scale is claimed, and general pair data involve of
order \(N^2K^2\) real variables. The thermodynamic limit is outside: changing-dimension limits and smooth, equivariant,
or functorial selection are open (\SG{open:cg-kron-classification-beyond-finite}).

\section{What is proved, and how solid it is}
\label{comp:sg2-scorecard}

\begin{scorecard}
\multicolumn{4}{@{}l}{\textit{G6. Score tangents and the Hermite spectrum}}\\
\SG{thm:rg-score-pushforward-defect} & Block-channel Fisher loss is an expected conditional variance & \SOLID & Standard in substance\\
\SG{prop:rg-score-block-lift} & Replication is a lift of norm \(\sqrt b\), not a kernel & \SOLID & Part (ii) a no-go\\
\SG{rem:ig-three-tangents} & Gaussian tangent meets bounded actions only at zero & \SOLID & Any nondegenerate Gaussian\\
\SG{thm:rg-gaussian-hermite-spectrum} & Normalized-sum spectrum \(b^{1-k/2}\) & \TOY & Scalar independent copies\\
\multicolumn{4}{@{}l}{\textit{G7. Aggregation and Kron reduction}}\\
\SG{hyp:gauss-global-interaction} & Precision is a spring-plus-pin network & \CHOICE & Null set for \(K,N\geq2\)\\
\SG{prop:cg-gaussian-aggregation} & Aggregation adds pins and cut weights, drops internal edges & \SOLID & Standard (Galerkin)\\
\SG{prop:cg-aggregation-conditioning} & Aggregated law conditions on zero internal disagreement & \SOLID & Linear statistic material\\
\SG{prop:ig-galerkin-pushforward} & Pushforward by \(P_\Lambda\); no kernel across precisions & \SOLID & Part (v) a no-go\\
\SG{prop:cg-normalized-schur-action} & Kron is the exact marginal & \SOLID & Standard\\
\SG{prop:cg-kron-leaves-family} & Unrestricted matrix Kron leaves the family & \SOLID & Refutation\\
\SG{thm:cg-independent-kron-classification} & Fixed-pin pair-tested families lie in one cone \(\mathcal W_H\) & \SOLID & SPD hidden pin\\
\SG{thm:cg-kron-boundary-classification} & Universal closure iff at most two terminals & \SOLID & Flat, independent, \(K\geq2\)\\
\SG{cor:cg-kron-anchored-agents} & With priors: iff one retained neighbor; else generic failure & \SOLID & Pins on hidden agents\\
\SG{prop:cg-kron-positive-pair-obstruction} & A marginal with no positive two-node form & \SOLID & Existence only\\
\SG{thm:cg-nonflat-series-kron} & Degree-two paths close, holonomy kept & \SOLID & Zero hidden pin\\
\multicolumn{4}{@{}l}{\textit{G8. Nonflat links, relaxation, and quotient volume}}\\
\SG{prop:cg-effective-support} & Twisted internal edge leaves a self term & \SOLID & Any links\\
\SG{prop:cg-gaussian-fixed-relaxation} & Exponential relaxation; exact natural gradient & \SOLID & Constant metric\\
\SG{prop:cg-gaussian-forward-kl-barycenter} & Forward-KL parent is the moment match & \SOLID & Standard\\
\SG{cor:cg-compact-holonomy-barycenter} & Invariant parent is the Haar-projected match & \COND & Compact holonomy\\
\SG{prop:hist-normal-barycenter-kernel} & No kernel realizes the two-normal barycenter & \SOLID & No-go\\
\SG{prop:cg-quotient-metric} & Membership Jacobian cancels in normalized laws & \SOLID & Any partition\\
\multicolumn{4}{@{}l}{\textit{G9. The operator arrow}}\\
\SG{def:grg-operator-arrow} & Aggregate, then rescale by declared \(\zeta_\ell\) & \CHOICE & No \(\zeta_\ell\) selected\\
\SG{prop:grg-aggregation-semigroup} & Positive, non-injective semigroup & \SOLID & Finite levels\\
\SG{prop:grg-interior-law-step} & Interior law step; kernel iff \(\zeta_\ell\leq1\) & \SOLID & Interior only\\
\SG{prop:grg-positive-additive-linear} & Biadditive family: closed, unique positive additive & \SOLID & Not attraction\\
\multicolumn{4}{@{}l}{\textit{G10. Reciprocal pairs, anchored stars, and loop holonomy}}\\
\SG{cor:obs-flat-fold-singular} & Flat unanchored fold is singular & \SOLID & One flat potential\\
\SG{cor:obs-holonomy-kernel-shrinkage} & Nonflat holonomy shrinks the kernel & \SOLID & Odd \(\operatorname{SO}(K)\) no-go\\
\SG{prop:obs-holonomy-determinant-factorization} & Links enter \(\det\Lambda\) via \(\det(I-\mathrm{Hol})\) & \SOLID & Two-residual pair\\
\SG{prop:obs-normalizer-link-dependence} & Normalizer varies with declared links & \SOLID & \(K=1\) witness\\
\SG{thm:obs-star-fixed-point-contraction} & Star schedule rate is a canonical correlation & \SOLID & Exact schedule\\
\SG{prop:obs-no-new-anchor-exit} & Strict flat coarse-graining forces trivial holonomy & \SOLID & No-go\\
\multicolumn{4}{@{}l}{\textit{GA. Decision metatheory}}\\
\SG{thm:cg-semialgebraic-kron-classification} & Finite semialgebraic Kron questions decidable & \SOLID & No complexity bound\\
\SG{prop:cg-kron-selector-boundary} & No continuous or equivariant selector in general & \SOLID & No-go\\
\SG{prop:cg-kron-enumerator-obstruction} & Enumerated families undecidable & \SOLID & No-go\\
\SG{thm:cg-variable-support-classification} & Singular laws decidable across ranks & \SOLID & Declared reference\\
\SG{prop:cg-rank-change-statistical-obstruction} & Rank change can break regularity & \SOLID & No-go\\
\end{scorecard}

Many SOLID entries are negative: every row marked a no-go or refutation, the necessity halves of the Kron
classifications, and the non-identifiability of folded holonomy marks. Each refutes a universal statement, so a small
witness suffices, and none is an open question.

\begin{hygiene}
Conditions that matter: membership in the interaction family, a modeling choice excluding almost every precision when
\(K,N\geq2\), on which all closure statements of G7 and G9 rest; \(K\geq2\) for Kron failure; the ``every independent
assignment on flat links'' quantifier of the graph criterion, so a specific correlated network may close outside it; a
positive-definite hidden pin for the fixed-anchor classification, degree two and zero pin for series closure, and pins
on all hidden agents for the anchored corollary; compact holonomy for the Haar-projected parent, which generic
\(\GL(K,\R)\) loops violate; \(\Lambda\succ0\) for law steps; fixed precision for Fisher contraction along aggregation; the
exact schedule for the star; \(\Gamma_b\supseteq\GL^+(K)\) for the rank-only label. Supplement G is real throughout.
Hygiene: standard Borel spaces, Lebesgue references, continuous conditional versions, domination for local Fisher
expansions, and encodings of algebraic coefficients.
\end{hygiene}

\section{Toy versus general, and the 512-agent target}
\label{comp:sg2-toy}

Most of the part is general in \(N\) and \(K\). Its small instances are counterexamples, which are not toy results: the
three-agent Kron witness refutes unrestricted closure at every size, and the graph theorem lifts the three-terminal
witness to every graph with a three-terminal hidden component. Only the Hermite spectrum is instance-bound: its
mechanism, that apparent growth comes from declared extensive normalization because exact arrows contract, rests on
general contraction statements, while \(b^{1-k/2}\) belongs to scalar independent copies. The one-dimensional
normalizer witness proves a negative whose mechanism, the determinant factorization, holds for the pair in every
dimension and instances the general kernel-holonomy isomorphism (\MT{prop:cg-kernel-holonomy}). For 512 or more agents
with \(\GL(K,\R)\) links the stated hypotheses give the following; the table adds no result.

{\small\renewcommand{\arraystretch}{1.15}
\begin{longtable}{@{}>{\raggedright\arraybackslash}p{0.34\textwidth}>{\raggedright\arraybackslash}p{0.60\textwidth}@{}}
\toprule
\textbf{Result} & \textbf{At \(N=512\) with \(\GL(K,\R)\) links}\\
\midrule\endhead
\bottomrule\endfoot
Aggregation, conditioning, Galerkin law & Holds for every hard partition, one step and finite depth\\
Twisted internal edge as self term & Holds for any invertible links\\
Biadditive closure, rank-only \(M\) & Holds, since \(\GL^+(K)\subseteq\GL(K,\R)\)\\
Kron closure, flat positive family & Fails generically for relays when every agent has a prior and \(K\geq2\)\\
Nonflat series closure & Holds on degree-two paths; branching can defeat every positive pair form\\
Holonomy-constrained parent & Empty off a null set of links; compact groups covered\\
Reciprocal pair & Generic nonflat pairs proper; odd \(\operatorname{SO}(K)\) always singular\\
Star contraction & Holds at any \(N\); read without observations, the rate tends to one as \(N\) grows\\
Hermite spectrum; decision metatheory & Not applicable; decidable in principle only\\
\end{longtable}}

Two gaps separate this part from a thermodynamic-limit reading of renormalization with Kadanoff blocks. It supplies
the exact one-step block map, the analogue of a recursion relation, and its finite compositions, but no hierarchical
identification, rescaling rule, or limiting protocol, all deferred to Supplement B. And its fibers are real:
\(\GL(K,\mathbb C)\) and \(\operatorname{SU}(N)\) enter only through real representations, and no statement here is
specialized to them.

\section{Open ends}
\label{comp:sg2-open}

The supplement's ledger is \SG{tab:gaussian-open}; every row bears on this part.

{\small\renewcommand{\arraystretch}{1.15}
\begin{longtable}{@{}>{\raggedright\arraybackslash}p{0.22\textwidth}>{\raggedright\arraybackslash}p{0.38\textwidth}>{\raggedright\arraybackslash}p{0.34\textwidth}@{}}
\toprule
\textbf{Item} & \textbf{What is unknown} & \textbf{What would settle it}\\
\midrule\endhead
\bottomrule\endfoot
\SG{open:grg-galerkin-monotonicity} & Monotonicity under varying precision, at the boundary, at unbounded depth & Lyapunov functional; boundary repair; projective limit\\
\SG{open:grg-stochastic-inverse} & Running aggregation backward & Priors for discarded weights and composable inverse kernels, or a no-go\\
\SG{open:grg-bayesian-bridge} & Relation to Bayesian renormalization \citep{BermanKlingerStapleton2023} & A common state space and a jointly monotone functional, or a counterexample\\
\SG{open:ig-model-space-fisher-transfer} & Model-space Fisher scale for a recognition precision & Coarse mean made generative, with nondegeneracy\\
\SG{open:exp-mass-sector} & Scaling of a retained mass coordinate & Induced mass law with an orbit theorem\\
\SG{open:exp-projection-commutation} & Information projection versus gauge or coarse maps & Intertwining theorem or counterexample\\
\SG{open:cg-kron-classification-beyond-finite} & Kron classification beyond finite encodings & Proof for a stated family, or finite reduction\\
\SG{open:obs-star-update-robustness} & Star under delayed, noisy, or asynchronous updates & Theorem with a rate, or counterexample\\
\end{longtable}}

Behind these sit a natural partition selector (\MT{open:cg-natural-selector}) and an intrinsic criterion for
recognition structure (\MT{open:restrict-intrinsic-criterion}). Fixed points, flows, and the thermodynamic limit are
declined and deferred to Supplement B. For large \(N\) the first three rows matter most: without them the exact step
says nothing about many steps.

\section{Verdict}
\label{comp:sg2-verdict}

\begin{verdictbox}{\DELIVERED}
\textbf{Claimed.} Closure of hard aggregation, identified with conditioning on zero within-cluster disagreement,
without selecting the family or a partition; refutation of unrestricted Kron closure, the congruence-diagonal and
fixed-anchor classifications, and the two-terminal criterion; an operator arrow that is a positive semigroup needing a
declared rescaling, with interior law steps and biadditive closure; reciprocal pairs, anchored stars, and loop
holonomy; a decision metatheory; and no fixed point or asymptotic scaling.

\textbf{Proved.} All of it, at the stated generality. G6 \DELIVERED{} (the spectrum is the scalar case it announces);
G7 \DELIVERED{} (explicit witnesses for every negative); G8 \DELIVERED{} (compactness is in the statement); G9
\DELIVERED{} (no flow asserted); G10 \DELIVERED{} (singularity confined to one flat potential); GA \DELIVERED{}
(decidability in principle).

\textbf{Gaps.} The part is narrower than the physicist's ``renormalization'' and says so: one exact step and its finite
compositions, hard partitions, a declared measure-zero family of real matrix-spring networks. Its positive content is
mostly standard mathematics in new typing; the new content is the Kron classification, the holonomy bookkeeping, and
the no-go results. For the 512-agent program, integrating out relays with matrix weights generically leaves the flat
positive family, holonomy-constrained Gaussian parents do not exist for generic \(\GL(K,\R)\) loops, and monotonicity
under varying precision, stochastic inversion, iteration, and limits remain open or deferred.
\end{verdictbox}

\chapter{Supplement B: Base Transport, Infinite Networks, and the Lattice Cluster}
\label{comp:sb}

\section{The question and the answer}
\label{comp:sb-question}

The main manuscript works at one base point \(c_*\): every frame is a point of one principal fiber,
the independent links on the edges of the agent network are the only transports, and the design is
\(\{c_*\}\). Supplement B collects what the theory would need if contexts varied over a base, if the
network became infinite, or if agents sat on a lattice. For each extension it asks which single-point
statements survive and exactly which additional data the extension consumes: a connection, a cover, a
design, a base measure, a scale identification, or a lattice realization (\SB{ch:base-introduction}).
Nothing in the main manuscript depends on it. This chapter reads it at overview depth with the scope
gate in force: it reports what the supplement proves and where it stops, and proposes nothing.

Part B-I answers the base question with kinematics. Relative frames, chosen connections, covariant
defects of cross maps, four independent notions of flatness, finite designs, a Fisher geometry pulled
back to the base, gauge-invariant durations for field histories, and a descent criterion are all
typed and proved, with counterexamples wherever an expected implication fails. No connection, base measure, or probability
law on fields is derived (\SB{open:prob-continuum-theory}), and graph links become base transports
only by declaration (\SB{hyp:geo-graph-base-transport}). Part B-II answers what Wilsonian vocabulary
requires: comparison maps that the coarse channels do not supply. A finite network at \(c_*\) carries
no asymptotics (\SB{prop:rg-finite-depth-no-asymptotics}), beta functions are scheme dependent, and on
Gaussian hierarchies attraction is proved only under a fixed primitive identification whose existence
is open (\SB{conj:grg-fixed-b-attraction}). No infinite agent network, exponent, or universality class
is constructed. Part B-III constructs one infinite-volume family completely: the square-lattice
\(\mathbb Z_2\) gauge--matter law at small plaquette coupling, with a coarse specification at every
exterior for every block size on a block-independent domain
(\SB{thm:rg-spatial-constrained-specification}), a differentiable renormalization map, and proved
volume limits of pressure, label selection, specific VFE, and Fisher densities
(\SB{cor:rg-thermo-native-label-limits}). Whether its sites are agents at \(c_*\) or base points is
open (\SB{open:rg-lattice-typing}).

The supplement does what its introduction says. Its title words promise more: no base transport of
beliefs is derived from inference, no infinite agent network is built, and the one constructed
infinite-volume family is untyped, abelian, and flows to the trivial fixed point
(\cref{comp:sb-verdict}).

\section{Objects and notation}
\label{comp:sb-notation}

The supplement adds its own symbol table for Parts B-I and B-II (\SB{app:base-backmatter}) and a local
table for the lattice part (\SB{tab:rg-lattice-notation}). The decoder covers what a reader meets first.

\begin{decoder}
Support \(\mathcal C_i\); design \(D=\{c_d\}\) & Contexts where agent \(i\) exists; a finite set of
sampled contexts, not a discretization of the base. & Local field support; probe points. \\
Relative frame \(h_i\) & Group element comparing belief and model frames, \(u_i^m=u_i^b\cdot h_i\). &
Bifundamental field of two gauge sectors. \\
Connections \(\omega_b,\omega_m\); transports \(\Omega_\gamma\) & Chosen rules for comparing fibers
along base curves. & Gauge potential and Wilson line. \\
Covariant defect \(\mathcal D\Phi\) & Failure of a belief-to-model cross map to commute with transport. &
\(\nabla\Phi\) for a Hom-valued field. \\
Graph flatness & Every closed walk of the agent graph has trivial link product. & Pure-gauge lattice
configuration. \\
Covariant jet \(D^\omega s\); pullback \(h^\omega_s\) & Rate of change of a belief field along the base
with frame rotation removed; its squared Fisher length. & \(D_\mu\phi\); gauged sigma-model metric. \\
Horizontal defect \(A_\Psi\) & Mismatch of fine and coarse connections under a coarse map over a base
map. & Mismatched gauge field under blocking. \\
Comparison data \(J_\ell\) & Declared identification of each blocked theory with a reference. &
Rescaling step of a Wilsonian RG map. \\
Relevant, marginal, irrelevant & Sign of a growth rate per unit log block scale, relative to declared
norms and normalizations. & Sign of an RG exponent \(y\). \\
\(z_e=\phi_x\psi_e\phi_{x+e_\iota}\) & Gauge-invariant link sign of the lattice model. & Unitary-gauge
link of the \(\mathbb Z_2\) Higgs model. \\
\(b=\tanh h\), \(t=\tanh j\) & Strengths of matter hopping \(h\) and plaquette coupling \(j\). &
High-temperature expansion variables. \\
Specification at every exterior & Coarse conditional laws defined and continuous at every boundary
condition, including null ones. & Gibbsianness of the renormalized measure. \\
Label \(\gamma\), bulk score \(\mathsf s_\gamma\) & A member of a finite catalog of lattice laws and its
extensive log-mass density. & Coexisting phases and their free energies. \\
\end{decoder}

Glyphs change meaning between parts. The letter \(b\) is the belief channel in Part B-I, a block ratio
\(b>1\) in Part B-II, and \(\tanh h\) on the lattice; \(t\) is a scale coordinate in
\SB{ch:wilsonian-limits} and \(\tanh j\) on the lattice; beta functions are \(\mathfrak b\), not the
attention rows \(\beta^x_{ij}\); on the lattice \(\rho=\tanh(h+2j)\) and \(\tau\) is a decay rate.
Words are narrower than a physicist expects. ``Pullback'' means the connection-relative pullback of
fiber tensors to the base, not Supplement G's restriction along a parameter map. ``Flat'' has four
meanings (\cref{comp:sb-flatness}). ``Connection'' can be a principal connection, the scale connection
of \SB{def:rg-scale-connection}, or a projected connection. ``Relevant'' is never a property of a flow
alone (\SB{def:rg-mode-exponents}), and ``fixed point'' is typed by tier: a fixed object at one tier need
not be fixed at another; six specific non-implications are proved in
\SB{prop:rg-fixed-object-nonimplication}.
``Declared'' means chosen rather than derived, and usually marks where modeling enters.

\section{Part B-I: base transport and multi-point designs}
\label{comp:sb-base}

\subsection{Frames, connections, and cross maps}
\label{comp:sb-frames}

Over a base, the belief and model channels of agent \(i\) may carry separately chosen local sections of
the one principal bundle, and the torsor structure makes their comparison canonical:
\[
u_i^m=u_i^b\cdot h_i,
\qquad
h_i\longmapsto g_i^bh_i(g_i^m)^{-1},
\qquad
\Theta_e^m=h_i^{-1}\Theta_e^bh_j\quad(e:j\to i).
\]
The relative frame \(h_i\) exists for any pair of representations, of any dimensions, and maps no
representation space into another (\SB{prop:geo-relative-frame}). It transforms on the left under
belief rechoices and on the right under model rechoices, like a bifundamental field. The two channels
carry one link on \(e\) exactly when the third equation holds (\SB{prop:geo-identified-links}); the main
manuscript's literal equality \(\Theta^b_e=\Theta^m_e\) is its gauge \(h_i=e_G\). The second atlas does
not enlarge the gauge group, whose active image at a context is still one copy of \(G\)
(\SB{prop:geo-two-atlas-active-image}); a genuine \(G_b\times G_m\) is an optional extension
(\SB{def:geo-product-extension}).

Connections are chosen data, with the textbook frame law \(A'=\operatorname{Ad}_gA-(\mathrm dg)g^{-1}\)
(\SB{prop:geo-connection-frame-laws}). The new object is the covariant defect of a cross map,
\[
(\mathcal D\Phi)_e(X)=T_e\Phi\bigl(H^b_eX\bigr)-H^m_{\Phi(e)}X,
\qquad\text{in the linear case}\qquad
D\Phi=\nabla^m\circ\Phi-\Phi\circ\nabla^b .
\]
The first term carries the belief connection's horizontal direction through \(\Phi\), the second is the
model connection's horizontal direction at the image. Their difference is vertical and vanishes exactly
when \(\Phi\) commutes with parallel transport along every base curve
(\SB{prop:geo-covariant-defect-parallelness}). At \(c_*\) there is no base vector and no defect. Across
scales, a coarse-graining map between principal bundles over a base map exists only when a
topological condition holds, as for the Hopf bundle, which no homomorphism \(z\mapsto z^n\) with
\(n\neq0\) maps into a trivial bundle (\SB{prop:rg-principal-scale-map-existence}).

\subsection{Four kinds of flatness}
\label{comp:sb-flatness}

At \(c_*\) every bundle is trivial and every connection vanishes, so ``flat'' can only mean graph
flatness. Over a base four conditions separate completely (\SB{prop:geo-four-notions}), with the first
decided by one Čech class common to both atlases (\SB{prop:geo-cech-class}).

\begin{center}
\small\renewcommand{\arraystretch}{1.15}
\begin{tabular}{@{}>{\raggedright\arraybackslash}p{0.2\textwidth}>{\raggedright\arraybackslash}p{0.3\textwidth}>{\raggedright\arraybackslash}p{0.42\textwidth}@{}}
\toprule
\textbf{Condition} & \textbf{Gauge-theory meaning} & \textbf{Separating witness} \\
\midrule
(a) Bundle triviality & A global gauge exists. & A flat \(S_3\)-bundle over a circle, glued by a
3-cycle, is nontrivial. \\
(b) Graph flatness & Agent-graph links are pure gauge. & Declared independently of (a), (c), (d). \\
(c) Connection flatness & \(F=0\). & \(x\mathrm dy\) on \(\R^2\): uniform field, not flat. \\
(d) Trivial holonomy on one loop & One Wilson loop is trivial. & \(x\mathrm dy\) around a circle
enclosing flux \(2\pi\); conversely an Aharonov--Bohm vortex is flat with nontrivial phase. \\
\bottomrule
\end{tabular}
\end{center}

For the program this means that flat links at \(c_*\), condition (b), say nothing about (a), (c), or
(d). The supplement also records that one principal class does not make the belief and model bundles
isomorphic (a cylinder and a Möbius band share one \(\{\pm1\}\)-bundle), and that a global section of an
associated bundle need not trivialize the principal bundle (\SB{prop:geo-cech-class}).

\subsection{When is a graph link a base transport?}
\label{comp:sb-graph-base}

When it is declared to be one. The standing hypothesis \SB{hyp:geo-graph-base-transport} chooses a
context per agent and a base curve per edge copy and requires
\(\widehat\rho_b(\Theta_e^b)=\Omega_{\gamma_e}\) and \(\widehat\rho_m(\Theta_e^m)=\widetilde\Omega_{\gamma_e}\).
Under it, graph loop products are base holonomies (\SB{prop:geo-graph-base-loop-transport}), a one-line
consequence of the homomorphism property. The supplement is explicit that the hypothesis is satisfiable
by definition, that no result realizes it for independently specified links
(\SB{open:rg-cover-graph-base-identification}), that recovering a connection from records is open
(\SB{open:rg-cover-holonomy-recovery}), and that even an exact realization concerns connection holonomy,
not bundle topology.

\begin{analogy}{structural}
At \(c_*\) the links are lattice gauge variables over a zero-dimensional base (\cref{comp:ch02}). The
hypothesis is Wilson's discretization read backwards: a lattice link is declared to be the path-ordered
exponential of a continuum potential along an assigned curve. In lattice gauge theory that is the
definition of the discretization, and here too it is a definition, not a theorem.
\end{analogy}

\subsection{Finite designs and shared links}
\label{comp:sb-designs}

A finite design is the multi-point version of \(\{c_*\}\) (\SB{def:prob-finite-design}). Under the
design-point product hypothesis, distinct design points are conditionally independent given the
structural data (\SB{hyp:gen-design-product}), and the free energy decomposes over the design
(\SB{cor:gen-design-factorwise-decomposition}); the hypothesis is substantive, since it excludes
correlation of an agent's states at two contexts beyond the structural data. If a submodel ties the
copies of one link across design points, each tied link ranging over all of \(G\), the surviving gauge
freedom is
\[
g_{d,i}^x=\bar g_i^xz_d,\qquad g_{d,j}^x=\bar g_j^xz_d,\qquad z_d\in Z(G),
\]
a context-independent rechoice times a central twist shared by both endpoints
(\SB{prop:gen-shared-link-stabilizer}). Tying a link across contexts turns local freedom along the design
into global freedom up to the center, which is \(\R^\times I\) for \(\GL(K,\R)\), \(\mathbb Z_N\) for
\(\mathrm{SU}(N)\), and trivial for \(\mathrm{SO}(3)\).

The finite construction supplies no probability law on a space of sections
(\SB{open:prob-continuum-theory}), and matching sections to finite laws at the design points determines
neither the joint nor the sections off the design (\SB{prop:prob-compatibility-nonidentifiability}). The
named obstruction is the one familiar from constructive field theory: reference measures are finite
products, and infinitely many non-probability factors give no measure, so a continuum theory must declare
a measure on fields directly. Until then every evidence-bound statement of the main manuscript and of
Supplement B is a statement about a finite design.

\subsection{Base-direction information geometry}
\label{comp:sb-pullback}

Once contexts vary, a belief field \(s\) changes along the base, and the connection separates the part of
that change that is only frame rotation:
\[
D^\omega s=\operatorname{ver}^\omega\circ Ts,
\qquad
h^\omega_s(X,Y)=g^F_{s(c)}\bigl(D^\omega s(X),D^\omega s(Y)\bigr)
\]
(\SB{def:pb-informational-pullbacks}). The first is the analogue of \(D_\mu\phi\); the second is the squared
Fisher speed of the belief along the base direction \(X\), whose square root is the Fisher--Rao distance per
unit displacement between the belief at \(c\) and the belief transported back from a neighbor. It is invariant under passive frame
changes (\SB{thm:pb-pullback-gauge-invariance}) but not under a change of connection: for
\(\omega'=\omega+\mathfrak d\),
\[
h^{\omega'}_s(X,Y)=h^\omega_s(X,Y)
+g^F\bigl(\mathsf J^s_{\mathfrak d}X,D^\omega sY\bigr)
+g^F\bigl(D^\omega sX,\mathsf J^s_{\mathfrak d}Y\bigr)
+g^F\bigl(\mathsf J^s_{\mathfrak d}X,\mathsf J^s_{\mathfrak d}Y\bigr),
\]
with \(\mathsf J^s_{\mathfrak d}\) the vertical field generated by the connection difference
(\SB{prop:pb-pullback-connection-change}). A belief constant over the base has zero pullback under the zero
connection and \(\mathrm c_1^2\mathrm dx^2\) under \(A'=\mathrm c_1\mathrm dx\), because the new connection
says the fiber is being shifted underneath it. This answers what a ``metric on the base built from beliefs'' is: a
metric relative to a chosen connection, and the documents choose none
(\SB{open:pb-canonical-pullback-geometry}). It is also only a semimetric, with radical the directions of
covariant constancy; constant rank gives a metric on a quotient vector bundle
(\SB{thm:pb-pullback-rank-quotient}), and a quotient manifold can fail even then, since the null directions
can form a contact distribution (\SB{prop:pb-contact-null-counterexample}). Combining the belief and model
channels needs declared weights (\SB{hyp:pb-weighted-product-geometry}).

Under a coarse map \(\Psi\) over a base map \(f\), fine and coarse base geometries obey the exact identity
\[
h^\omega_s-f^*\bar h^{\bar\omega}_{\bar s}=\delta_\Psi-\mathcal X_\Psi-\mathcal Q_\Psi
\]
(\SB{thm:pb-signed-base-comparison}). Here \(\delta_\Psi\) is the Fisher information lost through the fiber
channel, nonnegative for a Markov channel, the main manuscript's data-processing loss read along base
directions; \(\mathcal X_\Psi\) is a sign-indefinite cross term and \(\mathcal Q_\Psi\geq0\) a quadratic
term, both produced by the horizontal defect. With zero defect and a Markov channel the comparison is a
contraction and a
meta-agent's perceived base geometry is smaller than its children's
(\SB{cor:pb-base-fisher-contraction}, \SB{cor:pb-meta-perceived-geometry}); with nonzero defect it can be
negative at zero information loss (\SB{thm:pb-signed-positivity-criterion}). For maps induced by a principal
scale map the defect is the fundamental field of
\(\mathfrak A_{\mathcal P}=\mathcal P^*\bar\omega-\mathrm d\kappa\circ\omega\) and vanishes exactly when
\(\mathfrak A_{\mathcal P}\) takes values in the isotropy algebra of the coarse section value, with
\(\mathcal Q_\Psi\) the orbit Fisher form evaluated on \(\mathfrak A_{\mathcal P}\)
(\SB{thm:pb-isotropy-criterion}). Composed defects obey an ordered cocycle law
(\SB{thm:pb-anomaly-composition}), and the base Fisher defect composes up to an exact residual
(\SB{thm:pb-base-defect-cocycle}) whose relaxed pointwise problems are completely parameterized.

\begin{analogy}{structural}
The pullback is the kinetic metric \(\lvert D_\mu s\rvert^2\) of a gauged sigma model with the statistical
manifold as target, and its connection dependence is that of \(\lvert D_\mu\phi\rvert^2\) on \(A_\mu\). In
field theory the gauge field is dynamical or fixed by an equation of motion; the supplement supplies
neither. The isotropy criterion has the form of a Higgs mass term: a quadratic form in the gauge field
whose radical is the algebra fixing the vacuum value. The correspondence is of form only.
\end{analogy}

\subsection{Section-space histories}
\label{comp:sb-histories}

When an agent's state is a whole field, its inference history is a curve of sections, and it changes
pointwise vertically: at each context the belief moves within its fiber, and a base curve is not an
agent history (\SB{prop:hist-section-history-verticality}). On a strong Riemannian Hilbert manifold with
a free, proper, isometric gauge action, the quotient speed
\[
\lVert d\pi_Qv\rVert^2=\min_{\xi\in\mathfrak g}\mathsf g_Q(v-A_Q\xi,v-A_Q\xi)
\]
has a unique minimizer, so gauge-invariant durations are well defined
(\SB{thm:hist-strong-hilbert-quotient}); the minimization removes the velocity along gauge orbits, as a
Coulomb-type condition does. It fails for weak metrics: the \(L^2\) Fisher metric on \(H^1\) fields on a
circle gives a smooth proper quotient on which every tangent has zero speed
(\SB{prop:hist-proper-weak-fisher-quotient}), and that weak metric is itself the Fisher metric of a
natural observation family, so the obstruction is not an artifact of an exotic metric. Averaging a field
over the fibers of a base map loses a nonnegative amount of integrated Fisher information, given a Markov
fiber map, nonincreasing weights, and a jointly convex fiber chart; when the coarse metric is moreover
constant in the chart, the loss is exactly the sum of a channel loss, a weight gap, and a context variance
(\SB{thm:hist-averaging-defect}). Without joint convexity the conclusion fails: in the covariance sector of
the Gaussian chart the loss is negative on an explicit datum (\SB{prop:hist-averaging-convexity}).

\subsection{Covers, availability, and descent}
\label{comp:sb-covers}

\SB{ch:neighborhood-parents} is the part closest to the original picture of agents sprinkled over a base
with overlapping supports. Four agents on a square of contexts have strip supports that overlap in pairs
and triples with no common point (\SB{thm:rg-cover-nerve}), and each is present in a record with a smooth
availability probability \(w_i(c)\) that vanishes off its support (\SB{thm:rg-cover-availability}).
Attention then behaves as one would want over overlaps, with one qualification:
\(\beta^x_{ij}(c)=\eta^x_{ij}(c)/\alpha^x_i(c)\) is smooth where the receiver is present, while the event
law \(\eta\) is smooth everywhere and imposes no row where the receiver is absent
(\SB{prop:rg-cover-zero-occupancy}). The main manuscript's parent needs all its children present, so a
four-child parent here has empty domain; an availability parent weighting locally present children by
\(w_i/\sum w\) exists on the union of supports (\SB{thm:rg-cover-compatible-laws}). The weights are chosen,
and every compatibility statement holds by construction from one assumed ambient law
(\SB{hyp:rg-cover-ambient-law}); gluing independently supplied local laws is
\SB{open:rg-cover-local-gluing}.

The substantive result is the descent criterion: along a surjective submersion, a section descends
exactly when its image is constant on fibers, and a connection form exactly when its two pullbacks to the
fiber product agree (\SB{thm:rg-cover-descent-obstructions}). Around it, a coarse map with
context-dependent weights has defect \(O_{\mathrm c}(\mathrm dL)O_{\mathrm f}^{-1}\) even when all
connections agree (\SB{thm:rg-cover-nonflat-defect}); projecting a flat connection onto a moving quotient
creates curvature \(\bar{\mathsf g}^{-1}\mathrm{II}^\top\wedge\mathrm{II}\) from the second fundamental form,
and staged projections agree with the direct one when the intermediate quotient metric is propagated,
while resetting it can change the connection class (\SB{thm:rg-cover-projected-connection},
\SB{thm:rg-cover-quotient-composition}). Finally, a declared record
\(o_1=(\operatorname{Hol}_\iota-I)P_{\mathrm{pr}}Y+\varepsilon\) distinguishes two base-holonomy classes by a
frame-invariant statistic (\SB{thm:rg-cover-operational-holonomy}), the base companion of the main
manuscript's network-loop witness (\MT{thm:cg-loop-operational-holonomy}). The record reads the holonomy
by construction; the theorem shows that such a normalized experiment exists, not that the model's own
records are holonomy sensitive.

\begin{analogy}{exact}
The projected-connection curvature is the Gauss equation for an induced connection; with a flat ambient
connection and a moving rank-one image it is Berry curvature. The composition result says iterated
adiabatic projections agree with the direct one when each stage uses the metric induced by the last.
\end{analogy}

\subsection{How Part B-I connects to the base-manifold picture}
\label{comp:sb-vision}

Part B-I supplies the typed skeleton of the original picture: supports and Čech data, relative frames,
chosen connections, covariant defects, finite designs, a pulled-back base geometry, gauge-invariant
durations, availability-weighted attention over overlaps, and descent along base maps. It supplies no
dynamics for any base datum. The connection is chosen, not determined by inference; base measures and
weights are declared; the ambient law on a cover is assumed; and without a law on fields the evidence
bound over a base exists only on finite designs. The single-point graph links and the base transports
are related only by declaration. Part B-I is a precise list of what must be supplied when base transport
is taken up, not yet a theory of agents over a base.

\section{Part B-II: infinite networks and Wilsonian limits}
\label{comp:sb-wilson}

\subsection{What Wilsonian words require, and why finite depth has none}
\label{comp:sb-finite-depth}

The main manuscript coarse-grains finite networks by exact channels and keeps an effective action, a
closure residual, and finite-step multipliers (\cref{comp:ch11}). A Wilsonian step is ``block, then
identify the blocked theory with the original'', and the identification is extra structure: declared
comparison isomorphisms \(J_\ell\), giving the reference cocycle
\(\widehat{\mathfrak R}_\ell=J_{\ell+1}^{-1}\circ C_{\ell,\ell+1}\circ J_\ell\) (\SB{def:rg-comparison-data}).
Exponents are growth rates per unit of cumulative log block scale \(\varsigma_{n\leftarrow\ell}=\sum_k\log
b_k\),
\[
\chi_{n\leftarrow\ell}(u)=\frac{\log\lVert D_{n\leftarrow\ell}u\rVert_n-\log\lVert u\rVert_\ell}{\varsigma_{n\leftarrow\ell}},
\]
with \(D_{n\leftarrow\ell}\) the main manuscript's multipliers (\SB{def:rg-mode-exponents}). Where defined,
their canonical content is the forward Lyapunov filtration (\SB{prop:rg-lyapunov-filtration}); they are
invariant under tempered comparisons (\SB{thm:rg-tempered-comparison}), while a superexponential
comparison can change the apparent exponent (\SB{prop:rg-superexponential-distortion}); and no Oseledets splitting is
claimed, since a deterministic blocking sequence meets none of the hypotheses of the multiplicative
ergodic theorem \citep{Oseledets1968}. Even uniform projective contraction can leave an attracting
two-cycle rather than a fixed ray (\SB{prop:rg-contraction-no-fixed-ray}).

\SB{prop:rg-finite-depth-no-asymptotics} makes precise what every physicist knows about blocking a finite
system. With blocks of at least two agents the depth is at most \(\lfloor\log_2\lvert V_0\rvert\rfloor\), so
the limit defining an exponent does not exist; levels of different sizes admit no comparison isomorphism;
and exponent signs depend on the block schedule unless it is two-sided bounded: a mode doubled at each step has exponent 1 under
\(b_k=2\) and \(2/(n+1)\to0\) under \(b_k=2^{k+1}\). This does not conflict with renormalization at
\(c_*\) as a thermodynamic limit; it says that asymptotic vocabulary needs one, an infinite network or a
family in which size and depth grow together (\cref{comp:sb-universality}), whereas the main manuscript's
interaction multigraph is finite by definition (\MT{def:geo-graph-links}).

The exact step contributes no growth of its own: along an exact interaction orbit
\(\chi_{n\leftarrow\ell}(u)\leq\lvert V_n\rvert\log3/\varsigma_{n\leftarrow\ell}\), so no direction is
relevant in the native norms when \(\lvert V_n\rvert/\varsigma\to0\), and the same holds in the reference norm
when the comparisons are also tempered (\SB{cor:rg-interaction-no-relevant-mode}). The condition
\(\lvert V_n\rvert/\varsigma\to0\) is substantive: the bound grows with the number of coarse agents, and it
is vacuous for blocking by vertex ratios on a finite network. Relevance must come from a declared extensive normalization.
In the scalar Gaussian block each Hermite function of degree \(k\) has coefficient \(b^{1-k/2}\), so degree
one is relevant, two marginal, and higher degrees irrelevant, and the growth of degree one comes entirely
from the extensive replication normalization, of norm \(\sqrt b\), since the conditional expectation
contracts (\SB{def:rg-hermite-relevance}). Correlated coordinates
change this spectrum, a general \(\GL(d)\) action leaves the setting, and no link or holonomy enters
(\SB{prop:rg-hermite-scope}).

\begin{analogy}{structural}
In Wilson's renormalization group the blocking step, a conditional expectation, contracts; eigenvalues
above one come from the rescaling of fields that accompanies it \citep{wilson1974renormalization}. Part
B-II proves this division of labor for the exact step of the main manuscript. The Hermite spectrum is the
Gaussian fixed point of the central limit theorem read as a renormalization map, whose binary case is in
\citet{JonaLasinio2001}.
\end{analogy}

\subsection{Beta functions and fixed points}
\label{comp:sb-beta}

A beta function is a difference quotient on one reference space. For the action,
\(\mathfrak b^H_b[H;\varrho]=(\mathfrak R^H_b[H;\varrho]-H)/\log b\), and a change of reference
\(\varrho'=e^{-\Delta}\varrho\), \(H'=H-\Delta\), which leaves the weighted measure unchanged, shifts it by
the inhomogeneous term \(-\mathfrak b^H_b[\Delta;\varrho]\) (\SB{prop:rg-action-beta-reference-change}). A
beta function is thus a joint statement about the action and the reference trajectory. The retained beta
of a truncated family differs from the exact one by the transported closure residual over \(\log b\)
(\SB{prop:rg-retained-beta-residual}); a moving comparison frame gives even the identity step a nonzero
beta, and what survives smooth conjugacy are projective eigenvalue ratios at a fixed ray
(\SB{prop:rg-scheme-dependence}); and integer endpoints do not determine a continuous beta
(\SB{prop:rg-continuous-beta-underdetermined}).

A fixed theory is a measure pair, reference and weighted measure, fixed by the rescaled kernel. If the
kernel has a unique invariant law, every fixed pair has constant action
(\SB{prop:rg-fixed-pair-unique-ergodicity}), so a fixed theory with nontrivial action needs several
invariant laws; this depends on the typing choice that the reference is pushed by the same kernel. The
essential-spectrum theorem of the chapter has a hypothesis no exact arrow satisfies
(\SB{thm:rg-unital-essential-spectrum}), and one exact kernel, the circle doubling map, has spectral radius
1 on bounded functions and \(2^\alpha\) on Hölder functions (\SB{prop:rg-circle-norm-witness}), so a
spectrum belongs to a declared topology.

\subsection{The homogeneous Gaussian hierarchy}
\label{comp:sb-gauss-hierarchy}

On a homogeneous complete graph with equal blocks of size \(b\), common self term \(A\), and common
coupling \(W\), hard aggregation acts by
\[
A\longmapsto bA,\qquad W\longmapsto b^2W,
\]
informally, because a block has \(b\) self terms and two blocks have \(b^2\) cross couplings. On
\(\Sym^K\oplus\Sym^K\) the spectrum is \(b\) and \(b^2\), each with multiplicity \(K(K+1)/2\), so the ratio
outside the dominant sector is \(1/b\) (\SB{prop:rg-gaussian-dense-scaling}). Normalizing by \(b^2\) fixes the
coupling and contracts the self term toward a limiting ray with \(A=0\), a boundary operator whose
consensus zero modes make the Gaussian non-normalizable, so pinning, quotienting, or a small mass are
different infrared repairs. On general graphs the exponents follow the scaling of the cut: \(\xi=2\) for
the complete graph, \(\xi=0\) for a chain. This is mean-field power counting for Gaussian couplings.

It does not show attraction. Birkhoff's cone contraction cannot reach the boundary ray, since sum-only
aggregation is not primitive on the full cone and does not mix matrix directions
(\SB{prop:grg-full-cone-blocked}). Attraction is proved under one fixed primitive identification
\(B\otimes\operatorname{id}\) on edge indices:
\(r_B^{-\ell}(B^\ell W^{(0)})_f\to v_fM_\infty\) with \(M_\infty=\sum_eu_eW^{(0)}_e\) and \(u,v\) the Perron
vectors (\SB{conj:grg-fixed-b-attraction}). The spatial profile is attracted; the \(K\times K\) fiber matrix is
a Perron-weighted average of the initial couplings and is not selected. Whether any graph family admits such
an identification is \SB{open:rg-scalarized-attraction}, and nonautonomous identifications are not covered.
Attraction on an infinite hierarchy is therefore not established.

\subsection{Two-index limits, and what the absence of universality means}
\label{comp:sb-universality}

A thermodynamic scale statement has two indices, size \(n\) and depth \(\ell\), and three limits: fixed
depth after size, maximal depth, and a diagonal \(\ell(n)\) (\SB{def:rg-two-index-family}). They need not
agree: for the blocked path graph the fixed-depth limit of the spectral count is the arccos law, the
maximal-depth limit a single site (\SB{prop:rg-noncommuting-limits}). A countable index set carries a joint
law by Kolmogorov extension (\SB{thm:cg-projective-cylinder-law}), but consistency gives no ELBO limit and no
scale flow. That no universality class is identified (\SB{open:rg-two-index-limits}) means that the
supplement states only what one would need: a triangular family of finite agent networks at \(c_*\) with
consistency at fixed depth, uniform control of the coarse error, a diagonal scaling, and proved convergence
of a scheme-invariant observable (\SB{open:rg-two-index-limits}). Raw couplings and beta components are
scheme-dependent coordinates; a spectral exponent is a label only after convergence of the integrated
density of states is proved (\SB{def:rg-gaussian-spectral-labels}), and a finite graph has an atomic
spectrum. For the view of renormalization at \(c_*\) as a thermodynamic limit, the two-index family is the
typed home, and it is open.

\subsection{The numerical statements}
\label{comp:sb-numerical}

Supplement B has three NUMERICAL statements, all in Part B-II: a declared rescaling family on nested
six-agent cycles fails the semigroup law, with total variation 0.042 between direct and staged blocking
(\SB{rem:rg-compatibility-computation}); a frozen-attention projection scheme on a six-cycle and a
nine-cycle reaches fixed structures with vanishing pairwise block
(\SB{rem:rg-passive-blocking-computation}); and a row-regeneration scheme on the same instances keeps a
nonzero pairwise block (\SB{rem:rg-regenerated-attention-computation}). The schemes are declared
approximations, not the exact step, and the supplement says nothing here shows the exact pushforward
maintains interaction. The implementation has been retired and the values were not recomputed
(\SB{sec:rg-thermo-numerical-provenance}). The other rows of that index are exact rational checks of lattice
statements that rest on their proofs.

\section{Part B-III: the lattice cluster}
\label{comp:sb-lattice}

\subsection{The model and the flow}
\label{comp:sb-lattice-flow}

The source law lives on the edges of the square lattice, with matter signs \(\phi_x\), raw links \(\psi_e\),
and dressed signs \(z_e=\phi_x\psi_e\phi_{x+e_\iota}\). On a finite region with exterior \(\eta\) its weight
relative to counting measure is
\[
r^\eta_\Lambda(z)=\prod_{e\in\Lambda}(1+bz_e)\prod_{p\cap\Lambda\neq\varnothing}\bigl(1+t\chi_p(z\eta)\bigr),
\qquad b=\tanh h,\quad t=\tanh j
\]
(\SB{def:rg-spatial-source}). Since \(e^{hz}=\cosh h(1+bz)\) and \(e^{j\chi}=\cosh j(1+t\chi)\), this is
\(\exp(h\sum_ez_e+j\sum_p\chi_p)\), the \(\mathbb Z_2\) gauge theory of \citet{Wegner1971} coupled to a
\(\mathbb Z_2\) Higgs field as in \citet{FradkinShenker1979}, in unitary gauge: the first product is matter
hopping, a field on each link, and \(\chi_p\) is the plaquette holonomy. Small \(t\) is strong gauge
coupling, where \citet{OsterwalderSeiler1978} proved existence and analyticity. For \(0<b<1\) and \(0<t<1/6\)
the Gibbs law is unique and translation invariant, the matter signs are independent fair coins, and the
law depends on them only through \(z\) (\SB{thm:rg-spatial-source-law}).

The channel \(C_B\) keeps every \(B\)-th matter sign and replaces each coarse link by the product of the
\(B\) fine links along the straight path it covers; every fine link off the coarse grid lines is integrated
out. It is gauge equivariant, composes exactly, \(C_aC_B=C_{aB}\), and a coarse plaquette holonomy is the
product of the fine ones inside it (\SB{prop:rg-spatial-path-square}). The central theorem
(\SB{thm:rg-spatial-constrained-specification}) says that for \(0<b<1\), \(0<t<1/12\), and every \(B\), the
pushed law has a unique positive specification, defined and continuous at every coarse exterior including
null ones, whose dependence on the exterior beyond radius \(R\) decays like \(B^2(R+1)e^{-\tau(R+1)}\), and
that the domain \(12t<1\) does not depend on \(B\). The renormalized measure is Gibbsian, with a
renormalized Hamiltonian, at every block size on one fixed window, so the non-Gibbsianness catalogued by
\citet{VanEnterFernandezSokal1993} does not occur there. The generated interaction is complete with norm tending
to zero in \(B\) (\SB{thm:rg-spatial-generated-interaction}); for the integer factors \(a\in\{2,3,6\}\), on a
ball of character coefficients, the renormalization map \(\mathscr R_a\) is Fréchet \(C^1\) with contracting
linearization, so zero is its only fixed point there (\SB{thm:rg-spatial-local-operator}); and for a
compact box of rational source parameters and rational norm parameters, the source law enters that ball once
\(B\) exceeds an explicitly computable \(B_0\) (\SB{thm:rg-spatial-basin-entrance}).

The flow is to the trivial fixed point, and this should be read plainly. As \(B\) grows the blocked laws
converge on fixed cylinders to the fair product law, for every Gibbs law at any finite couplings
(\SB{prop:rg-spatial-cylinder-attraction}); informally, because a long Wilson-line parity multiplies many
signs whose conditional means stay away from \(\pm1\). Interacting fixed points, critical exponents, and a continuum
gauge-field measure lie outside the construction (\SB{open:rg-spatial-generality}). What is nontrivial is the
control: Gibbsianness uniform in block size, a complete interaction with explicit bounds, and a
differentiable map, all constructed. At each fixed depth the Fisher information per coarse cell is a
\(C^2\) density that is a regular metric on a nonempty open region (\SB{thm:rg-spatial-fisher-density}), and
\(\mathcal I_{aB}\preceq a^2\mathcal I_B\): per fine cell information decreases, the data-processing
inequality (\SB{prop:rg-spatial-fisher-contraction}).

\begin{analogy}{exact}
The source is the \(\mathbb Z_2\) gauge--Higgs model in unitary gauge, and the straight-path channel is
blocking of gauge links by straight Wilson lines along the coarse grid, a product map rather than
decimation.
\end{analogy}

\begin{analogy}{structural}
The open-ball construction is of the kind \citet{Israel1981} gave for decimation and Kadanoff maps of
Ising-spin interactions at high temperature; here the channel and the interaction space are different, and
the supplement separates its results from those maps.
\end{analogy}

\subsection{Pressure, label selection, and the specific VFE}
\label{comp:sb-pressure}

\SB{ch:geometry-pressure} places a finite catalog of such laws on one lattice; each label is a value of the
structural mark of the main manuscript's partition catalog (\MT{def:rg-geometry-ensemble},
\cref{comp:ch12}). Pressures exist with explicit surface rates, \(\lvert\mathfrak p-n^{-2}\log
Z_n\rvert\leq4\mathsf m_1/n\) (\SB{thm:rg-pressure-boundary}); existence is standard \citep{Georgii1988} and
the rates are the addition. Under global normalization the label with the largest bulk score takes all the
weight exponentially fast in the volume, and at a tie the outcome is decided by subextensive terms, not by
pressures (\SB{thm:rg-pressure-selection}). Forgetting the label can destroy Gibbsianness outright: a mixture
of two Bernoulli product laws has no continuous conditional at any exterior (\SB{thm:rg-geometry-erasure}).

Observation is handled for one record: every edge sensor reports \(+\)
(\SB{hyp:rg-pressure-positive-record}). It shifts each branch field from \(b_\gamma\) to
\((b_\gamma+\eta_\gamma)/(1+b_\gamma\eta_\gamma)\), and the specific VFE of a label-retaining recognition
satisfies \(\min_{\mathbb Q}\mathfrak f=-\bar e\), attained on the maximal posterior-score labels with the
shifted Gibbs laws (\SB{prop:rg-pressure-vfe-minimum}): the Gibbs variational principle in the
thermodynamic limit, with a finite label. Along a fixed blocking word
\(B^2\mathfrak f_{\mathrm{fine}}=\mathfrak f_{\mathrm{coarse,induced}}+\mathfrak d_{\mathrm{fib}}\) with
\(\mathfrak d_{\mathrm{fib}}\geq0\) (\SB{thm:rg-pressure-induced-vfe}), the volume-limit form of the main
manuscript's exact coarse free-energy identity. A general record produces random fields of both signs;
its quenched evidence density is \SB{open:rg-pressure-quenched-record}.

\subsection{Thermodynamic Fisher geometry}
\label{comp:sb-thermo-fisher}

On compact source boxes with \(h>0\) and \(0<j<\operatorname{atanh}(1/12)\), the native pressure is \(C^4\),
and since the native law is an exponential family in \(\mathsf g=(h,j)\) its Fisher density is the Hessian
of the log-partition density, \(\tfrac1NI\to\mathsf H=\nabla^2_{\mathsf g}\mathsf A\) in \(C^2\)
(\SB{thm:rg-thermo-native-cfour-pressure}, \SB{prop:rg-thermo-native-ctwo-metric}); pushes along a fixed word
have the same \(C^2\) density per coarse cell as the infinite-law marginals
(\SB{thm:rg-thermo-native-free-push-ctwo}). The main manuscript's finite-mixture asymptotics
(\MT{sec:cg-retained-label-asymptotics}) then hold with \(N=L^2\) cells
(\SB{cor:rg-thermo-native-label-limits}): with a strict winner the Fisher density converges to the winner's;
at a tie \(I/N^2\to\operatorname{Cov}_w(\mathsf v_\gamma)\), the covariance of the bulk-score slopes, which is
information about the label; and in a window of width \(1/N\) the label law is logistic.

\begin{analogy}{exact}
For an exponential family the Fisher metric is the Hessian of the log-partition function, the thermodynamic
metric of fluctuation theory \citep{Crooks2007}, per cell the susceptibility matrix.
\end{analogy}

\begin{analogy}{structural}
The tie and window statements have the form of finite-size rounding at a first-order transition, which
\citet{BorgsKotecky1990} derive from the lattice model itself. Here the phase mixture is prescribed through the
label weights, which makes the statement elementary; the correspondence is one of form.
\end{analogy}

\subsection{The typing question and the thermodynamic-limit reading}
\label{comp:sb-typing}

No lattice statement declares what a site is (\SB{open:rg-lattice-typing}). In the first reading the sites
are agents at \(c_*\), \(G=\mathbb Z_2\), the raw link is an independent link, the dressed sign is the per-edge
twist, and a plaquette is the holonomy of an agent four-cycle. In the second the sites are points of a
discretized base, the raw link is the transport of a \(\mathbb Z_2\) connection, a plaquette is a discrete
curvature, and the channel is a scale map over a base map. Every statement holds under either reading.

The first reading is the reader's own: links as lattice gauge theory over a zero-dimensional base, and
renormalization at \(c_*\) as a thermodynamic limit. Under it the lattice cluster is the one place in the
three documents where a thermodynamic limit of a link law on an agent graph is actually constructed. The
open problem
states what is missing before the reading can be adopted: the lattice law must be the link marginal of a
normalized joint with belief and model channels, a recognition law, and an ELBO on a countably infinite
agent set, whereas the main manuscript's multigraph is finite and the lattice family carries none of these.
The lattice part also fixes every blocking word before its volume limit; joint growth of depth and volume
is part of \SB{open:rg-spatial-generality}, and where depth and volume do grow together the limit is the
trivial one. Even under the first reading the group is finite and abelian, so the per-edge twist is gauge
invariant and nonabelian holonomy classes never appear, far from the target of 512 or more agents with
\(\GL(K,\R)\), \(\GL(K,\mathbb C)\), or \(\mathrm{SU}(N)\) links; the domain is a single-phase strong-coupling
region, with a unique Gibbs law and a smooth pressure; and observation enters through one constant record. The lattice cluster proves that the
coarse-channel calculus can be run to completion in infinite volume for one abelian family. It is not
evidence about formation in the target class.

\subsection{What is new, and the classical antecedents}
\label{comp:sb-antecedents}

The model and tools are classical: \citet{Wegner1971}, \citet{FradkinShenker1979}, and
\citet{OsterwalderSeiler1978} for the model; Dobrushin uniqueness \citep{Dobrushin1968} and boundary
comparison \citep{Georgii1988}; Kolmogorov extension; Kozlov's Gibbs representation \citep{Kozlov1974};
Israel's high-temperature renormalization maps \citep{Israel1981}; and \citet{VanEnterFernandezSokal1993}
for Gibbsian pathologies, of which \SB{thm:rg-geometry-erasure} is a Bernoulli case in sharper form. The
principal candidate for a new result is the block-size uniformity of \(12t<1\) in
\SB{thm:rg-spatial-constrained-specification}: the surveyed high-temperature results concern decimation,
Kadanoff, and averaging maps of Ising-spin interactions, not this gauge--matter path channel, and the proof
avoids a block estimate whose constant grows like \(Bt\). The literature check is limited to the works the
chapter cites, and priority is not claimed.

\section{What is proved, and how solid it is}
\label{comp:sb-scorecard}

Established no-go statements and separating counterexamples count as solid. The scope column says how far
each result's generality is from the program's target.

\begin{scorecard}
\SB{prop:geo-relative-frame}, \SB{prop:geo-identified-links} & A unique relative frame; one shared link
means \(\Theta^m_e=h_i^{-1}\Theta^b_eh_j\). & \SOLID & Standard in substance (torsor). \\
\SB{prop:geo-four-notions} & Bundle triviality, graph flatness, connection flatness, and one-loop holonomy
are independent. & \SOLID & \(\mathrm U(1)\) and \(S_3\) witnesses; at \(c_*\) (a), (c), (d) hold trivially and
only (b) has content. \\
\SB{hyp:geo-graph-base-transport} & A graph link is a base transport only by declaration. & \CHOICE &
Realization for independent links is open. \\
\SB{prop:gen-shared-link-stabilizer} & Tying a link across contexts leaves constant rechoices times a
central twist. & \SOLID & Admissible rechoices form a groupoid. \\
\SB{open:prob-continuum-theory} & No law on fields; evidence bounds are finite-design statements. & \OPEN &
No infinite product of reference measures. \\
\SB{thm:pb-pullback-gauge-invariance}, \SB{prop:pb-pullback-connection-change} & The base Fisher pullback is
passively covariant and changes with the connection by exact terms. & \SOLID & No connection-free base
metric; only a semimetric. \\
\SB{thm:pb-signed-base-comparison} & Fine minus coarse base Fisher is channel loss minus two anomaly
terms. & \SOLID & Exact; can be negative at zero loss. \\
\SB{thm:pb-isotropy-criterion} & The horizontal defect vanishes exactly when the defect form lies in the
isotropy algebra. & \SOLID & Pointwise; no compatible connection claimed. \\
\SB{thm:rg-cover-compatible-laws} & An ambient law gives compatible local data; availability parents glue.
& \COND & Holds by construction (functoriality) from one assumed ambient law; gluing from local data
open. \\
\SB{thm:rg-cover-descent-obstructions} & Sections and connections descend exactly under fiber constancy
and pullback agreement. & \SOLID & Standard in substance. \\
\SB{thm:rg-cover-operational-holonomy} & A declared record distinguishes two base-holonomy classes. &
\TOY & Base-loop instance of the MT witness, whose mechanism is general (\SOLID{} in \cref{comp:ch10}); one
loop, one connection family, Gaussian record; recovery beyond it is \SB{open:rg-cover-holonomy-recovery}. \\
\SB{prop:rg-finite-depth-no-asymptotics} & A finite network has bounded depth, no comparisons, and
schedule-dependent exponents. & \SOLID & A no-go for finite \(N\), not for \(N\to\infty\). \\
\SB{cor:rg-interaction-no-relevant-mode} & The exact step has no relevant direction; relevance needs an
extensive normalization. & \COND & Needs \(\lvert V_n\rvert/\varsigma\to0\): the bound is exponential in
\(\lvert V_n\rvert\), so it is silent when coarse levels keep many agents, and vacuous for finite vertex-ratio
blocking. Tempered comparisons are needed only for the reference-norm clause. \\
\SB{prop:rg-gaussian-dense-scaling} & Complete graph: \(A\mapsto bA\), \(W\mapsto b^2W\); dense
normalization contracts the self term by \(1/b\). & \SOLID & Exact; limiting ray is a boundary operator. \\
\SB{conj:grg-fixed-b-attraction} & Under a fixed primitive identification orbits converge; the fiber matrix
is not selected. & \COND & Standard in substance (Perron--Frobenius on a primitive \(B\)); existence of
the identification is open. \\
\SB{rem:rg-compatibility-computation}, \SB{rem:rg-passive-blocking-computation},
\SB{rem:rg-regenerated-attention-computation} & A rescaling family fails the semigroup law; frozen attention
kills the pairwise block, regeneration keeps it. & \NUM & Six- and nine-agent schemes; not
recomputed. \\
\SB{open:rg-two-index-limits} & Joint size and depth limits and universality for agent networks. & \OPEN &
Limits need not commute. \\
\SB{thm:rg-spatial-constrained-specification} & Every straight-path blocking of the \(\mathbb Z_2\)
gauge--matter law at \(12t<1\) is Gibbsian at every exterior, on a domain independent of block size. & \SOLID & Principal candidate for a new result; \(\mathbb Z_2\), strong
coupling. \\
\SB{thm:rg-spatial-generated-interaction}, \SB{thm:rg-spatial-local-operator} & Complete interaction tending
to zero; for \(a\in\{2,3,6\}\), a \(C^1\) map on a ball whose only fixed point is zero. & \SOLID & Israel-type
in form, different channel; trivial fixed point. \\
\SB{prop:rg-pressure-vfe-minimum} & Specific VFE minimum is minus the evidence density. & \COND & Only for the
constant all-positive record. \\
\SB{cor:rg-thermo-native-label-limits} & \(C^2\) Fisher densities; strict, tie (\(N^2\)), and \(1/N\)-window
limits. & \SOLID & Compact boxes; blocking word fixed before volume. \\
\SB{open:rg-lattice-typing} & Are lattice sites agents at \(c_*\) or base points? & \OPEN & No lattice
statement is yet about agents. \\
\end{scorecard}

\section{Conditions that matter}
\label{comp:sb-hygiene}

\begin{hygiene}
Substantive conditions, which could fail in the cases the program cares about: chosen connections
(\SB{def:geo-connections}), relative to which every base-direction statement is made; the declaration
\SB{hyp:geo-graph-base-transport}; conditional independence of design points (\SB{hyp:gen-design-product});
declared base measures, channel weights, and availability weights; constant rank of the covariant jet
(which can jump, as in the rank-jump witness of \SB{prop:pb-contact-null-counterexample}), integrable null directions, and basic tensors for quotient
geometries; chart joint convexity (\SB{hyp:hist-joint-convexity}),
which fails for Gaussian covariances; a strong metric on a declared manifold of fields
(\SB{hyp:hist-regular-section-space}), which the Fisher metric of the weak-metric witness is not; one
ambient law on a cover
(\SB{hyp:rg-cover-ambient-law}); declared comparison data, infinite blocking sequences, and tempered
comparisons for every asymptotic statement; one fixed primitive identification for attraction; and on the
lattice the group \(\mathbb Z_2\), small plaquette coupling with positive field, compact source boxes,
blocking words fixed before volume limits, controlled label growth (\SB{hyp:rg-pressure-label-growth}), and
the constant all-positive record (\SB{hyp:rg-pressure-positive-record}).

Hygiene a physicist can take for granted on a first reading: standard Borel spaces, measurability of section
responses (\SB{hyp:prob-context-section-response}), smoothness classes \(C^1\) to \(C^4\), paracompactness and
the constant-rank bookkeeping that restricts the defect classifications to regular strata (a technical
restriction, unlike the rank of the jet itself), integrability of the integrals, the submersion hypothesis of descent,
boundary conventions, and centering of squares in local limits.
\end{hygiene}

Toy content is limited and labeled: the four-agent cover, the holonomy witness, the scalar Hermite spectrum,
and the numerical six- and nine-agent schemes. The many small counterexamples of Part B-I are not toys;
each refutes a general implication and is a solid negative result. The lattice is infinite but special in
its group, abelian and finite, and in its single-phase coupling region. Its mechanism, keeping a renormalization map
Gibbsian uniformly in block size by a transport argument, is specific to strong coupling and unshown for
nonabelian or non-compact links; its constants (\(12\), \(\mathrm q\), \(\mathrm c_1\)) say nothing about the
target class.

\section{Open ends}
\label{comp:sb-open}

The supplement lists twenty open problems (\SB{tab:base-open}). For Part B-I, the unknowns are a probability
law on fields with continuum limits of the bound (\SB{open:prob-continuum-theory}); the orbit space of the
two-atlas coordinate action (\SB{open:gen-gauge-quotient}); a base geometry fixed by the model rather than a
chosen connection (\SB{open:pb-canonical-pullback-geometry}); whether declared recognition families have the
structure the history theorems use (\SB{open:hist-section-tier-admission}) and whether descending fields form
a smooth submanifold (\SB{open:hist-projectable-set}); gluing local laws on a cover
(\SB{open:rg-cover-local-gluing}); realizing independent graph links as base transports
(\SB{open:rg-cover-graph-base-identification}); recovering a connection from records
(\SB{open:rg-cover-holonomy-recovery}); quotients across rank changes (\SB{open:rg-cover-quotient-extension});
and informative averaging along a base map (\SB{open:rg-cover-base-aggregation}). Each is settled by a
construction for a declared class or by an obstruction.

For Part B-II, the unknowns are rescaling kernels on agent networks obeying the semigroup law
(\SB{open:rg-rescaling-compatibility}); a multiplicative ergodic structure for scale cocycles
(\SB{open:rg-multiplicative-ergodic}); relevance spectra of correlated or gauge-mixed Gaussian blocks
(\SB{open:rg-hermite-extensive-relevance}) and a basin for the Gaussian block law
(\SB{open:rg-hermite-nonlinear-attraction}); a canonical continuous beta (\SB{open:rg-continuous-beta}); joint
size and depth limits with a convergent scheme-invariant observable (\SB{open:rg-two-index-limits}); and a
fixed primitive identification for some graph family (\SB{open:rg-scalarized-attraction}). For Part B-III,
they are the typing (\SB{open:rg-lattice-typing}), settled by proving one reading; critical points,
depth-uniform Fisher floors, joint depth and volume limits, and self-consistent curvature actions
(\SB{open:rg-spatial-generality}); and the quenched evidence density for random records
(\SB{open:rg-pressure-quenched-record}), which needs every-exterior control of random fields of both signs.

Three gates carry most of the weight. Without \SB{open:prob-continuum-theory} there is no evidence bound over a
base. Without \SB{open:rg-cover-graph-base-identification} the agent graph's links and the base's connection
stay separate objects joined only by declaration. And \SB{open:rg-two-index-limits}, with
\SB{open:rg-lattice-typing}, gates renormalization as a thermodynamic limit: the typed statement exists, the
one constructed instance is an untyped abelian lattice at strong coupling, and no agent-network instance with
its recognition law exists.

\section{Verdict}
\label{comp:sb-verdict}

\begin{verdictbox}{\NARROWER}
\textbf{What it says it does.} It collects what the theory requires once contexts vary, networks become
infinite, and agents sit on a lattice, stating for each extension which single-point statements survive and
which additional data it consumes (\SB{ch:base-introduction}). It declines to identify the base with physical
space, graph links with base transports, or any continuum law, and declines asymptotic vocabulary for finite
networks.

\textbf{What the statements give.} Part B-I gives the kinematics of a base extension, with exact
identities and separating counterexamples: relative frames, covariant defects, four independent flatness
notions, the shared-link stabilizer, an exact fine-to-coarse base Fisher comparison with a sharp isotropy
criterion, durations on strong metrics, and descent. Part B-II gives a typing of Wilsonian vocabulary, no-go
and scheme-dependence results, exact Gaussian power counting, and one conditional attraction theorem. Part
B-III gives a complete constructive renormalization of one \(\mathbb Z_2\) family at strong coupling, with
pressure, label selection, specific VFE, and Fisher densities in the volume limit.

\textbf{The gap.} ``Base transport'' is transport by a chosen connection, not transport derived from
inference; base measures, weights, and ambient laws are declared; no law on fields exists; and graph links
become base transports only by declaration. ``Infinite networks'' names the apparatus an infinite network
would need; no infinite agent network with channels, recognition law, and ELBO is constructed, no exponent or
fixed point of the agent flow is established, and universality is a stated program. The lattice cluster is
untyped, abelian, confined to strong coupling, observed through one record, and flows to the trivial fixed
point. Each part does what its opening paragraph says; the supplement as a whole is a precise map of what base
transport, infinite networks, and spatial realization would require, not a theory of any of them.
\end{verdictbox}


\part{Back Matter}
\chapter*{Glossary}
\phantomsection\addcontentsline{toc}{chapter}{Glossary}
\markboth{Glossary}{}
\label{comp:glossary}

The technical documents use many words in senses of their own. This glossary gives each one a plain
meaning and, where one exists, the nearest physics analogue, with the companion chapter that explains
it. An analogue listed here is a pointer, not an identification: the chapter says whether the
correspondence is exact, structural, or loose.

\section*{Geometry and gauge at one base point}
\phantomsection\addcontentsline{toc}{section}{Geometry and gauge at one base point}

\begin{decoder}
Base point $c_*$ & The single context at which the whole main manuscript works (\cref{comp:ch01}). & A zero-dimensional base. \\
Frame $u_i$ & The agent's chosen point of the one principal fiber over $c_*$, used for both channels (\cref{comp:ch02}). & Local gauge choice; kinematically a group-valued Higgs field. \\
Passive re-expression $g_i$ & A per-agent change of frame; coordinates change, nothing physical does (\cref{comp:ch02}). & Local gauge transformation read as a change of trivialization. \\
Active gauge group & Automorphisms of the one shared fiber: a single copy of $G$ inside the product of per-agent groups (\cref{comp:ch02}). & A global symmetry. \\
Frame comparison $T_{ij}$ & The unique group element relating two agents' frames in the same fiber; no path involved (\cref{comp:ch02}). & Canonical pure-gauge link. \\
Link $\Theta^x_e$ & An independent group element on an oriented edge copy, per channel (\cref{comp:ch02}). & Lattice link variable. \\
Per-edge twist & A link written against one reference point; its conjugacy class is invariant (\cref{comp:ch02}). & A unitary-gauge (dressed) link. \\
Holonomy $\mathrm{Hol}^x(p)$ & Ordered product of links around a closed walk (\cref{comp:ch02}). & Wilson loop before the trace. \\
Graph flat; frame aligned & Every loop product is the identity; links equal to the frame comparisons (\cref{comp:ch02}). & Pure gauge; canonical pure-gauge background. \\
Law coherence & Agents' laws equal each other's transports on every internal edge, up to holonomy stabilizers (\cref{comp:ch02}). & Covariantly constant matter field. \\
Cross morphism $\Phi$ & A map between belief and model fibers; extra structure, not implied by the bundles (\cref{comp:ch02}). & Coupling between two matter species. \\
Tree transport, rooted holonomy & Link products along a spanning tree to a block root, and all loop products at the root (\cref{comp:ch10}). & Maximal-tree (axial) gauge fixing. \\
Block-dressed edge & A boundary link with tree transports attached at both ends (\cref{comp:ch10}). & Covariant link between block roots. \\
Connection Laplacian; fixed fiber & $\sum_e(z_i-\Theta_ez_j)^\top W_e(z_i-\Theta_ez_j)$ and its kernel (\cref{comp:ch10}). & Covariant lattice Laplacian and its zero modes. \\
Frame sector & The meta-agent's frame, a free orbit invisible to its members (\cref{comp:ch10}). & Pure gauge. \\
\end{decoder}

\section*{Probability, the joint law, and the evidence bound}
\phantomsection\addcontentsline{toc}{section}{Probability, the joint law, and the evidence bound}

\begin{decoder}
Record, record kernel $K_a$ & An observation or interaction factor attached to a set of agents, normalized as a conditional law (\cref{comp:ch03,comp:ch05}). & A bond or plaquette Boltzmann factor. \\
Evidence & Marginal likelihood of the records (\cref{comp:ch03}). & Partition function. \\
Recognition law; selected coupling & The population's approximate posterior, a joint law with the agents' local laws as marginals, chosen by a declared rule (\cref{comp:ch03,comp:ch05}). & Variational trial state, correlated in general. \\
Posterior slice & The exact posterior given the observation and the structural data (\cref{comp:ch01}). & Exact conditioned Gibbs state. \\
Law point & A belief treated as a latent random variable in its own right (\cref{comp:ch03}). & A one-body distribution treated as a field. \\
Evaluator & The map from a model presentation to a normalized generative kernel; a model is not itself a law (\cref{comp:ch03,comp:ch09}). & Coupling-to-Hamiltonian map. \\
Typing prohibition & No generative kernel may read a recognition law or the posterior (\cref{comp:ch04}). & A Hamiltonian independent of the trial state. \\
Extended evidence bound & $\log p(o)-\KL(\mathbb Q\Vert\boldsymbol\Pi)$, defined for every recognition law and equal to $-\infty$ on the singular branch (\cref{comp:ch05}). & Minus the Gibbs trial free energy. \\
Singular branch & The recognition law charges a set the posterior cannot see, so the bound is $-\infty$ (\cref{comp:ch05}). & A delta-function state on a continuum. \\
Total correlation & Relative entropy of a joint law from the product of its marginals (\cref{comp:ch05}). & Multi-information; the mean-field entropy error. \\
Local conditional free energy & A block's free energy with the rest of the population held fixed (\cref{comp:ch05}). & Finite-volume free energy with a boundary condition. \\
Exact block E-coordinate & Replacing a block's law by its exact local posterior (\cref{comp:ch05}). & Heat-bath update, in law. \\
Accepted M-coordinate & A parameter step that does not lower the bound at fixed recognition (\cref{comp:ch05}). & The M-step of expectation maximization. \\
Operational environment node & An environment rewritten as a node that emits records (\cref{comp:ch05}). & A boundary degree of freedom. \\
\end{decoder}

\section*{Attention and selective communication}
\phantomsection\addcontentsline{toc}{section}{Attention and selective communication}

\begin{decoder}
Source label (pointer) & A latent categorical variable naming which neighbor a receiver listens to (\cref{comp:ch06}). & A Potts-like hidden assignment. \\
Source prior $\pi_{ij}$; attention row $\beta_{ij}$ & Prior over sources, and the posterior (generative) or recognition row over them (\cref{comp:ch06}). & Prior weight; Boltzmann weight over sources. \\
Receiver occupancy $\alpha_i$; event law $\eta$ & Which receiver an event belongs to, and the joint receiver--source law $\eta=\alpha\beta$ (\cref{comp:ch06}). & Joint bond occupation. \\
Attention temperature $\tau_i$ & Scale of the record energy; in the Gaussian record only $W/\tau$ matters (\cref{comp:ch06}). & Canonical temperature. \\
Constant-row gap & The exact cost of using one attention row for every configuration (\cref{comp:ch06}). & Mean-field decoupling error. \\
Query and key maps $S^Q_e$, $S^K_e$ & Linear maps from the receiver's and the source's internal spaces into a common comparison space (\cref{comp:ch06}). & None standard. \\
Communication map $\mathsf L_e$ & Encoder, transmission, and decoder composed; may be singular, and it is not the link (\cref{comp:ch06}). & A general linear link field. \\
Value map & The matrix by which the exact receiver update moves toward a source: a Kalman gain times the key map (\cref{comp:ch06}). & Gain times observation matrix. \\
\end{decoder}

\section*{Information geometry and dynamics}
\phantomsection\addcontentsline{toc}{section}{Information geometry and dynamics}

\begin{decoder}
Fisher metric & Covariance of the score (\cref{comp:ch07}). & Susceptibility matrix. \\
History & An oriented curve of configurations with no time label (\cref{comp:ch07}). & A trajectory without its clock. \\
Fisher duration (information clock) & Fisher arc length along a history, invariant under reparameterization (\cref{comp:ch07}). & Thermodynamic length. \\
Oriented semiconjugacy & A coarse map carrying the fine flow to the coarse flow up to an increasing time change (\cref{comp:ch07}). & Exact dynamical coarse-graining. \\
Parity amplitude $\varkappa$ & Strength of a six-bit correlation invisible to every proper marginal (\cref{comp:ch07}). & A six-spin coupling. \\
Fast and slow blocks & A declared split of coordinates by rate (\cref{comp:ch08}). & Electrons and nuclei. \\
Reduced functional $F_{\mathrm{red}}$ & The free energy minimized over the fast block (\cref{comp:ch08}). & Born--Oppenheimer surface. \\
Instantaneous refresh & Exact resampling of the fine state given the parent at every step (\cref{comp:ch08}). & Adiabatic elimination at infinite separation. \\
Partial refresh & Resampling with probability below one, which makes the parent non-Markov (\cref{comp:ch08}). & Finite separation with a memory kernel. \\
\end{decoder}

\section*{Coarse-graining}
\phantomsection\addcontentsline{toc}{section}{Coarse-graining}

\begin{decoder}
Coarse channel & A Markov kernel on latent coordinates that reads neither recognition nor observation (\cref{comp:ch09}). & Kadanoff block-spin weight. \\
Energy precomposition; Galerkin restriction & Substituting a restricted configuration into an energy; compressing an operator (\cref{comp:ch09}). & Constrained configurations; the multigrid coarse operator. \\
Discarded conditional information & Expected relative entropy between recognition and posterior inside the blocks, given the coarse value (\cref{comp:ch09}). & Information lost inside blocks. \\
Pointwise parent datum & The pushforward of the generative, posterior, and recognition laws of a block through one channel (\cref{comp:ch09}). & The block-spin ensemble. \\
Extensive labeled family & A mixture over finitely many labels whose log masses grow with the number of agents (\cref{comp:ch09}). & A multiphase ensemble; ties are phase coexistence. \\
Natural selector & A relabeling-natural, gauge-compatible rule choosing a partition from the data; open (\cref{comp:ch09}). & A choice of blocking scheme. \\
Holonomy-conditioned score & Least weighted relative entropy from the transported marginals to a holonomy-invariant parent (\cref{comp:ch10}). & A twirled barycenter. \\
Effective action & Minus the log density of the pushed evidence measure (\cref{comp:ch11}). & Block Hamiltonian; integrating out. \\
Hoeffding coordinates & One interaction component per set of agents, centered in each agent (\cref{comp:ch11}). & Cluster or ANOVA expansion. \\
Hyperedge & An interaction component on three or more agents, generated by coarse-graining (\cref{comp:ch11}). & A multi-spin coupling from decimation. \\
Bounded tier & The regime of essentially bounded interactions in which the coupling calculus is proved (\cref{comp:ch11}). & Finite or compact spins. \\
Native growth bound & The multiplier norm is at most $3^{|V_n|}-1$ in the target's size (\cref{comp:ch11}). & No growth without rescaling. \\
Finite-network obstruction & Coupling spaces cannot be identified across a step that reduces the number of agents (\cref{comp:ch11}). & Finite-volume renormalization terminates. \\
Strong lumpability & A coarse process that is Markov for every initial law (\cref{comp:ch11}). & Kemeny--Snell lumpability. \\
Memory kernel; finite-memory defect & Exact non-Markov corrections, and the relative-entropy cost of a memory-$k$ approximation (\cref{comp:ch11}). & Nakajima--Zwanzig. \\
Membership kernel; structural mark & Soft assignment of an agent to a block, and the random partition a structural channel produces (\cref{comp:ch11,comp:ch12}). & Fuzzy block spin; random blocking. \\
Declared root & A parentless top latent with a supplied prior (\cref{comp:ch12}). & A boundary condition at the top. \\
Lifted joint & Parameters and record copies made latent so that delayed feedback is one joint law (\cref{comp:ch12}). & State augmentation in filtering. \\
\end{decoder}

\section*{Formation}
\phantomsection\addcontentsline{toc}{section}{Formation}

\begin{decoder}
Coherence episode & Coherence held by a named restoring interaction and observed on a horizon that separates it from equilibrium (\cref{comp:ch13}; rule R15). & A dynamical ordering event. \\
Restored equilibrium coherence & Coherence the stationary law already has with probability tending to one; not formation (\cref{comp:ch14}). & Equilibrium order. \\
Supplied-scale sequence & A second coalescence that happens regardless of the first stage (\cref{comp:ch13}; rule R14). & A hierarchy put into the Hamiltonian. \\
Frozen arm & The same model with learning switched off; the null for a selection claim (\cref{comp:ch14}; rule R13). & A control without feedback. \\
Energy strengths $\mathrm J_\ell=\epsilon^{-\mathrm Q_\ell}$ & The supplied stiffness of each level of the formation chain (\cref{comp:ch13}). & A coupling hierarchy. \\
Preparation & The chain's initial state: a displaced common mean with thermal spread (\cref{comp:ch13}). & A Gibbs state in an applied field switched off at time zero. \\
Marked parent & The pushforward onto the retained space, keeping every label and loop (\cref{comp:ch13}). & Exact decimation. \\
Source gating & An edge carries energy only while one of its endpoints attends to it (\cref{comp:ch13}). & Bonds slaved to labels. \\
Sheet & The slice of state space at one fixed attention assignment (\cref{comp:ch13}). & A sector. \\
Minimum edge cover & The fewest edges touching every agent; the stationary blocks of the gated family (\cref{comp:ch13}). & None. \\
Matching seed & An initial assignment that is a perfect matching, with the initial means in a small ball about one of two specified configurations $\pm x_\circ$ (\cref{comp:ch14}). & A favorable initial condition. \\
Source learning, mobility $\epsilon_{\mathrm L}$ & Prior rows relax toward the recognition marginal at a slow rate (\cref{comp:ch14}). & A slow variable. \\
Transient imprint & The time-integrated early asymmetry that learning amplifies (\cref{comp:ch14}). & A vanishing symmetry-breaking field. \\
Discarded conditional relative entropy & Mean relative entropy of the coordinates a parent drops, given the ones it keeps (\cref{comp:ch13,comp:ch14}). & Error of adiabatic elimination. \\
Activation clock & The time at which a parent is switched on (\cref{comp:ch14}). & Waiting for fast modes to thermalize. \\
\end{decoder}

\section*{The Gaussian realization}
\phantomsection\addcontentsline{toc}{section}{The Gaussian realization}

\begin{decoder}
Interaction family & Precision $\Lambda_{ii}=A_i+\sum_jW_{ij}$, $\Lambda_{ij}=-W_{ij}$ with positive semidefinite blocks (\cref{comp:sg1}). & Massive matrix-weighted lattice Laplacian. \\
Self term $A_i$ & The on-site block of the precision (\cref{comp:sg1}). & Mass or pinning term. \\
Gain absorption & Frames that turn every transition on a forest into the identity (\cref{comp:sg1}). & Axial gauge. \\
Orientation criterion & $(-1)^K\det\Lambda_{12}>0$: when a pair's coupling can be gauged to the family (\cref{comp:sg1}). & A ferromagnetic bond sign ($K=1$). \\
Determinant gap & The loss of the bound from factorizing a Gaussian by blocks (\cref{comp:sg1}). & Mean-field free-energy error. \\
Interaction pencil & The pair $(L,\Lambda)$ and its generalized eigenvalues (\cref{comp:sg1}). & A normal-mode problem. \\
Hard aggregation; Galerkin law & Identifying each cluster with one coordinate; the fine law conditioned on zero disagreement inside clusters (\cref{comp:sg2}). & Rigid-block Kadanoff block spin; the algebraic-multigrid operator. \\
Kron reduction & Integrating out hidden agents by a Schur complement (\cref{comp:sg2}). & Integrating out; star--mesh transform. \\
Internal residue & The term $(I-\Theta)^\top W(I-\Theta)$ a twisted internal edge leaves on a cluster (\cref{comp:sg2}). & Mass generated by a flux. \\
Operator arrow & Aggregation followed by a declared rescaling of the coarse precision (\cref{comp:sg2}). & Field rescaling of a renormalization step. \\
Declared repair & A pin, quotient, or retained mass at a singular member of the family (\cref{comp:sg1,comp:sg2}). & Infrared regulator. \\
Reciprocal fold & Two agents joined by opposite edges with no anchor; its precision is singular when flat (\cref{comp:sg2}). & Two masses on springs with no ground. \\
\end{decoder}

\section*{Base transport and the lattice cluster}
\phantomsection\addcontentsline{toc}{section}{Base transport and the lattice cluster}

\begin{decoder}
Relative frame $h_i$ & The group element comparing belief and model frames when they differ (\cref{comp:sb}). & A bifundamental field of two gauge sectors. \\
Covariant defect & The failure of a belief-to-model map to commute with parallel transport along the base (\cref{comp:sb}). & A covariant derivative of a Hom-valued field. \\
Four notions of flatness & Bundle triviality, graph flatness, connection flatness, trivial holonomy on a loop; at $c_*$ only graph flatness has content (\cref{comp:sb}). & Topological sector, pure gauge, zero field strength, trivial Wilson loop. \\
Fisher pullback & Squared Fisher speed of a belief along the base, relative to a chosen connection (\cref{comp:sb}). & Gauged sigma-model kinetic metric. \\
Comparison data & A declared identification of each blocked theory with a reference (\cref{comp:sb}). & The rescaling step of a Wilsonian flow. \\
Relevant, marginal, irrelevant & Sign of a growth rate per logarithmic block scale, relative to declared norms (\cref{comp:sb}). & Sign of a renormalization-group exponent. \\
Specification at every exterior & Coarse conditional laws continuous at every boundary condition (\cref{comp:sb}). & Gibbsianness of the renormalized measure. \\
Specific VFE & Variational free energy per cell in the volume limit (\cref{comp:sb}). & Free-energy density. \\
\end{decoder}

\chapter*{Open Ends of the Theory}
\phantomsection\addcontentsline{toc}{chapter}{Open Ends of the Theory}
\markboth{Open Ends of the Theory}{}
\label{comp:open-ends}

The main manuscript lists its open problems in one table (\MT{app:claim-ledger}), each with a
settling condition: a class of models and two outcomes, a proof or a counterexample. The individual
chapters of this companion explain them where they arise. This chapter groups them by the question
a physicist would ask, says in which class each one is posed, and ends with the questions the program
cares about that the technical documents do not pose at all.

One feature of the settling conditions shapes everything below. Rule R16 of the authoring
specification requires both outcomes to be nontrivial within a stated class, so the problems are
posed in classes narrow enough to be decidable. Several of the central ones are posed on the
four-cycle kinetics of \MT{def:rg-formation-cycle-experiment} or on switching diffusions of that
type. Settling such a problem settles it for that class; it would ground the question at small size
but would not answer it for a network of hundreds of agents.

\section*{The parent's own dynamics}
\phantomsection\addcontentsline{toc}{section}{The parent's own dynamics}

Exact coarse-graining produces the law of a parent at every instant
(\MT{thm:rg-pointwise-parent-datum}). Whether that law evolves as an autonomous process, a parent
with dynamics of its own, is \MT{open:rg-pointwise-parent-dynamics}, which the introduction calls the
defining open problem of the phase. It asks whether a time-homogeneous Markov parent, or one carrying
a finite filter from a declared regular class, reproduces the native path law of the retained
coordinates within a total variation that vanishes as $\epsilon\to0$, both on windows of fixed length
and over the complete history. It is posed for switching diffusions of the four-cycle type. The
four-cycle continuation results of \cref{comp:ch14} answer special cases on the positive side: a
parent activated at a suitable time, after a coherence prefix, has vanishing discarded information and
a retained path law close to the native one on short windows. They do not settle the problem, because
they fix the activation event and the horizon. \MT{open:obs-selective-formation} asks the same
question for a block that communicates through learned low-rank interfaces: is there a parent of fixed
lower dimension whose path law of the communicated coordinates is asymptotically exact on the learning
time scale?

\section*{Formation that is not supplied}
\phantomsection\addcontentsline{toc}{section}{Formation that is not supplied}

The formation theorem of \cref{comp:ch13} relaxes a prepared state through a hierarchy of interaction
strengths supplied by hand. Four open problems ask for more. \MT{open:rg-formation-regime-sharpness}
asks whether the doubling recursion $\mathrm Q_\ell=2\mathrm Q_{\ell+1}+20$ of the strengths can be
relaxed to inner base $4/3$, a separation that is still doubly exponential in the depth, with constants
independent of the depth. \MT{open:rg-formation-adaptive-family} asks, on the
four-cycle, for parents with vanishing conditional and dynamical errors together with a margin,
bounded below, by which learning or the loop coordinate of the initial mean changes the law of the
first formed block; or for a proof that for almost every initial mean both differences vanish.
\MT{open:rg-formation-general} asks for repeated endogenous formation on general finite graphs with
temperature-independent data: a second coalescence that depends on the first by a margin, is held by
a restoring interaction, is observed on dynamical horizons, and has well-behaved parents; or a proof
that none exists in that class. \MT{open:full-graph-hierarchy} asks for the full package: a lineage of
at least two levels formed with probability bounded below from every initial law with a density and a finite
second moment, an intrinsic
partition selector, and a slow manifold on which the upper level feeds back consistently. These are
the questions on which the program's claim that collectives form through coupled inference rests.
None of the four takes a limit of many agents; the last three are posed at fixed finite size in the
low-temperature limit, while the regime-sharpness problem asks for constants uniform in the depth,
and hence in the number of agents of the family.

\section*{Selecting structure}
\phantomsection\addcontentsline{toc}{section}{Selecting structure}

No functional analyzed in the documents selects a nontrivial partition into meta-agents: the costs
studied either prefer an endpoint (every agent separate, or everything merged) or need an external
coefficient. \MT{open:restrict-intrinsic-criterion} and \MT{open:cg-natural-selector} ask for a
functional of the laws, natural under relabeling and compatible with the gauge action, not monotone
under refinement, with a nondegenerate optimizer, or for an impossibility theorem in a stated class.
\MT{open:rg-invariant-retained-subspace} asks whether deterministic block channels generically admit
retained coordinates beyond the trivial families of \cref{comp:ch11}. Structural dynamics of the graph
itself is open too: \MT{open:obs-edge-creation} asks for a normalized joint in which edges are created
and deleted by a threshold rule while the collective free energy does not increase, and
\MT{open:obs-multihead-attention} asks for an exact recursion for multi-head attention at the level of
source labels.

\section*{Convergence and time scales}
\phantomsection\addcontentsline{toc}{section}{Convergence and time scales}

The alternating scheme of inference and learning is not known to converge
(\MT{open:elbo-alternating-convergence}). The reduction of fast beliefs onto a slow manifold is proved
for held assignments; tracking along assignment paths with jumps is \MT{open:obs-slow-manifold},
posed on the four-cycle. Delayed feedback across learning episodes has an exact joint law, and whether
the per-episode accounting reproduces the aggregate update to first order is
\MT{open:full-graph-episode-accounting}. How a meta-agent's information clock compares with its
constituents' is \MT{open:hist-meta-clock-ratio}.

\section*{Foundations}
\phantomsection\addcontentsline{toc}{section}{Foundations}

Five open problems concern the measure-theoretic and geometric footing: likelihood families with no
common dominating measure (\MT{open:prob-nondominated-likelihood}), a normalized prior over law points
(\MT{open:prob-law-point-prior}), equivariant posterior versions for noncompact stabilizers
(\MT{open:cg-equivariant-versions}), global Fisher quotients (\MT{open:hist-retained-fisher-quotient}),
and correlated lifts beyond the finite parity witness (\MT{open:hist-parity-lift-scope}). The
equivariant-versions problem is open only when a law's stabilizer is noncompact and its posterior has
no weakly continuous version (\MT{prop:cg-equivariant-bayes-versions},
\MT{cor:cg-continuous-version-equivariant}). Noncompactness of $\GL(K)$ alone does not put the
program's Gaussian laws in the open case: under a faithful linear representation a nondegenerate
Gaussian law has a compact stabilizer, conjugate into $\mathrm O(K)$, so the open case concerns
degenerate laws, nonfaithful actions, and non-Gaussian fibers.

\section*{Dynamic links and their coarse laws}
\phantomsection\addcontentsline{toc}{section}{Dynamic links and their coarse laws}

The existence of a normalized generator prior is settled for every Lie group
(\MT{prop:link-generator-prior}); no compactness or finite Haar volume is needed. What remains
open is a finite-dimensional family of link laws that survives coarse multiplication
(\MT{open:link-coarse-closure}). For $\GL^+(K)$, $K\geq2$, the stated family must contain
the exponential images of all nondegenerate Gaussian generator laws and be closed under
convolution, with a weakly continuous injective parameterization by a connected open subset of a
finite-dimensional Euclidean space. The exact coarse law already exists; products can leave the
real exponential image, and local chart generators need not be Gaussian
(\MT{prop:link-coarse-generator-law}; \cref{comp:ch12}). This open item concerns useful
closure, not the existence of a parent pushforward.

Supplement~G poses two further questions (\cref{comp:gaussian-latent-links}). For unanchored
undirected Gaussian matter on connected multigraphs of cycle rank at least two, positive-definite
weights, zero information vector, and a product nondegenerate Gaussian generator base, is the
joint normalizer finite at a stated $K\geq2$ (\SG{open:glink-unanchored-normalization})?
The scalar case is finite, while an unanchored single cycle diverges. For the unregularized
expected comparison-record energy at $K\geq4$ with negative-trace rank-one optimum, can a chart
link having a nonsemisimple negative eigenvalue attain the chart infimum
(\SG{open:glink-chart-infimum})? Nonattainment is already proved for $K=2,3$ and for the
semisimple-negative class. Neither question is the earlier law-point-prior problem.

\section*{Contact with the world}
\phantomsection\addcontentsline{toc}{section}{Contact with the world}

Two items concern testability rather than mathematics. \MT{claim:physical-law-identification} asks
for held-out evidence, under a protocol fixed in advance, that an object of the coarse-graining
matches a named target system beyond a baseline. \MT{open:phil-testability} asks for a test at one
base point that could falsify nontrivial loop holonomy or the formation of a block, or a proof that
the records cannot distinguish them. The operational holonomy theorem of \cref{comp:ch10} shows that
records sensitive to network-loop holonomy exist in principle.

\section*{The supplements}
\phantomsection\addcontentsline{toc}{section}{The supplements}

Supplement~G keeps its own open rows: the scaling of a retained mass sector
(\SG{open:exp-mass-sector}), commutation of the information projection with coarse maps
(\SG{open:exp-projection-commutation}), monotonicity of Galerkin aggregation beyond one fixed precision
(\SG{open:grg-galerkin-monotonicity}), robustness of the anchored-star update
(\SG{open:obs-star-update-robustness}), Kron classification beyond finite encodings
(\SG{open:cg-kron-classification-beyond-finite}), a bridge to Bayesian renormalization
(\SG{open:grg-bayesian-bridge}), transfer of a model-space Fisher scale
(\SG{open:ig-model-space-fisher-transfer}), and stochastic inversion of aggregation
(\SG{open:grg-stochastic-inverse}), and the two link-sector questions described above.
\Cref{comp:sg1,comp:sg2,comp:gaussian-latent-links} explain them. Supplement~B holds the
open problems of the deferred program, among them continuum law theory
(\SB{open:prob-continuum-theory}), graph-to-base identification
(\SB{open:rg-cover-graph-base-identification}), compatible rescaling kernels on agent networks
(\SB{open:rg-rescaling-compatibility}), two-index limits and universality
(\SB{open:rg-two-index-limits}), and the typing of lattice sites (\SB{open:rg-lattice-typing});
\cref{comp:sb} explains them.

\section*{Questions the documents do not pose}
\phantomsection\addcontentsline{toc}{section}{Questions the documents do not pose}

Some questions the program cares about appear in none of the open-problem tables, either because
they were deferred by decision or because no one has yet posed them in a decidable form. Listing them
is not a criticism of the documents, which state their scope; it is a map of where the theory stops.

None of the main manuscript's open problems takes the limit of many agents. The limits taken are low
temperature, slow learning, and small separation parameters; two problems concern size without a
limit (uniformity in the depth for the regime-sharpness problem, and the dependence of a clock ratio
on the number of agents in \MT{open:hist-meta-clock-ratio}). The main manuscript's only large-size
statements are the retained-label Fisher theorems of \cref{comp:ch09}. They assume that the log
masses of a finite labeled mixture grow in proportion to an extensive size such as the number of
agents, and derive the Fisher geometry from that assumption; near a tie between labels the natural
coordinate is $N(\theta-\theta_0)$ (\MT{thm:rg-thermo-label-crossover}), the form of the finite-size
rounding of a first-order transition, here in a retained label. They do not derive the extensive
growth from a network of interacting agents, and they say nothing about formation. A thermodynamic
limit of the agent network at one base point, with nonabelian and noncompact links, is therefore
neither proved nor posed there; Supplement~B holds the infinite-network and asymptotic apparatus.
For the same reason, phase transitions, condensation, and critical slowing down are neither ruled in
nor ruled out as formation mechanisms: at finite size, with a unique law-level dynamics, they cannot
occur in the sense of statistical mechanics, and the main manuscript contains no statement about
formation in the limit in which they could.

A normalized law on noncompact links is supplied by the Gaussian-reference generator construction
of \cref{comp:lie-links}. The Haar-reference law remains outside the noncompact scope of
\MT{ch:geometry}; the chosen generator metric and prior are substantive model data. The
holonomy-conditioned parent score of \cref{comp:ch10} is infinite almost surely on nonforest
connected blocks under absolutely continuous $\mathfrak{gl}(K)$ generator laws
(\MT{cor:link-holonomy-parent-almost-sure}). That score is one particular construction, a projection of
marginal laws onto a holonomy-fixed family, not the exact coarse channel, whose parent law always
exists. Using this fixed nondegenerate Gaussian parent family on the stated $\GL^+(K)$ blocks
requires restricted generator support or a different parent family. Compact-unitary links
are outside that negative result.

Base-manifold transport, the setting of the program's original picture of agents spread over a base
with attention varying smoothly over overlaps, is deferred by the author's decision until the
single-point theory is settled; its apparatus and open problems are in Supplement~B.


\bibliographystyle{plainnat}
\bibliography{references}
\end{document}
